\documentclass[11pt]{article}

\usepackage[margin=1.1in]{geometry}

\usepackage{amssymb,mathtools,natbib}

% Anchored formal environments for causalsmith papers.
% First mandatory argument = obj_id (machine anchor joining the paper to the
% Lean crosswalk; invisible in the rendered PDF). Optional argument = name.
% Each environment also defines \label{obj:<obj_id>} so prose can use
% \cref{obj:T-1}; cleveref derives the printed kind and number from the target.
\usepackage{amsmath}
\usepackage{mathrsfs}
\usepackage{amsthm}
\usepackage{booktabs}
% Long displays may break across pages; let TeX stretch a little harder before an
% overfull \hbox so wide formulas/tables are less likely to run off the page.
\allowdisplaybreaks
\setlength{\emergencystretch}{3em}
\newtheorem{theorem}{Theorem}
\newtheorem{lemma}{Lemma}
\newtheorem{assumption}{Assumption}
\newtheorem{definition}{Definition}
\newtheorem{proposition}{Proposition}
\newtheorem{remark}{Remark}
\newtheorem{citedresult}{Cited result}
% The amsthm counter is `algorithmbox` (not `algorithm`) so a later `\usepackage{algorithm}` cannot clash.
\newcounter{algorithmbox}
\renewcommand{\thealgorithmbox}{\arabic{algorithmbox}}
\NewDocumentEnvironment{theoremv}{m o s}{\begin{theorem}\label{obj:#1}{\normalfont[\texttt{\detokenize{#1}}]\IfValueT{#2}{\ (#2)}\IfBooleanT{#3}{\verificationfootnotemark}\ }}{\end{theorem}}
\NewDocumentEnvironment{lemmav}{m o s}{\begin{lemma}\label{obj:#1}{\normalfont[\texttt{\detokenize{#1}}]\IfValueT{#2}{\ (#2)}\IfBooleanT{#3}{\verificationfootnotemark}\ }}{\end{lemma}}
\NewDocumentEnvironment{assumptionv}{m o s}{\begin{assumption}\label{obj:#1}{\normalfont[\texttt{\detokenize{#1}}]\IfValueT{#2}{\ (#2)}\IfBooleanT{#3}{\verificationfootnotemark}\ }}{\end{assumption}}
\NewDocumentEnvironment{definitionv}{m o s}{\begin{definition}\label{obj:#1}{\normalfont[\texttt{\detokenize{#1}}]\IfValueT{#2}{\ (#2)}\IfBooleanT{#3}{\verificationfootnotemark}\ }}{\end{definition}}
% A `propositionv` is a proved result tiered as a Proposition (theorem-grade, machine-verified).
\NewDocumentEnvironment{propositionv}{m o s}{\begin{proposition}\label{obj:#1}{\normalfont[\texttt{\detokenize{#1}}]\IfValueT{#2}{\ (#2)}\IfBooleanT{#3}{\verificationfootnotemark}\ }}{\end{proposition}}
% A `remarkv` holds NON-load-bearing interpretive / open-question content (e.g. a causalsmith `oeq:`
% open question). It is neither proved nor assumed; no theorem may depend on it.
\NewDocumentEnvironment{remarkv}{m o s}{\begin{remark}\label{obj:#1}{\normalfont[\texttt{\detokenize{#1}}]\IfValueT{#2}{\ (#2)}\IfBooleanT{#3}{\verificationfootnotemark}\ }}{\end{remark}}
% An `algorithmv` is a DEFINITION-kind object presented as a boxed, numbered-step procedure (an
% estimator / selection rule built in stages). Same anchor contract as `definitionv`; its body is
% ordinary LaTeX whose steps are `\begin{enumerate}` items, so tex2html needs no extra machinery.
% Presented in the CONVENTIONAL algorithm style readers expect from `algorithm`+`algorithmic`:
% a rule above, an "Algorithm N (Title)" caption line, a rule under the caption, the numbered
% steps, and a closing rule. It is NOT a float — the anchor contract requires the object to stay
% where the outline places it, and floats drift. The obj-id tag is suppressed here (it is
% bookkeeping, and a caption line carrying it reads as debug output in a finished paper).
\NewDocumentEnvironment{algorithmv}{m o s}%
  {\par\addvspace{\medskipamount}\noindent\rule{\linewidth}{0.8pt}\par\nobreak\vspace{2pt}%
   \refstepcounter{algorithmbox}\label{obj:#1}%
   \noindent\textbf{Algorithm \thealgorithmbox}\IfValueT{#2}{ \textnormal{(#2)}}%
   \IfBooleanT{#3}{\verificationfootnotemark}\par\nobreak\vspace{2pt}%
   \noindent\rule{\linewidth}{0.4pt}\par\nobreak\vspace{2pt}\normalfont}%
  {\par\nobreak\vspace{2pt}\noindent\rule{\linewidth}{0.8pt}\par\addvspace{\medskipamount}}
% A `citedv` is an IMPORTED external result the paper relies on but does NOT prove
% (a causalsmith `gate_class:"cited"` node). It is machine-checked against the cited
% source at F2.5, never proved here; the fixed provenance note makes that explicit
% so the environment can never be read as a contribution of this paper. The \citep
% attribution is supplied by the surrounding prose (P2), keyed to references.bib.
\NewDocumentEnvironment{citedv}{m o s}{\begin{citedresult}\label{obj:#1}{\normalfont[\texttt{\detokenize{#1}}]\IfValueT{#2}{\ (#2)}\IfBooleanT{#3}{\verificationfootnotemark}\ \textit{(imported result; machine-checked against the cited source, not proved here.)}\ }}{\end{citedresult}}
% \leanref{obj_id}{display text}: an inline first-mention link to a from-note object's
% verified Lean code. In the PDF it renders the display text plainly (the Lean link is a
% web-only affordance); tex2html turns it into a drawer-opening span.
\newcommand{\leanref}[2]{#2}
% For a theorem/lemma whose Lean declaration accepts a source-matched published proposition as an
% input, the starred anchored environment puts the mark in its heading; the generated text remains
% outside the frozen mathematical environment. Keep the combined macro for older paper bundles.
\newcommand{\verificationfootnotemark}{\footnotemark}
\newcommand{\verificationfootnotetext}[1]{\footnotetext{#1}}
\newcommand{\verificationfootnote}[1]{\verificationfootnotemark\verificationfootnotetext{#1}}
% xcolor before hyperref (hyperref would otherwise pull in plain `color` for colorlinks, and a
% later xcolor load clashes). The page-1 identity stamp at the end of this file needs a grey.
\usepackage{xcolor}
\usepackage{graphicx}
\usepackage{eso-pic}
% hyperref follows natbib and the environment macros; cleveref must follow hyperref.
\usepackage[colorlinks=true,linkcolor=blue,citecolor=blue,urlcolor=blue]{hyperref}
\usepackage[capitalise,nameinlink,noabbrev]{cleveref}
% Pin the names explicitly: labels are placed inside the underlying amsthm environments, and these
% declarations make the PDF contract independent of cleveref's theorem-name inference. House style:
% reference names always print with a capital first letter ("Theorem 1", "Definition 2"), in any
% sentence position — hence `capitalise` above and capitalized \crefname forms below.
\crefname{theorem}{Theorem}{Theorems}
\Crefname{theorem}{Theorem}{Theorems}
\crefname{lemma}{Lemma}{Lemmas}
\Crefname{lemma}{Lemma}{Lemmas}
\crefname{assumption}{Assumption}{Assumptions}
\Crefname{assumption}{Assumption}{Assumptions}
\crefname{definition}{Definition}{Definitions}
\Crefname{definition}{Definition}{Definitions}
\crefname{proposition}{Proposition}{Propositions}
\Crefname{proposition}{Proposition}{Propositions}
\crefname{remark}{Remark}{Remarks}
\Crefname{remark}{Remark}{Remarks}
\crefname{citedresult}{Cited result}{Cited results}
\Crefname{citedresult}{Cited result}{Cited results}
\crefname{algorithmbox}{Algorithm}{Algorithms}
\Crefname{algorithmbox}{Algorithm}{Algorithms}
% ---------------------------------------------------------------------------
% Page-1 identity stamp (arXiv style) and the pinned title-block date.
%
% Split of responsibility: THIS file owns the FORMAT (placement, font, colour, punctuation),
% the generated `paper_stamp.tex` owns only the VALUES (number+version, area, revised date,
% DOI, link target). P4 refreshes this template on every emit, so a layout fix reaches every
% bundle from a recompile; and because `paper_stamp.tex` is not one of the P2 assembly
% digest's authored inputs, a DOI that arrives after publication needs only that one small
% file rewritten plus a recompile — never a redraft, never a reassembly.
%
% Defaults first, so a bundle WITHOUT a stamp file compiles exactly as it always did:
% `\cspaperdate` falls back to `\today` and no stamp is drawn.
\providecommand*{\cspaperdate}{\today}  % the \date{} target
\providecommand*{\csstampid}{}          % e.g. CSWP-2026-008v3 (empty ⇒ no stamp at all)
\providecommand*{\csstamparea}{}        % e.g. Stat
\providecommand*{\csstampdate}{}        % e.g. 27 Aug 2026
\providecommand*{\csstampdoi}{}         % e.g. 10.5281/zenodo.123 (empty ⇒ DOI part omitted)
\providecommand*{\csstamplink}{}        % href target (DOI url, else the paper's page; may be empty)
\IfFileExists{paper_stamp.tex}{% GENERATED by CausalSmith P4 — paper identity stamp values. Do not hand-edit.
%
% Field order on the stamp (owned by paper_macros.tex):
%   <number>v<version> · doi:<doi> · [<Area>] <date>   — the DOI is never the last token, so a
%   text extractor cannot run it into whatever follows. Without a DOI that part is omitted.
% Format lives in paper_macros.tex; only the values are here, and this file is NOT one of
% the P2 assembly digest's authored inputs. A DOI arriving after publication therefore needs
% this file rewritten plus a `latexmk` recompile — never a redraft or a reassembly.
\renewcommand*{\csstampid}{CSWP-2026-037v9}
\renewcommand*{\csstamparea}{Stat}
\renewcommand*{\csstampdate}{3 Oct 2026}
\renewcommand*{\csstampdoi}{}
\renewcommand*{\csstamplink}{https://causalsmith.org/papers/stat_annotation_rarearm_frontier_v1}
\renewcommand*{\cspaperdate}{3 October 2026}
}{}
\makeatletter
% Field order is deliberate: the DOI is NEVER the last token on the line. A trailing identifier
% runs into whatever a text extractor emits next, which makes it unsafe to match mechanically
% (the Zenodo stamp matcher reads exactly this line); the date is a safe terminator.
%   <number>v<version> · doi:<doi> · [<Area>] <date>      — with a DOI
%   <number>v<version> · [<Area>] <date>                  — without
\newcommand{\cs@stampsep}{\,\textperiodcentered\,}
\newcommand{\cs@stamptext}{%
  \csstampid
  \ifx\csstampdoi\@empty\else\cs@stampsep doi:\csstampdoi\fi
  \ifx\csstamparea\@empty\else\cs@stampsep[\csstamparea]\fi
  \ifx\csstampdate\@empty\else\quad\csstampdate\fi
}
\ifx\csstampid\@empty\else
  \definecolor{csstampgrey}{gray}{0.45}
  % Background shipout material: it occupies no space, so the text block and the \thanks
  % footnote are byte-identical to an unstamped compile. The starred form applies to the
  % NEXT page shipped only — i.e. page 1. Placed in the left margin (the text block starts
  % at 1.1in), reading bottom-to-top, in a Times-like face at 8pt.
  \AddToShipoutPictureBG*{%
    \AtPageLowerLeft{%
      \put(\LenToUnit{0.42in},\LenToUnit{1.1in}){%
        \rotatebox{90}{%
          \fontfamily{ptm}\fontsize{8}{10}\selectfont
          \ifx\csstamplink\@empty
            \textcolor{csstampgrey}{\cs@stamptext}%
          \else
            \expandafter\href\expandafter{\csstamplink}{\textcolor{csstampgrey}{\cs@stamptext}}%
          \fi
        }%
      }%
    }%
  }
\fi
\makeatother


\title{Treatment effect estimation with auxiliary records and diminishing overlap}

\author{CausalSmith\thanks{Machine-generated by the CausalSmith research pipeline.}}

\date{\cspaperdate}

\begin{document}

\maketitle

\begin{abstract}
We characterize minimax squared-error precision for a population-weighted treatment contrast with binary treatment and outcomes and finite discrete covariates. The experiment supplies independent samples from a common population: \(n\ge1\) complete records containing outcomes, treatment, and covariates, and \(m\ge0\) auxiliary treatment--covariate records. Covariates have \(d\ge2\) categories, and treatment propensities in occupied cells lie in \([\epsilon,1-\epsilon]\), where \(\epsilon\in(0,1/4]\) is a public overlap floor. Over arbitrary covariate probabilities and cellwise outcome means, the minimax risk is
\[
\asymp
\min\left\{
1,\frac{1}{n\epsilon}
+\left[
\frac{d}{(n+m)\epsilon\log(\exp(1)+n\epsilon)}
\right]^2
\right\},
\]
with numerical comparison constants independent of all four indices. A deterministic estimator on the original records attains this order with proof-oriented tuning and numerical constants, and a matching lower bound applies to all Borel randomized rules on those records. The comparison yields exact sequence criteria for consistency, for joint consistency and rare-label benchmark attainment, and for a vanishing risk ratio relative to complete records alone. At fixed complete-record size, alphabet size, and overlap floor, increasing auxiliary information converges to the exact-marginal minimax benchmark. Under potential-outcome consistency and conditional treatment exchangeability, the contrast equals the population average treatment effect.
\end{abstract}

\section{Introduction}

This paper establishes a uniform minimax precision frontier for estimating a population-weighted treatment contrast when outcome information and treatment--covariate information arrive through separate sampling channels. The model has binary treatment and outcomes and a finite covariate alphabet. Its unknown population law permits arbitrary covariate probabilities and unrestricted cellwise outcome means, subject to a public treatment-overlap restriction in every occupied covariate cell. The resulting risk characterization describes jointly how the complete-record budget, the auxiliary-record budget, the alphabet size, and a possibly diminishing overlap floor determine attainable precision.

The experiment contains \(n\ge1\), the number of complete outcome--treatment--covariate records, and \(m\ge0\), the number of independent auxiliary treatment--covariate records, sampled from the same population. The covariate alphabet has size \(d\ge2\), and \(\epsilon\in(0,1/4]\), the public overlap floor, bounds occupied-cell treatment propensities away from zero and one. Thus all \(n+m\) records carry information about population treatment--covariate frequencies, while the complete sample also carries outcomes. Conditional overlap accommodates arbitrarily small population covariate cells, making the distinction between these information budgets consequential when many cells are sparsely represented.

For this experiment, \cref{obj:thm:uniform-annotation-frontier} establishes that \(R(n,m,d,\epsilon)\), the minimax squared-error risk, satisfies
\[
R(n,m,d,\epsilon)
\asymp
\min\left\{
1,\frac{1}{n\epsilon}
+\left[
\frac{d}{(n+m)\epsilon\log(\exp(1)+n\epsilon)}
\right]^2
\right\}.
\]
The comparison constants are numerical and independent of every public index throughout the stated domain. The target is the population-covariate-weighted difference between the two conditional outcome means, defined in \cref{obj:def:target}. Every observable model law admits a compatible potential-outcome extension satisfying consistency and conditional treatment exchangeability, and the target equals the population average treatment effect under every such extension, by \cref{obj:lem:causal-realization}.

The two terms identify distinct sources of statistical difficulty. The inverse scale \(1/(n\epsilon)\) reflects information about rare-arm outcomes in the complete sample. The many-cell term reflects learning the unknown treatment--covariate distribution using both channels. A common-marginal testing construction in \cref{obj:thm:rare-label-floor} establishes a capped inverse-\(n\epsilon\) risk floor for every auxiliary budget, including the experiment supplied with the exact population marginal. The deterministic construction in \cref{obj:def:estimator-handle} attains the full frontier through independent observation pools, pilot localization, reciprocal-polynomial correction, clipping, and finite averaging over ordered prefixes. Its proof-oriented tuning establishes the statistical order uniformly over the model and public indices with comparison constants that absorb the zero-output threshold; \cref{sec:attaining-estimator} explains the mechanism.

The frontier also gives exact resource characterizations for arbitrary experiment sequences. Writing \(S_k=n_k\epsilon_k\) for rare-arm outcome information, \(N_k=n_k+m_k\) for the total marginal-information budget, and \(\ell_k=\log(\exp(1)+S_k)\) for the regularized logarithmic scale, \cref{obj:thm:annotation-resource-phases} shows that consistency is equivalent to
\[
S_k\longrightarrow\infty,
\qquad
d_k=o(N_k\epsilon_k\ell_k).
\]
Joint consistency and attainment of rare-label benchmark order are equivalent to
\[
S_k\longrightarrow\infty,
\qquad
d_k=O\!\left(\frac{N_k\epsilon_k\ell_k}{\sqrt{S_k}}\right).
\]
The same result characterizes a vanishing risk ratio relative to complete records alone, accounting for both supervised saturation and supervised risks that already vanish. These criteria quantify the marginal expansion needed to reach a precision order governed by complete-record outcomes. At fixed \(n,d,\epsilon\), \cref{obj:prop:benchmark-recovery} further establishes convergence of the minimax risk itself to the exact-marginal risk as \(m\to\infty\), with an explicit finite-sample comparison.

The closest supervised analysis is \citet[model (3), Theorems 1--2]{ZengBalakrishnanHanKennedy2026Discrete}, and the matching fixed-interior supervised characterization is credited to \citet[Theorem \texttt{thm:sharp-minimax-fixed-interior}]{InternalDiscreteATEPolynomialUpperMatch}. The fixed-overlap two-channel frontier is established by \citet[Theorem 4]{CausalSmith2026Annotation}. Relative to these comparisons, the present result characterizes the joint dependence on unequal information budgets and a possibly diminishing public overlap floor through constants uniform over all four indices. Reciprocal approximation in off-policy evaluation provides a methodological antecedent \citep[Section 3.3, Theorems 5--6]{MaZhuJiaoWainwright2022OPE}. Semi-supervised treatment-effect inference supplies complementary efficiency results under nuisance-estimation conditions \citep[Theorem 2.1 and Corollary 2.1]{ChakraborttyDai2024Semisupervised}, while the allocation results of \citet[Theorems 4.2 and 5.1 of the arXiv manuscript]{NwankwoGoldkindZhou2026Annotations} address annotation design under their adaptive sampling regime. Here the precision comparison concerns passive independent acquisition from a common population.

The results hold under the stated model and sampling assumptions.

The paper proceeds from related work to the experiment, target, and causal interpretation, followed by the uniform minimax comparison and benchmark recovery. It then explains the attaining estimator and develops the consistency and auxiliary-information resource characterizations before discussing their interpretation and limitations. The appendices present the causal and benchmark arguments, estimator risk bounds, testing and approximation lower bounds, resource derivations and verification note, and proofs of the main results.


\section{Related work}

The closest supervised comparison is the unrestricted discrete model of \citet[model (3), Theorems 1--2]{ZengBalakrishnanHanKennedy2026Discrete}. For a fixed overlap floor, their conventional estimators have squared-error scale \(n^{-1}+d^2/n^2\), while their converse has scale \(n^{-1}+d^2/(n^2\log^2 n)\) in its stated many-cell range, with floor-dependent constants. The public CausalSmith working paper \citet[Theorem \texttt{thm:sharp-minimax-fixed-interior}]{InternalDiscreteATEPolynomialUpperMatch} closes that gap: for each fixed \(0<\epsilon<1/2\), there are \(a_\epsilon,\rho_\epsilon,C_\epsilon>0\) and \(N_\epsilon<\infty\) such that, whenever \(n\ge N_\epsilon\) and \(d\le\rho_\epsilon n\log n\), the supervised minimax risk is bounded above and below by floor-dependent multiples of
\[
n^{-1}+\frac{d^2}{n^2\log^2 n}.
\]
Our comparison in \cref{obj:thm:uniform-annotation-frontier} treats the complete-record budget, auxiliary-record budget, alphabet size, and public overlap floor jointly, with numerical comparison constants independent of all four indices over their stated domains. This uniformity makes the dependence on diminishing overlap part of the risk characterization itself.

The nearest two-channel result is Theorem 4 of CausalSmith working paper CSWP-2026-022 \citep{CausalSmith2026Annotation}. At each fixed overlap floor it gives, with constants allowed to depend on that floor,
\[
\min\left\{1,\frac1n+\frac{d^2}{(n+m)^2[\log(\exp(1)n)]^2}\right\}
\]
as the risk order for the complete-plus-auxiliary experiment. The present \leanref{S-1}{passive two-channel experiment} likewise combines complete records with independent records containing treatment and covariate information, under \cref{obj:ass:data-product}. Within the overlap class in \cref{obj:def:model}, the two budgets may be unequal: complete records supply outcome information, while both channels supply information about the treatment--covariate distribution. The overlap-uniform comparison in \cref{obj:thm:uniform-annotation-frontier}, together with the rare-label benchmark in \cref{obj:thm:rare-label-floor}, gives these information sources separate roles with numerical constants uniform in the possibly diminishing floor. The benchmark specializations are collected in \cref{obj:prop:benchmark-recovery}. Their sequence implications in \cref{obj:thm:annotation-resource-phases} characterize minimax consistency, attainment of the rare-label benchmark order, and strict relative improvement from auxiliary records as the budgets, alphabet size, and overlap floor vary together.

Reciprocal-polynomial approximation has an important antecedent in \citet[Section 3.3, Theorems 5--6]{MaZhuJiaoWainwright2022OPE}. They study off-policy evaluation with supplied target-policy weights and partial knowledge of the behavior policy through a global lower bound on action probabilities. Their Chebyshev construction and approximation-based lower bound provide methodological precedents for estimating weighted outcome functionals with uncertain denominators. The upper-bound result uses a Poisson sampling model and a delocalization condition on the target policy, while the lower bound imposes its stated exploration-probability regime. Here the target uses unknown population covariate weights, and the overlap restriction controls treatment probabilities conditional on an occupied covariate cell. Consequently, small population cells remain admissible under fixed conditional overlap. The estimator in \cref{obj:def:estimator-handle} adapts reciprocal approximation to these unknown weights and to the distinct outcome and marginal-information budgets.

The passive experiment also belongs alongside semi-supervised treatment-effect inference. \citet[Theorem 2.1 and Corollary 2.1]{ChakraborttyDai2024Semisupervised} study labeled and outcome-unlabeled samples and obtain robustness and efficiency results under their nuisance-estimation conditions. \citet[Theorems 3.7, 3.13, and 5.1]{ZhangChakraborttyBradic2023DecayingOverlap} analyze missing-at-random labeling with decaying labeling probabilities, making the observation mechanism central to the available information. A complementary design question is addressed by \citet[Theorems 4.2 and 5.1 of the arXiv manuscript]{NwankwoGoldkindZhou2026Annotations}: their annotation-allocation result optimizes asymptotic variance, and their feasible procedure uses batch adaptation and cross-fitting. Those conclusions operate under their proportional budget and batch regime, causal identification and missing-at-random conditions, treatment and annotation positivity, and the stipulated moment, nuisance-consistency, marginal and product convergence-rate, function-class, and cross-batch dependence conditions. In our experiment, independent sampling from a common population fixes how records enter each channel. The squared-error comparison therefore characterizes the value of auxiliary records under passive acquisition, with treatment overlap and the two sample sizes specified as public experiment indices.

More broadly, limited-overlap precision and inference motivate attention to the information carried by rare treatment observations \citep{Rothe2017LimitedOverlap,ArmstrongKolesar2021ATEInference,MaWang2020RobustIPW}. Discrete functional estimation supplies complementary tools for aggregate targets when many cells have small counts, particularly polynomial approximation and moment matching \citep{JiaoVenkatHanWeissman2015Functionals,WuYang2016Entropy}. These strands meet in the paper's minimax squared-error question: how treatment overlap and the separate budgets for outcome and treatment--covariate information jointly determine the precision of a population-weighted treatment contrast.


\section{Experiment, target, and causal interpretation}

Throughout, logarithms are natural, norms of finite vectors are Euclidean, and numerical comparison constants are independent of the public experiment indices unless stated otherwise. The relations \(\lesssim\) and \(\asymp\) denote inequalities and two-sided comparisons up to such constants. The public indices are the complete-record sample size \leanref{sym:n}{\ensuremath{n}}\(\ge1\), the auxiliary sample size \leanref{sym:m}{\ensuremath{m}}\(\ge0\), the covariate alphabet size \leanref{sym:d}{\ensuremath{d}}\(\ge2\), and the overlap floor \leanref{sym:\epsilon}{\ensuremath{\epsilon}}\(\in(0,1/4]\). The covariate index \leanref{sym:j}{\ensuremath{j}} ranges from \(1\) to \(d\), and the treatment index \leanref{sym:a}{\ensuremath{a}} ranges over \(\{0,1\}\).

The two information budgets concern different parts of the same population distribution. A complete record contains the covariate coordinate \leanref{sym:X}{\ensuremath{X}}, binary treatment \leanref{sym:A}{\ensuremath{A}}, and binary observed outcome \leanref{sym:Y}{\ensuremath{Y}}; its probability law is \leanref{sym:P}{\ensuremath{P}}. We use the covariate probabilities \leanref{sym:p}{\ensuremath{p_j(P)}}, treatment--covariate probabilities \leanref{sym:s}{\ensuremath{s_{aj}(P)}}, success-marked probabilities \leanref{sym:q}{\ensuremath{q_{aj}(P)}}, treatment propensities \leanref{sym:e}{\ensuremath{e_j(P)}}, outcome regressions \leanref{sym:\mu}{\ensuremath{\mu_{aj}(P)}}, and treatment--covariate marginal \leanref{sym:P_{XA}}{\ensuremath{P_{XA}}}.

Explicitly, for \(a\in\{0,1\}\) and \(j\in\{1,\ldots,d\}\),
\[
\begin{aligned}
p_j(P)&=\sum_{a'\in\{0,1\}}\sum_{y\in\{0,1\}}P(X=j,A=a',Y=y),\\
s_{aj}(P)&=\sum_{y\in\{0,1\}}P(X=j,A=a,Y=y),\\
q_{aj}(P)&=P(X=j,A=a,Y=1),\\
\mu_{aj}(P)&=
\begin{cases}
q_{aj}(P)/s_{aj}(P),&s_{aj}(P)>0,\\
0,&s_{aj}(P)=0.
\end{cases}
\end{aligned}
\]
Thus \(s_{aj}(P)=P(X=j,A=a)\), while the zero convention makes every outcome regression total and leaves the target unchanged on null arm--cells.



These quantities distinguish conditional treatment overlap from population cell size. The propensity and marginal definitions used throughout are recorded explicitly below.

\begin{definitionv}{synth_2}[Cell-specific treatment propensity]
For a probability law \(P\) of \((X,A,Y)\) on \(\{1,\ldots,d\}\times\{0,1\}^2\), let \(p_j(P)=P(X=j)\). Define the treatment propensity in cell \(j\) by
\[
e_j(P)=
\begin{cases}
\Pr_P(A=1\mid X=j)
=\dfrac{P(X=j,A=1)}{p_j(P)},&p_j(P)>0,\\[6pt]
\dfrac12,&p_j(P)=0.
\end{cases}
\]
\end{definitionv}

\begin{assumptionv}{ass:overlap}[Public overlap restriction]
For every \(j\) with \(p_j(P)>0\),
\[
e_j(P)\in[\epsilon,1-\epsilon].
\]
\end{assumptionv}

We impose a public restriction on treatment assignment within every occupied covariate cell.

\begin{definitionv}{synth_1}[Covariate and treatment marginal]
The treatment–covariate marginal \(P_{XA}\) is the probability law obtained from \(P\), the probability law of complete records \((X,A,Y)\), by the projection \((X,A,Y)\mapsto(X,A)\). Here \(X\) takes values in a finite alphabet of size \(d\), for a nonnegative integer \(d\), and \(A,Y\in\{0,1\}\). Its cell-probability table is
\[
P_{XA}(j,a)
=
\sum_{y\in\{0,1\}} P(X=j,A=a,Y=y),
\]
for each covariate value \(j\) and treatment value \(a\in\{0,1\}\).
\end{definitionv}

The standard strong conditional overlap condition gives each treatment arm positive conditional probability in every occupied cell \citep[model (3)]{ZengBalakrishnanHanKennedy2026Discrete}; propensity adjustment traces to inverse-probability weighting \citep{Horvitz1952,Hirano2003}. Its public floor may vary across experiments, allowing treatment overlap to diminish along an experiment sequence. Covariate probabilities may be arbitrarily small, so the restriction accommodates populations with many sparsely represented cells. The null-cell convention in \cref{obj:synth_2} supplies a value for the propensity wherever the cell has zero population weight.

Sampling supplies complete records and auxiliary records independently from this population. Write \leanref{sym:\mathcal L}{\ensuremath{\mathcal L}} for the ordered complete sample, \leanref{sym:\mathcal V}{\ensuremath{\mathcal V}} for the ordered auxiliary sample, and \leanref{sym:\mathcal D}{\ensuremath{\mathcal D}} for their observed pair.



An independent seed \leanref{sym:U}{\ensuremath{U}} permits randomized estimation rules on these original records.

\begin{assumptionv}{ass:data-product}[Two-channel product sampling]
The observed data \(\mathcal D\) have distribution
\[
\mathcal D\sim P^{\otimes n}\otimes P_{XA}^{\otimes m}.
\]
\end{assumptionv}

The sampling restriction can now be stated as a product experiment.

\begin{definitionv}{synth_5}[Probability laws]
For a measurable space \(\mathcal Z\), let \(\mathsf{Prob}(\mathcal Z)\) denote the class of probability measures on \(\mathcal Z\).
\end{definitionv}

Independent complete and outcome-unlabeled samples from the same population are a standard sampling condition in semisupervised treatment-effect estimation \citep[Assumptions 2.1--2.2]{ChakraborttyDai2024Semisupervised}. Here the auxiliary records directly inform the same treatment--covariate table that governs the complete records. Consequently, \(n+m\) records carry marginal information, while \(n\) records also carry outcomes.

To specify the observable model, let probability laws have the following notation.

\begin{definitionv}{def:model}[Observable law class \(\mathcal M_{d,\epsilon}\)]
The observable law class is
\[
\mathcal M_{d,\epsilon}
=
\left\{
P\in\mathsf{Prob}\bigl(\{1,\ldots,d\}\times\{0,1\}^2\bigr):
P\text{ satisfies }\cref{obj:ass:overlap}
\right\},
\qquad d\ge2,\quad 0<\epsilon\le\tfrac14.
\]
This is model (3) of Zeng et al., restricted to the stated range of the public overlap floor \(\epsilon\).
\end{definitionv}

The observable law class \leanref{sym:\mathcal M}{\ensuremath{\mathcal M_{d,\epsilon}}} combines the finite record space with the public overlap restriction.

\begin{definitionv}{def:target}[Outcome-contrast estimand \(\tau(P)\)]
The estimand is
\[
\tau(P)
=
\sum_{j=1}^{d}p_j(P)\{\mu_{1j}(P)-\mu_{0j}(P)\},
\]
with summands for null covariate cells defined to equal zero.
\end{definitionv}

The correspondence with the discrete model in \citet[model (3)]{ZengBalakrishnanHanKennedy2026Discrete} fixes the population class for the two-channel experiment. Within that class, the target \leanref{sym:\tau}{\ensuremath{\tau(P)}} averages the conditional outcome contrast using population covariate weights.

\begin{definitionv}{synth_3}[Independent uniform randomization]
Let \(\mathrm{Uniform}[0,1]\) denote the probability law on \([0,1]\) assigning each Borel set its Lebesgue measure. The estimator randomization seed \(U\) has this law and is independent of the observed data \(\mathcal D=(\mathcal L,\mathcal V)\), where \(\mathcal L\) and \(\mathcal V\) are the ordered complete-record and auxiliary samples, respectively.
\end{definitionv}

Because the outcome regressions lie in \([0,1]\), the target lies in \([-1,1]\). This range motivates the following convention for admissible estimation rules \leanref{sym:T}{\ensuremath{T}}.



Clipping weakly reduces squared error for every target value in \([-1,1]\). The minimax risk \leanref{sym:R}{\ensuremath{R(n,m,d,\epsilon)}} therefore evaluates precision over all original-record rules with this bounded range.

\begin{definitionv}{def:risk}[Minimax risk $R(n,m,d,\epsilon)$]
The minimax squared-error risk is
\[
R(n,m,d,\epsilon)
=
\inf_T\sup_{P\in\mathcal M_{d,\epsilon}}
\mathbb E_{P^{\otimes n}\otimes P_{XA}^{\otimes m}\otimes\mathrm{Uniform}[0,1]}
\bigl[(T(\mathcal D,U)-\tau(P))^2\bigr].
\]
The infimum ranges over all Borel estimation rules on the original observed records, including rules randomized through the independent seed \(U\). Clipping under squared-error loss permits restriction to rules taking values in \([-1,1]\).
\end{definitionv}

The observable target also admits a potential-outcome interpretation. A \leanref{S-2}{compatible potential-outcome extension} has law \leanref{sym:H}{\ensuremath{H}} and carries a binary potential outcome \leanref{sym:Y_a}{\ensuremath{Y_a}} for each treatment value \(a\).



The causal interpretation rests on consistency and conditional treatment exchangeability under such an extension.

\begin{assumptionv}{ass:consistency}[Potential-outcome consistency condition]
Under the potential-outcome extension \(H\),
\[
Y=AY_1+(1-A)Y_0
\]
holds \(H\)-almost surely.
\end{assumptionv}

The standard consistency condition for binary treatment connects the observed outcome to the potential outcome under the received treatment \citep{SplawaNeyman1990,Rubin1974}. It supplies the link between the observed outcome regressions and the corresponding potential outcomes within each treatment arm. Conditional treatment exchangeability then connects those arm-specific observations to the potential-outcome distribution within the covariate cell.

\begin{assumptionv}{ass:exchangeability}[Conditional treatment exchangeability]
Under the potential-outcome extension \(H\),
\[
(Y_0,Y_1)\perp\!\!\!\perp A\mid X.
\]
\end{assumptionv}

This standard conditional treatment exchangeability condition requires treatment assignment to be independent of the pair of potential outcomes given the recorded covariate \citep{Rosenbaum1983,Imbens2015}. Together with consistency and overlap, it identifies the two conditional potential-outcome means from the observable outcome regressions.

\Cref{obj:lem:causal-realization} establishes both existence and identification: every law in \cref{obj:def:model} admits a compatible extension satisfying \cref{obj:ass:consistency,obj:ass:exchangeability}, and every compatible extension satisfying these conditions obeys
\[
\tau(P)=\mathbb E_H[Y_1-Y_0].
\]
Thus the observable functional in \cref{obj:def:target} is the average treatment effect throughout this class of compatible causal extensions. The observable minimax criterion consequently measures precision for that common causal target under the stated interpretation.

For comparison, the \leanref{S-3}{exact-marginal benchmark experiment} supplies the population treatment--covariate table directly, paralleling the role of known propensity weights in inverse-probability constructions \citep{Horvitz1952,Hirano2003}. Its estimation rule \leanref{sym:T^{\mathrm K}}{\ensuremath{T^{\mathrm K}}} receives this table jointly with the complete records and the seed.



The corresponding exact-marginal risk \leanref{sym:R^{\mathrm K}}{\ensuremath{R^{\mathrm K}(n,d,\epsilon)}} evaluates estimation with this population information available.

\begin{definitionv}{def:known-risk}[Exact-marginal risk $R^{\mathrm K}(n,d,\epsilon)$]
Define the minimax squared-error risk with the exact treatment--covariate marginal supplied by
\[
R^{\mathrm K}(n,d,\epsilon)
=
\inf_{T^{\mathrm K}}\sup_{P\in\mathcal M_{d,\epsilon}}
\mathbb E_{P^{\otimes n}\otimes\mathrm{Uniform}[0,1]}
\left[
\bigl(T^{\mathrm K}(\mathcal L,P_{XA},U)-\tau(P)\bigr)^2
\right].
\]
Here \(\mathcal L\) is the complete-record component of \(\mathcal D\), and the infimum ranges over Borel estimation rules receiving \(\mathcal L\), the exact marginal table \(P_{XA}\), and the independent seed \(U\).
\end{definitionv}

Supplying the exact table makes the population covariate weights and treatment propensities available to the rule, while the complete sample supplies outcome information. The main comparisons use the capped rare-label benchmark \leanref{sym:b}{\ensuremath{b(n,\epsilon)}} and the fixed-overlap annotation benchmark \leanref{sym:f}{\ensuremath{f(n,m,d)}}; their formulas and comparison properties are given in \cref{obj:prop:benchmark-recovery}.



The first function records the capped inverse scale of complete-record information under the public overlap floor. The second combines an inverse complete-sample term with a many-cell term governed by the total marginal-information budget. These numerical benchmarks provide a common scale for the exact-marginal and fixed-overlap comparisons in \cref{obj:prop:benchmark-recovery}.


\section{Minimax precision with two information budgets}

The complete records supply outcome information, while both complete and auxiliary records supply treatment--covariate information. The frontier proved in this section is \(R(n,m,d,\epsilon)\asymp r(n,m,d,\epsilon)\), where \(r(n,m,d,\epsilon)\) is the numerical risk rate defined next, with comparison constants uniform in all four public indices. Its two components have distinct origins: complete records estimate the rare-arm success masses, while all records estimate the treatment--covariate masses needed to convert those successes into population-weighted arm means. Sparse arm masses require a reciprocal correction; its usable polynomial degree is limited by rare-arm outcome information \(n\epsilon\), even when auxiliary records make the marginal sample much larger. We first define the resulting rate, then give the complete attaining rule required by the frozen verification layer, followed by the formal comparison and its benchmark consequences.

\begin{definitionv}{def:rate-handle}[Numerical rate $r(n,m,d,\epsilon)$]
Define
\[
r(n,m,d,\epsilon)
=
\min\left\{
1,\frac1S+\left(\frac d{N\epsilon\ell}\right)^2
\right\},
\qquad
N=n+m,\quad S=n\epsilon,\quad
\ell=\log(\exp(1)+S).
\]
All denominators are strictly positive on the public index range, and \(r(n,m,d,\epsilon)\) is a positive numerical function.
\end{definitionv}

\begin{definitionv}{synth_4}[Clipping to the target range]
For every \(x\in\mathbb R\), define
\[
\operatorname{clip}_{[-1,1]}(x)=\max\{-1,\min\{x,1\}\}.
\]
\end{definitionv}

\begin{algorithmv}{def:estimator-handle}[Prefix-average estimator $\widehat\tau$]
The estimator takes the ordered complete and auxiliary records \(\mathcal D\), the independent seed \(U\), and the public parameters \(n,m,d,\epsilon\) as inputs.
\begin{enumerate}
\item Set
\[
S=n\epsilon.
\]
If \(S<\exp(4096)\), output
\[
\widehat\tau(\mathcal D,U)=0.
\]
For \(S\ge\exp(4096)\), proceed through the remaining steps.

\item Define the pool lengths and tuning parameters by
\[
\begin{aligned}
h_0&=\lfloor n/3\rfloor,
&
h_p&=\lfloor(n-h_0)/2\rfloor,
&
h_f&=n-h_0-h_p,\\
M_p&=h_p+\lfloor m/2\rfloor,
&
M_f&=h_f+m-\lfloor m/2\rfloor,\\
u&=h_0/8,
&
t_p&=M_p/8,
&
t&=M_f/8,\\
L&=\lfloor\log S/1024\rfloor,
&
B&=\frac{2^{20}L}{\min\{t_p,t\}},
&
k_0&=\lfloor t_pB/4\rfloor.
\end{aligned}
\]
Use the first \(h_0\) complete records as the outcome pool. Form the pilot pool by concatenating the treatment--covariate projections of the next \(h_p\) complete records with the first \(\lfloor m/2\rfloor\) auxiliary records. Form the factorial pool by concatenating the remaining complete-record projections with the remaining auxiliary records. Preserve the stated order within each pool.

\item Define the Chebyshev polynomials and the reciprocal-approximation coefficients by
\[
\begin{aligned}
T_0(x)&=1,\qquad T_1(x)=x,\\
T_{h+1}(x)&=2xT_h(x)-T_{h-1}(x),\qquad h\ge1,\\
E_L(z)&=\frac{1-T_L(1-2z)}{2L^2z},\\
G_L(z)&=\frac{1-E_L(z)}z
=\sum_{h=0}^{L-2}g_hz^h.
\end{aligned}
\]
Use the removable extensions at zero in the definitions of \(E_L\) and \(G_L\).

\item For each prefix triple satisfying
\[
0\le h'\le h_0,\qquad
0\le h''\le M_p,\qquad
0\le h'''\le M_f,
\]
write the ordered outcome-pool observations as
\((X_i^{\mathrm o},A_i^{\mathrm o},Y_i^{\mathrm o})\), the pilot-pool observations as
\((X_i^{\mathrm p},A_i^{\mathrm p})\), and the factorial-pool observations as
\((X_i^{\mathrm f},A_i^{\mathrm f})\). For \(a\in\{0,1\}\) and \(j\in\{1,\ldots,d\}\), define
\[
\begin{aligned}
Z_{aj}
&=\sum_{i=1}^{h'}
\mathbf1\{X_i^{\mathrm o}=j,\ A_i^{\mathrm o}=a,\ Y_i^{\mathrm o}=1\},\\
J_{aj}
&=\sum_{i=1}^{h''}
\mathbf1\{X_i^{\mathrm p}=j,\ A_i^{\mathrm p}=a\},\\
K_{aj}
&=\sum_{i=1}^{h'''}
\mathbf1\{X_i^{\mathrm f}=j,\ A_i^{\mathrm f}=a\}.
\end{aligned}
\]
Define the falling factorials of an integer count \(x\) by
\[
(x)_0=1,\qquad
(x)_h=\prod_{r=0}^{h-1}(x-r),\qquad h\ge1.
\]

\item For each prefix triple, define
\[
\begin{aligned}
\widehat G_{aj}
&=\sum_{h=0}^{L-2}g_h\frac{(K_{aj})_h}{(tB)^h},\\
P_{aj}^{\mathrm{pol}}
&=\frac{Z_{aj}}u
\left(1+\frac{K_{1-a,j}\widehat G_{aj}}{tB}\right),\\
H_{aj}
&=\frac{Z_{aj}}u
\left(1+\frac{K_{1-a,j}}{K_{aj}+1}\right),\\
V_{aj}
&=\mathbf1\{J_{aj}\le k_0\}P_{aj}^{\mathrm{pol}}
+\mathbf1\{J_{aj}>k_0\}H_{aj},\\
F(h',h'',h''';\mathcal D)
&=\operatorname{clip}_{[-1,1]}
\left(\sum_{j=1}^d(V_{1j}-V_{0j})\right).
\end{aligned}
\]

\item Output the deterministic finite average
\[
\widehat\tau(\mathcal D,U)
=
\sum_{h'=0}^{h_0}\sum_{h''=0}^{M_p}\sum_{h'''=0}^{M_f}
\exp(-u-t_p-t)
\frac{u^{h'}}{h'!}
\frac{t_p^{h''}}{h''!}
\frac{t^{h'''}}{h'''!}
F(h',h'',h''';\mathcal D).
\]
Poisson prefix combinations exceeding a pool length have output zero. The estimator is independent of \(U\). Every denominator is a strictly positive public quantity or an integer count plus one, including for empty observed cells; all quantities in the construction are determined by the observed data and public parameters.
\end{enumerate}
\end{algorithmv}

The zero-output branch for \(S<\exp(4096)\) is a proof device: bounded loss on that branch is absorbed into the universal comparison constant because the numerical rate is at least \(\exp(-4096)\). Thus the construction certifies the uniform minimax order with proof-oriented constants; the threshold should not be read as calibrated moderate-sample tuning.





The rare-arm outcome scale \leanref{sym:S}{\ensuremath{S}} combines the complete-record sample size with the overlap floor. The total marginal-information budget \(N\) counts all records carrying treatment and covariate observations, and the regularized logarithm \leanref{sym:\ell}{\ensuremath{\ell}} depends on the outcome scale. Thus \(S^{-1}\) records the outcome-information component, while \(\{d/(N\epsilon\ell)\}^2\) records the many-cell component. The cap at one accommodates experiments with bounded information.

The rate separates the precision associated with observing rare-arm outcomes from the precision associated with learning the treatment--covariate distribution across many cells. Increasing the auxiliary sample enlarges \(N\), reducing the many-cell term, while the complete-record outcome scale remains \(S\). A smaller overlap floor raises both terms and also reduces the logarithmic factor. The following theorem compares this rate with the minimax risk in \cref{obj:def:risk}; its upper bound is attained by the deterministic prefix-average estimator \leanref{sym:\widehat\tau}{\ensuremath{\widehat\tau}} specified in \cref{obj:def:estimator-handle}. The causal clause connects the observable target to \cref{obj:lem:causal-realization}.

\begin{theoremv}{thm:uniform-annotation-frontier}[Uniform annotation frontier]
There exist numerical constants \(0<c\le C<\infty\), independent of all public indices, such that the following holds for every experiment with:
\begin{itemize}
\item \textbf{(Sample sizes.)} Integers \(n\ge1\) and \(m\ge0\).
\item \textbf{(Covariate alphabet.)} An integer \(d\ge2\).
\item \textbf{(Overlap floor.)} A real parameter \(0<\epsilon\le1/4\).
\end{itemize}
The deterministic estimator \(\widehat\tau(\mathcal D,U)\) constructed in \cref{obj:def:estimator-handle} is Borel measurable, ignores the randomization seed \(U\), and takes values in \([-1,1]\) for every data realization. With \(R(n,m,d,\epsilon)\) the minimax squared-error risk from \cref{obj:def:risk} and \(r(n,m,d,\epsilon)\) the rate from \cref{obj:def:rate-handle},
\[
c\,r(n,m,d,\epsilon)
\le R(n,m,d,\epsilon)
\le \sup_{P\in\mathcal M_{d,\epsilon}}
\mathbb E_P\!\left[
\bigl(\widehat\tau(\mathcal D,U)-\tau(P)\bigr)^2
\right]
\le C\,r(n,m,d,\epsilon).
\]
Here \(\mathcal M_{d,\epsilon}\) is the observable model class in \cref{obj:def:model}, \(\tau(P)\) is the observable target in \cref{obj:def:target}, and the expectation uses the experiment specified in \cref{obj:def:risk}.

Moreover, for every integer \(d\ge2\), every \(0<\epsilon\le1/4\), and every observable law \(P\in\mathcal M_{d,\epsilon}\), there exists a joint law \(H\) of the observed coordinates and binary potential outcomes \(Y_0,Y_1\) whose observed marginal is \(P\) and which satisfies \cref{obj:ass:consistency,obj:ass:exchangeability}. Every such law \(H\) satisfies
\[
\mathbb E_H[Y_1-Y_0]=\tau(P).
\]
\end{theoremv}

The outcome-information component has a direct lower benchmark that applies both to the two-channel risk and to the exact-marginal risk in \cref{obj:def:known-risk}. The next theorem establishes it using two laws with a common treatment--covariate marginal, so auxiliary records cannot distinguish them.

\begin{theoremv}{thm:rare-label-floor}[Uniform rare-label risk floors]
There is a numerical constant \(c>0\), independent of all public indices, such that, for every \(n,m,d\in\mathbb N\) satisfying
\begin{itemize}
\item \textbf{(Complete-record sample size.)} \(n\ge 1\).
\item \textbf{(Covariate alphabet.)} \(d\ge 2\).
\item \textbf{(Overlap parameter.)} \(0<\epsilon\le 1/4\).
\end{itemize}
the minimax risks in \cref{obj:def:risk,obj:def:known-risk} satisfy
\[
R(n,m,d,\epsilon)\ge c\,b(n,\epsilon),
\qquad
R^{\mathrm K}(n,d,\epsilon)\ge c\,b(n,\epsilon),
\qquad
b(n,\epsilon)=\min\{1,(n\epsilon)^{-1}\}.
\]
The first bound applies to the infimum over all clipped Borel randomized original-record rules specified in \cref{obj:def:risk}.

Moreover, for every \(d\in\mathbb N\), overlap parameter \(\epsilon\in\mathbb R\), and separation \(\delta\in\mathbb R\) satisfying
\begin{itemize}
\item \textbf{(Covariate alphabet.)} \(d\ge 2\).
\item \textbf{(Overlap parameter.)} \(0<\epsilon\le 1/4\).
\item \textbf{(Separation.)} \(0<\delta\le 1/8\).
\end{itemize}
the two complete-record laws \(P_+\) and \(P_-\) specified by
\[
\begin{aligned}
p_j(P_\pm)&=\frac12,
&
e_j(P_\pm)&=\epsilon,
&
\mu_{0j}(P_\pm)&=\frac12,
&
\mu_{1j}(P_\pm)&=\frac12\pm\delta,
&& j=1,2,\\
p_j(P_\pm)&=0,
&&&&&&&& j=3,\ldots,d
\end{aligned}
\]
belong to \(\mathcal M_{d,\epsilon}\) of \cref{obj:def:model}. Their targets, defined in \cref{obj:def:target}, and treatment–covariate marginals satisfy
\[
\tau(P_+)=\delta,
\qquad
\tau(P_-)=-\delta,
\qquad
(P_+)_{XA}=(P_-)_{XA}.
\]
At \(\delta=1/8\), their one-record Kullback--Leibler divergence is exactly
\[
D(P_+\Vert P_-)
=
\sum_{z:\,P_+(z)>0}
P_+(z)\log\frac{P_+(z)}{P_-(z)}
=
\frac{\epsilon}{4}\log\frac53,
\]
where \(z\) ranges over complete-record atoms and \(\log\) is the natural logarithm.
\end{theoremv}

The target separation is carried entirely by treated outcomes at the public overlap floor. The auxiliary observations and the supplied exact marginal have the same distribution under the two laws, so the floor persists for every auxiliary budget. We next compare the uniform rate with the rare-label benchmark, the fixed-overlap benchmark, and the exact-marginal experiment.

\begin{propositionv}{prop:benchmark-recovery}[Benchmark recovery and oracle comparison]
Let \(r(n,m,d,\epsilon)\), \(R(n,m,d,\epsilon)\), and \(R^{\mathrm K}(n,d,\epsilon)\) be as in \cref{obj:def:rate-handle,obj:def:risk,obj:def:known-risk}. Define the rare-label and fixed-overlap benchmark functions by
\[
b(n,\epsilon)=\min\left\{1,\frac{1}{n\epsilon}\right\},
\qquad
f(n,m,d)=\min\left\{1,\frac{1}{n}
+\frac{d^2}{(n+m)^2[\log(\exp(1)n)]^2}\right\}.
\]
The following statements hold.
\begin{enumerate}
\item For every fixed \(0<\epsilon\le 1/4\), there exist constants \(a,B_c>0\), depending only on \(\epsilon\), such that for all integers \(n\ge1\), \(m\ge0\), and \(d\ge2\),
\[
a f(n,m,d)\le r(n,m,d,\epsilon)\le B_c f(n,m,d).
\]

\item There exist numerical constants \(c,C>0\) such that for all integers \(n\ge1\), \(m\ge0\), and \(d\ge2\), and every \(0<\epsilon\le1/4\),
\[
c b(n,\epsilon)\le R(n,m,2,\epsilon)\le C b(n,\epsilon),
\qquad
c b(n,\epsilon)\le R^{\mathrm K}(n,d,\epsilon)\le C b(n,\epsilon).
\]

\item For all integers \(n\ge1\), \(m\ge1\), and \(d\ge2\), and every \(0<\epsilon\le1/4\),
\[
0\le R(n,m,d,\epsilon)-R^{\mathrm K}(n,d,\epsilon)
\le \min\left\{1,6(n+2)\sqrt{\frac{2d}{m}}\right\}.
\]

\item For every fixed integer \(n\ge1\), integer \(d\ge2\), and \(0<\epsilon\le1/4\), the limit along integer auxiliary sample sizes satisfies
\[
\lim_{m\to\infty}R(n,m,d,\epsilon)=R^{\mathrm K}(n,d,\epsilon).
\]

\item For every real \(M_0>0\), there exists \(\kappa>0\), depending only on \(M_0\), such that for all integers \(n\ge1\), \(m\ge0\), and \(d\ge2\), and every \(0<\epsilon\le1/4\) satisfying \(n\epsilon\le M_0\),
\[
R(n,m,d,\epsilon)\ge\kappa.
\]
\end{enumerate}
\end{propositionv}

At fixed overlap, the first clause recovers the benchmark combining inverse complete-sample precision with a many-cell term governed by the total sample size and a squared logarithmic factor; only these comparison constants may depend on the chosen floor. With two covariate cells, the risk has rare-label order for every auxiliary budget, and supplying the exact marginal gives that order for every allowed alphabet size. At fixed \(n,d,\epsilon\), the third and fourth clauses additionally give an explicit comparison and convergence of the minimax risk to the exact-marginal risk as \(m\to\infty\). \Cref{sec:attaining-estimator} explains the construction in \cref{obj:def:estimator-handle} and compares its guarantee with an inverse-count baseline. The technical many-cell lower-bound construction and its testing comparison are presented through \cref{obj:def:prior-handle,obj:lem:normalized-affine-testing}.


\section{An attaining estimator}\label{sec:attaining-estimator}

The estimator corrects a mismatch between the quantities observed with outcomes and the arm means entering the target. In arm \(a\) and cell \(j\), complete records identify the success mass \(q_{aj}=s_{aj}\mu_{aj}\), whereas the target contribution is \(p_j\mu_{aj}=q_{aj}p_j/s_{aj}\). Since \(p_j=s_{0j}+s_{1j}\), this requires estimating a reciprocal arm mass. The pilot pool localizes each arm--cell: populated cells use an inverse count, while sparse cells use a polynomial reciprocal whose monomials are estimated from factorial counts in an independent pool.

The degree calibration is tied to \(S=n\epsilon\), rather than to \(N=n+m\), because only the complete records carry the success marks \(q_{aj}\), and the least-favored arm supplies them at the overlap scale \(\epsilon\). Auxiliary records reduce marginal-count noise but cannot create outcome marks. The choice \(L=\lfloor\log S/1024\rfloor\) balances the sparse-cell approximation bias with the coefficient and outcome noise controlled by the complete sample. The formal clipping rule and complete prefix-average estimator appear in \cref{obj:synth_4,obj:def:estimator-handle}; this section explains their mechanism and compares their guarantee with a simpler inverse-count baseline.





A simpler comparison isolates the role of the polynomial correction. The inverse-count baseline \leanref{sym:T^{\mathrm B}}{\ensuremath{T^{\mathrm B}}} uses an outcome pool and a single marginal pool, applying the same count-plus-one correction throughout. It retains the clipping and finite-averaging construction, as specified next.





\begin{algorithmv}{def:baseline}[Inverse-count baseline $T^{\mathrm B}$]
The baseline takes the ordered complete and auxiliary records \(\mathcal D\) and the public parameters \(n,m,d,\epsilon\) as inputs.
\begin{enumerate}
\item If \(n<6\), output
\[
T^{\mathrm B}(\mathcal D)=0.
\]
For \(n\ge6\), proceed through the remaining steps.

\item Define the pool lengths and Poisson prefix means by
\[
h_{\mathrm o}=\lfloor n/2\rfloor,\qquad
h_{\mathrm m}=n-h_{\mathrm o}+m,\qquad
u_{\mathrm B}=h_{\mathrm o}/8,\qquad
t_{\mathrm B}=h_{\mathrm m}/8.
\]
Use the first \(h_{\mathrm o}\) complete records as the outcome pool. Form the marginal pool by concatenating the treatment--covariate projections of the remaining complete records with all auxiliary records, preserving their order.

\item Write the outcome-pool observations as
\((X_i^{\mathrm o},A_i^{\mathrm o},Y_i^{\mathrm o})\) and the marginal-pool observations as
\((X_i^{\mathrm m},A_i^{\mathrm m})\). For each prefix pair satisfying
\[
0\le h'\le h_{\mathrm o},\qquad
0\le h''\le h_{\mathrm m},
\]
and for \(a\in\{0,1\}\), \(j\in\{1,\ldots,d\}\), define
\[
\begin{aligned}
Z_{aj}^{\mathrm B}
&=\sum_{i=1}^{h'}
\mathbf1\{X_i^{\mathrm o}=j,\ A_i^{\mathrm o}=a,\ Y_i^{\mathrm o}=1\},\\
K_{aj}^{\mathrm B}
&=\sum_{i=1}^{h''}
\mathbf1\{X_i^{\mathrm m}=j,\ A_i^{\mathrm m}=a\},\\
H_{aj}^{\mathrm B}
&=\frac{Z_{aj}^{\mathrm B}}{u_{\mathrm B}}
\left(1+\frac{K_{1-a,j}^{\mathrm B}}{K_{aj}^{\mathrm B}+1}\right),\\
F_{\mathrm B}(h',h'';\mathcal D)
&=\operatorname{clip}_{[-1,1]}
\left(\sum_{j=1}^d(H_{1j}^{\mathrm B}-H_{0j}^{\mathrm B})\right).
\end{aligned}
\]

\item Output the deterministic finite average
\[
T^{\mathrm B}(\mathcal D)
=
\sum_{h'=0}^{h_{\mathrm o}}\sum_{h''=0}^{h_{\mathrm m}}
\exp(-u_{\mathrm B}-t_{\mathrm B})
\frac{u_{\mathrm B}^{h'}}{h'!}
\frac{t_{\mathrm B}^{h''}}{h''!}
F_{\mathrm B}(h',h'';\mathcal D).
\]
The weights are the probabilities of independent Poisson prefix lengths with means \(u_{\mathrm B}\) and \(t_{\mathrm B}\); prefix pairs exceeding either pool length have output zero.
\end{enumerate}
\end{algorithmv}

The denominator regularizes the marginal correction even when an observed arm--cell is empty. Outcome information enters through the success counts, while auxiliary records enlarge the pool estimating the treatment--covariate distribution. The following guarantee quantifies these two contributions uniformly over the observable model.

\begin{theoremv}{thm:uniform-baseline}[Uniform inverse-count risk guarantee]
There exists a universal numerical constant \(C>0\), independent of \(n,m,d,\epsilon\), such that, for every experiment satisfying
\begin{itemize}
\item \textbf{(Sample sizes.)} \(n,m\) are nonnegative integers with \(n\ge 1\).
\item \textbf{(Alphabet size.)} \(d\) is an integer with \(d\ge 2\).
\item \textbf{(Overlap parameter.)} \(\epsilon\) is real and \(0<\epsilon\le 1/4\).
\end{itemize}
the baseline \(T^{\mathrm B}\) of \cref{obj:def:baseline} is measurable and satisfies
\[
-1\le T^{\mathrm B}(\mathcal D)\le 1
\]
for every data realization \(\mathcal D\). For every observable law \(P\in\mathcal M_{d,\epsilon}\), its squared-error risk satisfies
\[
\mathbb E_P\!\left[
  \bigl(T^{\mathrm B}(\mathcal D)-\tau(P)\bigr)^2
\right]
\le
C\min\!\left\{
1,\,
\frac{1}{n\epsilon}
+
\left(\frac{d}{(n+m)\epsilon}\right)^2
\right\}.
\]
The expectation is over the two-channel data under \(P\); the rule ignores the independent uniform randomization seed.
\end{theoremv}

The bound reflects the different roles of the two pools. Rare-arm outcome information determines the term \(1/(n\epsilon)\), while the marginal pool controls the many-cell term through the total number of treatment--covariate observations. Increasing the auxiliary budget therefore sharpens the baseline guarantee through the second term. Comparing the two displayed worst-case upper guarantees, the attaining estimator replaces the baseline guarantee's many-cell term by one smaller by the squared regularized logarithmic factor \(\ell^2\), up to numerical constants, while preserving the inverse rare-arm information term. This comparison motivates the polynomial correction in sparse arm--cells and provides a transparent uniform benchmark for the three-pool construction.

The supporting moment bounds appear in \cref{obj:lem:poisson-inverse-moments,obj:lem:inverse-count-arm-risk}. The polynomial certificate and the localized arm-risk bound are given in \cref{obj:lem:chebyshev-factorial-certificate,obj:lem:uniform-hybrid-arm-risk}. Finally, \cref{obj:lem:finite-prefix-transfer} supplies the transfer argument connecting the Poisson calculations to the finite prefix averages on the original records.


\section{Consistency and the value of auxiliary records}

The risk comparison in \cref{obj:thm:uniform-annotation-frontier} answers three resource questions: which budgets support consistency, how much marginal information attains rare-label oracle order, and when auxiliary records yield a vanishing risk ratio relative to complete records alone. These questions allow the sample sizes, alphabet size, and overlap floor to vary jointly. We first specify the domain of experiment sequences.

\begin{definitionv}{synth_8}[Allowed experiment sequences]
An experiment sequence \(\mathbf v\) is an arbitrary sequence of public experiment indices, indexed by the nonnegative integers:
\[
\mathbf v=\bigl((n_k,m_k,d_k,\epsilon_k)\bigr)_{k=0}^{\infty}.
\]
For every \(k\), the complete-record size \(n_k\), auxiliary size \(m_k\), and alphabet size \(d_k\) are integers satisfying
\[
n_k\ge 1,\qquad m_k\ge 0,\qquad d_k\ge 2,
\]
and the public overlap floor \(\epsilon_k\) is a real number satisfying
\[
0<\epsilon_k\le \frac14.
\]
\end{definitionv}

An experiment sequence \leanref{sym:\mathbf v}{\ensuremath{\mathbf v}} permits arbitrary joint variation within these public ranges, including diminishing overlap and an expanding alphabet. The three resource criteria concern vanishing minimax risk, bounded risk relative to the rare-label benchmark, and vanishing risk relative to the supervised experiment. The benchmark retains its cap:
\[
b(n,\epsilon)=\min\{1,(n\epsilon)^{-1}\}.
\]
The following definition makes these comparisons precise and introduces the sequence-specific information scales.

\begin{definitionv}{def:phases}[Resource phases \(\mathcal C,\mathcal O,\mathcal I\)]
Define
\[
\begin{aligned}
\mathcal C
&=\left\{\mathbf v:
R(n_k,m_k,d_k,\epsilon_k)\longrightarrow 0\right\},\\
\mathcal O
&=\left\{\mathbf v:
n_k\epsilon_k\longrightarrow\infty,\quad
\limsup_{k\to\infty}
\frac{R(n_k,m_k,d_k,\epsilon_k)}{b(n_k,\epsilon_k)}
<\infty\right\},\\
\mathcal I
&=\left\{\mathbf v:
\frac{R(n_k,m_k,d_k,\epsilon_k)}
{R(n_k,0,d_k,\epsilon_k)}
\longrightarrow 0\right\}.
\end{aligned}
\]
Each set ranges over all allowed sequences of public experiment indices
\[
\mathbf v=((n_k,m_k,d_k,\epsilon_k))_{k\ge1}.
\]
For each such sequence, the rare-arm information scale and total marginal-information sample size are
\[
S_k=n_k\epsilon_k,\qquad N_k=n_k+m_k.
\]
Write \(\ell_k\) for the regularized logarithm \(\ell\) evaluated at the \(k\)-th experiment.
\end{definitionv}

Explicitly, the regularized logarithm is \(\ell_k=\log(\exp(1)+S_k)\). Sequences are indexed by \(k\ge0\) as in \cref{obj:synth_8}; the \(k\ge1\) display in \cref{obj:def:phases} denotes the same asymptotic tail and discards only the initial coordinate. The consistency phase \leanref{sym:\mathcal C}{\ensuremath{\mathcal C}} concerns absolute precision. The oracle phase \leanref{sym:\mathcal O}{\ensuremath{\mathcal O}} combines a diverging rare-arm information scale with attainment of the rare-label benchmark up to a constant factor. The improvement phase \leanref{sym:\mathcal I}{\ensuremath{\mathcal I}} measures a strict relative gain from auxiliary records. The next result characterizes all three criteria and supplies finite-experiment conditions for saturation and rare-label order.

\begin{theoremv}{thm:annotation-resource-phases}[Exact annotation resource phases]
Let \(\mathbf v=((n_k,m_k,d_k,\epsilon_k))_{k\ge0}\) be any sequence of allowed public experiment indices, and put
\[
S_k=n_k\epsilon_k,\qquad
N_k=n_k+m_k,\qquad
\ell_k=\log(\exp(1)+S_k).
\]
For the phases defined in \cref{obj:def:phases}, the following equivalences hold as \(k\to\infty\):
\[
\begin{aligned}
\mathbf v\in\mathcal C
&\Longleftrightarrow
S_k\longrightarrow\infty
\quad\text{and}\quad
d_k=o(N_k\epsilon_k\ell_k),\\
\mathbf v\in\mathcal O
&\Longleftrightarrow
S_k\longrightarrow\infty
\quad\text{and}\quad
d_k=O\!\left(\frac{N_k\epsilon_k\ell_k}{\sqrt{S_k}}\right),\\
\mathbf v\in\mathcal I
&\Longleftrightarrow
S_k\longrightarrow\infty,\quad
\frac{d_k}{\sqrt{S_k}\ell_k}\longrightarrow\infty,\quad
\frac{N_k/n_k}{\max\{1,d_k/(S_k\ell_k)\}}\longrightarrow\infty.
\end{aligned}
\]
If \(m_k=0\) for every \(k\), then
\[
\begin{aligned}
\mathbf v\in\mathcal C
&\Longleftrightarrow
S_k\longrightarrow\infty
\quad\text{and}\quad d_k=o(S_k\ell_k),\\
\mathbf v\in\mathcal O
&\Longleftrightarrow
S_k\longrightarrow\infty
\quad\text{and}\quad d_k=O(\sqrt{S_k}\ell_k),
\end{aligned}
\qquad
\mathbf v\notin\mathcal I.
\]
Whenever \(S_k\to\infty\), oracle membership also has the equivalent budget characterization
\[
\mathbf v\in\mathcal O
\quad\Longleftrightarrow\quad
\frac{d_k\sqrt{S_k}}{\epsilon_k\ell_k}=O(N_k).
\]

Write \(R\) for the minimax risk in \cref{obj:def:risk}. For every \(\mathbf v\in\mathcal O\), there exist constants \(c,C>0\) such that, for all sufficiently large \(k\),
\[
\frac{c}{S_k}
\le R(n_k,m_k,d_k,\epsilon_k)
\le \frac{C}{S_k}.
\]
Moreover, there exists \(c>0\) such that, for all sufficiently large \(k\) and every \(m'\in\mathbb N\),
\[
\frac{c}{S_k}\le R(n_k,m',d_k,\epsilon_k).
\]

For a single experiment, let \(n,m,d\in\mathbb N\) and \(\epsilon\in\mathbb R\) satisfy
\begin{itemize}
\item \textbf{(Complete-record budget.)} \(n\ge1\).
\item \textbf{(Alphabet size.)} \(d\ge2\).
\item \textbf{(Overlap range.)} \(0<\epsilon\le1/4\).
\end{itemize}
Put
\[
N=n+m,\qquad S=n\epsilon,\qquad
\ell=\log(\exp(1)+S),\qquad
b(n,\epsilon)=\min\{1,S^{-1}\}.
\]
There exists a universal constant \(c>0\) such that, for every such experiment,
\[
N\epsilon\ell\le d
\quad\Longrightarrow\quad
c\le R(n,m,d,\epsilon)\le1.
\]
There also exist universal constants \(c,C>0\) such that, for every such experiment,
\[
\frac{d\sqrt{n\epsilon}}{\epsilon\ell}\le N
\quad\Longrightarrow\quad
c\,b(n,\epsilon)\le R(n,m,d,\epsilon)\le C\,b(n,\epsilon).
\]
\end{theoremv}

The theorem separates the two resources directly. Complete records determine \(S_k\), so no auxiliary budget can overcome bounded rare-arm outcome information. Both channels determine \(N_k\), so auxiliary records can reduce the many-cell term once \(S_k\to\infty\). For a single experiment, the two useful endpoints are
\[
N\epsilon\ell\le d
\quad\Longrightarrow\quad R\asymp1,
\qquad
N\ge \frac{d\sqrt S}{\epsilon\ell}
\quad\Longrightarrow\quad
R\asymp \min\{1,S^{-1}\},
\]
with universal comparison constants. The first is marginal-information saturation; the second is the rare-label plateau.

\subsection{Three illustrative sequences}

Let \(t=k+4\), so every displayed public index is legal and \(\ell_k\asymp\log t\).
\begin{enumerate}
\item \textbf{Supervised rare-label order.} Take
\[
n_k=t^4,\qquad \epsilon_k=t^{-1},\qquad m_k=0,\qquad d_k=t.
\]
Then \(S_k=t^3\) and \(d_k/(\sqrt{S_k}\ell_k)\to0\). The supervised experiment is already in \(\mathcal O\), and \(R(n_k,0,d_k,\epsilon_k)\asymp t^{-3}\).

\item \textbf{Improvement between two consistent risks.} Take
\[
n_k=t^4,\qquad \epsilon_k=t^{-1},\qquad d_k=t^2,\qquad m_k=t^6.
\]
Without auxiliary records the many-cell term gives
\[
R(n_k,0,d_k,\epsilon_k)\asymp \frac{1}{t^2(\log t)^2},
\]
while the added marginal budget reduces the risk to \(R(n_k,m_k,d_k,\epsilon_k)\asymp t^{-3}\). Both risks vanish, but their ratio is of order \((\log t)^2/t\to0\), so the sequence lies in \(\mathcal I\).

\item \textbf{Expansion from saturation to consistency.} Take
\[
n_k=t^4,\qquad \epsilon_k=t^{-1},\qquad d_k=t^5,\qquad m_k=t^8.
\]
The supervised scale satisfies \(n_k\epsilon_k\ell_k\asymp t^3\log t\ll d_k\), so its risk stays of constant order. With auxiliary records,
\((n_k+m_k)\epsilon_k\ell_k\asymp t^7\log t\), making
\(d_k=o((n_k+m_k)\epsilon_k\ell_k)\); the two-channel risk is then of order \(t^{-3}\) and is consistent.
\end{enumerate}

These examples are illustrations of the order characterizations, rather than finite-budget calibrations. They also show why strict improvement can compare two vanishing risks or move an experiment out of saturation, while the rare-label plateau remains controlled by complete-record outcomes.


\section{Discussion and limitations}

The results provide a population benchmark for passive outcome annotation from a common distribution. Under the independent sampling condition in \cref{obj:ass:data-product}, complete records supply outcomes together with treatment and covariates, while auxiliary records supply treatment and covariates. The minimax comparison in \cref{obj:thm:uniform-annotation-frontier} characterizes the precision of the population contrast in \cref{obj:def:target} as both budgets, the covariate alphabet, and the public overlap floor vary. Under \cref{obj:ass:consistency,obj:ass:exchangeability}, this contrast is the population average treatment effect.

The marginal-budget plateau gives these information sources a distinct economic interpretation. Additional auxiliary records improve knowledge of treatment--covariate frequencies and reduce the many-cell component of the risk. Once that component is at or below the rare-label benchmark, the minimax risk has order \(\min\{1,(n\epsilon)^{-1}\}\), by \cref{obj:thm:uniform-annotation-frontier,obj:thm:rare-label-floor}. Thus auxiliary expansion can bring estimation to a precision scale governed by rare-arm outcomes in the complete sample. On the oracle sequences characterized in \cref{obj:thm:annotation-resource-phases}, this scale is eventually \((n_k\epsilon_k)^{-1}\). The plateau concerns risk order: it identifies the information budget governing precision once treatment--covariate information is sufficiently abundant.

The exact-marginal experiment sharpens this interpretation. Supplying the entire treatment--covariate probability table yields rare-label risk order for every allowed alphabet size, as established in \cref{obj:prop:benchmark-recovery}. Moreover, for each fixed complete-record sample size \(n\ge1\), fixed alphabet size \(d\ge2\), and fixed overlap floor \(0<\epsilon\le1/4\), the minimax risk converges to the exact-marginal risk as the integer auxiliary budget tends to infinity. This fixed-parameter limit describes what increasingly precise frequency information can deliver with a given outcome sample. The common-marginal alternatives in \cref{obj:thm:rare-label-floor} explain the remaining statistical difficulty: identical treatment--covariate frequencies accompany different rare-arm outcome distributions and different target values.

For annotation planning, these comparisons distinguish investment in population frequency information from investment in outcome information. They provide a benchmark for the precision attainable when acquisition follows the common-distribution experiment. Allocation procedures that choose which records receive outcomes introduce a further design problem. The asymptotic-variance allocation and feasible batch-adaptive procedure of \citet[Theorems 4.2 and 5.1 of the arXiv manuscript]{NwankwoGoldkindZhou2026Annotations} provide complementary context for that problem under their stated budget, identification, positivity, and regularity conditions.

\subsection{Limitations and future work}

The attaining estimator in \cref{obj:def:estimator-handle} outputs zero whenever \(n\epsilon<\exp(4096)\). Its uniform result therefore certifies minimax order with proof-oriented constants rather than providing calibrated finite-budget tuning: the comparison constant absorbs this threshold over all realistic values below it. The resource conditions in \cref{obj:thm:annotation-resource-phases} should accordingly be read as order-based planning benchmarks, not as calibrated recommendations for a particular annotation budget. Computationally practical tuning, stable evaluation, and implementation at moderate sample sizes require further analysis. Adaptive allocation would also change the observation mechanism in \cref{obj:ass:data-product}; extending the benchmark to allocation rules would require accounting for the dependence and selection induced by those rules.

The paper's loss is squared error for a population treatment contrast. Interval estimation requires additional distributional and coverage guarantees, while policy objectives require specifying the decision class and the welfare criterion. These questions have different statistical requirements from the minimax point-estimation comparison.

Finally, the overlap floor is a public restriction on population propensities in occupied cells. Empirical overlap assessments introduce uncertainty about that restriction, particularly when observed cell counts are small. Extensions to continuous covariates, multiple treatments, or outcomes beyond the binary case would require corresponding model and information-budget characterizations. The binary finite-alphabet experiment provides a reference point for studying these extensions.


\appendix

\section*{Appendices}

\section{Causal interpretation and benchmark proofs}

The causal interpretation and the benchmark comparisons concern the same observable class in \cref{obj:def:model}, defined by the public overlap restriction. The potential-outcome statement describes compatible extensions of each law in that class. The benchmark comparisons concern the information supplied to estimation rules under the sampling experiment in \cref{obj:ass:data-product}.

\subsection{Compatible extensions and identification}

The existence statement uses a conditional construction with independent binary treatment and potential outcomes. We first specify the probability law used in that construction.

\begin{definitionv}{synth_7}[Bernoulli law]
For \(p\in[0,1]\), \(\operatorname{Bernoulli}(p)\) is the probability law on \(\{0,1\}\) assigning probability \(p\) to \(1\) and probability \(1-p\) to \(0\).
\end{definitionv}

This law expresses binary conditional distributions directly in terms of their success probabilities. In the following construction, the observed propensity and outcome regressions determine those probabilities within each occupied covariate cell.

\begin{lemmav}{lem:causal-realization}[Causal realization and identification]
Let \(d\) be a natural number, \(\epsilon\) a real number, and \(P\) a probability law of the complete observation \((X,A,Y)\). Suppose:
\begin{itemize}
\item \textbf{(Covariate alphabet.)} \(d\ge 2\).
\item \textbf{(Overlap range.)} \(0<\epsilon\le 1/4\).
\item \textbf{(Model membership.)} \(P\in\mathcal M_{d,\epsilon}\), as defined in \cref{obj:def:model}.
\end{itemize}
Then \(P\) admits a joint probability law \(H\) of \((X,A,Y,Y_0,Y_1)\), with binary potential outcomes \(Y_0,Y_1\), whose \((X,A,Y)\) marginal equals \(P\) and which satisfies \cref{obj:ass:consistency,obj:ass:exchangeability}. Call precisely these extensions compatible.

One compatible extension is given by
\[
\Pr_H(X=j)=p_j(P),
\]
and, for every occupied cell \(j\), by the conditional law
\[
\mathcal L_H(A,Y_0,Y_1\mid X=j)
=
\operatorname{Bernoulli}\!\bigl(e_j(P)\bigr)
\otimes
\operatorname{Bernoulli}\!\bigl(\mu_{0j}(P)\bigr)
\otimes
\operatorname{Bernoulli}\!\bigl(\mu_{1j}(P)\bigr),
\qquad Y=Y_A.
\]
For every compatible extension \(H\), the causal contrast equals the target defined in \cref{obj:def:target}:
\[
\mathbb E_H[Y_1-Y_0]=\tau(P).
\]
\end{lemmav}

\begin{proof}[Proof of \cref{obj:lem:causal-realization}]
\begin{enumerate}
\item
Throughout, \(j\in\{1,\ldots,d\}\), and all treatment and outcome indices range over \(\{0,1\}\). Recall the cell quantities
\[
\begin{aligned}
p_j(P)&=P(X=j),\\
s_{aj}(P)&=P(X=j,A=a),\\
q_{aj}(P)&=P(X=j,A=a,Y=1),\\
\mu_{aj}(P)&=
\begin{cases}
q_{aj}(P)/s_{aj}(P),&s_{aj}(P)>0,\\
0,&s_{aj}(P)=0.
\end{cases}
\end{aligned}
\]
The propensity convention in \cref{obj:synth_2} is
\[
e_j(P)=
\begin{cases}
s_{1j}(P)/p_j(P),&p_j(P)>0,\\
1/2,&p_j(P)=0.
\end{cases}
\]
Since these quantities are obtained by summing probabilities,
\[
p_j(P)\ge0,\qquad
\sum_{j=1}^{d}p_j(P)=1,\qquad
s_{0j}(P)+s_{1j}(P)=p_j(P),\qquad
0\le q_{aj}(P)\le s_{aj}(P).
\]
Thus \(e_j(P)\in[0,1]\): on an occupied cell this follows from
\(0\le s_{1j}(P)\le p_j(P)\), and on a null cell it follows from the
value \(1/2\). Similarly, division by \(s_{aj}(P)>0\), together with
the stipulated value on a null arm, gives
\[
0\le\mu_{aj}(P)\le1.
\]

Define the treatment and potential-outcome Bernoulli masses by
\[
\beta^A_j(a)=
\begin{cases}
1-e_j(P),&a=0,\\
e_j(P),&a=1,
\end{cases}
\qquad
\beta^Y_{aj}(z)=
\begin{cases}
1-\mu_{aj}(P),&z=0,\\
\mu_{aj}(P),&z=1.
\end{cases}
\]
These definitions apply to every cell and arm, including those with
zero mass. By \cref{obj:synth_7}, they are nonnegative and satisfy
\[
\sum_a\beta^A_j(a)=1,\qquad
\sum_z\beta^Y_{aj}(z)=1.
\]
For a binary pair \((z_0,z_1)\), write
\[
z_a=
\begin{cases}
z_0,&a=0,\\
z_1,&a=1.
\end{cases}
\]
Construct \(H_{\mathrm{ind}}\) through its atom masses:
\[
\Pr_{H_{\mathrm{ind}}}
 (X=j,A=a,Y=y,Y_0=z_0,Y_1=z_1)
=
p_j(P)\beta^A_j(a)
\beta^Y_{0j}(z_0)\beta^Y_{1j}(z_1)
\mathbf1\{y=z_a\}.
\]
Every entry is nonnegative. Summing the factual-outcome indicator
first and then the remaining binary coordinates yields
\[
\begin{aligned}
&\sum_{j=1}^{d}\sum_{a,y,z_0,z_1}
\Pr_{H_{\mathrm{ind}}}
 (X=j,A=a,Y=y,Y_0=z_0,Y_1=z_1)\\
&\quad=
\sum_{j=1}^{d}p_j(P)
 \left(\sum_a\beta^A_j(a)\right)
 \left(\sum_{z_0}\beta^Y_{0j}(z_0)\right)
 \left(\sum_{z_1}\beta^Y_{1j}(z_1)\right)
=1.
\end{aligned}
\]
Hence the displayed table defines a probability law; normalization
preserves its atom masses.
% lean: CausalSmith.Stat.AnnotationRarearmFrontier.independentExtension_fullMass

\item
Summing out both potential outcomes gives
\[
\begin{aligned}
&\Pr_{H_{\mathrm{ind}}}(X=j,A=a,Y=y)\\
&\quad=
\sum_{z_0,z_1}
p_j(P)\beta^A_j(a)
\beta^Y_{0j}(z_0)\beta^Y_{1j}(z_1)
\mathbf1\{y=z_a\}\\
&\quad=p_j(P)\beta^A_j(a)\beta^Y_{aj}(y).
\end{aligned}
\]
Indeed, for \(a=0\) the indicator selects \(z_0=y\) and the
\(z_1\)-sum equals one; for \(a=1\) it selects \(z_1=y\) and the
\(z_0\)-sum equals one.

The propensity weights reconstruct the treatment--covariate masses:
\[
p_j(P)\beta^A_j(a)=s_{aj}(P).
\]
For \(p_j(P)=0\), nonnegativity and
\(s_{0j}(P)+s_{1j}(P)=0\) give both arm masses zero. For
\(p_j(P)>0\), the two identities are
\[
p_j(P)e_j(P)=s_{1j}(P),\qquad
p_j(P)\{1-e_j(P)\}=p_j(P)-s_{1j}(P)=s_{0j}(P).
\]
Likewise,
\[
s_{aj}(P)\beta^Y_{aj}(y)=P(X=j,A=a,Y=y).
\]
If \(s_{aj}(P)=0\), both observed outcome atoms vanish. If
\(s_{aj}(P)>0\), the identities for success and failure are
\begingroup\small
\[
s_{aj}(P)\mu_{aj}(P)=q_{aj}(P),\qquad
s_{aj}(P)\{1-\mu_{aj}(P)\}
=s_{aj}(P)-q_{aj}(P)
=P(X=j,A=a,Y=0).
\]
\endgroup
Combining these calculations proves
\[
\Pr_{H_{\mathrm{ind}}}(X=j,A=a,Y=y)
=P(X=j,A=a,Y=y).
\]
Equality at every atom proves equality of the observed marginal with
\(P\), and in particular
\[
\Pr_{H_{\mathrm{ind}}}(X=j)=p_j(P).
\]
% lean: CausalSmith.Stat.AnnotationRarearmFrontier.independentExtension_observedMarginal

\item
The defining indicator gives
\[
\Pr_{H_{\mathrm{ind}}}
 (X=j,A=a,Y=y,Y_0=z_0,Y_1=z_1)=0
\qquad\text{whenever }y\ne z_a.
\]
Consequently,
\[
Y=AY_1+(1-A)Y_0
\qquad H_{\mathrm{ind}}\text{-almost surely},
\]
which is \cref{obj:ass:consistency}.
% lean: CausalSmith.Stat.AnnotationRarearmFrontier.independentExtension_consistency

\item
Summing over the factual outcome, and then using the unit sums of the
Bernoulli factors, gives
\[
\begin{aligned}
&\Pr_{H_{\mathrm{ind}}}(X=j,A=a,Y_0=z_0,Y_1=z_1)
=p_j(P)\beta^A_j(a)\beta^Y_{0j}(z_0)\beta^Y_{1j}(z_1),\\
&\Pr_{H_{\mathrm{ind}}}(X=j,A=a)=p_j(P)\beta^A_j(a),\\
&\Pr_{H_{\mathrm{ind}}}(X=j,Y_0=z_0,Y_1=z_1)
=p_j(P)\beta^Y_{0j}(z_0)\beta^Y_{1j}(z_1).
\end{aligned}
\]
Therefore, for every cell,
\[
\begin{aligned}
&\Pr_{H_{\mathrm{ind}}}(X=j,A=a,Y_0=z_0,Y_1=z_1)\,p_j(P)\\
&\quad=
\Pr_{H_{\mathrm{ind}}}(X=j,A=a)\,
\Pr_{H_{\mathrm{ind}}}(X=j,Y_0=z_0,Y_1=z_1).
\end{aligned}
\]
On an occupied cell, division by \(p_j(P)^2\) proves conditional
independence of \(A\) and \((Y_0,Y_1)\). More explicitly,
\[
\Pr_{H_{\mathrm{ind}}}
 (A=a,Y_0=z_0,Y_1=z_1\mid X=j)
=\beta^A_j(a)\beta^Y_{0j}(z_0)\beta^Y_{1j}(z_1).
\]
This is the stated product of three Bernoulli laws. Null cells have
zero probability, so the calculation establishes
\cref{obj:ass:exchangeability}. Together with the preceding steps,
\(H_{\mathrm{ind}}\) is a compatible extension.
% lean: CausalSmith.Stat.AnnotationRarearmFrontier.independentExtension_exchangeability

\item
Now fix any compatible extension \(H\). Define its potential-outcome
atom masses and its cellwise potential-success masses by
\[
\begin{aligned}
\rho^H_j(a,z_0,z_1)
&=\Pr_H(X=j,A=a,Y_0=z_0,Y_1=z_1)\\
&=\sum_y\Pr_H(X=j,A=a,Y=y,Y_0=z_0,Y_1=z_1),\\
\eta^H_{aj}
&=\sum_{a',z_0,z_1}\rho^H_j(a',z_0,z_1)\,z_a.
\end{aligned}
\]
These definitions apply to every cell and both arms. Expanding the
finite expectation, summing out \(Y\), and distributing the
difference gives
\[
\begin{aligned}
\mathbb E_H[Y_1-Y_0]
&=\sum_{j=1}^{d}\sum_{a,y,z_0,z_1}
 \Pr_H(X=j,A=a,Y=y,Y_0=z_0,Y_1=z_1)(z_1-z_0)\\
&=\sum_{j=1}^{d}\bigl(\eta^H_{1j}-\eta^H_{0j}\bigr).
\end{aligned}
\]
We identify the two success masses in each summand.
% lean: CausalSmith.Stat.AnnotationRarearmFrontier.poContrast_eq_sum_potentialSuccessMass

\item
Since the observed marginal of \(H\) equals \(P\), summing its atom
masses gives
\[
p_j(P)=\sum_{a',z_0,z_1}\rho^H_j(a',z_0,z_1),
\qquad
s_{aj}(P)=\sum_{z_0,z_1}\rho^H_j(a,z_0,z_1).
\]
By \cref{obj:ass:consistency},
\[
\Pr_H(X=j,A=a,Y=1-z_a,Y_0=z_0,Y_1=z_1)=0.
\]
Thus summing the factual-success atoms selects precisely those
potential-outcome pairs with \(z_a=1\), yielding
\[
q_{aj}(P)=\sum_{z_0,z_1}\rho^H_j(a,z_0,z_1)\,z_a.
\]
Conditional exchangeability in
\cref{obj:ass:exchangeability} supplies the atom identity
\[
\rho^H_j(a,z_0,z_1)\,p_j(P)
=
s_{aj}(P)\sum_{a'}\rho^H_j(a',z_0,z_1).
\]
For \(p_j(P)>0\), this follows by multiplying the conditional
factorization by \(p_j(P)^2\). For \(p_j(P)=0\), all the corresponding
atom masses are zero, so the same identity holds.

Multiplying this identity by \(z_a\) and summing over the potential
outcomes proves
\[
\begin{aligned}
q_{aj}(P)p_j(P)
&=\sum_{z_0,z_1}
 \rho^H_j(a,z_0,z_1)p_j(P)\,z_a\\
&=s_{aj}(P)\sum_{z_0,z_1}\sum_{a'}
 \rho^H_j(a',z_0,z_1)\,z_a\\
&=s_{aj}(P)\eta^H_{aj}.
\end{aligned}
\]
% lean: CausalSmith.Stat.AnnotationRarearmFrontier.compatible_success_identity

\item
Fix \(j\) and \(a\). First suppose \(p_j(P)=0\).
Nonnegativity of the atom masses and \(0\le z_a\le1\) give
\[
0\le\eta^H_{aj}
\le\sum_{a',z_0,z_1}\rho^H_j(a',z_0,z_1)
=p_j(P)=0.
\]
Hence
\[
\eta^H_{aj}=0=p_j(P)\mu_{aj}(P).
\]

Next suppose \(p_j(P)>0\). Model membership in
\cref{obj:def:model} supplies \cref{obj:ass:overlap}, so
\[
s_{1j}(P)=p_j(P)e_j(P)\ge\epsilon p_j(P)>0,
\qquad
s_{0j}(P)=p_j(P)\{1-e_j(P)\}\ge\epsilon p_j(P)>0.
\]
Dividing the preceding success identity by the positive arm mass
therefore gives
\[
\eta^H_{aj}
=p_j(P)\frac{q_{aj}(P)}{s_{aj}(P)}
=p_j(P)\mu_{aj}(P).
\]
The two cases establish this identity for every cell and both arms.
% lean: CausalSmith.Stat.AnnotationRarearmFrontier.compatible_cell_mean

\item
Substituting first the treatment success mass and then the control
success mass into the contrast expansion yields
\[
\begin{aligned}
\mathbb E_H[Y_1-Y_0]
&=\sum_{j=1}^{d}
 \bigl\{p_j(P)\mu_{1j}(P)-p_j(P)\mu_{0j}(P)\bigr\}\\
&=\sum_{j=1}^{d}p_j(P)
 \bigl\{\mu_{1j}(P)-\mu_{0j}(P)\bigr\}
=\tau(P),
\end{aligned}
\]
where the last equality is \cref{obj:def:target}. Since \(H\) was
arbitrary, this identifies the contrast for every compatible
extension.
% lean: CausalSmith.Stat.AnnotationRarearmFrontier.compatible_poContrast
\end{enumerate}
\end{proof}

The product construction supplies an existence witness for every observable law in the stated class. It chooses conditional independence between the two potential outcomes as well as between treatment and the potential-outcome pair. Compatibility itself requires the observed marginal, consistency, and conditional treatment exchangeability; the identifying equality applies to every extension satisfying those requirements. Thus the causal contrast is invariant across compatible choices of the joint potential-outcome distribution.

The public overlap restriction gives both treatment arms positive probability in every occupied covariate cell. Under a compatible extension, consistency and conditional exchangeability identify each conditional potential-outcome mean with its corresponding observed outcome regression. The target in \cref{obj:def:target} aggregates their contrast using the population covariate weights. Null covariate cells have zero weight under both the observable law and its extensions. These roles keep the observable restriction in \cref{obj:ass:overlap} distinct from the identification conditions imposed on the extension in \cref{obj:ass:consistency,obj:ass:exchangeability}.

\subsection{Two-cell and exact-marginal comparisons}

The benchmark functions \(b(n,\epsilon)\) and \(f(n,m,d)\) are defined in \cref{obj:prop:benchmark-recovery}: the former is the capped inverse rare-arm outcome scale, and the latter combines inverse complete-sample precision with a many-cell term at fixed overlap. The numerical rate \(r(n,m,d,\epsilon)\) is defined in \cref{obj:def:rate-handle}. These definitions distinguish the numerical comparisons from the minimax risks in \cref{obj:def:risk,obj:def:known-risk}.

For a two-cell covariate alphabet, \cref{obj:prop:benchmark-recovery} gives rare-label order uniformly over the auxiliary budget and the allowed overlap range. The uniform risk comparison in \cref{obj:thm:uniform-annotation-frontier} supplies the connection between the numerical rate and attainable precision. The matching lower bound in \cref{obj:thm:rare-label-floor} locates the persistent difficulty in the treated outcomes of the complete records.

The exact-marginal comparison has the same rare-label order for every allowed alphabet size. In this experiment, the population covariate weights and treatment propensities are supplied through the marginal table, as specified in \cref{obj:def:known-risk}. The common-marginal alternatives in \cref{obj:thm:rare-label-floor} retain their target separation when this table is supplied. Together, the two comparisons in \cref{obj:prop:benchmark-recovery} characterize experiments whose minimax precision is governed by complete-record outcome information, with numerical comparison constants uniform over their stated public indices.

At each fixed positive overlap floor, the comparison with \(f(n,m,d)\) in \cref{obj:prop:benchmark-recovery} expresses the rate in the conventional complete-sample and total-sample scales. Its constants may depend on that fixed floor. The two-cell and exact-marginal comparisons retain uniform constants as the floor varies within the stated range.

\subsection{The fixed-parameter experiment limit}

The exact-marginal risk provides the limiting benchmark for increasing auxiliary information. For \(n\ge1\), \(m\ge1\), \(d\ge2\), and \(0<\epsilon\le1/4\), \cref{obj:prop:benchmark-recovery} bounds the nonnegative difference between the two-channel and exact-marginal risks by
\[
\min\left\{1,6(n+2)\sqrt{\frac{2d}{m}}\right\}.
\]
This comparison concerns the minimax risk itself. At fixed \(n,d,\epsilon\), its dependence on the auxiliary sample size quantifies convergence to the exact-marginal experiment.

The same proposition establishes that bounded rare-arm outcome information gives a uniformly positive risk floor over all auxiliary budgets and allowed alphabet sizes. The experiment limit and this floor therefore describe complementary aspects of auxiliary information: increasing auxiliary records approaches the precision available with the population marginal supplied, while the rare-label benchmark characterizes the outcome information carried by the fixed complete sample.


\section{Estimator risk bounds}

The upper risk comparisons combine three ingredients: moment bounds for smoothed inverse counts, polynomial estimation within a localized range of arm-cell probabilities, and finite averaging over ordered prefixes. The auxiliary Poisson experiments separate the relevant counts, while the prefix transfer returns their risk bounds to the original fixed-size experiment.

\subsection{Poisson counts and inverse-count contributions}

We first collect the probability laws, count variables, and arm-level quantities used in the inverse-count calculations.



Adding one to the denominator makes the contribution well defined at an empty arm count. The following moment bounds quantify the resulting smoothing and its dispersion over the full range of Poisson intensities.

\begin{lemmav}{lem:poisson-inverse-moments}[Inverse Poisson count moments]
Let \(\lambda\ge 0\), let \(K\sim\operatorname{Poi}(\lambda)\) have the Poisson law with mean \(\lambda\), and define
\[
g(K)=\frac{1}{K+1},
\qquad C_0=2^{16}.
\]
Then
\[
\mathbb E[g(K)]
=
\begin{cases}
1,&\lambda=0,\\
\dfrac{1-e^{-\lambda}}{\lambda},&\lambda>0,
\end{cases}
\qquad
\mathbb E[g(K)^2]\le \frac{C_0}{(1+\lambda)^2},
\qquad
\operatorname{Var}(g(K))\le \frac{C_0\lambda}{(1+\lambda)^4}.
\]
\end{lemmav}

\begin{proof}[Proof of \cref{obj:lem:poisson-inverse-moments}]
\begin{enumerate}
\item For \(k\in\mathbb Z_{\ge0}\), write
\[
p_\lambda(k)=\Pr(K=k)=e^{-\lambda}\frac{\lambda^k}{k!},
\qquad
g(k)=\frac1{k+1},
\qquad
r_{\mathrm{inv}}(k)=\frac1{(k+1)(k+2)},
\]
with \(0^0=1\). For any real sequence, define its forward increment by
\[
\Delta f_{\mathrm{aux}}(k)
=f_{\mathrm{aux}}(k+1)-f_{\mathrm{aux}}(k),
\qquad
f_{\mathrm{aux}}\colon\mathbb Z_{\ge0}\to\mathbb R.
\]
Then
\[
0<g(k)\le1,\qquad 0<r_{\mathrm{inv}}(k)\le1,
\qquad
\Delta g(k)
=\frac1{k+2}-\frac1{k+1}
=-r_{\mathrm{inv}}(k).
\]
Thus \(g(K)\), \(r_{\mathrm{inv}}(K)\), and \(\Delta g(K)\) are square-integrable. % lean: Causalean.Mathlib.Probability.Poisson.shifted_reciprocal_memLp_two
% lean: Causalean.Mathlib.Probability.Poisson.shifted_reciprocal_pair_memLp_two

\item We first establish the forward-increment variance inequality for a sequence \(f_{\mathrm{aux}}\) of finite support. Expanding the square and using \(\sum_{k\ge0}p_\lambda(k)=1\) gives
\[
\begin{aligned}
\operatorname{Var}(f_{\mathrm{aux}}(K))
&=\frac12\sum_{i,j\ge0}p_\lambda(i)p_\lambda(j)
  \bigl(f_{\mathrm{aux}}(j)-f_{\mathrm{aux}}(i)\bigr)^2\\
&=\sum_{0\le i<j}p_\lambda(i)p_\lambda(j)
  \bigl(f_{\mathrm{aux}}(j)-f_{\mathrm{aux}}(i)\bigr)^2.
\end{aligned}
\]
For \(0\le i<j\), telescoping and Cauchy--Schwarz yield
\[
f_{\mathrm{aux}}(j)-f_{\mathrm{aux}}(i)
=\sum_{k=i}^{j-1}\Delta f_{\mathrm{aux}}(k),
\qquad
\bigl(f_{\mathrm{aux}}(j)-f_{\mathrm{aux}}(i)\bigr)^2
\le (j-i)\sum_{k=i}^{j-1}\bigl(\Delta f_{\mathrm{aux}}(k)\bigr)^2.
\]

The Poisson mass recurrence is
\[
(k+1)p_\lambda(k+1)=\lambda p_\lambda(k)
\qquad(k\ge0).
\]
Consequently, for every \(k\ge0\),
\[
\sum_{i=0}^{k}i\,p_\lambda(i)
=\lambda\sum_{i<k}p_\lambda(i),
\qquad
\sum_{j=k+1}^{\infty}j\,p_\lambda(j)
=\lambda\sum_{j=k}^{\infty}p_\lambda(j).
\]
Here and below the sum over \(i<k\) is empty when \(k=0\). These identities imply
\[
\begin{aligned}
&\sum_{i=0}^{k}\sum_{j=k+1}^{\infty}
p_\lambda(i)p_\lambda(j)(j-i)\\
&\quad=\lambda\left[
\left(\sum_{i=0}^{k}p_\lambda(i)\right)
\left(\sum_{j=k}^{\infty}p_\lambda(j)\right)
-
\left(\sum_{i<k}p_\lambda(i)\right)
\left(\sum_{j=k+1}^{\infty}p_\lambda(j)\right)
\right]\\
&\quad=\lambda p_\lambda(k)
\left[
\sum_{i<k}p_\lambda(i)+p_\lambda(k)
+\sum_{j=k+1}^{\infty}p_\lambda(j)
\right]
=\lambda p_\lambda(k).
\end{aligned}
\]
All these identities include \(\lambda=0\). Interchanging the nonnegative path sums now gives
\[
\begin{aligned}
\operatorname{Var}(f_{\mathrm{aux}}(K))
&\le
\sum_{0\le i<j}p_\lambda(i)p_\lambda(j)(j-i)
\sum_{k=i}^{j-1}\bigl(\Delta f_{\mathrm{aux}}(k)\bigr)^2\\
&=
\sum_{k\ge0}\bigl(\Delta f_{\mathrm{aux}}(k)\bigr)^2
\sum_{i=0}^{k}\sum_{j=k+1}^{\infty}
p_\lambda(i)p_\lambda(j)(j-i)\\
&=\lambda\,\mathbb E\bigl[(\Delta f_{\mathrm{aux}}(K))^2\bigr].
\end{aligned}
\]
The last sum is finite because \(\Delta f_{\mathrm{aux}}\) has finite support. % lean: Causalean.Mathlib.Probability.PoissonAddOnePoincare.poisson_addOne_poincare_finiteSupport

\item To extend this inequality, suppose that \(f_{\mathrm{aux}}(K)\) and \(\Delta f_{\mathrm{aux}}(K)\) are square-integrable. For \(N_{\mathrm{cut}}\in\mathbb Z_{\ge0}\), define
\[
f_{\mathrm{aux},N_{\mathrm{cut}}}(k)
=
\begin{cases}
f_{\mathrm{aux}}(k),&k<N_{\mathrm{cut}},\\
0,&k\ge N_{\mathrm{cut}},
\end{cases}
\]
and the errors
\[
\begin{aligned}
e^{\mathrm{val}}_{N_{\mathrm{cut}}}(k)
&=f_{\mathrm{aux},N_{\mathrm{cut}}}(k)-f_{\mathrm{aux}}(k),\\
e^{\mathrm{inc}}_{N_{\mathrm{cut}}}(k)
&=\Delta f_{\mathrm{aux},N_{\mathrm{cut}}}(k)
  -\Delta f_{\mathrm{aux}}(k).
\end{aligned}
\]
The value error satisfies
\[
\mathbb E\bigl[(e^{\mathrm{val}}_{N_{\mathrm{cut}}}(K))^2\bigr]
=\sum_{k\ge N_{\mathrm{cut}}}p_\lambda(k)f_{\mathrm{aux}}(k)^2
\longrightarrow0.
\]
For the increment error, the identity
\[
f_{\mathrm{aux}}(k+1)
=\Delta f_{\mathrm{aux}}(k)+f_{\mathrm{aux}}(k)
\]
implies
\[
f_{\mathrm{aux}}(k+1)^2
\le2\bigl((\Delta f_{\mathrm{aux}}(k))^2+f_{\mathrm{aux}}(k)^2\bigr).
\]
Thus the shifted sequence is square-integrable. The truncation definition also gives
\[
\bigl(e^{\mathrm{inc}}_{N_{\mathrm{cut}}}(k)\bigr)^2
\le2\bigl(f_{\mathrm{aux}}(k+1)^2+f_{\mathrm{aux}}(k)^2\bigr).
\]
For each fixed \(k\), this error is zero once \(N_{\mathrm{cut}}\ge k+2\). Dominated convergence therefore yields
\[
\mathbb E\bigl[(e^{\mathrm{inc}}_{N_{\mathrm{cut}}}(K))^2\bigr]
\longrightarrow0.
\]

Cauchy--Schwarz and expansion of the squares give
\[
\left|
\mathbb E[f_{\mathrm{aux},N_{\mathrm{cut}}}(K)]
-\mathbb E[f_{\mathrm{aux}}(K)]
\right|
\le
\sqrt{\mathbb E\bigl[(e^{\mathrm{val}}_{N_{\mathrm{cut}}}(K))^2\bigr]},
\]
\[
\begin{aligned}
&\left|
\mathbb E[f_{\mathrm{aux},N_{\mathrm{cut}}}(K)^2]
-\mathbb E[f_{\mathrm{aux}}(K)^2]
\right|\\
&\quad\le
\mathbb E\bigl[(e^{\mathrm{val}}_{N_{\mathrm{cut}}}(K))^2\bigr]
+2\sqrt{
\mathbb E[f_{\mathrm{aux}}(K)^2]\,
\mathbb E\bigl[(e^{\mathrm{val}}_{N_{\mathrm{cut}}}(K))^2\bigr]},
\end{aligned}
\]
and
\[
\begin{aligned}
&\left|
\mathbb E[(\Delta f_{\mathrm{aux},N_{\mathrm{cut}}}(K))^2]
-\mathbb E[(\Delta f_{\mathrm{aux}}(K))^2]
\right|\\
&\quad\le
\mathbb E\bigl[(e^{\mathrm{inc}}_{N_{\mathrm{cut}}}(K))^2\bigr]
+2\sqrt{
\mathbb E[(\Delta f_{\mathrm{aux}}(K))^2]\,
\mathbb E\bigl[(e^{\mathrm{inc}}_{N_{\mathrm{cut}}}(K))^2\bigr]}.
\end{aligned}
\]
Hence both the variance and the expected squared increment of the truncations converge to those of \(f_{\mathrm{aux}}\). Passing to the limit in the finite-support inequality proves
\[
\operatorname{Var}(f_{\mathrm{aux}}(K))
\le\lambda\,\mathbb E[(\Delta f_{\mathrm{aux}}(K))^2].
\]
Applying this to \(g\), with the square-integrability established above, gives
\[
\operatorname{Var}(g(K))
\le\lambda\,\mathbb E[r_{\mathrm{inv}}(K)^2].
\] % lean: Causalean.Mathlib.Probability.PoissonAddOnePoincare.poisson_addOne_poincare

\item We next bound the second moment of \(g(K)\). Since \(0\le g(k)^2\le1\), for \(0\le\lambda\le1\),
\[
(1+\lambda)^2\mathbb E[g(K)^2]\le4\le8.
\]
For \(\lambda>1\), the pointwise comparison
\[
g(k)^2=\frac1{(k+1)^2}
\le\frac2{(k+1)(k+2)}
=2r_{\mathrm{inv}}(k)
\]
follows from \(k+2\le2(k+1)\). Applying the Poisson recurrence twice gives
\[
\lambda^2p_\lambda(k)r_{\mathrm{inv}}(k)=p_\lambda(k+2).
\]
The reciprocal-weighted series is summable because its terms lie between zero and \(p_\lambda(k)\). Summing the identity therefore yields
\[
\begin{aligned}
\lambda^2\mathbb E[r_{\mathrm{inv}}(K)]
&=\sum_{k\ge0}p_\lambda(k+2)\\
&=1-p_\lambda(0)-p_\lambda(1)
=1-e^{-\lambda}(1+\lambda).
\end{aligned}
\]
It follows that
\[
\lambda^2\mathbb E[g(K)^2]
\le2\lambda^2\mathbb E[r_{\mathrm{inv}}(K)]
=2\bigl(1-e^{-\lambda}(1+\lambda)\bigr)
\le2.
\]
Since \((1+\lambda)^2\le4\lambda^2\) for \(\lambda>1\), we obtain
\[
(1+\lambda)^2\mathbb E[g(K)^2]\le8.
\]
Together, the two cases establish this bound for every \(\lambda\ge0\). % lean: Causalean.Mathlib.Probability.Poisson.shifted_reciprocal_factorial_second_moment
% lean: Causalean.Mathlib.Probability.Poisson.shifted_reciprocal_sq_rate_bound
% lean: Causalean.Mathlib.Probability.Poisson.shifted_reciprocal_sq_bound

\item For the squared reciprocal pair, define
\[
r_{\mathrm{inv},4}(k)
=\frac1{(k+1)(k+2)(k+3)(k+4)}
\qquad(k\in\mathbb Z_{\ge0}).
\]
The calculation
\[
6(k+1)(k+2)-(k+3)(k+4)=5k^2+11k\ge0
\]
shows that
\[
r_{\mathrm{inv}}(k)^2\le6r_{\mathrm{inv},4}(k),
\qquad
\mathbb E[r_{\mathrm{inv}}(K)^2]
\le6\mathbb E[r_{\mathrm{inv},4}(K)].
\]
Also \(0\le r_{\mathrm{inv}}(k)^2\le1\). Thus, for \(0\le\lambda\le1\),
\[
(1+\lambda)^4\mathbb E[r_{\mathrm{inv}}(K)^2]
\le16\le96.
\]

For \(\lambda>1\), four applications of the Poisson recurrence give
\[
\lambda^4p_\lambda(k)r_{\mathrm{inv},4}(k)=p_\lambda(k+4).
\]
The reciprocal-weighted series is again dominated by the Poisson mass series, so summation gives
\[
\begin{aligned}
\lambda^4\mathbb E[r_{\mathrm{inv},4}(K)]
&=\sum_{k\ge0}p_\lambda(k+4)\\
&=1-\sum_{k=0}^{3}p_\lambda(k)\\
&=1-e^{-\lambda}
\left(1+\lambda+\frac{\lambda^2}{2}+\frac{\lambda^3}{6}\right).
\end{aligned}
\]
Consequently,
\[
\begin{aligned}
\lambda^4\mathbb E[r_{\mathrm{inv}}(K)^2]
&\le6\lambda^4\mathbb E[r_{\mathrm{inv},4}(K)]\\
&=6\left[
1-e^{-\lambda}
\left(1+\lambda+\frac{\lambda^2}{2}+\frac{\lambda^3}{6}\right)
\right]
\le6.
\end{aligned}
\]
Using \((1+\lambda)^4\le16\lambda^4\) now yields
\[
(1+\lambda)^4\mathbb E[r_{\mathrm{inv}}(K)^2]\le96.
\]
This bound therefore holds for every \(\lambda\ge0\). % lean: Causalean.Mathlib.Probability.Poisson.shifted_reciprocal_factorial_fourth_moment
% lean: Causalean.Mathlib.Probability.Poisson.shifted_reciprocal_pair_sq_bound

\item To evaluate the first moment, first take \(\lambda=0\). The Poisson masses then satisfy
\[
p_0(0)=1,\qquad p_0(k)=0\quad(k\ge1),
\]
and hence
\[
\mathbb E[g(K)]=g(0)=1.
\]
For \(\lambda>0\), the recurrence gives
\[
\lambda p_\lambda(k)g(k)=p_\lambda(k+1).
\]
Because \(0\le p_\lambda(k)g(k)\le p_\lambda(k)\), this identity can be summed:
\[
\lambda\mathbb E[g(K)]
=\sum_{k\ge0}p_\lambda(k+1)
=1-p_\lambda(0)
=1-e^{-\lambda}.
\]
Dividing by \(\lambda>0\) proves the asserted positive-rate formula. % lean: Causalean.Mathlib.Probability.Poisson.shifted_reciprocal_first_moment

\item Finally, \(1+\lambda>0\) and \(8\le2^{16}=C_0\), so
\[
\mathbb E[g(K)^2]
\le\frac8{(1+\lambda)^2}
\le\frac{C_0}{(1+\lambda)^2}.
\]
Multiplying the variance inequality by the positive quantity \((1+\lambda)^4\), and then using the reciprocal-pair bound and \(\lambda\ge0\), gives
\[
\begin{aligned}
(1+\lambda)^4\operatorname{Var}(g(K))
&\le\lambda(1+\lambda)^4\mathbb E[r_{\mathrm{inv}}(K)^2]\\
&\le96\lambda
\le C_0\lambda.
\end{aligned}
\]
Division by \((1+\lambda)^4>0\) proves
\[
\operatorname{Var}(g(K))
\le\frac{C_0\lambda}{(1+\lambda)^4},
\]
including \(\lambda=0\). % lean: CausalSmith.Stat.AnnotationRarearmFrontier.poisson_inverse_moments
\end{enumerate}
\end{proof}

The expectation identity records the smoothing error, while the second-moment and variance bounds control its contribution to sampling error. Applied to the independent outcome-success and marginal counts in \cref{obj:lem:inverse-count-arm-risk}, these controls give an arm-level bound that is uniform over the overlap class.

\begin{lemmav}{lem:inverse-count-arm-risk}[Inverse-count arm risk bounds]
There exists a universal numerical constant \(C>0\) such that the following holds.
\begin{itemize}
\item \textbf{(Parameters.)} Let \(d\ge 2\) be an integer, \(0<\epsilon\le 1/4\), \(a\in\{0,1\}\), and \(u,t>0\).
\item \textbf{(Population.)} Let \(P\in\mathcal M_{d,\epsilon}\), as defined in \cref{obj:def:model}. For \(j=1,\ldots,d\) and \(b\in\{0,1\}\), write
\[
p_j(P)=P(X=j),\qquad
s_{bj}(P)=P(X=j,A=b),\qquad
q_{bj}(P)=P(X=j,A=b,Y=1),
\]
and define the outcome-success conditional mean by
\[
\mu_{bj}(P)=
\begin{cases}
q_{bj}(P)/s_{bj}(P),&s_{bj}(P)>0,\\
0,&s_{bj}(P)=0.
\end{cases}
\]
\item \textbf{(Poisson counts.)} On a common probability space, let \(Z_j,K_j,W_j\) be measurable counts satisfying
\[
Z_j\sim\operatorname{Poi}\!\bigl(uq_{aj}(P)\bigr),\qquad
K_j\sim\operatorname{Poi}\!\bigl(ts_{aj}(P)\bigr),\qquad
W_j\sim\operatorname{Poi}\!\bigl(ts_{1-a,j}(P)\bigr),
\]
where \(\operatorname{Poi}(\lambda)\) denotes the Poisson law with mean \(\lambda\). The entire collection of \(3d\) counts is mutually independent.
\end{itemize}
Define the cell contributions and the covariate-weighted arm mean by
\[
H_j=\frac{Z_j}{u}\left(1+\frac{W_j}{K_j+1}\right),
\qquad
\psi_a=\sum_{j=1}^{d}p_j(P)\mu_{aj}(P).
\]
Then
\[
\left|\mathbb E\!\left[\sum_{j=1}^{d}H_j\right]-\psi_a\right|
\le
\min\left\{1,\frac{d}{\exp(1)t\epsilon}\right\},
\qquad
\operatorname{Var}\!\left(\sum_{j=1}^{d}H_j\right)
\le
C\left(\frac{1}{u\epsilon}+\frac{1}{t\epsilon}\right).
\]
\end{lemmav}

\begin{proof}[Proof of \cref{obj:lem:inverse-count-arm-risk}]
Take \(C=2^{20}>0\), and fix the parameters, law, and counts in the statement. Throughout, suppress the argument \(P\) in the population quantities. For every cell,
\[
p_j=s_{aj}+s_{1-a,j},\qquad
0\le q_{aj}\le s_{aj},\qquad
0\le\mu_{aj}\le1,\qquad
\sum_{j=1}^{d}p_j=1.
\]
Indeed, the arm probabilities partition the covariate-cell probability, and the outcome-success event is contained in the corresponding arm-cell event. Dividing \(0\le q_{aj}\le s_{aj}\) by \(s_{aj}>0\) gives the conditional-mean bounds; when \(s_{aj}=0\), both \(q_{aj}\) and \(\mu_{aj}\) are zero.

\begin{enumerate}
\item \emph{Exact expectation.}
Define the nonnegative Poisson means
\[
\lambda_j=t s_{aj},\qquad \rho_j=t s_{1-a,j}.
\]
The Poisson first moments give
\[
\mathbb E Z_j=u q_{aj},\qquad \mathbb E W_j=\rho_j.
\]
By \cref{obj:lem:poisson-inverse-moments},
\[
\mathbb E\!\left[\frac1{K_j+1}\right]
=
\begin{cases}
1,&\lambda_j=0,\\[2pt]
\dfrac{1-e^{-\lambda_j}}{\lambda_j},&\lambda_j>0.
\end{cases}
\]
Mutual independence permits factorization of the product expectation:
\[
\begin{aligned}
\mathbb E H_j
&=\frac1u\left(
\mathbb E Z_j+
\mathbb E Z_j\,
\mathbb E\!\left[\frac1{K_j+1}\right]\,
\mathbb E W_j\right)\\
&=q_{aj}\left(1+\rho_j
\begin{cases}
1,&\lambda_j=0,\\[2pt]
\dfrac{1-e^{-\lambda_j}}{\lambda_j},&\lambda_j>0.
\end{cases}\right).
\end{aligned}
\]
The count factors are integrable, the reciprocal factor is bounded by one, and independence makes their products integrable. Thus finite-sum linearity yields
\[
\mathbb E\!\left[\sum_{j=1}^{d}H_j\right]
=
\sum_{j=1}^{d}q_{aj}\left(1+\rho_j
\begin{cases}
1,&\lambda_j=0,\\[2pt]
\dfrac{1-e^{-\lambda_j}}{\lambda_j},&\lambda_j>0.
\end{cases}\right).
\]
% lean: CausalSmith.Stat.AnnotationRarearmFrontier.inverse_count_poisson_arm_mean

\item \emph{Bias bound.}
First suppose \(s_{aj}=0\). Then \(q_{aj}=\mu_{aj}=0\), so
\[
\mathbb E H_j
=p_j\mu_{aj}-s_{1-a,j}\mu_{aj}e^{-t s_{aj}}
=0.
\]
If \(s_{aj}>0\), then \(\lambda_j>0\), and the preceding expectation simplifies to
\[
\begin{aligned}
\mathbb E H_j
&=q_{aj}
+s_{1-a,j}\frac{q_{aj}}{s_{aj}}
  \bigl(1-e^{-t s_{aj}}\bigr)\\
&=(s_{aj}+s_{1-a,j})\mu_{aj}
-s_{1-a,j}\mu_{aj}e^{-t s_{aj}}\\
&=p_j\mu_{aj}-s_{1-a,j}\mu_{aj}e^{-t s_{aj}}.
\end{aligned}
\]
Consequently, in all cells,
\[
\mathbb E\!\left[\sum_{j=1}^{d}H_j\right]-\psi_a
=-\sum_{j=1}^{d}s_{1-a,j}\mu_{aj}e^{-t s_{aj}}.
\]
% lean: CausalSmith.Stat.AnnotationRarearmFrontier.inverse_count_cell_mean_identity

For \(p_j>0\), \cref{obj:ass:overlap} gives
\(s_{1j}\ge\epsilon p_j\) and \(s_{0j}\ge\epsilon p_j\):
multiply its lower propensity bound and its upper propensity bound by
\(p_j\), respectively. For \(p_j=0\), the same inequalities follow from
nonnegativity. Hence, for every \(j\),
\[
s_{aj}\ge\epsilon p_j,\qquad
0\le s_{1-a,j}\mu_{aj}\le s_{1-a,j}\le p_j.
\]
Since every summand in the bias expression is nonnegative,
\[
\left|\mathbb E\!\left[\sum_{j=1}^{d}H_j\right]-\psi_a\right|
=\sum_{j=1}^{d}s_{1-a,j}\mu_{aj}e^{-t s_{aj}}
\le\sum_{j=1}^{d}p_j e^{-t\epsilon p_j}.
\]
The last sum is at most one, because
\[
\sum_{j=1}^{d}p_j e^{-t\epsilon p_j}
\le\sum_{j=1}^{d}p_j=1.
\]
The elementary exponential bound gives
\[
(t\epsilon p_j)e^{-t\epsilon p_j}\le e^{-1},
\qquad
p_j e^{-t\epsilon p_j}\le\frac1{\exp(1)t\epsilon}.
\]
Summing this second bound proves
\[
\left|\mathbb E\!\left[\sum_{j=1}^{d}H_j\right]-\psi_a\right|
\le
\min\left\{1,\frac{d}{\exp(1)t\epsilon}\right\}.
\]
% lean: CausalSmith.Stat.AnnotationRarearmFrontier.inverse_count_mean_bias_bound

\item \emph{Weight moments.}
Use the reciprocal function and numerical constant from
\cref{obj:lem:poisson-inverse-moments}, and define the nonnegative real
random variables
\[
g(K_j)=\frac1{K_j+1},\qquad
C_0=2^{16},\qquad
\omega_j=W_jg(K_j),\qquad
\beta_j=1+\omega_j.
\]
The cited result supplies
\[
\mathbb E[g(K_j)^2]\le\frac{C_0}{(1+\lambda_j)^2},
\qquad
\operatorname{Var}(g(K_j))
\le\frac{C_0\lambda_j}{(1+\lambda_j)^4}.
\]
The usual Poisson moments are
\[
\mathbb E W_j=\rho_j,\qquad
\mathbb E W_j^2=\rho_j^2+\rho_j,\qquad
\operatorname{Var}(W_j)=\rho_j.
\]
Independence of \(W_j\) and \(K_j\) therefore gives
\[
\mathbb E\omega_j^2
=(\rho_j^2+\rho_j)\mathbb E[g(K_j)^2].
\]
In particular, \(\omega_j\) and \(\beta_j\) are square integrable. Since
\((1+\omega_j)^2\le2+2\omega_j^2\), by
\((\omega_j-1)^2\ge0\),
\[
\mathbb E\beta_j^2
\le2+2\mathbb E\omega_j^2
\le2+2C_0\frac{\rho_j^2+\rho_j}{(1+\lambda_j)^2}.
\]
Expanding variance and factoring the independent expectations yields
\[
\begin{aligned}
\operatorname{Var}(\beta_j)
&=\operatorname{Var}(\omega_j)\\
&=(\rho_j^2+\rho_j)\mathbb E[g(K_j)^2]
-\rho_j^2\bigl(\mathbb E g(K_j)\bigr)^2\\
&=\rho_j\mathbb E[g(K_j)^2]
+\rho_j^2\operatorname{Var}(g(K_j))\\
&\le C_0\left\{
\frac{\rho_j}{(1+\lambda_j)^2}
+\frac{\rho_j^2\lambda_j}{(1+\lambda_j)^4}
\right\}.
\end{aligned}
\]
These calculations also cover \(\lambda_j=0\) and \(\rho_j=0\).
% lean: CausalSmith.Stat.AnnotationRarearmFrontier.inverse_count_weight_moments

\item \emph{Cell variance.}
The Poisson moments of the outcome-success count give
\[
\mathbb E(Z_j/u)=q_{aj},\qquad
\mathbb E[(Z_j/u)^2]=q_{aj}^2+\frac{q_{aj}}u,
\qquad
\operatorname{Var}(Z_j/u)=\frac{q_{aj}}u.
\]
Moreover, \(Z_j/u\) is independent of \(\beta_j\), because \(Z_j\) is
independent of the pair \((W_j,K_j)\). Their product is square integrable:
its second moment factors into the two finite second moments. Thus
\[
\begin{aligned}
\operatorname{Var}(H_j)
&=\mathbb E[(Z_j/u)^2]\mathbb E\beta_j^2
-\bigl(\mathbb E(Z_j/u)\bigr)^2
 \bigl(\mathbb E\beta_j\bigr)^2\\
&=\frac{q_{aj}}u\,\mathbb E\beta_j^2
+q_{aj}^2\operatorname{Var}(\beta_j).
\end{aligned}
\]
Define the deterministic cell envelope by
\[
\mathcal V_{\mathrm{cell},j}
=
\frac{q_{aj}}u\left(
2+2C_0\frac{\rho_j^2+\rho_j}{(1+\lambda_j)^2}
\right)
+C_0q_{aj}^2\left(
\frac{\rho_j}{(1+\lambda_j)^2}
+\frac{\rho_j^2\lambda_j}{(1+\lambda_j)^4}
\right).
\]
Multiplying the weight bounds by the nonnegative coefficients
\(q_{aj}/u\) and \(q_{aj}^2\) gives
\[
\operatorname{Var}(H_j)\le\mathcal V_{\mathrm{cell},j}.
\]
% lean: CausalSmith.Stat.AnnotationRarearmFrontier.inverse_count_poisson_cell_variance

\item \emph{Variance of the sum.}
For \(j\ne k\), the count blocks use disjoint members of the mutually
independent collection, so
\[
(Z_j,K_j,W_j)\ \perp\!\!\!\perp\ (Z_k,K_k,W_k),
\qquad
H_j\ \perp\!\!\!\perp\ H_k.
\]
The second assertion follows by applying the measurable function defining
each contribution to its block. Pairwise independence and square
integrability now give
\[
\operatorname{Var}\!\left(\sum_{j=1}^{d}H_j\right)
=\sum_{j=1}^{d}\operatorname{Var}(H_j)
\le\sum_{j=1}^{d}\mathcal V_{\mathrm{cell},j}.
\]
% lean: CausalSmith.Stat.AnnotationRarearmFrontier.inverse_count_cell_blocks_independent

\item \emph{Deterministic cell envelope.}
If \(s_{aj}=0\), then \(q_{aj}=0\), and
\(\epsilon(s_{aj}+s_{1-a,j})\le s_{aj}\) forces
\(s_{1-a,j}=0\). Thus \(p_j=0\) and
\(\mathcal V_{\mathrm{cell},j}=0\).

Suppose henceforth that \(s_{aj}>0\). Overlap implies
\[
p_j\le\frac{s_{aj}}\epsilon\le\frac{p_j}\epsilon,
\qquad
s_{1-a,j}\le\frac{p_j}\epsilon,
\qquad
q_{aj}\le\frac{p_j}\epsilon.
\]
It also gives \(\epsilon s_{1-a,j}\le s_{aj}\), whence
\[
\epsilon s_{1-a,j}^2
\le s_{aj}s_{1-a,j}
\le s_{aj}p_j,
\qquad
\frac{s_{1-a,j}^2}{s_{aj}}\le\frac{p_j}\epsilon.
\]
% lean: CausalSmith.Stat.AnnotationRarearmFrontier.inverse_count_overlap_ratio_bound

The rational factors satisfy
\[
0\le\frac{\lambda_j}{1+\lambda_j}\le1,
\qquad
\lambda_j\le(1+\lambda_j)^2.
\]
Consequently,
\[
\begin{aligned}
\frac{s_{aj}\rho_j^2}{(1+\lambda_j)^2}
&=\frac{s_{1-a,j}^2}{s_{aj}}
  \left(\frac{\lambda_j}{1+\lambda_j}\right)^2
\le\frac{s_{1-a,j}^2}{s_{aj}},\\
\frac{s_{aj}\rho_j}{(1+\lambda_j)^2}
&=s_{1-a,j}\frac{\lambda_j}{(1+\lambda_j)^2}
\le s_{1-a,j},\\
\frac{s_{aj}^2\rho_j}{(1+\lambda_j)^2}
&=\frac{s_{1-a,j}}t
  \left(\frac{\lambda_j}{1+\lambda_j}\right)^2
\le\frac{s_{1-a,j}}t,\\
\frac{s_{aj}^2\rho_j^2\lambda_j}{(1+\lambda_j)^4}
&=\frac{s_{1-a,j}^2}{t s_{aj}}
  \left(\frac{\lambda_j}{1+\lambda_j}\right)^4
\le\frac{s_{1-a,j}^2}{t s_{aj}}.
\end{aligned}
\]
% lean: CausalSmith.Stat.AnnotationRarearmFrontier.inverse_count_rational_factors

Using \(q_{aj}\le s_{aj}\) in the first two bounds yields
\[
\begin{aligned}
q_{aj}\frac{\rho_j^2+\rho_j}{(1+\lambda_j)^2}
&\le
\frac{s_{aj}\rho_j^2}{(1+\lambda_j)^2}
+\frac{s_{aj}\rho_j}{(1+\lambda_j)^2}\\
&\le\frac{s_{1-a,j}^2}{s_{aj}}+s_{1-a,j}
\le\frac{2p_j}\epsilon.
\end{aligned}
\]
Similarly, \(q_{aj}^2\le s_{aj}^2\) and the last two bounds give
\[
\begin{aligned}
q_{aj}^2\left(
\frac{\rho_j}{(1+\lambda_j)^2}
+\frac{\rho_j^2\lambda_j}{(1+\lambda_j)^4}
\right)
&\le
\frac{s_{aj}^2\rho_j}{(1+\lambda_j)^2}
+\frac{s_{aj}^2\rho_j^2\lambda_j}{(1+\lambda_j)^4}\\
&\le\frac{s_{1-a,j}}t
+\frac{s_{1-a,j}^2}{t s_{aj}}
\le\frac{2p_j}{t\epsilon}.
\end{aligned}
\]
Substitution into the cell envelope proves
\[
\begin{aligned}
\mathcal V_{\mathrm{cell},j}
&\le
\frac2u\frac{p_j}\epsilon
+\frac{2C_0}u\frac{2p_j}\epsilon
+C_0\frac{2p_j}{t\epsilon}\\
&=(2+4C_0)\frac{p_j}{u\epsilon}
+2C_0\frac{p_j}{t\epsilon}\\
&\le
(2+4C_0)p_j
\left(\frac1{u\epsilon}+\frac1{t\epsilon}\right).
\end{aligned}
\]
The last inequality adds the nonnegative quantity
\((2+2C_0)p_j/(t\epsilon)\). Together with the zero-mass case, this
establishes the envelope for every cell.
% lean: CausalSmith.Stat.AnnotationRarearmFrontier.inverse_count_cell_variance_envelope

\item \emph{Summation and constant calibration.}
Summing the cell envelopes and using \(\sum_jp_j=1\) gives
\[
\begin{aligned}
\operatorname{Var}\!\left(\sum_{j=1}^{d}H_j\right)
&\le\sum_{j=1}^{d}\mathcal V_{\mathrm{cell},j}\\
&\le(2+4C_0)
\left(\frac1{u\epsilon}+\frac1{t\epsilon}\right).
\end{aligned}
\]
Finally,
\[
2+4C_0=262146\le1048576=2^{20}=C.
\]
Since the reciprocal-intensity sum is nonnegative, this proves
\[
\operatorname{Var}\!\left(\sum_{j=1}^{d}H_j\right)
\le C\left(\frac1{u\epsilon}+\frac1{t\epsilon}\right),
\]
as required.
% lean: CausalSmith.Stat.AnnotationRarearmFrontier.inverse_count_arm_variance_envelope
\end{enumerate}
\end{proof}

The bias bound captures the accumulation of smoothing error across cells. The variance bound separates the outcome-success intensity from the marginal-count intensity, with the overlap floor controlling both contributions. These are the armwise risk inputs for the inverse-count baseline in \cref{obj:def:baseline}; finite averaging supplies the corresponding rule on the observed records.

\subsection{Transfer to finite ordered pools}

The transfer argument applies to any bounded statistic of independent ordered prefixes. Its notation is local: the risk-bound symbol \(A\) below is distinct from the treatment coordinate.



A finite pool supports every prefix up to its observed length. Averaging over those prefixes, with zero output on overflow, yields the deterministic transfer rule in the next result.

\begin{lemmav}{lem:finite-prefix-transfer}[Finite prefix risk transfer]
Let \(I\) be a finite index set. Suppose the following conditions hold.
\begin{itemize}
\item \textbf{(Pools.)} For each \(i\in I\), let \(P_i\) be a probability law on an arbitrary measurable observation space, and let \(h_i\geq 1\) be an integer pool length. Observations within pool \(i\) are iid with law \(P_i\), and the pools are mutually independent.
\item \textbf{(Statistic.)} Let \(T\) be a measurable real-valued statistic on families of finite ordered prefixes, including empty prefixes, such that
\[
-1\leq T(s)\leq 1
\]
for every prefix family \(s\).
\item \textbf{(Target and risk.)} Let \(\vartheta\in[-1,1]\) be the target and \(A\in\mathbb R\) the risk bound. Let \(J_i\sim\operatorname{Poi}(h_i/8)\) be mutually independent prefix lengths, independent of all observations. For independent infinite iid observation sequences with respective laws \(P_i\), suppose
\[
\mathbb E\!\left[
\left(T\bigl((s_{i,1:J_i})_{i\in I}\bigr)-\vartheta\right)^2
\right]\leq A.
\]
\end{itemize}
For finite ordered pools \(s=(s_{i,1:h_i})_{i\in I}\), define the prefix-averaged statistic by
\[
T_{\mathrm{fin}}(s)
=
\sum_{\substack{(j_i)_{i\in I}\\0\leq j_i\leq h_i\ \forall i}}
\left(
\prod_{i\in I}
e^{-h_i/8}\frac{(h_i/8)^{j_i}}{j_i!}
\right)
T\bigl((s_{i,1:j_i})_{i\in I}\bigr).
\]
This is the conditional average over the prefix lengths with output zero whenever some \(J_i>h_i\). It defines a total deterministic measurable rule, takes values in \([-1,1]\) for every finite pool family, and satisfies
\[
\mathbb E\!\left[
\left(T_{\mathrm{fin}}\bigl((s_{i,1:h_i})_{i\in I}\bigr)
-\vartheta\right)^2
\right]
\leq A+\sum_{i\in I}e^{-h_i},
\]
where the expectation is under the independent finite iid pools. For observation spaces equipped with their Borel sigma-algebras, the rule is Borel measurable.
\end{lemmav}

\begin{proof}[Proof of \cref{obj:lem:finite-prefix-transfer}]
\begin{enumerate}
\item For a length vector \(\mathbf j=(j_i)_{i\in I}\in\mathbb N_0^I\), define
\begingroup\small
\[
\omega(\mathbf j)
=
\prod_{i\in I}
e^{-h_i/8}\frac{(h_i/8)^{j_i}}{j_i!},
\qquad
\mathsf{Good}
=
\{(s,\mathbf j):j_i\le h_i\text{ for every }i\in I\},
\qquad
\mathsf{Overflow}=\mathsf{Good}^{\,c}.
\]
\endgroup
Here \(s\) ranges over the finite ordered pool families in the statement. Define, for every such \(s\) and every \(\mathbf j\in\mathbb N_0^I\),
\[
T_{\mathrm{cap}}(s,\mathbf j)
=
\begin{cases}
T\bigl((s_{i,1:j_i})_{i\in I}\bigr),
   &(s,\mathbf j)\in\mathsf{Good},\\
0,&(s,\mathbf j)\in\mathsf{Overflow}.
\end{cases}
\]
On \(\mathsf{Good}\), the assumed range of \(T\) gives
\(\lvert T_{\mathrm{cap}}(s,\mathbf j)\rvert\le1\); on
\(\mathsf{Overflow}\), this follows from its zero value. Thus
\[
\lvert T_{\mathrm{cap}}(s,\mathbf j)\rvert\le1
\quad\text{for every }(s,\mathbf j).
\]
Writing \(\mathbb E_J\) for expectation over the independent prefix lengths, we have
\[
\begin{aligned}
\mathbb E_J T_{\mathrm{cap}}(s,J)
&=
\sum_{\mathbf j\in\mathbb N_0^I}
\omega(\mathbf j)T_{\mathrm{cap}}(s,\mathbf j)\\
&=
\sum_{\substack{\mathbf j\in\mathbb N_0^I\\j_i\le h_i\ \forall i}}
\omega(\mathbf j)
T\bigl((s_{i,1:j_i})_{i\in I}\bigr)
=
T_{\mathrm{fin}}(s).
\end{aligned}
\]
Every prefix projection in the last sum is measurable, and the sum is finite. Consequently \(T_{\mathrm{fin}}\) is a total deterministic measurable rule, with
\[
\lvert T_{\mathrm{fin}}(s)\rvert
\le
\mathbb E_J\lvert T_{\mathrm{cap}}(s,J)\rvert
\le1.
\]
For Borel observation spaces, these compositions and their finite sum are Borel measurable. % lean: Causalean.Stat.FiniteRaoBlackwell.IndependentPoissonPrefix.FiniteFamily.finite_prefix_transfer

\item Write \(\mathbb E_s\) for expectation under the independent finite iid pools and
\[
\mathbb E_{s,J}[\cdot]=\mathbb E_s\mathbb E_J[\cdot].
\]
The bounds on \(T_{\mathrm{cap}}\) and \(\vartheta\) imply
\[
\lvert T_{\mathrm{cap}}(s,J)-\vartheta\rvert\le2,
\qquad
0\le(T_{\mathrm{cap}}(s,J)-\vartheta)^2\le4,
\]
so all the following losses are integrable. For each fixed \(s\), Jensen's inequality for the square function gives
\[
\begin{aligned}
(T_{\mathrm{fin}}(s)-\vartheta)^2
&=
\left(\mathbb E_J[T_{\mathrm{cap}}(s,J)-\vartheta]\right)^2\\
&\le
\mathbb E_J[(T_{\mathrm{cap}}(s,J)-\vartheta)^2].
\end{aligned}
\]
Integrating over the finite pools yields
\[
\mathbb E_s[(T_{\mathrm{fin}}(s)-\vartheta)^2]
\le
\mathbb E_{s,J}[(T_{\mathrm{cap}}(s,J)-\vartheta)^2].
\] % lean: Causalean.Stat.FiniteRaoBlackwell.IndependentPoissonPrefix.FiniteFamily.sqRisk_prefixRaoBlackwellStatistic_le_capped

\item Let \(\boldsymbol s^{\mathrm{Poi}}\) denote the untruncated random prefix family from the hypothesis. Its conditional law is
\[
\mathcal L(\boldsymbol s^{\mathrm{Poi}}\mid J=\mathbf j)
=
\bigotimes_{i\in I}P_i^{\otimes j_i},
\qquad
\mathbf j\in\mathbb N_0^I,
\]
where
\[
P_i^{\otimes0}=\delta_{()},
\]
the point mass on the empty ordered prefix. For every admissible \(\mathbf j\), independence and iid sampling give the same conditional prefix law for the finite pools:
\[
\mathcal L\bigl((s_{i,1:j_i})_{i\in I}\mid J=\mathbf j\bigr)
=
\bigotimes_{i\in I}P_i^{\otimes j_i},
\qquad j_i\le h_i\ \forall i.
\]
Indeed, projection onto the first \(j_i\) coordinates sends
\(P_i^{\otimes h_i}\) to \(P_i^{\otimes j_i}\), and the coordinatewise laws multiply across the independent pools. Summing these equalities over the admissible lengths therefore gives
\[
\begin{aligned}
&\mathbb E_{s,J}\!\left[
(T_{\mathrm{cap}}(s,J)-\vartheta)^2
\mathbf1_{\mathsf{Good}}(s,J)\right]\\
&\quad=
\sum_{\substack{\mathbf j\in\mathbb N_0^I\\j_i\le h_i\ \forall i}}
\omega(\mathbf j)
\int
(T(s')-\vartheta)^2\,
d\!\left(\bigotimes_{i\in I}P_i^{\otimes j_i}\right)(s')\\
&\quad=
\mathbb E\!\left[
(T(\boldsymbol s^{\mathrm{Poi}})-\vartheta)^2
\mathbf1_{\{J_i\le h_i\ \forall i\}}\right].
\end{aligned}
\]
In the integral, \(s'\) ranges over prefix families with respective lengths \(j_i\). % lean: Causalean.Stat.FiniteRaoBlackwell.IndependentPoissonPrefix.FiniteFamily.sqRisk_cappedPrefixStatistic_restrict_nonoverflow_eq

Splitting the capped loss between the two events, and using its zero output on overflow, now yields
\[
\begin{aligned}
\mathbb E_{s,J}[(T_{\mathrm{cap}}(s,J)-\vartheta)^2]
&=
\mathbb E\!\left[
(T(\boldsymbol s^{\mathrm{Poi}})-\vartheta)^2
\mathbf1_{\{J_i\le h_i\ \forall i\}}\right]\\
&\qquad+
\vartheta^2\Pr(\exists i\in I:J_i>h_i)\\
&\le
\mathbb E[(T(\boldsymbol s^{\mathrm{Poi}})-\vartheta)^2]
+\Pr(\exists i\in I:J_i>h_i)\\
&\le
\mathbb E[(T(\boldsymbol s^{\mathrm{Poi}})-\vartheta)^2]
+\sum_{i\in I}\Pr(J_i>h_i).
\end{aligned}
\]
The first inequality discards a nonnegative complementary loss and uses
\(\vartheta^2\le1\). The last follows by writing the overflow event as
\(\bigcup_{i\in I}\{J_i>h_i\}\) and applying the finite union bound. % lean: Causalean.Stat.FiniteRaoBlackwell.IndependentPoissonPrefix.FiniteFamily.zero_overflow_prefix_risk

\item Fix \(i\in I\). The Poisson exponential moment follows directly from its probability masses and the exponential series:
\[
\begin{aligned}
\mathbb E[8^{J_i}]
&=
e^{-h_i/8}
\sum_{k=0}^{\infty}
\frac{(h_i/8)^k8^k}{k!}\\
&=
e^{-h_i/8}
\sum_{k=0}^{\infty}\frac{h_i^k}{k!}
=
e^{7h_i/8}.
\end{aligned}
\]
Since \(\{J_i>h_i\}\subseteq\{J_i\ge h_i\}\), exponential Markov gives
\[
\Pr(J_i>h_i)e^{-7h_i/8}
\le
e^{-h_i\log8}.
\]
The elementary numerical lower bound \(\log2>5/8\) implies
\[
\log8=3\log2\ge\frac{15}{8}.
\]
As \(h_i\ge0\), we obtain
\[
\Pr(J_i>h_i)e^{-7h_i/8}
\le
e^{-h_i\log8}
\le
e^{-15h_i/8}
=
e^{-h_i}e^{-7h_i/8}.
\]
Cancelling the strictly positive factor \(e^{-7h_i/8}\) proves
\[
\Pr(J_i>h_i)\le e^{-h_i}.
\] % lean: Causalean.Stat.FiniteRaoBlackwell.IndependentPoissonPrefix.FiniteFamily.eighth_pool_poisson_overflow

\item Combining the preceding bounds with the assumed Poisson-prefix risk gives
\[
\begin{aligned}
\mathbb E_s[(T_{\mathrm{fin}}(s)-\vartheta)^2]
&\le
\mathbb E[(T(\boldsymbol s^{\mathrm{Poi}})-\vartheta)^2]
+\sum_{i\in I}\Pr(J_i>h_i)\\
&\le
A+\sum_{i\in I}e^{-h_i}.
\end{aligned}
\]
These calculations include \(I=\varnothing\): empty products equal \(1\), empty sums equal \(0\), the pool and prefix families are the unique empty families, and the overflow event is empty. % lean: CausalSmith.Stat.AnnotationRarearmFrontier.finite_prefix_transfer
\end{enumerate}
\end{proof}

The transfer preserves the prefix risk bound up to an explicit exponentially small remainder determined by the finite pool lengths. It also provides measurability and boundedness on every observed pool family. For the baseline, \cref{obj:lem:inverse-count-arm-risk,obj:lem:finite-prefix-transfer} supply the count-risk and finite-experiment ingredients of \cref{obj:thm:uniform-baseline}.

\subsection{A localized polynomial certificate}

The attaining construction in \cref{obj:def:estimator-handle} uses a polynomial correction on the branch selected by a small pilot count. The scalar certificate below collects its approximation and moment controls. We introduce the certificate notation before stating those bounds.



The residual bound controls reciprocal approximation on the unit interval. The same certificate gives unbiased factorial estimation of the polynomial and a second-moment bound for every nonnegative argument.

\begin{lemmav}{lem:chebyshev-factorial-certificate}[Chebyshev factorial moment certificate]
Let \(L\) be an integer with \(L\ge2\). Let \(T_L\) be the Chebyshev polynomial specified by
\[
T_0(x)=1,\qquad T_1(x)=x,\qquad
T_{h+1}(x)=2xT_h(x)-T_{h-1}(x)\quad(h\ge1),
\]
and define
\[
C_L(z)=T_L(1-2z).
\]
Define the polynomials \(E_L\) and \(G_L\) by the identities
\[
E_L(z)=\frac{1-C_L(z)}{2L^2z},
\qquad
G_L(z)=\frac{1-E_L(z)}{z}
\quad(z\ne0),
\]
with their removable values at zero. Then
\[
0\le E_L(z)\le\min\{1,(L^2z)^{-1}\}
\quad(0<z\le1),
\qquad E_L(0)=1.
\]
Moreover,
\[
\deg G_L\le L-2,\qquad
G_L(z)=\sum_{h=0}^{L-2}g_hz^h,\qquad
\sum_{h=0}^{L-2}|g_h|\le14^L,
\]
where \(g_h\) is the coefficient of \(z^h\) in \(G_L\).

For real parameters \(t,B,z\), define
\[
D=tB,
\]
and suppose that
\begin{itemize}
\item \textbf{(Positive scales.)} \(t>0\) and \(B>0\).
\item \textbf{(Nonnegative argument.)} \(z\ge0\).
\item \textbf{(Degree calibration.)} \(D\ge L\).
\end{itemize}
Let \(K\sim\operatorname{Poi}(tzB)\), the Poisson law with mean \(tzB\), and define
\[
(K)_0=1,\qquad
(K)_h=\prod_{i=0}^{h-1}(K-i)\quad(h\ge1),
\qquad
\widehat G=\sum_{h=0}^{L-2}g_h\frac{(K)_h}{D^h}.
\]
Then
\[
\mathbb E[\widehat G]=G_L(z),
\qquad
\mathbb E[\widehat G^2]\le392^L(1+z)^{2L}.
\]
\end{lemmav}

\begin{proof}[Proof of \cref{obj:lem:chebyshev-factorial-certificate}]
\begin{enumerate}
\item The Chebyshev recurrence gives
\[
T_h(1)=1,\qquad \deg T_h\le h,\qquad
T_h(\cos\theta)=\cos(h\theta)
\quad(h\ge0,\ 0\le\theta\le\pi).
\]
Indeed, the initial values and the cosine addition formula establish the last identity by induction. Consequently,
\[
|T_L(x)|\le1\qquad(-1\le x\le1).
\]
Applying \cref{obj:lem:markov-polynomial} with polynomial \(T_L\), degree bound \(L\), and uniform bound \(1\), we obtain
\[
|T_L'(x)|\le L^2\qquad(-1\le x\le1).
\]
For \(0\le z\le1\), the points \(1-2z\) and \(1\) both belong to this interval. The derivative bound therefore gives
\[
|T_L(1-2z)-T_L(1)|\le L^2|(1-2z)-1|=2L^2z.
\]
Combining this with the range bound and \(T_L(1)=1\) yields
\[
0\le1-C_L(z),\qquad
1-C_L(z)\le2,\qquad
1-C_L(z)\le2L^2z.
\]

The numerator \(1-C_L\) has zero constant coefficient. Polynomial division by \(z\) thus gives the identity
\[
zE_L(z)=\frac{1-C_L(z)}{2L^2}.
\]
For \(z\ne0\), this implies
\[
E_L(z)=\frac{1-C_L(z)}{2L^2z}.
\]
When \(0<z\le1\), the denominator is positive. Dividing the three numerator bounds by it gives
\[
0\le E_L(z),\qquad
E_L(z)\le1,\qquad
E_L(z)\le\frac{1}{L^2z},
\]
and hence the required minimum bound.
% lean: CausalSmith.Stat.AnnotationRarearmFrontier.chebE_approximation_bounds

\item Differentiating the polynomial identity from the preceding step and evaluating at zero gives
\[
E_L(0)=-\frac{C_L'(0)}{2L^2}
      =\frac{T_L'(1)}{L^2}.
\]
The endpoint derivative identity is
\[
T_L'(1)=L^2.
\]
It follows directly by differentiating the Chebyshev recurrence: the initial derivatives are
\[
T_0'(1)=0,\qquad T_1'(1)=1,
\]
and, using \(T_h(1)=1\),
\[
T_{h+2}'(1)=2+2T_{h+1}'(1)-T_h'(1).
\]
Induction gives \(T_h'(1)=h^2\). Since \(L\ge2\), division by \(L^2\) is valid, and therefore
\[
E_L(0)=1.
\]
Thus \(1-E_L\) also has zero constant coefficient. Its polynomial division by \(z\) gives
\[
zG_L(z)=1-E_L(z).
\]
For every \(z\ne0\), the two polynomial identities consequently yield
\[
E_L(z)=\frac{1-C_L(z)}{2L^2z},
\qquad
G_L(z)=\frac{1-E_L(z)}{z}.
\]
These identities specify the polynomial continuations at zero as well.
% lean: CausalSmith.Stat.AnnotationRarearmFrontier.cheb_continuation_eval

\item Composition with the affine polynomial \(1-2z\) gives
\[
\deg C_L\le L,\qquad \deg(1-C_L)\le L.
\]
Scalar multiplication preserves this degree bound, and polynomial division by \(z\) lowers it by one. Applying the same argument to \(1-E_L\) gives
\[
\deg E_L\le L-1,\qquad
\deg(1-E_L)\le L-1,\qquad
\deg G_L\le L-2.
\]
The inequalities also hold for a zero polynomial, with the usual degree convention. Therefore the coefficient expansion is
\[
G_L(z)=\sum_{h=0}^{L-2}g_hz^h.
\]
% lean: CausalSmith.Stat.AnnotationRarearmFrontier.chebG_natDegree_le

\item For any real polynomial \(p_{\mathrm{aux}}\), use coefficient notation and define its coefficient norm by
\[
p_{\mathrm{aux}}(z)
 =\sum_{h\ge0}\bigl([z^h]p_{\mathrm{aux}}\bigr)z^h,
\qquad
\|p_{\mathrm{aux}}\|_{\mathrm{coef},1}
 :=\sum_{h\ge0}|[z^h]p_{\mathrm{aux}}|.
\]
Both sums have finite support. The coefficient triangle inequality and convolution formula give, for real polynomials \(p_{\mathrm{aux}},q_{\mathrm{aux}}\) and \(a_{\mathrm{sc}}\in\mathbb R\),
\[
\begin{aligned}
\|p_{\mathrm{aux}}-q_{\mathrm{aux}}\|_{\mathrm{coef},1}
 &\le \|p_{\mathrm{aux}}\|_{\mathrm{coef},1}
      +\|q_{\mathrm{aux}}\|_{\mathrm{coef},1},\\
\|p_{\mathrm{aux}}q_{\mathrm{aux}}\|_{\mathrm{coef},1}
 &\le \|p_{\mathrm{aux}}\|_{\mathrm{coef},1}
       \|q_{\mathrm{aux}}\|_{\mathrm{coef},1},\\
\|a_{\mathrm{sc}}p_{\mathrm{aux}}\|_{\mathrm{coef},1}
 &=|a_{\mathrm{sc}}|\|p_{\mathrm{aux}}\|_{\mathrm{coef},1}.
\end{aligned}
\]
Multiplication by \(z\) preserves this norm:
\[
\|zp_{\mathrm{aux}}\|_{\mathrm{coef},1}
 =\|p_{\mathrm{aux}}\|_{\mathrm{coef},1}.
\]
For the zero polynomial both sides are zero; otherwise this follows by shifting the coefficients, since
\[
[z^0](zp_{\mathrm{aux}})=0,\qquad
[z^{h+1}](zp_{\mathrm{aux}})=[z^h]p_{\mathrm{aux}}.
\]

Define the shifted family
\[
C_{\mathrm{sh},h}(z):=T_h(1-2z)
\quad(h\ge0,\ z\in\mathbb R),
\qquad C_{\mathrm{sh},L}=C_L.
\]
Its initial coefficient norms and recurrence are
\[
\|C_{\mathrm{sh},0}\|_{\mathrm{coef},1}=1,\qquad
\|C_{\mathrm{sh},1}\|_{\mathrm{coef},1}=3,
\]
\[
C_{\mathrm{sh},h+2}
 =2(1-2z)C_{\mathrm{sh},h+1}-C_{\mathrm{sh},h}.
\]
Since
\[
\|2(1-2z)\|_{\mathrm{coef},1}=|2|+|-4|=6,
\]
the norm inequalities imply
\[
\|C_{\mathrm{sh},h+2}\|_{\mathrm{coef},1}
 \le6\|C_{\mathrm{sh},h+1}\|_{\mathrm{coef},1}
    +\|C_{\mathrm{sh},h}\|_{\mathrm{coef},1}.
\]
Starting from \(1\le7^0\) and \(3\le7^1\), two-step induction gives
\[
\|C_{\mathrm{sh},h}\|_{\mathrm{coef},1}\le7^h,
\]
because
\[
6\cdot7^{h+1}+7^h
 =43\cdot7^h
 \le49\cdot7^h
 =7^{h+2}.
\]

Applying the norm identities to \(zE_L=(1-C_L)/(2L^2)\), we obtain
\[
\|E_L\|_{\mathrm{coef},1}
 =\frac{\|1-C_L\|_{\mathrm{coef},1}}{2L^2}
 \le\frac{1+7^L}{2L^2}
 \le7^L.
\]
The last inequality uses \(L^2\ge1\) and \(7^L\ge1\). Applying them next to \(zG_L=1-E_L\) yields
\[
\begin{aligned}
\sum_{h=0}^{L-2}|g_h|
 &=\|G_L\|_{\mathrm{coef},1}\\
 &=\|1-E_L\|_{\mathrm{coef},1}\\
 &\le1+\|E_L\|_{\mathrm{coef},1}\\
 &\le1+7^L
 \le2\cdot7^L
 \le2^L7^L
 =14^L.
\end{aligned}
\]
Here the first equality uses the degree bound, and \(2\le2^L\) follows from \(L\ge2\).
% lean: CausalSmith.Stat.AnnotationRarearmFrontier.chebG_coeff_sum_le

\item Fix \(t,B>0\), \(z\ge0\), and \(D=tB\ge L\), and define the Poisson mean
\[
\lambda_{\mathrm P}:=Dz=tzB\ge0.
\]
The law of \(K\) is
\[
\mathbb P(K=k)=e^{-\lambda_{\mathrm P}}
              \frac{\lambda_{\mathrm P}^{\,k}}{k!}
\qquad(k\in\mathbb N_0).
\]
We use \(0^0=1\) and the empty-product value \(1\).

For every integer \(h\ge0\), the terms with \(k<h\) have \((k)_h=0\), while \((k)_h=k!/(k-h)!\) for \(k\ge h\). Shifting the Poisson series therefore gives
\[
\begin{aligned}
\mathbb E[(K)_h]
 &=e^{-\lambda_{\mathrm P}}
   \sum_{k=h}^{\infty}
   \frac{\lambda_{\mathrm P}^{\,k}}{(k-h)!}\\
 &=\lambda_{\mathrm P}^{\,h}e^{-\lambda_{\mathrm P}}
   \sum_{k=0}^{\infty}\frac{\lambda_{\mathrm P}^{\,k}}{k!}
 =\lambda_{\mathrm P}^{\,h}.
\end{aligned}
\]
The exponential series converges, so these factorial moments are finite. This computation includes \(\lambda_{\mathrm P}=0\) and \(h=0\).

Since \(D>0\), each summand of the estimator satisfies
\[
\mathbb E\!\left[g_h\frac{(K)_h}{D^h}\right]
 =g_h\frac{(Dz)^h}{D^h}
 =g_hz^h.
\]
Integrating the finite sum term by term and using the polynomial expansion gives
\[
\mathbb E[\widehat G]
 =\sum_{h=0}^{L-2}g_hz^h
 =G_L(z).
\]
% lean: CausalSmith.Stat.AnnotationRarearmFrontier.chebG_factorial_unbiased

\item To obtain the mixed moments, define the falling-factorial polynomials
\[
F_{\mathrm{fac},h}(x):=\prod_{i=0}^{h-1}(x-i)
\quad(h\in\mathbb N_0,\ x\in\mathbb R),
\qquad F_{\mathrm{fac},0}(x)=1.
\]
For an integer \(k\ge0\), they satisfy
\[
F_{\mathrm{fac},h}(k)=(k)_h.
\]

Let \(h,h'\in\mathbb N_0\). First suppose \(h\le h'\). The binomial addition identity for falling-factorial polynomials, applied to \(x=(x-h)+h\), gives
\[
F_{\mathrm{fac},h'}(x)
 =\sum_{j_{\mathrm{ov}}=0}^{h'}
   \binom{h'}{j_{\mathrm{ov}}}
   F_{\mathrm{fac},j_{\mathrm{ov}}}(h)
   F_{\mathrm{fac},h'-j_{\mathrm{ov}}}(x-h).
\]
The product definition supplies
\[
F_{\mathrm{fac},h}(x)
F_{\mathrm{fac},h'-j_{\mathrm{ov}}}(x-h)
 =F_{\mathrm{fac},h+h'-j_{\mathrm{ov}}}(x).
\]
Moreover,
\[
F_{\mathrm{fac},j_{\mathrm{ov}}}(h)=0
 \quad(j_{\mathrm{ov}}>h),
\qquad
F_{\mathrm{fac},j_{\mathrm{ov}}}(h)
 =j_{\mathrm{ov}}!\binom{h}{j_{\mathrm{ov}}}
 \quad(0\le j_{\mathrm{ov}}\le h).
\]
Multiplying the addition identity by \(F_{\mathrm{fac},h}(x)\) and removing its zero terms thus gives
\[
F_{\mathrm{fac},h}(x)F_{\mathrm{fac},h'}(x)
 =\sum_{j_{\mathrm{ov}}=0}^{h}
   \binom{h}{j_{\mathrm{ov}}}
   \binom{h'}{j_{\mathrm{ov}}}
   j_{\mathrm{ov}}!\,
   F_{\mathrm{fac},h+h'-j_{\mathrm{ov}}}(x).
\]
If \(h'>h\) fails, exchange the two orders; the product and the coefficients are symmetric. Evaluating the resulting identity at any integer \(k\ge0\) yields
\[
(k)_h(k)_{h'}
 =\sum_{j_{\mathrm{ov}}=0}^{\min(h,h')}
   \binom{h}{j_{\mathrm{ov}}}
   \binom{h'}{j_{\mathrm{ov}}}
   j_{\mathrm{ov}}!\,(k)_{h+h'-j_{\mathrm{ov}}}.
\]
% lean: Causalean.Stat.Concentration.Poisson.descFactorial_mul

\item Define, for \(h,h'\in\mathbb N_0\),
\[
M_{\mathrm{fac}}(h,h')
 :=\mathbb E\!\left[
       \frac{(K)_h}{D^h}\frac{(K)_{h'}}{D^{h'}}
     \right].
\]
The product expansion and the finite factorial moments justify termwise expectation. They give
\[
\begin{aligned}
M_{\mathrm{fac}}(h,h')
 &=\sum_{j_{\mathrm{ov}}=0}^{\min(h,h')}
   \binom{h}{j_{\mathrm{ov}}}
   \binom{h'}{j_{\mathrm{ov}}}
   j_{\mathrm{ov}}!\,
   \frac{(Dz)^{h+h'-j_{\mathrm{ov}}}}{D^{h+h'}}\\
 &=\sum_{j_{\mathrm{ov}}=0}^{\min(h,h')}
   \binom{h}{j_{\mathrm{ov}}}
   \left(
     \frac{\binom{h'}{j_{\mathrm{ov}}}j_{\mathrm{ov}}!}
          {D^{j_{\mathrm{ov}}}}
   \right)
   z^{h+h'-j_{\mathrm{ov}}}.
\end{aligned}
\]
All terms are nonnegative, so \(M_{\mathrm{fac}}(h,h')\ge0\).

Now suppose \(h,h'\le L\). Since \(h'\le L\le D\), the overlap coefficient satisfies
\[
0\le
\frac{\binom{h'}{j_{\mathrm{ov}}}j_{\mathrm{ov}}!}
     {D^{j_{\mathrm{ov}}}}
 =\frac{(h')_{j_{\mathrm{ov}}}}{D^{j_{\mathrm{ov}}}}
 \le\frac{(h')^{j_{\mathrm{ov}}}}{D^{j_{\mathrm{ov}}}}
 \le1.
\]
The falling-factorial bound follows by bounding each factor by \(h'\); when the order exceeds \(h'\), the falling factorial is zero. The order-zero case has value \(1\).

Also, for every overlap index in the displayed sum,
\[
z^{h+h'-j_{\mathrm{ov}}}
 \le(1+z)^{h+h'-j_{\mathrm{ov}}}
 \le(1+z)^{2L},
\]
because \(z\ge0\), \(1+z\ge1\), and \(h+h'-j_{\mathrm{ov}}\le2L\). These inequalities include \(z=0\) under the stated power convention. Hence
\[
\begin{aligned}
M_{\mathrm{fac}}(h,h')
 &\le\sum_{j_{\mathrm{ov}}=0}^{\min(h,h')}
       \binom{h}{j_{\mathrm{ov}}}(1+z)^{2L}\\
 &\le\sum_{j_{\mathrm{ov}}=0}^{h}
       \binom{h}{j_{\mathrm{ov}}}(1+z)^{2L}\\
 &=2^h(1+z)^{2L}\\
 &\le2^L(1+z)^{2L}.
\end{aligned}
\]
The second inequality adds nonnegative terms, and the equality is the binomial theorem.
% lean: CausalSmith.Stat.AnnotationRarearmFrontier.normalized_factorial_mixed_bound

\item Expanding the square and integrating the finite sum gives
\[
\mathbb E[\widehat G^2]
 =\sum_{h=0}^{L-2}\sum_{h'=0}^{L-2}
   g_hg_{h'}M_{\mathrm{fac}}(h,h').
\]
Every product is integrable by its finite factorial expansion. For each pair of indices,
\[
g_hg_{h'}\le|g_h||g_{h'}|,
\qquad
0\le M_{\mathrm{fac}}(h,h')
 \le2^L(1+z)^{2L},
\]
where the moment bound applies because \(h,h'\le L-2\le L\). Consequently,
\[
\begin{aligned}
\mathbb E[\widehat G^2]
 &\le\sum_{h=0}^{L-2}\sum_{h'=0}^{L-2}
       |g_h||g_{h'}|\,2^L(1+z)^{2L}\\
 &=\left(\sum_{h=0}^{L-2}|g_h|\right)^2
      2^L(1+z)^{2L}\\
 &\le(14^L)^2\,2^L(1+z)^{2L}\\
 &=392^L(1+z)^{2L}.
\end{aligned}
\]
The coefficient sum is nonnegative, so its bound may be squared. All index ranges remain valid at \(L=2\), when the estimator sum consists of its order-zero term.
% lean: CausalSmith.Stat.AnnotationRarearmFrontier.chebG_factorial_second_moment
\end{enumerate}
\end{proof}

The certificate links polynomial approximation to a statistic built from observed counts. Its second-moment bound depends explicitly on both the degree and the argument, making localization part of the risk calculation. The independent pilot count supplies the branch selection, while the inverse-count moment bounds control the contribution on the complementary branch. The following result assembles these ingredients at the calibration used by the attaining estimator.

\begin{lemmav}{lem:uniform-hybrid-arm-risk}[Uniform hybrid arm risk]
There exists a numerical constant \(C>0\) such that the following holds for every \(n,m,d\in\mathbb N\), overlap parameter \(\epsilon\in\mathbb R\), observable law \(P\), arm \(a\in\{0,1\}\), and pool intensities \(u,t_p,t\in\mathbb R\).
\begin{itemize}
\item \textbf{(Experiment.)} The indices satisfy
\[
n\ge1,\qquad d\ge2,\qquad 0<\epsilon\le\frac14,\qquad
S=n\epsilon\ge\exp(4096),\qquad N=n+m.
\]
\item \textbf{(Model.)} The law satisfies \(P\in\mathcal M_{d,\epsilon}\), as defined in \cref{obj:def:model}.
\item \textbf{(Intensities.)} The pool intensities satisfy
\[
u\ge\frac n{32},\qquad t_p\ge\frac N{64},\qquad
t\ge\frac N{64},\qquad \frac13\le\frac t{t_p}\le3.
\]
\item \textbf{(Calibration.)} Set
\[
L=\left\lfloor\frac{\log S}{1024}\right\rfloor,\qquad
H_0=2^{20},\qquad B=\frac{H_0L}{\min\{t_p,t\}},\qquad
k_0=\left\lfloor\frac{t_pB}{4}\right\rfloor.
\]
\item \textbf{(Poisson pools.)} On any probability space, let
\(\{Z_{bj},J_{bj},K_{bj}:b\in\{0,1\},\,1\le j\le d\}\)
be jointly independent measurable counts with
\[
Z_{bj}\sim\operatorname{Poi}\!\left(uq_{bj}(P)\right),\qquad
J_{bj}\sim\operatorname{Poi}\!\left(t_ps_{bj}(P)\right),\qquad
K_{bj}\sim\operatorname{Poi}\!\left(ts_{bj}(P)\right).
\]
Here \(\operatorname{Poi}(\lambda)\) is the Poisson law with mean \(\lambda\), and the population quantities are
\[
q_{bj}(P)=P(X=j,A=b,Y=1),\qquad
s_{bj}(P)=\sum_{y\in\{0,1\}}P(X=j,A=b,Y=y),
\]
\[
p_j(P)=\sum_{b\in\{0,1\}}s_{bj}(P),\qquad
\mu_{bj}(P)=
\begin{cases}
q_{bj}(P)/s_{bj}(P),&s_{bj}(P)>0,\\
0,&s_{bj}(P)=0.
\end{cases}
\]
\end{itemize}
Using the monomial coefficients \(g_h\) of the reciprocal-approximation polynomial \(G_L\) in the estimator construction, form
\[
\widehat G_{aj}
=\sum_{h=0}^{L-2}g_h\frac{(K_{aj})_h}{(tB)^h},
\qquad
(K)_0=1,\qquad
(K)_h=\prod_{r=0}^{h-1}(K-r)\quad(h\ge1),
\]
and
\[
V_{aj}=
\begin{cases}
\displaystyle
\frac{Z_{aj}}u
\left(1+\frac{K_{1-a,j}\widehat G_{aj}}{tB}\right),
&J_{aj}\le k_0,\\[6pt]
\displaystyle
\frac{Z_{aj}}u
\left(1+\frac{K_{1-a,j}}{K_{aj}+1}\right),
&J_{aj}>k_0.
\end{cases}
\]
Then
\[
\mathbb E\!\left[
\left(\sum_{j=1}^dV_{aj}
-\sum_{j=1}^dp_j(P)\mu_{aj}(P)\right)^2
\right]
\le
C\left\{\frac1S+
\left(\frac d{N\epsilon L}\right)^2\right\}.
\]
The same constant \(C\) applies to every choice of indices, law, arm, intensities, and probability space satisfying these conditions.
\end{lemmav}

\begin{proof}[Proof of \cref{obj:lem:uniform-hybrid-arm-risk}]
Fix parameters and counts satisfying the hypotheses. Write \(\mathbb P\) for their probability measure on the given sample space \(\Omega\).

\begin{enumerate}
\item \textit{Identify the joint count law.}
Define the count space and the selected count vector by
\[
\mathbb N_0=\{0,1,2,\ldots\},\qquad
\mathcal E=(\mathbb N_0^4)^d,\qquad
\mathscr T(\omega)
=
\bigl(J_{aj}(\omega),Z_{aj}(\omega),
K_{aj}(\omega),K_{1-a,j}(\omega)\bigr)_{j=1}^d.
\]
For \(j=1,\ldots,d\), define probability measures
\[
\begin{aligned}
\nu_j={}&
\operatorname{Poi}\!\bigl(t_ps_{aj}(P)\bigr)
\otimes\operatorname{Poi}\!\bigl(uq_{aj}(P)\bigr)\\
&\otimes\operatorname{Poi}\!\bigl(ts_{aj}(P)\bigr)
\otimes\operatorname{Poi}\!\bigl(ts_{1-a,j}(P)\bigr),
\qquad
\nu=\bigotimes_{j=1}^d\nu_j,
\end{aligned}
\]
with coordinates ordered as in \(\mathscr T\).

The selected counts form a subfamily of the stipulated independent family: their indices are distinct, including the two factorial-pool indices because \(a\ne1-a\). Thus, for every
\[
\boldsymbol c=(c_j)_{j=1}^d\in\mathcal E,
\qquad
c_j=(c_j^{\mathrm p},c_j^{\mathrm o},
c_j^{\mathrm f},c_j^{\mathrm{opp}}),
\]
independence gives
\[
\begin{aligned}
\mathbb P\{\mathscr T=\boldsymbol c\}
=\prod_{j=1}^d\bigl[
&\mathbb P\{J_{aj}=c_j^{\mathrm p}\}
\mathbb P\{Z_{aj}=c_j^{\mathrm o}\}\\
&\mathbb P\{K_{aj}=c_j^{\mathrm f}\}
\mathbb P\{K_{1-a,j}=c_j^{\mathrm{opp}}\}
\bigr]
=\nu\{\boldsymbol c\}.
\end{aligned}
\]
Indeed, the event on the left is the intersection of these singleton-coordinate events. Equality on singletons determines measures on the countable space \(\mathcal E\), so
\[
\mathbb P\circ\mathscr T^{-1}=\nu.
\]
% lean: CausalSmith.Stat.AnnotationRarearmFrontier.hybrid_four_count_product_law

\item \textit{Establish the population inputs.}
For the fixed law and arm, set
\begingroup\footnotesize
\[
s_j=s_{aj}(P),\qquad
v_j=s_{1-a,j}(P),\qquad
q_j=q_{aj}(P),\qquad
p_j=p_j(P),\qquad
\bar\mu_j=\mu_{aj}(P),
\qquad
\psi_a=\sum_{j=1}^dp_j\bar\mu_j.
\]
\endgroup
The probability definitions imply
\[
s_j,v_j\ge0,\qquad
p_j=s_j+v_j,\qquad
0\le q_j\le s_j,\qquad
0\le\bar\mu_j\le1,\qquad
\sum_{j=1}^dp_j=1.
\]
If \(s_j=0\), the bound \(0\le q_j\le s_j\) gives \(q_j=0\); if \(s_j>0\), the defining quotient for \(\bar\mu_j\) gives the same identity
\[
s_j\bar\mu_j=q_j
\qquad(j=1,\ldots,d).
\]
By \cref{obj:def:model}, the law satisfies \cref{obj:ass:overlap}. On occupied cells, each arm proportion lies in \([\epsilon,1-\epsilon]\), hence
\[
\epsilon p_j\le s_j.
\]
On a null covariate cell, both arm masses vanish and this inequality also holds. In particular,
\[
s_j=0\quad\Longrightarrow\quad
p_j=v_j=q_j=\bar\mu_j=0.
\]
The corresponding outcome count has mean zero, so both branch contributions vanish almost surely.

The public parameters satisfy
\begingroup\small
\[
S=n\epsilon,\qquad N=n+m,\qquad
0<\epsilon\le1,\qquad
S\le N\epsilon,\qquad
\frac S{32}\le u\epsilon,\qquad
\frac N{64}\le t_p,t,\qquad
\frac13\le\frac t{t_p}.
\]
\endgroup
Consequently \(S,N,u,t_p,t\) are strictly positive.

On \(\mathcal E\), define the coordinate counts by
\[
\bigl(\mathsf J_j(\boldsymbol c),\mathsf Z_j(\boldsymbol c),
\mathsf K_j(\boldsymbol c),\mathsf W_j(\boldsymbol c)\bigr)=c_j.
\]
Define real-valued functions on this entire space by
\[
\begin{aligned}
\mathcal G_j
&=\sum_{h=0}^{L-2}g_h\frac{(\mathsf K_j)_h}{(tB)^h},\\
\mathcal P_j
&=\frac{\mathsf Z_j}{u}
\left(1+\frac{\mathsf W_j\mathcal G_j}{tB}\right),\\
\mathcal H_j
&=\frac{\mathsf Z_j}{u}
\left(1+\frac{\mathsf W_j}{\mathsf K_j+1}\right),\\
\mathcal I_j&=\mathbf1\{\mathsf J_j\le k_0\},\qquad
\mathcal V_j=\mathcal I_j\mathcal P_j+(1-\mathcal I_j)\mathcal H_j.
\end{aligned}
\]
Under \(\nu\), the four coordinate counts in cell \(j\) have independent Poisson laws with means
\[
t_ps_j,\qquad uq_j=us_j\bar\mu_j,\qquad ts_j,\qquad tv_j,
\]
respectively. We write \(\mathbb E_\nu\) and \(\operatorname{Var}_\nu\) for expectation and variance under \(\nu\).
% lean: CausalSmith.Stat.AnnotationRarearmFrontier.uniform_hybrid_arm_risk

\item \textit{Calibrate the degree and localization scale.}
Introduce
\[
\eta_S=\log S,\qquad
F_{\mathrm{grow}}=(2^{24})^L,\qquad
\rho_{\mathrm{pil}}=S^{-20}.
\]
Since \(\eta_S\ge4096\), the definition of \(L\) gives
\[
L\ge4,\qquad
L\le\frac{\eta_S}{1024}<L+1,\qquad
\eta_S<1280L.
\]
Using \(\log2\le1\) and \(\log L\le L\),
\[
\log(F_{\mathrm{grow}}L^3)
=24L\log2+3\log L
\le27L
\le\frac{27}{1024}\eta_S
\le\frac{\eta_S}{2}.
\]
Therefore
\[
F_{\mathrm{grow}}L^3\le\sqrt S,
\qquad
e^{-200000L}\le e^{-20\eta_S}=\rho_{\mathrm{pil}}.
\]
Moreover, \(B>0\), and the minimum in its definition yields
\[
t_pB
=H_0L\frac{t_p}{\min\{t_p,t\}}\ge H_0L,
\qquad
tB
=H_0L\frac{t}{\min\{t_p,t\}}\ge H_0L\ge L.
\]

Define the light and heavy index sets, their light mass, and their light count by
\[
\mathcal J_{\mathrm{light}}=\{j:1\le j\le d,\ s_j\le B\},
\qquad
\mathcal J_{\mathrm{heavy}}=\{j:1\le j\le d,\ s_j>B\},
\]
\[
M_{\mathrm{light}}=\sum_{j\in\mathcal J_{\mathrm{light}}}p_j,
\qquad
c_{\mathrm{light}}=|\mathcal J_{\mathrm{light}}|.
\]
For every light cell, overlap gives
\[
p_j\le\frac{s_j}{\epsilon}\le\frac B\epsilon.
\]
Summing this inequality and using total mass one proves
\[
0\le M_{\mathrm{light}}
\le\min\left\{1,\frac{c_{\mathrm{light}}B}{\epsilon}\right\}
\le\frac{dB}{\epsilon},
\qquad
0\le c_{\mathrm{light}}\le d.
\]
These formulas also apply when either index set is empty, with its sum equal to zero.
% lean: CausalSmith.Stat.AnnotationRarearmFrontier.hybrid_degree_calibration
% lean: CausalSmith.Stat.AnnotationRarearmFrontier.light_cell_mass_bound

\item \textit{Bound the inverse-branch variance.}
Use the reciprocal function and numerical constant supplied by
\cref{obj:lem:poisson-inverse-moments}:
\[
g(k)=\frac1{k+1}\quad(k\in\mathbb N_0),
\qquad
C_0=2^{16},
\qquad
C_{\mathrm{inv}}=2+4C_0=262146.
\]
For each cell define the inverse multiplier
\[
\mathcal W_{\mathrm{inv},j}
=1+\mathsf W_jg(\mathsf K_j).
\]
Independence, the Poisson first and second moments, and
\cref{obj:lem:poisson-inverse-moments} give
\[
\begin{aligned}
\mathbb E_\nu[\mathcal W_{\mathrm{inv},j}^2]
&\le
2+2\bigl((tv_j)^2+tv_j\bigr)
\mathbb E_\nu[g(\mathsf K_j)^2]\\
&\le
2+2C_0\frac{(tv_j)^2+tv_j}{(1+ts_j)^2},
\end{aligned}
\]
where the first inequality uses \((1+x)^2\le2+2x^2\). The variance of the independent product is
\[
\begin{aligned}
\operatorname{Var}_\nu(\mathcal W_{\mathrm{inv},j})
&=tv_j\,\mathbb E_\nu[g(\mathsf K_j)^2]
+(tv_j)^2\operatorname{Var}_\nu(g(\mathsf K_j))\\
&\le C_0\left\{
\frac{tv_j}{(1+ts_j)^2}
+\frac{(tv_j)^2ts_j}{(1+ts_j)^4}
\right\}.
\end{aligned}
\]
Also,
\[
\mathbb E_\nu\!\left[\frac{\mathsf Z_j}{u}\right]=q_j,
\qquad
\mathbb E_\nu\!\left[\left(\frac{\mathsf Z_j}{u}\right)^2\right]
=\frac{q_j}{u}+q_j^2.
\]
Factoring its independent first and second moments therefore gives
\[
\operatorname{Var}_\nu(\mathcal H_j)
=
\frac{q_j}{u}\mathbb E_\nu[\mathcal W_{\mathrm{inv},j}^2]
+q_j^2\operatorname{Var}_\nu(\mathcal W_{\mathrm{inv},j}).
\]

For \(s_j>0\), overlap implies
\[
\frac{v_j^2}{s_j}\le\frac{p_j}{\epsilon},
\qquad
q_j,v_j\le\frac{p_j}{\epsilon}.
\]
Indeed, \(\epsilon v_j\le s_j\), and multiplying by \(v_j\le p_j\) proves the first bound. The elementary inequalities
\[
0\le\frac{ts_j}{1+ts_j}\le1,
\qquad
ts_j\le(1+ts_j)^2
\]
give the four rational bounds
\[
\begin{aligned}
\frac{s_j(tv_j)^2}{(1+ts_j)^2}
&\le\frac{v_j^2}{s_j},&
\frac{s_jtv_j}{(1+ts_j)^2}
&\le v_j,\\
\frac{s_j^2tv_j}{(1+ts_j)^2}
&\le\frac{v_j}{t},&
\frac{s_j^2(tv_j)^2ts_j}{(1+ts_j)^4}
&\le\frac{v_j^2}{ts_j}.
\end{aligned}
\]
For example, the first, third, and fourth left-hand sides equal their respective right-hand sides multiplied by
\((ts_j/(1+ts_j))^2\), \((ts_j/(1+ts_j))^2\), and
\((ts_j/(1+ts_j))^4\). Hence \(q_j\le s_j\) gives
\[
q_j\frac{(tv_j)^2+tv_j}{(1+ts_j)^2}
\le\frac{v_j^2}{s_j}+v_j
\le\frac{2p_j}{\epsilon},
\]
\[
q_j^2\left\{
\frac{tv_j}{(1+ts_j)^2}
+\frac{(tv_j)^2ts_j}{(1+ts_j)^4}
\right\}
\le\frac{v_j}{t}+\frac{v_j^2}{ts_j}
\le\frac{2p_j}{t\epsilon}.
\]
Substitution into the variance identity proves
\[
\operatorname{Var}_\nu(\mathcal H_j)
\le C_{\mathrm{inv}}p_j
\left(\frac1{u\epsilon}+\frac1{t\epsilon}\right).
\]
For \(s_j=0\), both sides vanish by the null-cell identities already established. Summing over all cells gives
\[
\sum_{j=1}^d\operatorname{Var}_\nu(\mathcal H_j)
\le C_{\mathrm{inv}}
\left(\frac1{u\epsilon}+\frac1{t\epsilon}\right).
\]
% lean: CausalSmith.Stat.AnnotationRarearmFrontier.hybrid_canonical_inverse_variance
% lean: CausalSmith.Stat.AnnotationRarearmFrontier.hybrid_canonical_inverse_variance_sum

\item \textit{Bound the light polynomial variances.}
Define, for every cell,
\[
\zeta_j=\frac{s_j}{B},
\qquad
\mathcal W_{\mathrm{pol},j}
=1+\frac{\mathsf W_j\mathcal G_j}{tB}.
\]
The hypotheses of \cref{obj:lem:chebyshev-factorial-certificate} hold because \(L\ge4\), \(t,B>0\), \(\zeta_j\ge0\), and \(tB\ge L\). For a light cell it supplies
\[
\mathbb E_\nu[\mathcal G_j^2]
\le392^L(1+\zeta_j)^{2L}
\le1568^L
\le F_{\mathrm{grow}}.
\]
Independence of the two factorial counts and
\(\mathbb E_\nu[\mathsf W_j^2]=(tv_j)^2+tv_j\) then imply
\[
\begin{aligned}
\mathbb E_\nu[\mathcal W_{\mathrm{pol},j}^2]
&\le
2+2\mathbb E_\nu[\mathcal G_j^2]
\left(\frac{v_j^2}{B^2}+\frac{v_j}{tB^2}\right)\\
&\le
2F_{\mathrm{grow}}
\left(1+\frac{v_j^2}{B^2}+\frac{v_j}{tB^2}\right).
\end{aligned}
\]
The outcome count is independent of this multiplier, so
\[
\mathbb E_\nu[\mathcal P_j^2]
=
\left(\frac{q_j}{u}+q_j^2\right)
\mathbb E_\nu[\mathcal W_{\mathrm{pol},j}^2].
\]

For \(s_j\le B\), the following bounds retain each mixed monomial:
\[
\begin{aligned}
\frac{s_jv_j^2}{B^2}&\le\frac{p_j}{\epsilon},&
\frac{s_jv_j}{tB^2}&\le v_j,&
\frac{s_j^2v_j}{tB^2}&\le\frac{p_j}{t},\\
\frac{s_j^2v_j^2}{B^2}&\le p_j^2,&
s_j^2&\le p_j^2,&
p_j&\le\frac B\epsilon.
\end{aligned}
\]
When \(s_j>0\), the first follows from
\(s_jv_j^2/B^2\le v_j^2/s_j\le p_j/\epsilon\).
The second follows from \(s_j\le B\le tB^2\), since \(tB\ge1\).
The next two use \(s_j^2\le B^2\) and \(v_j\le p_j\).
When \(s_j=0\), all these inequalities follow from \(v_j=p_j=0\).

Expanding the product and using \(q_j\le s_j\) gives
\[
\begin{aligned}
&\left(\frac{q_j}{u}+q_j^2\right)
\left(1+\frac{v_j^2}{B^2}+\frac{v_j}{tB^2}\right)\\
&\quad\le
\frac{s_j}{u}
+\frac{s_jv_j^2}{uB^2}
+\frac{s_jv_j}{utB^2}
+s_j^2+\frac{s_j^2v_j^2}{B^2}
+\frac{s_j^2v_j}{tB^2}\\
&\quad\le
\frac{3p_j}{u\epsilon}+\frac{p_j}{t}+2p_j^2
\le
\frac{3p_j}{u\epsilon}+\frac{p_j}{t}
+\frac{2B^2}{\epsilon^2}.
\end{aligned}
\]
Here \(s_j,v_j\le p_j/\epsilon\) uses \(\epsilon\le1\).
Thus
\[
\operatorname{Var}_\nu(\mathcal P_j)
\le\mathbb E_\nu[\mathcal P_j^2]
\le
6F_{\mathrm{grow}}
\left(\frac{p_j}{u\epsilon}+\frac{p_j}{t}
+\frac{B^2}{\epsilon^2}\right).
\]
Summation yields
\[
\sum_{j\in\mathcal J_{\mathrm{light}}}
\operatorname{Var}_\nu(\mathcal P_j)
\le
6F_{\mathrm{grow}}
\left(
\frac{M_{\mathrm{light}}}{u\epsilon}
+\frac{M_{\mathrm{light}}}{t}
+\frac{c_{\mathrm{light}}B^2}{\epsilon^2}
\right).
\]
% lean: CausalSmith.Stat.AnnotationRarearmFrontier.light_polynomial_variance_sum

\item \textit{Assemble the light mixture variances.}
Define
\[
\pi_j=\nu\{\mathsf J_j\le k_0\}\in[0,1],
\qquad
\delta_j=\mathbb E_\nu[\mathcal P_j]
-\mathbb E_\nu[\mathcal H_j].
\]
The factorial unbiasedness in
\cref{obj:lem:chebyshev-factorial-certificate} gives
\[
\mathbb E_\nu[\mathcal G_j]=G_L(\zeta_j).
\]
For \(s_j>0\), independence and the reciprocal expectation from
\cref{obj:lem:poisson-inverse-moments} therefore give
\[
\begin{aligned}
\mathbb E_\nu[\mathcal P_j]
&=s_j\bar\mu_j
\left(1+\frac{v_j}{B}G_L(\zeta_j)\right)
=p_j\bar\mu_j-v_j\bar\mu_jE_L(\zeta_j),\\
\mathbb E_\nu[\mathcal H_j]
&=s_j\bar\mu_j
\left(1+tv_j\frac{1-e^{-ts_j}}{ts_j}\right)
=p_j\bar\mu_j-v_j\bar\mu_je^{-ts_j}.
\end{aligned}
\]
The polynomial identity used here is
\[
\zeta_jG_L(\zeta_j)=1-E_L(\zeta_j),
\]
as specified in \cref{obj:def:estimator-handle}. The two branch-mean formulas also hold when \(s_j=0\), because \(p_j=v_j=0\).

The pilot is independent of both branches. Consequently,
\[
\begin{aligned}
\mathbb E_\nu[\mathcal V_j]
&=\pi_j\mathbb E_\nu[\mathcal P_j]
+(1-\pi_j)\mathbb E_\nu[\mathcal H_j],\\
\mathbb E_\nu[\mathcal V_j^2]
&=\pi_j\mathbb E_\nu[\mathcal P_j^2]
+(1-\pi_j)\mathbb E_\nu[\mathcal H_j^2].
\end{aligned}
\]
Subtracting the square of the first identity from the second gives
\[
\operatorname{Var}_\nu(\mathcal V_j)
=
\pi_j\operatorname{Var}_\nu(\mathcal P_j)
+(1-\pi_j)\operatorname{Var}_\nu(\mathcal H_j)
+\pi_j(1-\pi_j)\delta_j^2.
\]

For \(0<s_j\le B\), the residual bound in
\cref{obj:lem:chebyshev-factorial-certificate} yields
\[
0\le v_j\bar\mu_jE_L(\zeta_j)
\le\frac{v_jB}{L^2s_j}
\le\frac B{\epsilon L^2}.
\]
The last inequality follows from \(\epsilon v_j\le s_j\). Using the branch means and \((x-y)^2\le2x^2+2y^2\), we obtain
\[
\pi_j(1-\pi_j)\delta_j^2
\le
2\left(\frac B{\epsilon L^2}\right)^2
+2\bigl(p_je^{-ts_j}\bigr)^2.
\]
For \(s_j=0\), the branch means vanish and the same bound applies.

Overlap and the elementary bound \(xe^{-x}\le e^{-1}\) for \(x\ge0\) give
\[
\bigl(p_je^{-ts_j}\bigr)^2
\le p_j^2e^{-2t\epsilon p_j}
\le\frac{p_j}{2e\,t\epsilon}.
\]
Thus
\[
\sum_{j\in\mathcal J_{\mathrm{light}}}
\pi_j(1-\pi_j)\delta_j^2
\le
2c_{\mathrm{light}}
\left(\frac B{\epsilon L^2}\right)^2
+\frac1{e\,t\epsilon}.
\]
Combining this with the polynomial variance bound and
\(0\le\pi_j\le1\) proves
\[
\begin{aligned}
\sum_{j\in\mathcal J_{\mathrm{light}}}
\operatorname{Var}_\nu(\mathcal V_j)
\le{}&
\sum_{j\in\mathcal J_{\mathrm{light}}}
\operatorname{Var}_\nu(\mathcal H_j)\\
&+6F_{\mathrm{grow}}
\left(
\frac{M_{\mathrm{light}}}{u\epsilon}
+\frac{M_{\mathrm{light}}}{t}
+\frac{c_{\mathrm{light}}B^2}{\epsilon^2}
\right)\\
&+2c_{\mathrm{light}}
\left(\frac B{\epsilon L^2}\right)^2
+\frac1{e\,t\epsilon}.
\end{aligned}
\]
% lean: CausalSmith.Stat.AnnotationRarearmFrontier.hybrid_branch_means
% lean: CausalSmith.Stat.AnnotationRarearmFrontier.light_hybrid_variance_sum

\item \textit{Control the heavy mixture variances.}
For a heavy cell define its pilot mean by
\[
\lambda_{\mathrm p,j}=t_ps_j>0.
\]
The cutoff satisfies \(k_0<\lambda_{\mathrm p,j}/4\). The Poisson lower-tail Bernstein inequality states
\[
\nu\left\{
\lambda_{\mathrm p,j}-\mathsf J_j
>\sqrt{2\lambda_{\mathrm p,j}\xi}
\right\}\le e^{-\xi}
\qquad(\xi\ge0).
\]
Since
\[
\sqrt{2\lambda_{\mathrm p,j}(\lambda_{\mathrm p,j}/4)}
=\frac{\lambda_{\mathrm p,j}}{\sqrt2}
<\frac{3\lambda_{\mathrm p,j}}4,
\]
the event \(\{\mathsf J_j\le k_0\}\) is contained in the displayed deviation event with \(\xi=\lambda_{\mathrm p,j}/4\). Hence
\[
\pi_j\le e^{-t_ps_j/4}.
\]

For every \(\zeta_j\ge0\),
\cref{obj:lem:chebyshev-factorial-certificate} supplies
\[
\mathbb E_\nu[\mathcal G_j^2]
\le392^L(1+\zeta_j)^{2L}
\le F_{\mathrm{grow}}(1+\zeta_j)^{2L}.
\]
The multiplier calculation used for light cells now gives
\[
\mathbb E_\nu[\mathcal W_{\mathrm{pol},j}^2]
\le
2F_{\mathrm{grow}}(1+\zeta_j)^{2L}
\left(1+\frac{v_j^2}{B^2}+\frac{v_j}{tB^2}\right).
\]
For \(s_j>B\), one has \(ts_j\ge tB\ge1\). Applying the mixed-monomial calculation with localization scale \(s_j\), while retaining the squared cell mass, gives
\[
\left(\frac{q_j}{u}+q_j^2\right)
\left(1+\frac{v_j^2}{s_j^2}+\frac{v_j}{ts_j^2}\right)
\le
\frac{3p_j}{u\epsilon}+\frac{p_j}{t}+2p_j^2.
\]
Also,
\[
\begin{aligned}
1+\frac{v_j^2}{B^2}+\frac{v_j}{tB^2}
&=
1+\zeta_j^2
\left(\frac{v_j^2}{s_j^2}+\frac{v_j}{ts_j^2}\right)\\
&\le
(1+\zeta_j)^2
\left(1+\frac{v_j^2}{s_j^2}+\frac{v_j}{ts_j^2}\right).
\end{aligned}
\]
Factoring the independent outcome second moment therefore proves
\[
\mathbb E_\nu[\mathcal P_j^2]
\le
6F_{\mathrm{grow}}(1+\zeta_j)^{2L+2}
\left(\frac{p_j}{u\epsilon}+p_j^2+\frac{p_j}{t}\right).
\]

The pilot tail absorbs this growth. Indeed, \(\zeta_j>1\),
\(t_ps_j\ge H_0L\zeta_j\), \(\log(1+\zeta_j)\le\zeta_j\), and
\(2L+2\le(5/2)L\). Thus
\[
\begin{aligned}
&\log\left(
e^{-t_ps_j/4}F_{\mathrm{grow}}(1+\zeta_j)^{2L+2}
\right)\\
&\qquad\le
-\frac{H_0L\zeta_j}{4}
+24L+\frac52L\zeta_j\\
&\qquad\le
-\left(\frac{H_0}{4}-24-\frac52\right)L\zeta_j
\le-200000L.
\end{aligned}
\]
Consequently,
\[
\pi_jF_{\mathrm{grow}}(1+\zeta_j)^{2L+2}
\le e^{-200000L}\le\rho_{\mathrm{pil}},
\qquad
\pi_j\le\rho_{\mathrm{pil}}.
\]
Since \(0\le p_j\le1\) and \(\sum_jp_j=1\),
\[
\sum_{j\in\mathcal J_{\mathrm{heavy}}}
\pi_j\mathbb E_\nu[\mathcal P_j^2]
\le
6\rho_{\mathrm{pil}}
\left(\frac1{u\epsilon}+1+\frac1t\right).
\]

The inverse branch mean lies in \([0,p_j]\), by its explicit formula. Moreover,
\((\mathbb E_\nu[\mathcal P_j])^2\le
\mathbb E_\nu[\mathcal P_j^2]\). Hence
\[
\pi_j(1-\pi_j)\delta_j^2
\le
2\pi_j\mathbb E_\nu[\mathcal P_j^2]+2\pi_jp_j^2,
\]
and summing gives
\[
\sum_{j\in\mathcal J_{\mathrm{heavy}}}
\pi_j(1-\pi_j)\delta_j^2
\le
14\rho_{\mathrm{pil}}
\left(\frac1{u\epsilon}+1+\frac1t\right).
\]
Using the mixture variance identity and
\(\operatorname{Var}_\nu(\mathcal P_j)\le
\mathbb E_\nu[\mathcal P_j^2]\), we conclude
\[
\begin{aligned}
\sum_{j\in\mathcal J_{\mathrm{heavy}}}
\operatorname{Var}_\nu(\mathcal V_j)
\le{}&
\sum_{j\in\mathcal J_{\mathrm{heavy}}}
\operatorname{Var}_\nu(\mathcal H_j)\\
&+20\rho_{\mathrm{pil}}
\left(\frac1{u\epsilon}+1+\frac1t\right).
\end{aligned}
\]
Adding the light and heavy bounds proves
\[
\begin{aligned}
\sum_{j=1}^d\operatorname{Var}_\nu(\mathcal V_j)
\le{}&
\sum_{j=1}^d\operatorname{Var}_\nu(\mathcal H_j)\\
&+6F_{\mathrm{grow}}
\left(
\frac{M_{\mathrm{light}}}{u\epsilon}
+\frac{M_{\mathrm{light}}}{t}
+\frac{c_{\mathrm{light}}B^2}{\epsilon^2}
\right)\\
&+2c_{\mathrm{light}}
\left(\frac B{\epsilon L^2}\right)^2
+\frac1{e\,t\epsilon}\\
&+20\rho_{\mathrm{pil}}
\left(\frac1{u\epsilon}+1+\frac1t\right).
\end{aligned}
\]
% lean: CausalSmith.Stat.AnnotationRarearmFrontier.false_light_absorption
% lean: CausalSmith.Stat.AnnotationRarearmFrontier.heavy_hybrid_variance_sum
% lean: CausalSmith.Stat.AnnotationRarearmFrontier.hybrid_canonical_variance_sum

\item \textit{Bound the selected bias.}
Define the exponential pilot parameter by
\[
\theta_{\mathrm{pil}}=\log\left(1+\frac t{t_p}\right)>0.
\]
The Poisson exponential moment and exponential Markov inequality give
\[
\begin{aligned}
(1-\pi_j)e^{-ts_j}
&\le
e^{-\theta_{\mathrm{pil}}k_0}
\mathbb E_\nu[e^{\theta_{\mathrm{pil}}\mathsf J_j}]
e^{-ts_j}\\
&=
e^{-\theta_{\mathrm{pil}}k_0}
\exp\left\{t_ps_j(e^{\theta_{\mathrm{pil}}}-1)-ts_j\right\}\\
&=e^{-\theta_{\mathrm{pil}}k_0}.
\end{aligned}
\]
This calculation applies to every \(s_j\ge0\). The ratio condition and
\(\log(1+x)\ge2x/(2+x)\) for \(x\ge0\) imply
\[
\theta_{\mathrm{pil}}\ge\log(4/3)\ge\frac27.
\]
The floor inequality \(k_0+1>t_pB/4\), together with \(L\ge4\), gives
\[
k_0\ge\frac{H_0L}{8}.
\]
It follows that
\[
(1-\pi_j)e^{-ts_j}
\le
\exp\left(-\frac{H_0L}{28}\right)
\le e^{-30000L}
\le e^{-20\eta_S}
=\rho_{\mathrm{pil}}.
\]

The branch means yield, for \(s_j>0\),
\[
\mathbb E_\nu[\mathcal V_j]-p_j\bar\mu_j
=
-v_j\bar\mu_j
\left\{\pi_jE_L(\zeta_j)+(1-\pi_j)e^{-ts_j}\right\}.
\]
For \(s_j\le B\), the light residual estimate therefore gives
\[
\left|\mathbb E_\nu[\mathcal V_j]-p_j\bar\mu_j\right|
\le\frac B{\epsilon L^2}+p_j\rho_{\mathrm{pil}}.
\]

For heavy cells, the coefficient bound in
\cref{obj:lem:chebyshev-factorial-certificate} implies
\[
|G_L(\zeta_j)|
\le
\sum_{h=0}^{L-2}|g_h|\zeta_j^h
\le14^L(1+\zeta_j)^L.
\]
The defining polynomial identity gives
\[
\begin{aligned}
|E_L(\zeta_j)|
&=|1-\zeta_jG_L(\zeta_j)|\\
&\le1+\zeta_j14^L(1+\zeta_j)^L\\
&\le2F_{\mathrm{grow}}(1+\zeta_j)^{2L+2}.
\end{aligned}
\]
Here each of the two terms in the preceding sum is bounded by
\(F_{\mathrm{grow}}(1+\zeta_j)^{2L+2}\).
Using \(v_j\bar\mu_j\le p_j\) and the heavy pilot absorption proves
\[
\pi_j|v_j\bar\mu_jE_L(\zeta_j)|
\le2p_j\rho_{\mathrm{pil}}.
\]
Together with the opposite pilot bound, this gives
\[
\left|\mathbb E_\nu[\mathcal V_j]-p_j\bar\mu_j\right|
\le3p_j\rho_{\mathrm{pil}}
\qquad(s_j>B).
\]
For \(s_j=0\), the bias is zero. Applying the triangle inequality cell by cell and summing \(p_j\) proves
\[
\left|
\sum_{j=1}^d
\bigl(\mathbb E_\nu[\mathcal V_j]-p_j\bar\mu_j\bigr)
\right|
\le\frac{dB}{\epsilon L^2}+3\rho_{\mathrm{pil}}.
\]
Since \(S\ge1\), one has \(0\le\rho_{\mathrm{pil}}\le1\), and hence
\[
\begin{aligned}
\left(
\sum_{j=1}^d
\bigl(\mathbb E_\nu[\mathcal V_j]-p_j\bar\mu_j\bigr)
\right)^2
&\le
\frac{2d^2B^2}{\epsilon^2L^4}
+18\rho_{\mathrm{pil}}^2\\
&\le
\frac{2d^2B^2}{\epsilon^2L^4}
+18\rho_{\mathrm{pil}}.
\end{aligned}
\]
% lean: CausalSmith.Stat.AnnotationRarearmFrontier.pilot_false_heavy_absorption
% lean: CausalSmith.Stat.AnnotationRarearmFrontier.hybrid_selected_squared_bias_sum

\item \textit{Form the canonical mean-squared-error envelope.}
Poisson counts have moments of every order, and the inverse denominators are at least one, so all the branch and hybrid statistics have finite second moments. Under the product measure \(\nu\), distinct cells are independent. Therefore the variance and squared-bias decomposition is
\[
\begin{aligned}
\mathbb E_\nu\left[
\left(\sum_{j=1}^d\mathcal V_j-\psi_a\right)^2
\right]
={}&
\sum_{j=1}^d\operatorname{Var}_\nu(\mathcal V_j)\\
&+
\left(
\sum_{j=1}^d
\bigl(\mathbb E_\nu[\mathcal V_j]-p_j\bar\mu_j\bigr)
\right)^2.
\end{aligned}
\]
Substituting the variance and squared-bias bounds gives
\[
\begin{aligned}
\mathbb E_\nu\left[
\left(\sum_{j=1}^d\mathcal V_j-\psi_a\right)^2
\right]
\le{}&
\sum_{j=1}^d\operatorname{Var}_\nu(\mathcal H_j)\\
&+6F_{\mathrm{grow}}
\left(
\frac{M_{\mathrm{light}}}{u\epsilon}
+\frac{M_{\mathrm{light}}}{t}
+\frac{c_{\mathrm{light}}B^2}{\epsilon^2}
\right)\\
&+2c_{\mathrm{light}}
\left(\frac B{\epsilon L^2}\right)^2
+\frac1{e\,t\epsilon}\\
&+20\rho_{\mathrm{pil}}
\left(\frac1{u\epsilon}+1+\frac1t\right)\\
&+\frac{2d^2B^2}{\epsilon^2L^4}
+18\rho_{\mathrm{pil}}.
\end{aligned}
\]
% lean: CausalSmith.Stat.AnnotationRarearmFrontier.hybrid_canonical_mse_decomposition
% lean: CausalSmith.Stat.AnnotationRarearmFrontier.hybrid_canonical_mse_envelope

\item \textit{Normalize the retained terms.}
Define the nonnegative deterministic envelope
\[
\begin{aligned}
\mathscr E_{\mathrm{rate}}
={}&
\frac1S+\frac{d^2B^2}{\epsilon^2L^4}\\
&+F_{\mathrm{grow}}
\left(
\frac{M_{\mathrm{light}}}{S}
+\frac{M_{\mathrm{light}}}{t}
+\frac{dB^2}{\epsilon^2}
\right)
+\frac{dB^2}{\epsilon^2L^4}
+\rho_{\mathrm{pil}}.
\end{aligned}
\]
The intensity bounds, \(S\le N\epsilon\), and \(\epsilon\le1\) imply the first three bounds below; the experiment hypothesis \(S\ge\exp(4096)\ge1\) gives the fourth:
\[
\frac1{u\epsilon}\le\frac{32}{S},
\qquad
\frac1{t\epsilon}\le\frac{64}{S},
\qquad
\frac1t\le\frac{64}{S},
\qquad
\frac1S\le1.
\]
Using \(c_{\mathrm{light}}\le d\), each term of the preceding mean-squared-error envelope satisfies
\[
\begin{aligned}
\sum_j\operatorname{Var}_\nu(\mathcal H_j)
&\le96C_{\mathrm{inv}}\mathscr E_{\mathrm{rate}},\\
6F_{\mathrm{grow}}
\left(
\frac{M_{\mathrm{light}}}{u\epsilon}
+\frac{M_{\mathrm{light}}}{t}
+\frac{c_{\mathrm{light}}B^2}{\epsilon^2}
\right)
&\le192\mathscr E_{\mathrm{rate}},\\
2c_{\mathrm{light}}
\left(\frac B{\epsilon L^2}\right)^2
&\le2\mathscr E_{\mathrm{rate}},\\
\frac1{e\,t\epsilon}
&\le64\mathscr E_{\mathrm{rate}},\\
20\rho_{\mathrm{pil}}
\left(\frac1{u\epsilon}+1+\frac1t\right)
&\le1940\mathscr E_{\mathrm{rate}},\\
\frac{2d^2B^2}{\epsilon^2L^4}
&\le2\mathscr E_{\mathrm{rate}},\\
18\rho_{\mathrm{pil}}
&\le18\mathscr E_{\mathrm{rate}}.
\end{aligned}
\]
For the polynomial term, the first reciprocal bound multiplies
\(M_{\mathrm{light}}/S\) by \(32\), and the other two nonnegative summands are also bounded by \(32\) times their counterparts in \(\mathscr E_{\mathrm{rate}}\).
For the pilot term, the parentheses are at most \(32+1+64=97\).
Adding the seven inequalities, with \(C_{\mathrm{inv}}=262146\), yields
\[
\mathbb E_\nu\left[
\left(\sum_j\mathcal V_j-\psi_a\right)^2
\right]
\le25168234\,\mathscr E_{\mathrm{rate}}.
\]
% lean: CausalSmith.Stat.AnnotationRarearmFrontier.hybrid_canonical_envelope_rate

\item \textit{Absorb the deterministic envelope.}
Define
\[
A_{\mathrm{bw}}=64H_0,
\qquad
\zeta_{\mathrm{rate}}=\frac{d}{N\epsilon L}.
\]
Since \(\min\{t_p,t\}\ge N/64\),
\[
B\le\frac{A_{\mathrm{bw}}L}{N}.
\]
Substituting this bound into the light mass and bandwidth terms gives
\[
M_{\mathrm{light}}
\le A_{\mathrm{bw}}\zeta_{\mathrm{rate}}L^2,
\]
\[
\frac{dB^2}{\epsilon^2}
\le
\frac{A_{\mathrm{bw}}^2\zeta_{\mathrm{rate}}L^3}{N\epsilon}
\le
\frac{A_{\mathrm{bw}}^2\zeta_{\mathrm{rate}}L^3}{S},
\]
\[
\frac{d^2B^2}{\epsilon^2L^4}
\le A_{\mathrm{bw}}^2\zeta_{\mathrm{rate}}^2,
\qquad
\frac{dB^2}{\epsilon^2L^4}
\le\frac{d^2B^2}{\epsilon^2L^4}
\le A_{\mathrm{bw}}^2\zeta_{\mathrm{rate}}^2,
\]
where the last comparison uses \(d\ge1\).

The calibration \(F_{\mathrm{grow}}L^3\le\sqrt S\) and \(L\ge1\) imply
\[
F_{\mathrm{grow}}M_{\mathrm{light}}
\le A_{\mathrm{bw}}\zeta_{\mathrm{rate}}\sqrt S.
\]
Together with \(1/t\le64/S\), the preceding bandwidth bounds give
\[
\begin{aligned}
\frac{F_{\mathrm{grow}}M_{\mathrm{light}}}{S}
&\le\frac{A_{\mathrm{bw}}\zeta_{\mathrm{rate}}}{\sqrt S},\\
\frac{F_{\mathrm{grow}}M_{\mathrm{light}}}{t}
&\le\frac{64A_{\mathrm{bw}}\zeta_{\mathrm{rate}}}{\sqrt S},\\
F_{\mathrm{grow}}\frac{dB^2}{\epsilon^2}
&\le\frac{A_{\mathrm{bw}}^2\zeta_{\mathrm{rate}}}{\sqrt S}.
\end{aligned}
\]
Consequently,
\[
F_{\mathrm{grow}}
\left(
\frac{M_{\mathrm{light}}}{S}
+\frac{M_{\mathrm{light}}}{t}
+\frac{dB^2}{\epsilon^2}
\right)
\le
(65A_{\mathrm{bw}}+A_{\mathrm{bw}}^2)
\frac{\zeta_{\mathrm{rate}}}{\sqrt S}.
\]
The nonnegative square
\((\zeta_{\mathrm{rate}}-1/\sqrt S)^2\) gives
\[
\frac{2\zeta_{\mathrm{rate}}}{\sqrt S}
\le\zeta_{\mathrm{rate}}^2+\frac1S.
\]
Also \(S\ge1\) implies \(\rho_{\mathrm{pil}}=S^{-20}\le S^{-1}\).
Define the numerical constants
\[
C_{\mathrm{abs}}
=
2+2A_{\mathrm{bw}}^2
+\frac{65A_{\mathrm{bw}}+A_{\mathrm{bw}}^2}{2},
\qquad
C=25168234\,C_{\mathrm{abs}}>0.
\]
All the retained terms are therefore bounded as
\[
\mathscr E_{\mathrm{rate}}
\le
C_{\mathrm{abs}}
\left(\frac1S+\zeta_{\mathrm{rate}}^2\right).
\]
We have proved the canonical estimate
\[
\mathbb E_\nu\left[
\left(\sum_{j=1}^d\mathcal V_j-\psi_a\right)^2
\right]
\le
C\left\{
\frac1S+\left(\frac{d}{N\epsilon L}\right)^2
\right\}.
\]
The displayed construction makes \(C\) a fixed numerical constant.
% lean: CausalSmith.Stat.AnnotationRarearmFrontier.hybrid_intermediate_envelope_rate
% lean: CausalSmith.Stat.AnnotationRarearmFrontier.hybrid_canonical_arm_rate

\item \textit{Transfer the canonical estimate to the given counts.}
Define the squared-error function on the count space by
\[
\mathcal R:\mathcal E\longrightarrow[0,\infty),
\qquad
\mathcal R(\boldsymbol c)
=
\left(
\sum_{j=1}^d\mathcal V_j(\boldsymbol c)-\psi_a
\right)^2.
\]
The definitions of the branches give, for every \(\omega\in\Omega\),
\[
\mathcal V_j(\mathscr T(\omega))=V_{aj}(\omega).
\]
The count map is measurable, and every real-valued function on the countable count space is measurable. Using the joint law established at the outset, change of variables therefore gives
\[
\begin{aligned}
\mathbb E\left[
\left(\sum_{j=1}^dV_{aj}-\psi_a\right)^2
\right]
&=
\int_\Omega\mathcal R(\mathscr T(\omega))\,d\mathbb P(\omega)\\
&=
\int_{\mathcal E}\mathcal R(\boldsymbol c)\,d\nu(\boldsymbol c)\\
&\le
C\left\{
\frac1S+\left(\frac{d}{N\epsilon L}\right)^2
\right\}.
\end{aligned}
\]
Finally, \(\psi_a=\sum_jp_j(P)\mu_{aj}(P)\). The same numerical \(C\) thus applies to every admissible choice in the statement.
% lean: CausalSmith.Stat.AnnotationRarearmFrontier.uniform_hybrid_arm_risk
\end{enumerate}
\end{proof}

The arm-risk bound separates the rare-arm outcome scale from the many-cell approximation term. Under the stated logarithmic degree calibration, the polynomial correction improves the dependence of the latter term on total marginal information. Pilot localization and the count-moment controls keep the combined contribution within the displayed squared-error bound, uniformly over the observable laws and balanced pool intensities in the statement.

For the large-information branch of \cref{obj:def:estimator-handle}, \cref{obj:lem:uniform-hybrid-arm-risk} supplies the risk bound for each arm. Their clipped contrast has the same risk order for the target in \cref{obj:def:target}. The prescribed pool allocation satisfies the intensity and balance conditions, and \cref{obj:lem:finite-prefix-transfer} carries this bound to the deterministic finite prefix average on the original complete and auxiliary records. The exponentially small transfer remainder is absorbed by the inverse rare-arm information term, while the degree is comparable to the regularized logarithm in \cref{obj:def:rate-handle}. Together with the bounded-target control for the zero-output branch, these ingredients give the attaining-estimator upper comparison in \cref{obj:thm:uniform-annotation-frontier}.


\section{Testing and approximation lower bounds}

The lower comparison in \cref{obj:thm:uniform-annotation-frontier} combines rare-label testing with a many-cell construction based on polynomial approximation. The first component isolates the information supplied by complete records about rare-arm outcomes. The second connects a reciprocal approximation gap to priors on complete-record tables, with a common distribution of treatment--covariate marginals. We first fix the testing distances and signed-measure norm used in these arguments.



The probability distance measures separation between observation laws, whereas the signed-measure norm fixes the normalization of an approximation certificate. The next three results supply, respectively, the finite-alphabet Pinsker inequality \citep[Lemma 2.5]{Tsybakov2009Nonparametric}, Markov's polynomial inequality \citep[Theorem 1.1]{Pierzchala2016Polynomial}, and the finite alternation representation used in polynomial moment matching \citep[Remark 4 and Appendix E, identity (34)]{WuYang2016Entropy}. These are standard tools; the new statistical construction begins with their overlap-scaled specialization below.

\begingroup\footnotesize
\begin{lemmav}{lem:pinsker-finite}[Pinsker inequality on finite alphabets]
Let \(\nu,\nu'\) be probability measures on a finite set. Define their total variation distance and natural-logarithm Kullback--Leibler divergence by
\[
\operatorname{TV}(\nu,\nu')
=\frac12\sum_z|\nu(z)-\nu'(z)|,
\qquad
D(\nu\Vert\nu')
=
\begin{cases}
\displaystyle\sum_z \nu(z)\log\frac{\nu(z)}{\nu'(z)},
& \nu'(z)=0\Longrightarrow \nu(z)=0\ \text{for every }z,\\
+\infty,& \text{otherwise}.
\end{cases}
\]
Each summand with \(\nu(z)=0\) is defined to be zero. Then, with the inequality interpreted in the extended nonnegative reals,
\[
\operatorname{TV}(\nu,\nu')
\le \sqrt{\frac{D(\nu\Vert\nu')}{2}}.
\]
\end{lemmav}

\begin{proof}[Proof of \cref{obj:lem:pinsker-finite}]
\begin{enumerate}
\item Let \(\mathcal Z_{\mathrm{fin}}\) denote the finite alphabet. If some \(z\in\mathcal Z_{\mathrm{fin}}\) satisfies \(\nu'(z)=0\) and \(\nu(z)>0\), then
\[
D(\nu\Vert\nu')=+\infty,
\qquad
\sqrt{D(\nu\Vert\nu')/2}=+\infty,
\]
which proves the inequality in this case. Henceforth assume
\[
\nu'(z)=0\Longrightarrow \nu(z)=0
\qquad(z\in\mathcal Z_{\mathrm{fin}}).
\]
% lean: CausalSmith.Stat.AnnotationRarearmFrontier.PinskerFinite

\item Equip the alphabet with its full power-set sigma-algebra and construct the discrete probability measures
\[
\bar\nu(E)=\sum_{z\in E}\nu(z),
\qquad
\bar\nu'(E)=\sum_{z\in E}\nu'(z)
\qquad(E\subseteq\mathcal Z_{\mathrm{fin}}).
\]
In particular,
\[
\bar\nu(\{z\})=\nu(z),
\qquad
\bar\nu'(\{z\})=\nu'(z).
\]
If \(\bar\nu'(E)=0\), nonnegativity forces \(\nu'(z)=0\) for every \(z\in E\). The support condition then gives
\[
\bar\nu(E)=\sum_{z\in E}\nu(z)=0.
\]
Thus \(\bar\nu\ll\bar\nu'\).

Define the density and log-likelihood functions on the whole alphabet by
\[
\rho_{\mathrm{fin}}(z)=
\begin{cases}
\nu(z)/\nu'(z),&\nu'(z)>0,\\
0,&\nu'(z)=0,
\end{cases}
\qquad
\lambda_{\mathrm{fin}}(z)=
\begin{cases}
\log\rho_{\mathrm{fin}}(z),&\nu(z)>0,\\
0,&\nu(z)=0.
\end{cases}
\]
The singleton density identity is
\[
\rho_{\mathrm{fin}}(z)\bar\nu'(\{z\})
=\bar\nu(\{z\}),
\]
so \(\rho_{\mathrm{fin}}\) is a density of \(\bar\nu\) with respect to \(\bar\nu'\). At every positive \(\nu\)-mass, both masses in the ratio are positive. Consequently, the log-likelihood is integrable on this finite space, and the measure divergence is finite. Writing its real value as
\[
\mathcal K_{\mathrm{fin}}
:=\int\lambda_{\mathrm{fin}}\,d\bar\nu,
\]
evaluation on singletons gives
\[
\mathcal K_{\mathrm{fin}}
=\sum_{z\in\mathcal Z_{\mathrm{fin}}}
  \nu(z)\lambda_{\mathrm{fin}}(z)
=\sum_{z\in\mathcal Z_{\mathrm{fin}}}
  \nu(z)\log\frac{\nu(z)}{\nu'(z)}
=D(\nu\Vert\nu')<\infty.
\]
Here, as in the statement, every zero-\(\nu\)-mass summand is zero.
% lean: Causalean.Mathlib.InformationTheory.klDiv_toReal_eq_sum_measureReal

\item Define the event-supremum distance and the positive-gap event by
\[
\mathsf V_{\mathrm{fin}}
:=\sup_{E\subseteq\mathcal Z_{\mathrm{fin}}}
  |\bar\nu(E)-\bar\nu'(E)|,
\qquad
\mathcal E_{\mathrm{fin},+}
:=\{z\in\mathcal Z_{\mathrm{fin}}:\nu(z)-\nu'(z)\ge0\}.
\]
Normalization yields
\[
\sum_z\bigl(\nu(z)-\nu'(z)\bigr)=0.
\]
Splitting this sum and its absolute-value counterpart over
\(\mathcal E_{\mathrm{fin},+}\) and its complement gives
\[
\begin{aligned}
0
&=\sum_{z\in\mathcal E_{\mathrm{fin},+}}(\nu(z)-\nu'(z))
 +\sum_{z\notin\mathcal E_{\mathrm{fin},+}}(\nu(z)-\nu'(z)),\\
\sum_z|\nu(z)-\nu'(z)|
&=\sum_{z\in\mathcal E_{\mathrm{fin},+}}(\nu(z)-\nu'(z))
 -\sum_{z\notin\mathcal E_{\mathrm{fin},+}}(\nu(z)-\nu'(z))\\
&=2\sum_{z\in\mathcal E_{\mathrm{fin},+}}(\nu(z)-\nu'(z)).
\end{aligned}
\]
The signs defining the event justify the second line. Finite additivity on the disjoint singletons also gives
\[
\bar\nu(\mathcal E_{\mathrm{fin},+})
-\bar\nu'(\mathcal E_{\mathrm{fin},+})
=\sum_{z\in\mathcal E_{\mathrm{fin},+}}(\nu(z)-\nu'(z)).
\]
Every event gap is bounded by the defining supremum; its terms lie in \([0,1]\). Therefore
\[
\operatorname{TV}(\nu,\nu')
=\sum_{z\in\mathcal E_{\mathrm{fin},+}}(\nu(z)-\nu'(z))
\le
\left|\sum_{z\in\mathcal E_{\mathrm{fin},+}}(\nu(z)-\nu'(z))\right|
\le\mathsf V_{\mathrm{fin}}.
\]
% lean: CausalSmith.Stat.AnnotationRarearmFrontier.PinskerFinite.hTV

\item We next bound the absolute density deviation. Define, for \(t_{\mathrm{fin}}\in[0,\infty)\),
\[
\kappa_{\mathrm{fin}}(t_{\mathrm{fin}})
:=t_{\mathrm{fin}}\log t_{\mathrm{fin}}+1-t_{\mathrm{fin}},
\qquad
\Phi_{\mathrm{fin}}(t_{\mathrm{fin}})
:=(t_{\mathrm{fin}}+2)\kappa_{\mathrm{fin}}(t_{\mathrm{fin}})
-\frac32(t_{\mathrm{fin}}-1)^2,
\]
with \(0\log0=0\). The density identity and normalization imply
\[
\int\rho_{\mathrm{fin}}\,d\bar\nu'=1,
\qquad
\mathcal K_{\mathrm{fin}}
=\int\kappa_{\mathrm{fin}}(\rho_{\mathrm{fin}})\,d\bar\nu'.
\]
Indeed, the density identity converts the logarithmic term into
\(\int\lambda_{\mathrm{fin}}\,d\bar\nu\), and the remaining integral is \(1-1=0\).

For \(t_{\mathrm{fin}}>0\), differentiation gives
\[
\begin{aligned}
\Phi_{\mathrm{fin}}'(t_{\mathrm{fin}})
&=2(t_{\mathrm{fin}}+1)\log t_{\mathrm{fin}}
  -4(t_{\mathrm{fin}}-1),\\
\Phi_{\mathrm{fin}}''(t_{\mathrm{fin}})
&=2\bigl(\log t_{\mathrm{fin}}+t_{\mathrm{fin}}^{-1}-1\bigr)
\ge0.
\end{aligned}
\]
The last inequality follows from
\[
-\log t_{\mathrm{fin}}
=\log(t_{\mathrm{fin}}^{-1})
\le t_{\mathrm{fin}}^{-1}-1.
\]
Hence \(\Phi_{\mathrm{fin}}'\) is increasing. Since
\[
\Phi_{\mathrm{fin}}(1)=\Phi_{\mathrm{fin}}'(1)=0,
\]
its derivative is nonpositive on \((0,1]\) and nonnegative on
\([1,\infty)\). Thus \(\Phi_{\mathrm{fin}}\) decreases toward zero on
\((0,1]\) and increases from zero on \([1,\infty)\). At the endpoint,
\(\Phi_{\mathrm{fin}}(0)=1/2\). Dividing its nonnegativity by the positive denominator yields
\[
0\le
\frac{(t_{\mathrm{fin}}-1)^2}{t_{\mathrm{fin}}+2}
\le\frac23\kappa_{\mathrm{fin}}(t_{\mathrm{fin}})
\qquad(t_{\mathrm{fin}}\ge0).
\]

Define the following functions on \(\mathcal Z_{\mathrm{fin}}\):
\[
\begin{aligned}
g_{\mathrm{fin}}(z)
&:=\frac{(\rho_{\mathrm{fin}}(z)-1)^2}{\rho_{\mathrm{fin}}(z)+2},\\
f_{\mathrm{fin},1}(z)
&:=\frac{|\rho_{\mathrm{fin}}(z)-1|}
        {\sqrt{\rho_{\mathrm{fin}}(z)+2}},\\
f_{\mathrm{fin},2}(z)
&:=\sqrt{\rho_{\mathrm{fin}}(z)+2}.
\end{aligned}
\]
All these functions are integrable, together with their squares, on the finite space. The scalar bound gives
\[
0\le\int g_{\mathrm{fin}}\,d\bar\nu'
\le\frac23\mathcal K_{\mathrm{fin}},
\qquad
\mathcal K_{\mathrm{fin}}\ge0.
\]
Moreover,
\[
f_{\mathrm{fin},1}^2=g_{\mathrm{fin}},
\qquad
f_{\mathrm{fin},2}^2=\rho_{\mathrm{fin}}+2,
\qquad
f_{\mathrm{fin},1}f_{\mathrm{fin},2}
=|\rho_{\mathrm{fin}}-1|,
\]
and therefore
\[
\int f_{\mathrm{fin},2}^2\,d\bar\nu'=1+2=3.
\]
Cauchy--Schwarz now yields
\[
\begin{aligned}
\int|\rho_{\mathrm{fin}}-1|\,d\bar\nu'
&\le
\sqrt{\int f_{\mathrm{fin},1}^2\,d\bar\nu'}
\sqrt{\int f_{\mathrm{fin},2}^2\,d\bar\nu'}\\
&=\sqrt{\int g_{\mathrm{fin}}\,d\bar\nu'}\sqrt3\\
&\le\sqrt{\frac23\mathcal K_{\mathrm{fin}}}\sqrt3
=\sqrt{2\mathcal K_{\mathrm{fin}}}.
\end{aligned}
\]
% lean: Causalean.Stat.pinskerBound_of_ac_of_ne_top

\item To relate the event supremum to this integral, define
\[
\delta_{\mathrm{fin}}(z):=\rho_{\mathrm{fin}}(z)-1
\qquad(z\in\mathcal Z_{\mathrm{fin}}).
\]
Then
\[
\int\delta_{\mathrm{fin}}\,d\bar\nu'=0,
\qquad
\bar\nu(E)-\bar\nu'(E)
=\int_E\delta_{\mathrm{fin}}\,d\bar\nu'
\qquad(E\subseteq\mathcal Z_{\mathrm{fin}}).
\]
The first identity follows from normalization; the second follows by integrating the density and subtracting the integral of \(1\). Splitting the zero integral gives
\[
\int_{E^c}\delta_{\mathrm{fin}}\,d\bar\nu'
=-\int_E\delta_{\mathrm{fin}}\,d\bar\nu'.
\]
On each of \(E\) and \(E^c\), integration of
\(\delta_{\mathrm{fin}}\le|\delta_{\mathrm{fin}}|\) and
\(-\delta_{\mathrm{fin}}\le|\delta_{\mathrm{fin}}|\) implies
\[
\left|\int_E\delta_{\mathrm{fin}}\,d\bar\nu'\right|
\le\int_E|\delta_{\mathrm{fin}}|\,d\bar\nu',
\qquad
\left|\int_E\delta_{\mathrm{fin}}\,d\bar\nu'\right|
\le\int_{E^c}|\delta_{\mathrm{fin}}|\,d\bar\nu'.
\]
Adding these inequalities proves
\[
|\bar\nu(E)-\bar\nu'(E)|
\le\frac12\int|\delta_{\mathrm{fin}}|\,d\bar\nu'.
\]
Taking the supremum over events and using the preceding bound, we obtain
\[
\mathsf V_{\mathrm{fin}}
\le\frac12\int|\rho_{\mathrm{fin}}-1|\,d\bar\nu'
\le\frac12\sqrt{2\mathcal K_{\mathrm{fin}}}
=\sqrt{\frac{\mathcal K_{\mathrm{fin}}}{2}}.
\]
The final equality follows by squaring the two nonnegative quantities.
% lean: Causalean.Stat.tvDist_le_half_integral_abs_rnDeriv

\item Combining the finite-alphabet comparison with this bound and the divergence identity gives
\[
\operatorname{TV}(\nu,\nu')
\le\mathsf V_{\mathrm{fin}}
\le\sqrt{\frac{\mathcal K_{\mathrm{fin}}}{2}}
=\sqrt{\frac{D(\nu\Vert\nu')}{2}}.
\]
In the present case, the divergence is finite and nonnegative, so the real inequality carries directly to the extended nonnegative reals, where the square root is the power \(1/2\). Together with the infinite-divergence case, this proves the claim.
% lean: CausalSmith.Stat.AnnotationRarearmFrontier.PinskerFinite
\end{enumerate}
\end{proof}
\endgroup

\Cref{obj:lem:pinsker-finite} expresses testing separation in terms of the natural-logarithm divergence. For the approximation component, the relevant control instead concerns how rapidly a bounded polynomial can vary over its interval.

\begin{lemmav}{lem:markov-polynomial}[Markov inequality for polynomials]
Let \(L\) be a nonnegative integer, \(g\) a real polynomial of degree at most \(L\), and \(M_b\in\mathbb R\). Suppose that
\[
|g(x)|\le M_b \qquad \text{for every } x\in[-1,1].
\]
Then
\[
|g'(x)|\le L^2 M_b \qquad \text{for every } x\in[-1,1].
\]
\end{lemmav}

\begin{proof}[Proof of \cref{obj:lem:markov-polynomial}]
\begin{enumerate}
\item For a continuous real function \(f\) on \([-1,1]\), define
\[
\|f\|_{\infty}
  :=\sup\{|f(x)|:x\in[-1,1]\}.
\]
The set in this definition is nonempty and bounded, by compactness and continuity. Consequently, for every \(c\in\mathbb R\),
\[
\|f\|_{\infty}\le c
\quad\Longleftrightarrow\quad
|f(x)|\le c\quad\text{for every }x\in[-1,1].
\]
Indeed, each value is bounded by the supremum, and any common upper bound bounds the supremum. We will establish
\[
\|g'\|_{\infty}\le L^2\|g\|_{\infty}
\]
and then apply this equivalence to the hypothesis and conclusion.
% lean: Causalean.Mathlib.Analysis.FinitePolynomialAlternationDuality.intervalSupNorm_le_iff

\item If \(L=0\), the polynomial \(g\) is constant, so
\[
g'=0,\qquad \|g'\|_{\infty}=0=L^2\|g\|_{\infty}.
\]
Suppose henceforth that \(L>0\). Define the Chebyshev polynomials, viewed over either \(\mathbb R\) or \(\mathbb C\), by
\[
T_0(z)=1,\qquad T_1(z)=z,\qquad
T_{h+1}(z)=2zT_h(z)-T_{h-1}(z)\quad(h\ge1),
\]
and define their extrema and root angles by
\[
\xi_i:=\cos\frac{i\pi}{L}\quad(0\le i\le L),\qquad
\alpha_k:=\frac{(2k+1)\pi}{2L},\qquad
\rho_k:=\cos\alpha_k\quad(0\le k<L).
\]
The extrema are distinct and decreasing, and
\[
\xi_i\in[-1,1],\qquad
T_L(\xi_i)=(-1)^i.
\]
In particular,
\[
|g(\xi_i)|\le\|g\|_{\infty},
\qquad
0\le |g(\xi_0)|\le\|g\|_{\infty}.
\]
For \(0\le i\le L\), define the cardinal polynomials
\[
\lambda_i(z):=
\prod_{\substack{0\le j\le L\\j\ne i}}
\frac{z-\xi_j}{\xi_i-\xi_j}.
\]
Lagrange interpolation and differentiation give, for every real \(x\),
\[
g(z)=\sum_{i=0}^{L}g(\xi_i)\lambda_i(z),
\qquad
g'(x)=\sum_{i=0}^{L}g(\xi_i)\lambda_i'(x).
\]
Thus, for \(x\in[-1,1]\),
\[
|g'(x)|
\le\sum_{i=0}^{L}|g(\xi_i)|\,|\lambda_i'(x)|
\le\|g\|_{\infty}\sum_{i=0}^{L}|\lambda_i'(x)|.
\]
It remains to bound the final sum by \(L^2\).
% lean: Causalean.Mathlib.Analysis.FinitePolynomialAlternationDuality.duffinSchaeffer_derivative_le

\item Fix \(x\in[-1,1]\), and define
\[
\varepsilon_{x,i}:=
\begin{cases}
-1,&\lambda_i'(x)<0,\\
1,&\lambda_i'(x)\ge0,
\end{cases}
\qquad
g_{\mathrm{sgn},x}(z):=\sum_{i=0}^{L}\varepsilon_{x,i}\lambda_i(z).
\]
These definitions include the case \(\lambda_i'(x)=0\). Interpolation gives
\[
\deg g_{\mathrm{sgn},x}\le L,\qquad
g_{\mathrm{sgn},x}(\xi_i)=\varepsilon_{x,i},
\qquad
|g_{\mathrm{sgn},x}(\xi_i)|\le1.
\]
We first establish its derivative bound at the roots of \(T_L\).

Define
\[
\kappa_L:=L\,2^{L-1},\qquad
\mathcal N_L(z):=\prod_{j=0}^{L}(z-\xi_j),\qquad
\eta_i:=
\left(\prod_{\substack{0\le j\le L\\j\ne i}}
(\xi_i-\xi_j)\right)^{-1}.
\]
The following polynomial identities will be used:
\[
\kappa_L\mathcal N_L(z)=(z^2-1)T_L'(z),
\qquad
(1-z^2)T_L''(z)-zT_L'(z)+L^2T_L(z)=0.
\]
For the first identity, both sides have degree \(L+1\), leading coefficient \(\kappa_L\), and vanish at every \(\xi_i\): the endpoints annul \(z^2-1\), and the interior extrema annul \(T_L'\). The second is the Chebyshev differential identity, obtained by differentiating
\(T_L(\cos\theta)=\cos(L\theta)\).

Let \(\zeta\in\mathbb R\) satisfy \(T_L(\zeta)=0\). The explicit roots \(\rho_k\) give
\[
|\zeta|<1,\qquad \zeta\ne\xi_i\quad(0\le i\le L).
\]
Hence \(\mathcal N_L(\zeta)\ne0\), and the nodal identity implies
\(T_L'(\zeta)\ne0\). At this root, the differential identity yields
\[
(1-\zeta^2)T_L''(\zeta)=\zeta T_L'(\zeta).
\]
Differentiating
\[
\kappa_L\prod_{j\ne i}(z-\xi_j)
=\frac{(z^2-1)T_L'(z)}{z-\xi_i}
\]
at \(\zeta\), and using the preceding identity, gives
\[
\lambda_i'(\zeta)
=
\eta_i\,
\frac{T_L'(\zeta)(1-\zeta\xi_i)}
     {\kappa_L(\zeta-\xi_i)^2}.
\]
Since the extrema decrease,
\[
(-1)^i\eta_i>0.
\]
Also \(1-\zeta\xi_i>0\), because \(|\zeta|<1\) and \(|\xi_i|\le1\). Therefore
\[
(-1)^i\lambda_i'(\zeta)T_L'(\zeta)>0.
\]
Interpolation of \(T_L\) gives
\[
T_L'(\zeta)=\sum_{i=0}^{L}(-1)^i\lambda_i'(\zeta).
\]
If \(T_L'(\zeta)>0\), all summands on the right are positive; if \(T_L'(\zeta)<0\), all are negative. In either case,
\[
\sum_{i=0}^{L}|\lambda_i'(\zeta)|=|T_L'(\zeta)|.
\]
Consequently,
\[
\begin{aligned}
|g_{\mathrm{sgn},x}'(\zeta)|
&=\left|\sum_{i=0}^{L}
g_{\mathrm{sgn},x}(\xi_i)\lambda_i'(\zeta)\right|\\
&\le\sum_{i=0}^{L}
|g_{\mathrm{sgn},x}(\xi_i)|\,|\lambda_i'(\zeta)|\\
&\le\sum_{i=0}^{L}|\lambda_i'(\zeta)|
=|T_L'(\zeta)|.
\end{aligned}
\]
% lean: Causalean.Mathlib.Analysis.FinitePolynomialAlternationDuality.abs_eval_derivative_le_chebyshev_at_root

\item At the fixed point \(x\), differentiation of the signed interpolant gives
\[
g_{\mathrm{sgn},x}'(x)
=\sum_{i=0}^{L}\varepsilon_{x,i}\lambda_i'(x)
=\sum_{i=0}^{L}|\lambda_i'(x)|.
\]
Indeed, a negative weight is multiplied by \(-1\), and a nonnegative weight by \(1\). We next propagate the root bounds from the preceding step to this derivative value.
% lean: Causalean.Mathlib.Analysis.FinitePolynomialAlternationDuality.sum_abs_chebyshevDerivativeWeight_le_sq

\item Define the cardinal polynomials at the roots and a reflected-root polynomial by
\[
\lambda_{\mathrm{root},k}(z):=
\prod_{\substack{0\le j<L\\j\ne k}}
\frac{z-\rho_j}{\rho_k-\rho_j}
\quad(0\le k<L),
\qquad
\mathcal R_x(z):=
2^{L-1}\prod_{k=0}^{L-1}(z+|x-\rho_k|).
\]
Empty products equal \(1\), so these definitions also apply when \(L=1\).
Since
\[
T_L(z)=2^{L-1}\prod_{k=0}^{L-1}(z-\rho_k),
\]
we have
\[
T_L'(\rho_k)\lambda_{\mathrm{root},k}(x)
=
2^{L-1}\prod_{\substack{0\le j<L\\j\ne k}}(x-\rho_j).
\]
The derivative \(g_{\mathrm{sgn},x}'\) has degree at most \(L-1\). Interpolating it at the \(L\) roots and using the root bounds therefore gives
\[
\begin{aligned}
|g_{\mathrm{sgn},x}'(x)|
&\le\sum_{k=0}^{L-1}
|g_{\mathrm{sgn},x}'(\rho_k)|\,
|\lambda_{\mathrm{root},k}(x)|\\
&\le\sum_{k=0}^{L-1}
|T_L'(\rho_k)\lambda_{\mathrm{root},k}(x)|\\
&=2^{L-1}\sum_{k=0}^{L-1}
\prod_{\substack{0\le j<L\\j\ne k}}|x-\rho_j|
=|\mathcal R_x'(0)|.
\end{aligned}
\]
The last equality follows by differentiating the product defining
\(\mathcal R_x\); all its derivative summands at zero are nonnegative.
Moreover, for every \(y\in\mathbb R\),
\[
|\mathcal R_x(iy)|
=2^{L-1}\prod_{k=0}^{L-1}
\sqrt{(x-\rho_k)^2+y^2}
=|T_L(x+iy)|,
\qquad
\deg\mathcal R_x\le L.
\]
% lean: Causalean.Mathlib.Analysis.FinitePolynomialAlternationDuality.exists_chebyshevRootReflection

\item We now establish the vertical comparison
\[
|T_L(x+iy)|\le |T_L(1+iy)|
\qquad(x\in[-1,1],\ y\in\mathbb R).
\]
Define the squared chord function and the list of paired chord squares by
\[
d_{\mathrm{ch}}(\beta):=2-2\cos\beta
\quad(\beta\in\mathbb R),
\qquad
\mathscr D_L(\theta):=
\bigl(d_{\mathrm{ch}}(\theta+\alpha_k),
      d_{\mathrm{ch}}(\theta-\alpha_k)\bigr)_{k=0}^{L-1},
\]
where successive pairs are concatenated into a list of length \(2L\).
For an even list of nonnegative entries, define
\[
\mathcal P_t(v_1,\ldots,v_{2h})
:=\prod_{j=1}^{h}(v_{2j-1}v_{2j}+t)
\qquad(h\in\mathbb N,\ t\ge0),
\]
with value \(1\) for the empty list.

Sorting a list in decreasing order and pairing adjacent entries maximizes
\(\mathcal P_t\). To see this, for
\(v_1\ge v_2\ge v_3\ge v_4\ge0\), the two exchanges satisfy
\[
\begin{aligned}
&(v_1v_2+t)(v_3v_4+t)
 -(v_1v_3+t)(v_2v_4+t)
=t(v_1-v_4)(v_2-v_3)\ge0,\\
&(v_1v_2+t)(v_3v_4+t)
 -(v_1v_4+t)(v_2v_3+t)
=t(v_1-v_3)(v_2-v_4)\ge0.
\end{aligned}
\]
Thus the two largest entries can be paired together without decreasing the product. Removing that pair and inducting proves the assertion; the other factors are nonnegative.

For the fixed \(x\), define
\[
\theta_x:=\arccos x,\qquad
\delta_{\mathrm{ang}}:=\frac{\pi}{2L}.
\]
Choose \(k_{\mathrm{shift}}\in\mathbb Z\) nearest to \(L\theta_x/\pi\), and put
\[
\phi_x:=\theta_x-\frac{k_{\mathrm{shift}}\pi}{L},
\qquad
|\phi_x|\le\delta_{\mathrm{ang}}.
\]
The two lists \(\mathscr D_L(\theta_x)\) and
\(\mathscr D_L(\phi_x)\) have the same multiset of entries: their angles run through the odd multiples of \(\pi/(2L)\), shifted by the respective angle, and adding an integer multiple of \(\pi/L\) permutes these angles modulo \(2\pi\).

Define
\[
q_{k,+}:=d_{\mathrm{ch}}(\alpha_k+|\phi_x|),
\qquad
q_{k,-}:=d_{\mathrm{ch}}(\alpha_k-|\phi_x|)
\quad(0\le k<L).
\]
All the arguments lie in \([0,\pi]\). Since \(d_{\mathrm{ch}}\) increases on this interval and consecutive root angles differ by \(2\delta_{\mathrm{ang}}\),
\[
q_{k,-}\le q_{k,+},
\qquad
q_{k,+}\le q_{k+1,-}\quad(0\le k<L-1).
\]
Hence reversing the root blocks and placing \(q_{k,+}\) before \(q_{k,-}\) gives a decreasing list. Replacing \(\phi_x\) by \(|\phi_x|\) merely exchanges the two entries within a block, and reversing blocks preserves their paired product. The sorting inequality therefore implies
\[
\mathcal P_{4y^2}\bigl(\mathscr D_L(\theta_x)\bigr)
\le
\mathcal P_{4y^2}\bigl(\mathscr D_L(\phi_x)\bigr).
\]
The elementary chord identity
\[
d_{\mathrm{ch}}(\theta+\alpha_k)
d_{\mathrm{ch}}(\theta-\alpha_k)
=4(\cos\theta-\rho_k)^2
\]
turns this into
\[
4^L\left(\prod_{k=0}^{L-1}|x-\rho_k+iy|\right)^2
\le
4^L\left(\prod_{k=0}^{L-1}|\cos\phi_x-\rho_k+iy|\right)^2.
\]
The products are nonnegative, so
\[
\prod_{k=0}^{L-1}|x-\rho_k+iy|
\le
\prod_{k=0}^{L-1}|\cos\phi_x-\rho_k+iy|.
\]
Also
\[
\rho_k\le\cos\delta_{\mathrm{ang}}
\le\cos\phi_x\le1.
\]
Thus each factor on the right is bounded by its endpoint counterpart:
\[
|\cos\phi_x-\rho_k+iy|^2
=(\cos\phi_x-\rho_k)^2+y^2
\le(1-\rho_k)^2+y^2
=|1-\rho_k+iy|^2.
\]
Multiplying these bounds and then multiplying by \(2^{L-1}\) proves the required vertical comparison.
% lean: Causalean.Mathlib.Analysis.FinitePolynomialAlternationDuality.norm_eval_chebyshev_le_endpoint_vertical

\item Define
\[
\mathcal S(z):=T_L(1+z).
\]
Its degree and derivative at zero are
\[
\deg\mathcal S=L,\qquad \mathcal S'(0)=T_L'(1)=L^2>0,
\]
and its roots are \(\rho_k-1<0\). The preceding steps give
\[
|\mathcal R_x(iy)|=|T_L(x+iy)|
\le|T_L(1+iy)|=|\mathcal S(iy)|
\qquad(y\in\mathbb R).
\]
On the closed right half-plane, define
\[
\mathcal F_x(z):=\frac{\mathcal R_x(z)}{\mathcal S(z)}
\qquad(\operatorname{Re}z\ge0).
\]
The denominator has no zero there. This quotient is holomorphic in the open right half-plane, continuous on its closure, and, since
\(\deg\mathcal R_x\le\deg\mathcal S\),
\[
\mathcal F_x(z)=O(1)\quad\text{as }|z|\to\infty.
\]
In particular it has the growth bound required by the
Phragm\'en--Lindel\"of principle and is bounded along the positive real axis.
Its boundary values satisfy
\[
|\mathcal F_x(iy)|\le1.
\]
The right-half-plane Phragm\'en--Lindel\"of principle therefore gives
\[
|\mathcal F_x(z)|\le1
\qquad(\operatorname{Re}z\ge0).
\]

Suppose, for a contradiction, that
\(|\mathcal R_x'(0)|>|\mathcal S'(0)|\), and define
\[
\lambda_{\mathrm{der}}:=
\frac{\mathcal R_x'(0)}{\mathcal S'(0)},
\qquad
\mathcal E_x(z):=\mathcal R_x(z)-\lambda_{\mathrm{der}}\mathcal S(z).
\]
Then
\[
|\lambda_{\mathrm{der}}|>1,\qquad \mathcal E_x'(0)=0.
\]
For \(\operatorname{Re}z\ge0\), the equality \(\mathcal E_x(z)=0\) would imply
\(\mathcal F_x(z)=\lambda_{\mathrm{der}}\), contradicting the quotient bound.
Thus every root of \(\mathcal E_x\) lies in the open left half-plane.

If \(\mathcal E_x\) is nonconstant, the Gauss--Lucas theorem places every root of its derivative in the convex hull of its roots, hence in the open left half-plane. This contradicts \(\mathcal E_x'(0)=0\).
If \(\mathcal E_x\) is constant, define
\[
c_{\mathcal E}:=\mathcal E_x(0).
\]
Then
\[
\mathcal F_x(iy)
=\lambda_{\mathrm{der}}+\frac{c_{\mathcal E}}{\mathcal S(iy)}
\longrightarrow\lambda_{\mathrm{der}}
\quad\text{as }y\to+\infty,
\]
because \(\deg\mathcal S=L>0\). The boundary bound implies
\(|\lambda_{\mathrm{der}}|\le1\), again a contradiction. We conclude that
\[
|\mathcal R_x'(0)|\le|\mathcal S'(0)|=L^2.
\]
% lean: Causalean.Mathlib.Analysis.FinitePolynomialAlternationDuality.norm_eval_derivative_zero_le_chebyshev_endpoint_of_vertical

\item Combining the signed-interpolation identity, root reflection, and derivative comparison yields
\[
\sum_{i=0}^{L}|\lambda_i'(x)|
=g_{\mathrm{sgn},x}'(x)
\le |g_{\mathrm{sgn},x}'(x)|
\le|\mathcal R_x'(0)|
\le L^2.
\]
Returning to the interpolation bound for \(g\), we obtain
\[
|g'(x)|\le L^2\|g\|_{\infty}
\qquad(x\in[-1,1]),
\]
and hence
\[
\|g'\|_{\infty}\le L^2\|g\|_{\infty}.
\]
The hypothesis and the supremum characterization in the first step give
\[
\|g\|_{\infty}\le M_b.
\]
Since \(L^2\ge0\), multiplication and transitivity give
\[
\|g'\|_{\infty}
\le L^2\|g\|_{\infty}
\le L^2M_b.
\]
Finally, each derivative value is bounded by its supremum, so
\[
|g'(x)|\le L^2M_b
\qquad\text{for every }x\in[-1,1].
\]
% lean: CausalSmith.Stat.AnnotationRarearmFrontier.MarkovPolynomial
\end{enumerate}
\end{proof}

The degree-squared derivative bound identifies the interval scale relevant to the reciprocal function below. A complementary duality result represents its approximation error by a signed measure that annihilates all moments through the prescribed degree.

\begingroup\footnotesize
\begin{lemmav}{lem:finite-polynomial-duality}[Finite polynomial approximation duality]
Let \(v_-<v_+\), let \(g:[v_-,v_+]\to\mathbb R\) be continuous, and let \(L\ge0\) be an integer. There exist nodes
\[
v_-\le v_1<\cdots<v_{L+2}\le v_+
\]
and real weights \(w_1,\ldots,w_{L+2}\) such that
\[
\sum_{i=1}^{L+2}|w_i|=1,
\qquad
\sum_{i=1}^{L+2}w_i v_i^h=0
\quad\text{for every integer }0\le h\le L,
\]
and
\[
\sum_{i=1}^{L+2}w_i g(v_i)
=
\inf_{Q:\,\deg Q\le L}
\sup_{v\in[v_-,v_+]}|g(v)-Q(v)|,
\]
where the infimum ranges over real polynomials of degree at most \(L\). Equivalently, the finite signed measure
\[
\sigma=\sum_{i=1}^{L+2}w_i\delta_{v_i}
\]
has support on at most \(L+2\) points and satisfies
\[
\|\sigma\|_{\mathrm{TV}}=1,
\qquad
\int v^h\,d\sigma(v)=0
\quad\text{for every integer }0\le h\le L,
\qquad
\int g\,d\sigma
=
\inf_{Q:\,\deg Q\le L}
\sup_{v\in[v_-,v_+]}|g(v)-Q(v)|.
\]
Here \(\delta_{v_i}\) is the unit point mass at \(v_i\), and \(\|\sigma\|_{\mathrm{TV}}=\sum_{i=1}^{L+2}|w_i|\) is the total mass of the variation measure of \(\sigma\).
\end{lemmav}

\begin{proof}[Proof of \cref{obj:lem:finite-polynomial-duality}]
\begin{enumerate}
\item For a continuous function \(f:[v_-,v_+]\to\mathbb R\), write
\[
\|f\|_\infty=\sup_{v\in[v_-,v_+]}|f(v)|,
\qquad
\mathscr P_L=\{Q:\ Q\text{ is a real polynomial of degree at most }L\}.
\]
The restrictions of the polynomials in \(\mathscr P_L\) form a finite-dimensional normed space. Hence the ball
\[
\mathscr B_{\mathrm{app}}
=
\{Q\in\mathscr P_L:\ \|Q\|_\infty\le 2\|g\|_\infty+1\}
\]
is compact, and the continuous function \(Q\mapsto\|g-Q\|_\infty\) attains its minimum there. Choose a minimizer \(Q_*\). Since the zero polynomial belongs to this ball,
\[
\|g-Q_*\|_\infty\le \|g\|_\infty.
\]
For every \(Q\in\mathscr P_L\setminus\mathscr B_{\mathrm{app}}\), the triangle inequality gives
\[
\|g-Q\|_\infty
\ge \|Q\|_\infty-\|g\|_\infty
> \|g\|_\infty+1
> \|g-Q_*\|_\infty.
\]
Thus \(Q_*\) is a global minimizer. Define the approximation error and residual by
\[
\mathcal E_{\mathrm{app}}
=
\inf_{Q\in\mathscr P_L}\|g-Q\|_\infty
=
\|g-Q_*\|_\infty\ge0,
\qquad
\rho(v)=g(v)-Q_*(v)
\quad(v\in[v_-,v_+]).
\]
In particular, \(\rho\) is continuous and
\[
|\rho(v)|\le\mathcal E_{\mathrm{app}}
\quad(v\in[v_-,v_+]).
\]
% lean: Causalean.Mathlib.Analysis.FinitePolynomialAlternationDuality.exists_bestPolynomial

\item If \(\mathcal E_{\mathrm{app}}=0\), define
\[
v_i=v_-+\frac{i-1}{L+1}(v_+-v_-)
\quad(1\le i\le L+2),
\qquad
\eta=1.
\]
The denominator is positive for every \(L\ge0\), so these nodes are strictly increasing and belong to \([v_-,v_+]\). Moreover, the residual bound implies \(\rho(v)=0\) throughout the interval. Consequently,
\[
\rho(v_i)=\eta(-1)^{i-1}\mathcal E_{\mathrm{app}}
\quad(1\le i\le L+2).
\]
We next obtain the same identity when \(\mathcal E_{\mathrm{app}}>0\).
% lean: Causalean.Mathlib.Analysis.FinitePolynomialAlternationDuality.exists_equioscillationWitness

\item Suppose \(\mathcal E_{\mathrm{app}}>0\), and assume for contradiction that there are no \(L+2\) strictly increasing interval points at which \(\rho\) alternates between its two extreme values. Define
\[
\mathcal K_{\mathrm{ext}}
=
\{v\in[v_-,v_+]:|\rho(v)|=\mathcal E_{\mathrm{app}}\},
\qquad
\mathcal K_\pm
=
\{v\in\mathcal K_{\mathrm{ext}}:
\rho(v)=\pm\mathcal E_{\mathrm{app}}\}.
\]
These sets are compact; \(\mathcal K_{\mathrm{ext}}\) is nonempty because \(|\rho|\) attains its maximum, and \(\mathcal K_+\) and \(\mathcal K_-\) are disjoint.

We construct a polynomial \(Q_{\mathrm{sgn}}\in\mathscr P_L\) satisfying
\[
\rho(v)Q_{\mathrm{sgn}}(v)>0
\quad(v\in\mathcal K_{\mathrm{ext}}).
\]
If \(\mathcal K_+=\varnothing\), take \(Q_{\mathrm{sgn}}(v)=-1\); then the product equals \(\mathcal E_{\mathrm{app}}>0\) on \(\mathcal K_{\mathrm{ext}}\). If \(\mathcal K_+\ne\varnothing\) and \(\mathcal K_-=\varnothing\), take \(Q_{\mathrm{sgn}}(v)=1\), with the same conclusion.

It remains to consider the case in which both sign sets are nonempty. Their compactness and disjointness provide a number \(\delta_{\mathrm{ext}}>0\) such that
\[
|v-v'|>\delta_{\mathrm{ext}}
\quad(v\in\mathcal K_+,\ v'\in\mathcal K_-).
\]
Set
\[
\varepsilon_{\mathrm{ext}}=\frac{\delta_{\mathrm{ext}}}{5}.
\]
Choose a finite subcover of \(\mathcal K_{\mathrm{ext}}\) by open balls of radius \(\varepsilon_{\mathrm{ext}}\) centered in \(\mathcal K_{\mathrm{ext}}\), and list the distinct centers in increasing order:
\[
c_1<\cdots<c_{N_{\mathrm{cov}}},
\qquad
c_j\in\mathcal K_{\mathrm{ext}},
\qquad
\mathcal K_{\mathrm{ext}}
\subseteq
\bigcup_{j=1}^{N_{\mathrm{cov}}}
(c_j-\varepsilon_{\mathrm{ext}},c_j+\varepsilon_{\mathrm{ext}}),
\qquad
N_{\mathrm{cov}}\ge1.
\]
Define their signs and the number of adjacent sign changes by
\[
s_j=\frac{\rho(c_j)}{\mathcal E_{\mathrm{app}}}\in\{-1,1\},
\qquad
d_{\mathrm{flip}}
=
\#\{j:\ 1\le j<N_{\mathrm{cov}},\ s_j\ne s_{j+1}\}.
\]
If \(d_{\mathrm{flip}}>L\), selecting one center from each of the first \(L+2\) successive constant-sign runs gives \(L+2\) increasing alternating extrema, contrary to the supposition. Therefore
\[
d_{\mathrm{flip}}\le L.
\]

Define polynomials backwards along the list of centers:
\begingroup\footnotesize
\[
Q_{N_{\mathrm{cov}}}(v)=s_{N_{\mathrm{cov}}},
\qquad
Q_j(v)=
\begin{cases}
Q_{j+1}(v),&s_j=s_{j+1},\\[2mm]
\displaystyle
\left(v-\frac{c_j+c_{j+1}}2\right)Q_{j+1}(v),
&s_j\ne s_{j+1},
\end{cases}
\quad(1\le j<N_{\mathrm{cov}}),
\qquad
Q_{\mathrm{sgn}}=Q_1.
\]
\endgroup
Each sign change contributes one linear factor, so
\[
\deg Q_{\mathrm{sgn}}\le d_{\mathrm{flip}}\le L.
\]
To check its sign, descending induction gives
\[
s_jQ_j(v)>0
\quad\text{if }v<c_j+\varepsilon_{\mathrm{ext}},
\]
and
\[
s_kQ_j(v)>0
\quad\text{if }j\le k\le N_{\mathrm{cov}}
\text{ and }|v-c_k|<\varepsilon_{\mathrm{ext}}.
\]
Indeed, the terminal polynomial is the constant \(s_{N_{\mathrm{cov}}}\). At an equal-sign pair the induction carries over unchanged. At an opposite-sign pair,
\[
c_{j+1}-c_j>\delta_{\mathrm{ext}}>4\varepsilon_{\mathrm{ext}}.
\]
The new linear factor is negative for \(v<c_j+\varepsilon_{\mathrm{ext}}\) and positive within \(\varepsilon_{\mathrm{ext}}\) of every later center. These signs, together with \(s_j=-s_{j+1}\), establish the induction.

For any \(v\in\mathcal K_{\mathrm{ext}}\), choose a covering center \(c_j\). Opposite signs at \(v\) and \(c_j\) would imply
\[
|v-c_j|>\delta_{\mathrm{ext}},
\]
contradicting \(|v-c_j|<\varepsilon_{\mathrm{ext}}\). Hence
\[
\rho(v)=s_j\mathcal E_{\mathrm{app}},
\qquad
\rho(v)Q_{\mathrm{sgn}}(v)
=
\mathcal E_{\mathrm{app}}\,s_jQ_{\mathrm{sgn}}(v)>0.
\]
This completes the sign-polynomial construction.
% lean: Causalean.Mathlib.Analysis.FinitePolynomialAlternationDuality.exists_signPolynomial_of_no_alternatingExtrema

\item We now turn that strict sign agreement into a uniform improvement. Define
\[
p_{\mathrm{sgn}}(v)=\rho(v)Q_{\mathrm{sgn}}(v),
\qquad
\mathcal B_{\mathrm{sgn}}
=
\{v\in[v_-,v_+]:p_{\mathrm{sgn}}(v)\le0\}.
\]
This set is compact and contains no point of \(\mathcal K_{\mathrm{ext}}\). Define
\[
c_{\mathrm{imp}}
=
\begin{cases}
\displaystyle
\frac{\max_{v\in\mathcal B_{\mathrm{sgn}}}|\rho(v)|
+\mathcal E_{\mathrm{app}}}{2},
&\mathcal B_{\mathrm{sgn}}\ne\varnothing,\\[3mm]
\mathcal E_{\mathrm{app}}/2,
&\mathcal B_{\mathrm{sgn}}=\varnothing.
\end{cases}
\]
In the first case the maximum is strictly below \(\mathcal E_{\mathrm{app}}\), since otherwise a maximizing point would belong to \(\mathcal K_{\mathrm{ext}}\). Thus in either case
\[
c_{\mathrm{imp}}<\mathcal E_{\mathrm{app}},
\qquad
|\rho(v)|\ge c_{\mathrm{imp}}
\ \Longrightarrow\
p_{\mathrm{sgn}}(v)>0
\quad(v\in[v_-,v_+]).
\]
Define
\begingroup\small
\[
\mathcal H_{\mathrm{imp}}
=
\{v\in[v_-,v_+]:|\rho(v)|\ge c_{\mathrm{imp}}\},
\qquad
m_{\mathrm{imp}}
=
\min_{v\in\mathcal H_{\mathrm{imp}}}p_{\mathrm{sgn}}(v),
\qquad
M_{\mathrm{imp}}
=
\max_{v\in[v_-,v_+]}|Q_{\mathrm{sgn}}(v)|.
\]
\endgroup
The set \(\mathcal H_{\mathrm{imp}}\) is compact and contains \(\mathcal K_{\mathrm{ext}}\). Therefore \(m_{\mathrm{imp}}>0\). Also \(M_{\mathrm{imp}}>0\), because \(Q_{\mathrm{sgn}}\) is nonzero at every extremum. Set
\[
t_{\mathrm{imp}}
=
\min\left\{
\frac{m_{\mathrm{imp}}}{M_{\mathrm{imp}}^2},
\frac{\mathcal E_{\mathrm{app}}-c_{\mathrm{imp}}}{2M_{\mathrm{imp}}}
\right\}>0.
\]
For \(v\in\mathcal H_{\mathrm{imp}}\),
\[
t_{\mathrm{imp}}Q_{\mathrm{sgn}}(v)^2
\le t_{\mathrm{imp}}M_{\mathrm{imp}}^2
\le m_{\mathrm{imp}}
\le p_{\mathrm{sgn}}(v)
<2p_{\mathrm{sgn}}(v).
\]
Consequently,
\[
\begin{aligned}
\bigl(\rho(v)-t_{\mathrm{imp}}Q_{\mathrm{sgn}}(v)\bigr)^2
&=\rho(v)^2-2t_{\mathrm{imp}}p_{\mathrm{sgn}}(v)
+t_{\mathrm{imp}}^2Q_{\mathrm{sgn}}(v)^2\\
&<\rho(v)^2
\le\mathcal E_{\mathrm{app}}^2.
\end{aligned}
\]
For \(v\in[v_-,v_+]\setminus\mathcal H_{\mathrm{imp}}\),
\[
\begin{aligned}
|\rho(v)-t_{\mathrm{imp}}Q_{\mathrm{sgn}}(v)|
&\le|\rho(v)|+t_{\mathrm{imp}}|Q_{\mathrm{sgn}}(v)|\\
&<c_{\mathrm{imp}}
+\frac{\mathcal E_{\mathrm{app}}-c_{\mathrm{imp}}}{2}
<\mathcal E_{\mathrm{app}}.
\end{aligned}
\]
The perturbed residual is continuous on the compact interval, so these pointwise strict inequalities imply
\[
\|\rho-t_{\mathrm{imp}}Q_{\mathrm{sgn}}\|_\infty
<\mathcal E_{\mathrm{app}}.
\]
But the polynomial
\[
Q_{\mathrm{imp}}=Q_*+t_{\mathrm{imp}}Q_{\mathrm{sgn}}
\]
has degree at most \(L\), and hence
\[
\mathcal E_{\mathrm{app}}
\le\|g-Q_{\mathrm{imp}}\|_\infty
=\|\rho-t_{\mathrm{imp}}Q_{\mathrm{sgn}}\|_\infty
<\mathcal E_{\mathrm{app}},
\]
a contradiction.

Thus, when \(\mathcal E_{\mathrm{app}}>0\), there exist increasing interval nodes \(v_1,\ldots,v_{L+2}\) with alternating residual values. Define their orientation by
\[
\eta=\frac{\rho(v_1)}{\mathcal E_{\mathrm{app}}}\in\{-1,1\}.
\]
Together with the zero-error construction, we have in both cases
\[
v_-\le v_1<\cdots<v_{L+2}\le v_+,
\qquad
\eta\in\{-1,1\},
\qquad
\rho(v_i)=\eta(-1)^{i-1}\mathcal E_{\mathrm{app}}.
\]
% lean: Causalean.Mathlib.Analysis.FinitePolynomialAlternationDuality.exists_strict_uniformImprovement
% lean: Causalean.Mathlib.Analysis.FinitePolynomialAlternationDuality.exists_equioscillationWitness

\item Define the nodal weights, their normalizing sum, and the normalized weights by
\[
\beta_i
=
\prod_{\substack{1\le j\le L+2\\j\ne i}}(v_i-v_j)^{-1},
\qquad
Z_{\mathrm{norm}}=\sum_{i=1}^{L+2}|\beta_i|,
\qquad
\widetilde w_i=\frac{\beta_i}{Z_{\mathrm{norm}}}
\quad(1\le i\le L+2).
\]
Distinctness of the nodes makes every factor nonzero, so \(Z_{\mathrm{norm}}>0\). Hence
\[
\sum_{i=1}^{L+2}|\widetilde w_i|
=
\frac{\sum_{i=1}^{L+2}|\beta_i|}{Z_{\mathrm{norm}}}
=1.
\]
There are exactly \(L+2-i\) negative factors in \(\beta_i\). Therefore
\[
\beta_i=(-1)^{L+2-i}|\beta_i|,
\qquad
\widetilde w_i=(-1)^{L+2-i}|\widetilde w_i|.
\]
% lean: Causalean.Mathlib.Analysis.FinitePolynomialAlternationDuality.sum_abs_normalizedLagrangeWeight_eq_one
% lean: Causalean.Mathlib.Analysis.FinitePolynomialAlternationDuality.normalizedLagrangeWeight_alternates

\item For every real polynomial \(Q\) of degree at most \(L\), Lagrange interpolation at these \(L+2\) nodes gives
\[
Q(v)
=
\sum_{i=1}^{L+2}
Q(v_i)
\prod_{\substack{1\le j\le L+2\\j\ne i}}
\frac{v-v_j}{v_i-v_j}
\quad(v\in\mathbb R).
\]
The coefficient of \(v^{L+1}\) on the left is zero, whereas on the right it is \(\sum_i\beta_iQ(v_i)\). Dividing by \(Z_{\mathrm{norm}}\) yields
\[
\sum_{i=1}^{L+2}\widetilde w_iQ(v_i)=0.
\]
In particular,
\[
\sum_{i=1}^{L+2}\widetilde w_i v_i^h=0
\quad(0\le h\le L),
\qquad
\sum_{i=1}^{L+2}\widetilde w_iQ_*(v_i)=0.
\]
% lean: Causalean.Mathlib.Analysis.FinitePolynomialAlternationDuality.sum_normalizedLagrangeWeight_mul_pow_eq_zero
% lean: Causalean.Mathlib.Analysis.FinitePolynomialAlternationDuality.sum_normalizedLagrangeWeight_mul_eval_eq_zero

\item Define the weighted target value by
\[
a_{\mathrm{dual}}
=
\sum_{i=1}^{L+2}\widetilde w_i g(v_i).
\]
Polynomial cancellation and the alternating residual identity give
\[
\begin{aligned}
a_{\mathrm{dual}}
&=
\sum_{i=1}^{L+2}\widetilde w_i\bigl(g(v_i)-Q_*(v_i)\bigr)
+\sum_{i=1}^{L+2}\widetilde w_iQ_*(v_i)\\
&=
\sum_{i=1}^{L+2}
(-1)^{L+2-i}|\widetilde w_i|\,
\eta(-1)^{i-1}\mathcal E_{\mathrm{app}}\\
&=
\eta(-1)^{L+1}\mathcal E_{\mathrm{app}}
\sum_{i=1}^{L+2}|\widetilde w_i|\\
&=\eta(-1)^{L+1}\mathcal E_{\mathrm{app}}.
\end{aligned}
\]
Since \(\eta\in\{-1,1\}\) and \(\mathcal E_{\mathrm{app}}\ge0\), this proves
\[
|a_{\mathrm{dual}}|=\mathcal E_{\mathrm{app}}.
\]
% lean: Causalean.Mathlib.Analysis.FinitePolynomialAlternationDuality.exists_alternationDualCertificate
% lean: Causalean.Mathlib.Analysis.FinitePolynomialAlternationDuality.AlternationDualCertificate.toFiniteMomentDual
% lean: Causalean.Mathlib.Analysis.FinitePolynomialAlternationDuality.exists_finiteMomentDual

\item Finally, choose the orientation of the weights by
\[
w_i=
\begin{cases}
\widetilde w_i,&a_{\mathrm{dual}}\ge0,\\
-\widetilde w_i,&a_{\mathrm{dual}}<0,
\end{cases}
\quad(1\le i\le L+2).
\]
If \(a_{\mathrm{dual}}\ge0\), the preceding identities directly give
\[
\sum_{i=1}^{L+2}|w_i|=1,
\qquad
\sum_{i=1}^{L+2}w_iv_i^h=0
\quad(0\le h\le L),
\qquad
\sum_{i=1}^{L+2}w_ig(v_i)
=a_{\mathrm{dual}}
=|a_{\mathrm{dual}}|
=\mathcal E_{\mathrm{app}}.
\]
If \(a_{\mathrm{dual}}<0\), negating every weight preserves its absolute value and reverses every weighted sum:
\[
\sum_{i=1}^{L+2}|w_i|
=\sum_{i=1}^{L+2}|-\widetilde w_i|=1,
\qquad
\sum_{i=1}^{L+2}w_iv_i^h
=-\sum_{i=1}^{L+2}\widetilde w_iv_i^h=0,
\]
and
\[
\sum_{i=1}^{L+2}w_ig(v_i)
=-a_{\mathrm{dual}}
=|a_{\mathrm{dual}}|
=\mathcal E_{\mathrm{app}}.
\]
Thus both cases establish the required weights, including the case of zero approximation error.

For the signed-measure formulation, define
\[
\sigma=\sum_{i=1}^{L+2}w_i\delta_{v_i}.
\]
Because the nodes are distinct, for every real-valued function \(f\) on the interval,
\[
\int f(v)\,d\sigma(v)=\sum_{i=1}^{L+2}w_if(v_i),
\qquad
\|\sigma\|_{\mathrm{TV}}=\sum_{i=1}^{L+2}|w_i|=1,
\qquad
\operatorname{supp}\sigma\subseteq\{v_1,\ldots,v_{L+2}\}.
\]
Applying the integral identity to \(f(v)=v^h\) and to \(f=g\) gives the asserted moments and target value.
% lean: CausalSmith.Stat.AnnotationRarearmFrontier.FinitePolynomialDuality
\end{enumerate}
\end{proof}
\endgroup

The finite support in \cref{obj:lem:finite-polynomial-duality} makes the certificate suitable for a finite latent distribution. The following specialization retains a fixed reciprocal gap while placing its scale at the interval length divided by the squared degree.

\begin{lemmav}{lem:shrinking-cone-dual}[Finite reciprocal-gap certificate]
Let the parameters satisfy:
\begin{itemize}
\item \textbf{(Degree.)} \(L\) is an integer with \(L\ge 2\).
\item \textbf{(Interval.)} \(B>0\).
\item \textbf{(Overlap.)} \(0<\epsilon\le 1/4\).
\end{itemize}
Set
\[
\alpha=\frac{B}{100L^2}.
\]
There exist strictly increasing nodes \(0\le v_0<\cdots<v_{L+1}\le B\) and real weights \(c_0,\ldots,c_{L+1}\) defining the finite signed measure
\[
\sigma=\sum_{i=0}^{L+1}c_i\delta_{v_i},
\]
where \(\delta_{v_i}\) is the unit point mass at \(v_i\), such that
\[
\|\sigma\|_{\mathrm{TV}}
=\sum_{i=0}^{L+1}|c_i|=1,
\qquad
\int v^h\,d\sigma(v)=0
\quad\text{for every integer }0\le h\le L,
\]
and
\[
\int\frac{\alpha}{\alpha+(1-\epsilon)v}\,d\sigma(v)
\ge\frac18.
\]
Here \(\|\sigma\|_{\mathrm{TV}}\) denotes the total mass of the variation measure of \(\sigma\).
\end{lemmav}

\begin{proof}[Proof of \cref{obj:lem:shrinking-cone-dual}]
\begin{enumerate}
\item Define the reciprocal target and its slope parameter by
\[
c_\epsilon:=1-\epsilon,
\qquad
g_{\mathrm{rec}}(v):=\frac{\alpha}{\alpha+c_\epsilon v}
\quad (v\in[0,B]).
\]
The hypotheses give
\[
\alpha>0,
\qquad
c_\epsilon\ge\frac34>0,
\qquad
\alpha+c_\epsilon v\ge\alpha>0
\quad (v\in[0,B]).
\]
Thus \(g_{\mathrm{rec}}\) is continuous on \([0,B]\). % lean: CausalSmith.Stat.AnnotationRarearmFrontier.cone_reciprocal_continuous

\item We first select a polynomial attaining the best uniform approximation error. For a continuous real function \(\varphi\) on a nonempty compact interval, write
\[
\|\varphi\|_{\infty,[a_-,a_+]}
:=
\sup_{v\in[a_-,a_+]}|\varphi(v)|.
\]
Define
\[
\mathcal V_{\mathrm{app}}
:=
\left\{
Q|_{[0,B]}:
Q\text{ is a real polynomial of degree at most }L
\right\},
\]
\[
\rho_{\mathrm{app}}
:=
2\|g_{\mathrm{rec}}\|_{\infty,[0,B]}+1,
\qquad
\mathcal K_{\mathrm{app}}
:=
\left\{
\varphi\in\mathcal V_{\mathrm{app}}:
\|\varphi\|_{\infty,[0,B]}\le\rho_{\mathrm{app}}
\right\}.
\]
The space \(\mathcal V_{\mathrm{app}}\) is finite dimensional, being spanned by the restrictions of \(1,v,\ldots,v^L\). Hence \(\mathcal K_{\mathrm{app}}\) is compact and contains zero. The continuous error functional therefore attains its minimum there at the restriction of a polynomial \(Q_{\mathrm{app}}\) of degree at most \(L\). Comparison with zero gives
\[
\|g_{\mathrm{rec}}-Q_{\mathrm{app}}\|_{\infty,[0,B]}
\le
\|g_{\mathrm{rec}}\|_{\infty,[0,B]}.
\]
For \(\varphi\in\mathcal K_{\mathrm{app}}\), minimality gives
\[
\|g_{\mathrm{rec}}-Q_{\mathrm{app}}\|_{\infty,[0,B]}
\le
\|g_{\mathrm{rec}}-\varphi\|_{\infty,[0,B]}.
\]
For \(\varphi\in\mathcal V_{\mathrm{app}}\setminus\mathcal K_{\mathrm{app}}\), the reverse triangle inequality gives
\[
\begin{aligned}
\|g_{\mathrm{rec}}-\varphi\|_{\infty,[0,B]}
&\ge
\|\varphi\|_{\infty,[0,B]}
-\|g_{\mathrm{rec}}\|_{\infty,[0,B]}\\
&>
\|g_{\mathrm{rec}}\|_{\infty,[0,B]}+1
\ge
\|g_{\mathrm{rec}}-Q_{\mathrm{app}}\|_{\infty,[0,B]}.
\end{aligned}
\]
Consequently, defining the best approximation error by
\[
E_{\mathrm{app}}
:=
\inf_{\substack{Q\text{ real polynomial}\\ \deg Q\le L}}
\|g_{\mathrm{rec}}-Q\|_{\infty,[0,B]},
\]
we have
\[
E_{\mathrm{app}}
=
\|g_{\mathrm{rec}}-Q_{\mathrm{app}}\|_{\infty,[0,B]}.
\]
% lean: Causalean.Mathlib.Analysis.FinitePolynomialAlternationDuality.exists_bestPolynomial

\item Suppose, for contradiction, that
\[
\|g_{\mathrm{rec}}-Q_{\mathrm{app}}\|_{\infty,[0,B]}<\frac18.
\]
Then, for every \(v\in[0,B]\),
\[
|g_{\mathrm{rec}}(v)-Q_{\mathrm{app}}(v)|<\frac18,
\qquad
0\le g_{\mathrm{rec}}(v)\le1.
\]
The triangle inequality therefore yields
\[
|Q_{\mathrm{app}}(v)|
\le
|Q_{\mathrm{app}}(v)-g_{\mathrm{rec}}(v)|
+g_{\mathrm{rec}}(v)
<\frac98,
\qquad
\|Q_{\mathrm{app}}\|_{\infty,[0,B]}\le\frac98.
\]
Define the comparison point by
\[
z_{\mathrm{rec}}:=\frac{\alpha}{c_\epsilon}.
\]
Since \(L^2\ge4\) and \(c_\epsilon\ge3/4\),
\[
0<z_{\mathrm{rec}}
=
\frac{B}{100L^2c_\epsilon}
\le\frac{B}{300}\le B.
\]
Direct evaluation gives
\[
g_{\mathrm{rec}}(0)=1,
\qquad
g_{\mathrm{rec}}(z_{\mathrm{rec}})=\frac12.
\]
The approximation bounds at these two points imply
\[
Q_{\mathrm{app}}(0)>\frac78,
\qquad
Q_{\mathrm{app}}(z_{\mathrm{rec}})<\frac58,
\qquad
Q_{\mathrm{app}}(0)-Q_{\mathrm{app}}(z_{\mathrm{rec}})>\frac14.
\]
By the mean value theorem, there exists \(v_{\mathrm{mvt}}\in(0,z_{\mathrm{rec}})\) such that
\[
Q_{\mathrm{app}}'(v_{\mathrm{mvt}})
=
\frac{Q_{\mathrm{app}}(z_{\mathrm{rec}})-Q_{\mathrm{app}}(0)}
{z_{\mathrm{rec}}},
\]
and hence
\[
|Q_{\mathrm{app}}(z_{\mathrm{rec}})-Q_{\mathrm{app}}(0)|
=
|Q_{\mathrm{app}}'(v_{\mathrm{mvt}})|\,z_{\mathrm{rec}}.
\]
% lean: CausalSmith.Stat.AnnotationRarearmFrontier.cone_reciprocal_polynomial_gap

\item To bound this derivative, define the polynomial pulled back to the unit interval by
\[
Q_{\mathrm{pull}}(x_{\mathrm{unit}})
:=
Q_{\mathrm{app}}\!\left(\frac B2x_{\mathrm{unit}}+\frac B2\right)
\quad (x_{\mathrm{unit}}\in\mathbb R).
\]
It has degree at most \(L\), and the affine map sends \([-1,1]\) onto \([0,B]\). Thus
\[
\|Q_{\mathrm{pull}}\|_{\infty,[-1,1]}
=
\|Q_{\mathrm{app}}\|_{\infty,[0,B]},
\qquad
Q_{\mathrm{pull}}'(x_{\mathrm{unit}})
=
\frac B2
Q_{\mathrm{app}}'\!\left(\frac B2x_{\mathrm{unit}}+\frac B2\right).
\]
Applying \cref{obj:lem:markov-polynomial} with bound
\(\|Q_{\mathrm{app}}\|_{\infty,[0,B]}\) gives
\[
|Q_{\mathrm{pull}}'(x_{\mathrm{unit}})|
\le
L^2\|Q_{\mathrm{app}}\|_{\infty,[0,B]}
\quad (x_{\mathrm{unit}}\in[-1,1]).
\]
For every \(v\in[0,B]\), the inverse coordinate
\[
x_{\mathrm{unit}}(v):=\frac{2v-B}{B}
\]
belongs to \([-1,1]\). Substituting it and dividing by \(B/2>0\) yields
\[
\|Q_{\mathrm{app}}'\|_{\infty,[0,B]}
\le
\frac{2L^2}{B}\|Q_{\mathrm{app}}\|_{\infty,[0,B]}
\le
\frac{9L^2}{4B}.
\]
% lean: Causalean.Mathlib.Analysis.FinitePolynomialAlternationDuality.markov_derivative_Icc

In particular,
\[
|Q_{\mathrm{app}}'(v_{\mathrm{mvt}})|\,z_{\mathrm{rec}}
\le
\frac{9L^2}{4B}\,z_{\mathrm{rec}}
=
\frac{9}{400c_\epsilon}
\le
\frac{3}{100}
<
\frac14.
\]
This contradicts the chain
\[
\frac14
<
Q_{\mathrm{app}}(0)-Q_{\mathrm{app}}(z_{\mathrm{rec}})
\le
|Q_{\mathrm{app}}(z_{\mathrm{rec}})-Q_{\mathrm{app}}(0)|
=
|Q_{\mathrm{app}}'(v_{\mathrm{mvt}})|\,z_{\mathrm{rec}}.
\]
Therefore
\[
\|g_{\mathrm{rec}}-Q_{\mathrm{app}}\|_{\infty,[0,B]}
\ge\frac18,
\qquad
E_{\mathrm{app}}\ge\frac18.
\]
% lean: CausalSmith.Stat.AnnotationRarearmFrontier.cone_reciprocal_polynomial_gap

\item Apply \cref{obj:lem:finite-polynomial-duality} to the continuous function \(g_{\mathrm{rec}}\) on the nondegenerate interval \([0,B]\). Indexing its nodes from zero, it supplies
\[
0\le v_0<\cdots<v_{L+1}\le B,
\qquad
c_0,\ldots,c_{L+1}\in\mathbb R,
\]
with
\[
\sum_{i=0}^{L+1}|c_i|=1,
\qquad
\sum_{i=0}^{L+1}c_i v_i^h=0
\quad (h=0,\ldots,L),
\qquad
\sum_{i=0}^{L+1}c_i g_{\mathrm{rec}}(v_i)=E_{\mathrm{app}}.
\]
Define
\[
\sigma:=\sum_{i=0}^{L+1}c_i\delta_{v_i}.
\]
Because the nodes are distinct, its total variation and moments satisfy
\[
\|\sigma\|_{\mathrm{TV}}
=
\sum_{i=0}^{L+1}|c_i|=1,
\qquad
\int v^h\,d\sigma(v)
=
\sum_{i=0}^{L+1}c_i v_i^h=0
\quad (h=0,\ldots,L).
\]
Finally,
\[
\int\frac{\alpha}{\alpha+(1-\epsilon)v}\,d\sigma(v)
=
\sum_{i=0}^{L+1}c_i g_{\mathrm{rec}}(v_i)
=
E_{\mathrm{app}}
\ge\frac18,
\]
which establishes the certificate. % lean: CausalSmith.Stat.AnnotationRarearmFrontier.shrinking_cone_dual
\end{enumerate}
\end{proof}

The certificate separates two mathematical roles: its vanishing moments support the observation comparison, and its reciprocal integral supplies the target separation. Both properties hold on the same finite support under the degree, interval, and overlap conditions in \cref{obj:lem:shrinking-cone-dual}.

The testing priors use a two-cell construction for the rare-label component and a certificate-based construction for the many-cell component. We first specify the two-cell tables and the variation measure needed to turn signed certificate weights into latent probabilities.

\begingroup\small
\begin{definitionv}{synth_9}[Rare-label probability tables]
For \(d\ge2\), \(0<\epsilon\le1/4\), and \(0<\delta\le1/8\), define \(P_1\) and \(P_0\) on \(\{1,\ldots,d\}\times\{0,1\}^2\) by
\[
\begin{aligned}
P_1(X=j,A=1,Y=1)&=\frac{\epsilon}{2}\left(\frac12+\delta\right),&
P_1(X=j,A=1,Y=0)&=\frac{\epsilon}{2}\left(\frac12-\delta\right),\\
P_0(X=j,A=1,Y=1)&=\frac{\epsilon}{2}\left(\frac12-\delta\right),&
P_0(X=j,A=1,Y=0)&=\frac{\epsilon}{2}\left(\frac12+\delta\right),
\end{aligned}
\qquad j\in\{1,2\}.
\]
For both tables, set
\[
P_1(X=j,A=0,Y=y)=P_0(X=j,A=0,Y=y)=\frac{1-\epsilon}{4},
\qquad j\in\{1,2\},\quad y\in\{0,1\},
\]
and assign zero probability to every record with \(j>2\).
Thus \(P_1\) is the positive-sign law and \(P_0\) the negative-sign law of \(\cref{obj:thm:rare-label-floor}\). In the rare-label branch of \(\cref{obj:def:prior-handle}\), use these tables with
\[
\delta=\frac{1}{16\sqrt{\max\{1,n\epsilon\}}},
\qquad
\Pi_1=\delta_{P_1},\qquad \Pi_0=\delta_{P_0},
\]
where \(\delta_P\) denotes the Dirac probability measure at the table \(P\).
\end{definitionv}
\endgroup

These tables have the same treatment--covariate marginal \leanref{sym:P_{XA}}{\ensuremath{P_{XA}}}, while their treated outcome means carry opposite perturbations. They supply the rare-label comparison in \cref{obj:thm:rare-label-floor}, including the benchmark with the exact marginal table supplied. For the many-cell construction, the variation measure assigns nonnegative mass to the certificate's atoms and leaves their signs available for the outcome assignment.

\begin{definitionv}{synth_6}[Total-variation measure]
For the selected finite signed measure \(\sigma\), let \(|\sigma|\) denote its total-variation measure. For every measurable set \(E\),
\[
|\sigma|(E)=\sup_{\{E_1,\ldots,E_r\}}\sum_{i=1}^{r}|\sigma(E_i)|,
\]
where the supremum ranges over all finite measurable partitions of \(E\). In particular, if \(\sigma=\sum_{i=1}^{r}c_i\delta_{v_i}\) with distinct atoms \(v_i\), then
\[
|\sigma|=\sum_{i=1}^{r}|c_i|\delta_{v_i}.
\]
\end{definitionv}

The construction below chooses its branch using the public information scales. In the affine branch, a common latent distribution determines the arm masses, a core branch carries the signed outcome assignment, and a reservoir supplies the remaining raw mass before normalization.

\begingroup\scriptsize\setlength{\arraycolsep}{2pt}%
\advance\hsize by 8pt\advance\linewidth by 8pt
\begin{algorithmv}{def:prior-handle}[Two testing prior construction]
The testing priors \(\Pi_1,\Pi_0\) are measures on complete-record probability tables, defined for public parameters \(n,m,d\in\mathbb N\) and \(\epsilon\in\mathbb R\) satisfying
\begin{itemize}
\item \textbf{(Complete-record budget.)} \(n\ge1\).
\item \textbf{(Alphabet size.)} \(d\ge2\).
\item \textbf{(Overlap range.)} \(0<\epsilon\le1/4\).
\end{itemize}
Their construction is as follows.
\begin{enumerate}
\item Put
\[
S=n\epsilon,\qquad
\ell=\log(\exp(1)+S),\qquad
N=n+m,\qquad
x=\frac d{N\epsilon\ell}.
\]
If \(S<\exp(4096)\) or \(x^2\le S^{-1}\), set
\[
\delta=\frac1{16\sqrt{\max\{1,S\}}}.
\]
For \(r\in\{0,1\}\), define the two-cell table by
\[
P_r(X=j,A=a,Y=y)=
\begin{cases}
\displaystyle
\frac12\epsilon^a(1-\epsilon)^{1-a}
\left(\frac12+(2r-1)a\delta\right)^y
\left(\frac12-(2r-1)a\delta\right)^{1-y},
&j\in\{1,2\},\\[4pt]
0,&3\le j\le d,
\end{cases}
\qquad a,y\in\{0,1\}.
\]
Output
\[
\Pi_1=\delta_{P_1},\qquad \Pi_0=\delta_{P_0}.
\]
Otherwise continue with the following steps.

\item Set
\[
\begin{aligned}
u&=128n,\qquad w=128N,\qquad M=u+w,\\
L&=\lceil8\ell\rceil,\qquad
B=\frac L{1000M},\qquad
\alpha=\frac B{100L^2},\\
b_0&=\frac\alpha\epsilon,\qquad
K_*=\min\left\{d-1,\left\lfloor\frac1{100b_0}\right\rfloor\right\}.
\end{aligned}
\]

\item Fix a witness supplied by \cref{obj:lem:shrinking-cone-dual}, represented as
\[
\sigma=\sum_{i=1}^{L+2}\omega_i\delta_{v_i},
\]
with defining properties
\[
\begin{gathered}
0\le v_1<\cdots<v_{L+2}\le B,\qquad
\sum_{i=1}^{L+2}|\omega_i|=1,\\
\sum_{i=1}^{L+2}\omega_i v_i^h=0
\quad(h=0,\ldots,L),\\
\sum_{i=1}^{L+2}\omega_i
\frac{\alpha}{\alpha+(1-\epsilon)v_i}\ge\frac18.
\end{gathered}
\]
At each atom of positive variation mass, define the polar sign
\[
h(v)=\frac{\sigma(\{v\})}{|\sigma|(\{v\})}\in\{-1,1\}.
\]

\item Define the raw covariate and arm masses by
\[
w_*(v)=b_0+v,\qquad
S_1(v)=\alpha+(1-\epsilon)v,\qquad
S_0(v)=(1-\epsilon)b_0+\epsilon v,
\qquad 0\le v\le B.
\]
For each rare cell \(i=1,\ldots,K_*\), independently draw a latent coordinate \(V_i\) and a branch flag \(c_i\in\{0,1\}\), with \(c_i=1\) denoting the core branch, according to
\[
\begin{aligned}
\Pr(V_i\in dv,\ c_i=1)
&=\frac{\alpha}{S_1(v)}|\sigma|(dv),\\
\Pr(V_i=0,\ c_i=0)
&=1-\int\frac{\alpha}{S_1(v)}|\sigma|(dv).
\end{aligned}
\]
The filler and core branches remain distinct when the core measure has an atom at zero. For \(r\in\{0,1\}\), set
\[
\mu_1^{(r)}(v,c)=
\begin{cases}
\bigl(1+(2r-1)h(v)\bigr)/2,&c=1,\\
0,&c=0,
\end{cases}
\qquad
\mu_0^{(r)}(v,c)=0.
\]
Both priors use raw arm masses \(S_1(v),S_0(v)\).

\item Let \(V\) have the common latent-coordinate marginal just defined, and put
\[
\bar w=\mathbb E[w_*(V)],\qquad
p_*=1-K_*\bar w,\qquad
Q=p_*+\sum_{i=1}^{K_*}w_*(V_i).
\]
Append a reservoir cell with raw covariate mass \(p_*\), propensity \(1/2\), and both outcome means zero. Set the remaining cells null.

\item For \(r\in\{0,1\}\), define the normalized complete-record table by
\[
P_r(X=j,A=a,Y=y)=
\begin{cases}
\displaystyle
\frac{S_a(V_j)}Q
\left[
y\mu_a^{(r)}(V_j,c_j)
+(1-y)\bigl(1-\mu_a^{(r)}(V_j,c_j)\bigr)
\right],
&1\le j\le K_*,\\[6pt]
\displaystyle\frac{p_*}{2Q}\mathbf1\{y=0\},
&j=K_*+1,\\[6pt]
0,&K_*+1<j\le d,
\end{cases}
\qquad a,y\in\{0,1\}.
\]
Output the pushforward priors
\[
\Pi_1=\operatorname{Law}(P_1),\qquad
\Pi_0=\operatorname{Law}(P_0)
\]
of the common finite product latent distribution under the respective table maps.
\end{enumerate}
\end{algorithmv}
\endgroup

The testing priors \leanref{sym:\Pi_a}{\ensuremath{\Pi_1,\Pi_0}} in \cref{obj:def:prior-handle} are distributions over normalized probability tables. In the affine branch, the reservoir mass \(p_*\) complements the expected total rare-cell mass, while \(Q\) is the realized total raw mass. Normalizing by \(Q\) therefore retains the variation in rare-cell masses across latent draws. The core and filler flags specify distinct outcome assignments even when they share a latent coordinate at zero.

The reciprocal certificate supplies the separation ingredient for the covariate-weighted target in \cref{obj:def:target}. Its moment conditions supply the approximation ingredient for comparing observations under the two hypotheses. The next result assembles these roles under the large-information and affine-regime conditions, including the overlap validity of every normalized table.

\begin{lemmav}{lem:normalized-affine-testing}[Normalized affine testing priors]
There exists a numerical constant \(c>0\) such that the following holds for every \(n,m,d\in\mathbb N\) and \(\epsilon\in\mathbb R\). Put
\[
S=n\epsilon,\qquad N=n+m,\qquad
\ell=\log(\exp(1)+S),\qquad x=\frac{d}{N\epsilon\ell},
\]
where \(S\) is the rare-arm information scale, \(N\) is the total marginal-information sample size, \(\ell\) is the regularized logarithmic scale, and \(x\) is the normalized many-cell difficulty. Suppose that:
\begin{itemize}
\item \textbf{(Public indices.)} \(n\ge1\) and \(d\ge2\).
\item \textbf{(Overlap parameter.)} \(0<\epsilon\le1/4\).
\item \textbf{(Information scale.)} \(S\ge\exp(4096)\).
\item \textbf{(Affine regime.)} \(x^2>S^{-1}\).
\end{itemize}
For every signed certificate \(\sigma\) satisfying the certificate conditions at the affine tuning parameters in \cref{obj:def:prior-handle}, and every latent configuration in that construction, the normalized laws under both hypotheses belong to \(\mathcal M_{d,\epsilon}\) of \cref{obj:def:model}, and their treatment--covariate marginal tables \(P_{XA}\) agree. For each such \(\sigma\), the two product priors obtained by pushing the common product latent distribution through the normalized-law construction also induce exactly the same distribution of \(P_{XA}\).

The selected priors \(\Pi_1\) and \(\Pi_0\) in \cref{obj:def:prior-handle} induce exactly the same distribution of \(P_{XA}\), and
\[
R(n,m,d,\epsilon)\ge c\min\{1,x^2\}.
\]
Here \(R(n,m,d,\epsilon)\) is the minimax squared-error risk in \cref{obj:def:risk}, with the infimum over all clipped Borel randomized rules in the original fixed-size two-channel experiment.
\end{lemmav}

\begin{proof}[Proof of \cref{obj:lem:normalized-affine-testing}]
Set
\[
c=\frac{10^{-20}}{128}>0,
\]
and fix public parameters satisfying the hypotheses. Throughout the proof, the quantities
\[
S,\ N,\ \ell,\ x,\ u,\ w,\ M,\ L,\ B,\ \alpha,\ b_0,\ K_*
\]
have exactly the values in the affine branch of \cref{obj:def:prior-handle}.

\begin{enumerate}
\item
We first establish support and the common marginal for an arbitrary admissible certificate \(\sigma\). The tuning formulas give
\[
128N\le M\le256N,\qquad
8\ell\le L\le9\ell,\qquad
B,\alpha,b_0>0,\qquad K_*b_0\le\frac1{100}.
\]
Indeed, \(M=128(n+N)\) and \(n\le N\). Also \(\ell>1\), so
\[
8\ell\le\lceil8\ell\rceil<8\ell+1\le9\ell.
\]
The last bound follows by multiplying
\[
K_*\le\left\lfloor\frac1{100b_0}\right\rfloor
\le\frac1{100b_0}
\]
by \(b_0\). In particular, \(L\ge2\), and \(K_*\le d-1\).

Write the certificate and its variation measure as
\[
\sigma=\sum_{i=1}^{L+2}\omega_i\delta_{v_i},
\qquad
\eta=|\sigma|=\sum_{i=1}^{L+2}|\omega_i|\delta_{v_i}.
\]
For configurations involving a zero-weight node, extend the polar sign by
\[
h(v_i)=
\begin{cases}
1,&\omega_i\ge0,\\
-1,&\omega_i<0.
\end{cases}
\]
Thus \(|\omega_i|h(v_i)=\omega_i\), including when \(\omega_i=0\).

The one-coordinate latent space and its probability law are
\[
\Lambda=\{(0,0)\}\cup\{(v_i,1):1\le i\le L+2\},
\]
\[
\begin{aligned}
\nu_{\mathrm{cell}}\{(v_i,1)\}
&=\frac{\alpha}{S_1(v_i)}|\omega_i|,\\
\nu_{\mathrm{cell}}\{(0,0)\}
&=1-\sum_{i=1}^{L+2}\frac{\alpha}{S_1(v_i)}|\omega_i|.
\end{aligned}
\]
The branch flag distinguishes the filler state \((0,0)\) from a core state at the node zero. Since \(S_1(v)\ge\alpha>0\),
\[
0\le\sum_i\frac{\alpha}{S_1(v_i)}|\omega_i|
\le\sum_i|\omega_i|=1,
\]
so these masses are nonnegative and sum to one.

For a coordinate \((V,c_i)\) with this law,
\[
\mathbb EV
=\int v\frac{\alpha}{S_1(v)}\,d\eta(v)
\le\frac{\alpha}{1-\epsilon},
\]
because \(v\alpha/S_1(v)\le\alpha/(1-\epsilon)\). Consequently,
\[
\bar w=b_0+\mathbb EV
\le b_0+\frac{\alpha}{1-\epsilon}
\le\frac43b_0,
\]
where \(\alpha=\epsilon b_0\) and \(\epsilon\le1/4\). Hence, for every latent configuration,
\[
p_*=1-K_*\bar w
\ge1-\frac1{75}>\frac12,
\qquad
Q=p_*+\sum_{j=1}^{K_*}(b_0+V_j)>\frac12.
\]
These assertions also cover an empty rare-cell set, in which case \(p_*=Q=1\).

For every \(v\in[0,B]\), the raw arm masses satisfy
\[
\begin{aligned}
S_1(v)+S_0(v)&=b_0+v=w_*(v),\\
S_1(v)-\epsilon w_*(v)&=(1-2\epsilon)v\ge0,\\
(1-\epsilon)w_*(v)-S_1(v)&=(1-2\epsilon)b_0\ge0.
\end{aligned}
\]
Both raw arms are positive. The assigned conditional means belong to \([0,1]\), and summing the raw complete-record table over all atoms gives
\[
p_*+\sum_{j=1}^{K_*}w_*(V_j)=Q.
\]
Division by \(Q\) therefore produces probability tables. Their occupied rare cells satisfy the overlap inequalities above; the reservoir has propensity \(1/2\); and the remaining cells have mass zero. Thus
\[
P_\iota(\boldsymbol\lambda)\in\mathcal M_{d,\epsilon},
\qquad
\iota\in\{0,1\},\quad
\boldsymbol\lambda=((V_j,c_j))_{j=1}^{K_*}\in\Lambda^{K_*},
\]
by \cref{obj:def:model}.

Summing out the outcome coordinate gives, under either hypothesis,
\[
(P_\iota(\boldsymbol\lambda))_{XA}(j,a)=
\begin{cases}
S_a(V_j)/Q,&1\le j\le K_*,\\
p_*/(2Q),&j=K_*+1,\\
0,&j>K_*+1.
\end{cases}
\]
In particular,
\[
(P_1(\boldsymbol\lambda))_{XA}
=(P_0(\boldsymbol\lambda))_{XA}
\quad\text{for every }\boldsymbol\lambda.
\]
Define the common product latent law and its pushforward priors by
\[
\nu_{\mathrm{lat}}=\nu_{\mathrm{cell}}^{\otimes K_*},
\qquad
\Pi_\iota^\sigma=\operatorname{Law}_{\nu_{\mathrm{lat}}}
\bigl(P_\iota(\boldsymbol\lambda)\bigr).
\]
Pushing the pointwise equality through this common latent law yields
\[
\operatorname{Law}_{\Pi_1^\sigma}(P_{XA})
=\operatorname{Law}_{\Pi_0^\sigma}(P_{XA}).
\]
% lean: CausalSmith.Stat.AnnotationRarearmFrontier.affine_prior_support_and_common_marginal

\item
The assumed inequalities
\[
S\ge\exp(4096),\qquad x^2>S^{-1}
\]
select the affine branch of \cref{obj:def:prior-handle}. The preceding equality applies to its selected certificate, and hence
\[
\operatorname{Law}_{\Pi_1}(P_{XA})
=\operatorname{Law}_{\Pi_0}(P_{XA}).
\]
For the risk bound, fix an admissible certificate. Such a certificate exists because \(L\ge2\), \(B>0\), and \(0<\epsilon\le1/4\): \cref{obj:lem:shrinking-cone-dual} supplies the strictly ordered nodes, variation mass one, vanishing moments through degree \(L\), and reciprocal gap
\[
\int\frac{\alpha}{\alpha+(1-\epsilon)v}\,d\sigma(v)\ge\frac18.
\]
% lean: CausalSmith.Stat.AnnotationRarearmFrontier.affine_selectedPrior_common_marginal
% lean: CausalSmith.Stat.AnnotationRarearmFrontier.affine_schedule_dual

\item
Define the raw targets, their means, their gap, and the normalized targets by
\[
\begin{aligned}
Z_\iota^{\mathrm{tar}}(\boldsymbol\lambda)
&=\sum_{j=1}^{K_*}w_*(V_j)\mu_1^{(\iota)}(V_j,c_j),\\
t_\iota^{\mathrm{tar}}
&=\int Z_\iota^{\mathrm{tar}}\,d\nu_{\mathrm{lat}},\\
\Delta_{\mathrm{tar}}&=t_1^{\mathrm{tar}}-t_0^{\mathrm{tar}},\\
\tau_\iota(\boldsymbol\lambda)
&=\tau(P_\iota(\boldsymbol\lambda))
=\frac{Z_\iota^{\mathrm{tar}}(\boldsymbol\lambda)}{Q(\boldsymbol\lambda)}.
\end{aligned}
\]
The last identity follows from \cref{obj:def:target}: the control means and the reservoir outcome means are zero. Pointwise,
\[
0\le Z_\iota^{\mathrm{tar}}
\le\sum_{j=1}^{K_*}w_*(V_j)\le Q,
\qquad
0\le\tau_\iota\le1.
\]

The difference of the one-cell target means is
\[
\mathbb E_{\nu_{\mathrm{cell}}}
\left[w_*(V)\{\mu_1^{(1)}(V,c_i)-\mu_1^{(0)}(V,c_i)\}\right]
=\int w_*(v)\frac{\alpha}{S_1(v)}\,d\sigma(v).
\]
Here the filler contributes zero, while multiplying the core probability by the polar sign restores the signed weight. The identity
\[
w_*(v)\frac{\alpha}{S_1(v)}
=\frac{\alpha}{1-\epsilon}
+\frac{b_0(1-2\epsilon)}{1-\epsilon}
  \frac{\alpha}{S_1(v)}
\]
and the zero-th moment of \(\sigma\) give
\[
\int w_*(v)\frac{\alpha}{S_1(v)}\,d\sigma(v)
=\frac{b_0(1-2\epsilon)}{1-\epsilon}
 \int\frac{\alpha}{S_1(v)}\,d\sigma(v)
\ge\frac{b_0}{12},
\]
since \((1-2\epsilon)/(1-\epsilon)\ge2/3\). Additivity of expectation now gives
\[
\Delta_{\mathrm{tar}}\ge\frac{K_*b_0}{12}.
\]

We next calibrate this gap in terms of \(x\). The tuning formulas imply
\[
b_0=\frac1{100000M\epsilon L}
\ge\frac1{230400000N\epsilon\ell}.
\]
As \(S\ge1\), we have \(M\epsilon\ge128S\) and \(L\ge8\), so \(b_0\le1/100\). Thus
\[
\left\lfloor\frac1{100b_0}\right\rfloor\ge1.
\]
The defining floor inequality then gives
\[
\frac1{100b_0}
<
\left\lfloor\frac1{100b_0}\right\rfloor+1
\le2\left\lfloor\frac1{100b_0}\right\rfloor,
\qquad
\left\lfloor\frac1{100b_0}\right\rfloor b_0\ge\frac1{200}.
\]
Together with \(d-1\ge d/2\), this yields
\[
(d-1)b_0\ge\frac{x}{460800000}.
\]
If the alphabet cap determines \(K_*\), this is the required mass bound. If the floor cap determines \(K_*\), the bound \(1/200\) applies. Therefore
\[
K_*b_0\ge\frac{\min\{x,1\}}{460800000},
\qquad
\Delta_{\mathrm{tar}}\ge10^{-10}\min\{x,1\}>0.
\]
In particular, \(K_*\ge1\) in the present regime.
% lean: CausalSmith.Stat.AnnotationRarearmFrontier.affine_raw_target_mean_gap_uniform

\item
Consider the following raw Poisson experiment. Conditional on
\(\boldsymbol\lambda\) and hypothesis \(\iota\), the complete-record atom counts and auxiliary atom counts are mutually independent, with respective means
\[
uQ\,P_\iota(\boldsymbol\lambda)(j,a,y),
\qquad
wQ\,(P_\iota(\boldsymbol\lambda))_{XA}(j,a).
\]
We use the Poisson convention
\[
\operatorname{Poi}(\lambda_{\mathrm P})\{k\}
=e^{-\lambda_{\mathrm P}}\frac{\lambda_{\mathrm P}^k}{k!},
\qquad \lambda_{\mathrm P}\ge0,\quad k\in\mathbb N,
\]
with \(0^0=1\).

In one rare cell, combine the two control-channel counts. The resulting four-count mixture is
\[
\begin{aligned}
\nu_\iota^{\mathrm{cell}}
=\int_{\Lambda}
&\operatorname{Poi}\bigl(uS_1(v)\mu_1^{(\iota)}(v,c_i)\bigr)
\otimes
\operatorname{Poi}\bigl(uS_1(v)(1-\mu_1^{(\iota)}(v,c_i))\bigr)\\
&\otimes\operatorname{Poi}\bigl(wS_1(v)\bigr)
\otimes\operatorname{Poi}\bigl(MS_0(v)\bigr)
\,d\nu_{\mathrm{cell}}(v,c_i).
\end{aligned}
\]
The first two coordinates count treated complete outcomes equal to one and zero, respectively.

When both complete treated counts are zero, the two mixture probabilities agree. When both are positive, both mixture probabilities vanish. For
\(\rho,\gamma,\kappa\in\mathbb N\), define the signed discrepancy
\[
D_{\rho,\gamma,\kappa}^{\mathrm{cnt}}
=
\int e^{-M(b_0+v)}
\frac{(uS_1(v))^{\rho+1}(wS_1(v))^\gamma(MS_0(v))^\kappa}
     {(\rho+1)!\gamma!\kappa!}
\frac{\alpha}{S_1(v)}\,d\sigma(v).
\]
Expanding the conditional Poisson probabilities and the core mixture gives
\[
\begin{aligned}
\nu_1^{\mathrm{cell}}\{(\rho+1,0,\gamma,\kappa)\}
-\nu_0^{\mathrm{cell}}\{(\rho+1,0,\gamma,\kappa)\}
&=D_{\rho,\gamma,\kappa}^{\mathrm{cnt}},\\
\nu_1^{\mathrm{cell}}\{(0,\rho+1,\gamma,\kappa)\}
-\nu_0^{\mathrm{cell}}\{(0,\rho+1,\gamma,\kappa)\}
&=-D_{\rho,\gamma,\kappa}^{\mathrm{cnt}}.
\end{aligned}
\]
Indeed, the filler has treated mean zero under both hypotheses. On the core branch, the difference between the two deterministic outcome assignments is the polar sign, and
\(|\omega_i|h(v_i)=\omega_i\). These are all the possible nonzero discrepancies.
% lean: CausalSmith.Stat.AnnotationRarearmFrontier.affineCellPoissonPredictive_equation_six

\item
We bound the sum of these discrepancies by moment cancellation. Define
\[
\begin{aligned}
\varphi_{\rho,\gamma,\kappa}(v)
&=S_1(v)^{\rho+\gamma}S_0(v)^\kappa,\\
\beta_{\rho,\gamma,\kappa}^{\mathrm{cnt}}
&=e^{-Mb_0}
  \frac{u^\rho w^\gamma M^\kappa}{(\rho+1)!\gamma!\kappa!},\\
a_j^{\mathrm{cnt}}
&=\sum_{\rho,\gamma,\kappa\ge0}
  \beta_{\rho,\gamma,\kappa}^{\mathrm{cnt}}
  [v^j]\varphi_{\rho,\gamma,\kappa}(v),
\qquad j\in\mathbb N.
\end{aligned}
\]
Here \([v^j]\varphi(v)=\varphi^{(j)}(0)/j!\) denotes the coefficient of \(v^j\). Factoring one treated intensity gives
\[
D_{\rho,\gamma,\kappa}^{\mathrm{cnt}}
=u\alpha\,\beta_{\rho,\gamma,\kappa}^{\mathrm{cnt}}
\int e^{-Mv}\varphi_{\rho,\gamma,\kappa}(v)\,d\sigma(v).
\]
Every coefficient of \(\varphi_{\rho,\gamma,\kappa}\) is nonnegative.

For \(j\le L\), the matched moments imply
\[
\int e^{-Mv}v^j\,d\sigma(v)
=
\int\left(e^{-Mv}
-\sum_{k=0}^{L-j}\frac{(-Mv)^k}{k!}\right)v^j\,d\sigma(v).
\]
Bounding the exponential remainder on \([0,B]\), and using
\(\|\sigma\|_{\mathrm{TV}}=1\), gives
\[
\left|\int e^{-Mv}v^j\,d\sigma(v)\right|
\le B^j\sum_{\substack{k\ge0\\k+j>L}}\frac{(MB)^k}{k!}.
\]
For \(j>L\), the full absolute exponential series gives the same inequality, with every \(k\ge0\) included. Expanding the polynomial therefore yields
\[
\left|\int e^{-Mv}\varphi_{\rho,\gamma,\kappa}(v)\,d\sigma(v)\right|
\le
\sum_{j\ge0}[v^j]\varphi_{\rho,\gamma,\kappa}(v)\,
B^j\sum_{\substack{k\ge0\\k+j>L}}\frac{(MB)^k}{k!}.
\]

To sum over the counts, use \(1/(\rho+1)!\le1/\rho!\). Coefficientwise, the nonnegative count series obey
\[
\begin{aligned}
\sum_{\rho\ge0}\frac{u^\rho S_1(v)^\rho}{(\rho+1)!}
&\preceq e^{uS_1(v)},\\
\sum_{\gamma\ge0}\frac{w^\gamma S_1(v)^\gamma}{\gamma!}
&=e^{wS_1(v)},\\
\sum_{\kappa\ge0}\frac{M^\kappa S_0(v)^\kappa}{\kappa!}
&=e^{MS_0(v)},
\end{aligned}
\]
where \(\preceq\) means that each coefficient on the left is at most the corresponding coefficient on the right. Thus
\[
0\le a_j^{\mathrm{cnt}}
\le e^{-Mb_0}[v^j]e^{M(S_1(v)+S_0(v))}
=[v^j]e^{Mv}
=\frac{M^j}{j!}.
\]
This calculation cancels the common baseline using
\(S_1(v)+S_0(v)=b_0+v\).
The dominating series is finite:
\[
\sum_{j\ge0}\frac{M^jB^j}{j!}
\sum_{k\ge0}\frac{(MB)^k}{k!}
=e^{2MB}.
\]
Consequently the nonnegative coefficient sums can be interchanged, and the signed discrepancies are absolutely summable. We obtain
\[
\begin{aligned}
\sum_{\rho,\gamma,\kappa\ge0}
|D_{\rho,\gamma,\kappa}^{\mathrm{cnt}}|
&\le
u\alpha\sum_{j\ge0}a_j^{\mathrm{cnt}}B^j
\sum_{\substack{k\ge0\\k+j>L}}\frac{(MB)^k}{k!}\\
&\le
u\alpha\sum_{\substack{j,k\ge0\\j+k>L}}
\frac{(MB)^{j+k}}{j!k!}\\
&=
u\alpha\sum_{k_{\mathrm{tail}}>L}\frac{(2MB)^{k_{\mathrm{tail}}}}{k_{\mathrm{tail}}!}.
\end{aligned}
\]
The last equality uses
\[
\sum_{j=0}^{k_{\mathrm{tail}}}\frac1{j!(k_{\mathrm{tail}}-j)!}=\frac{2^{k_{\mathrm{tail}}}}{k_{\mathrm{tail}}!}.
\]
% lean: CausalSmith.Stat.AnnotationRarearmFrontier.affine_signed_raw_count_tail_bound

\item
For probability measures, use the total variation convention
\[
\operatorname{TV}(\nu,\nu')
=\sup_{\mathscr E\ \mathrm{measurable}}
|\nu(\mathscr E)-\nu'(\mathscr E)|.
\]
On a countable space this equals one half of the sum of absolute singleton discrepancies. The two disjoint patterns in the preceding calculation therefore give
\[
\operatorname{TV}(\nu_1^{\mathrm{cell}},\nu_0^{\mathrm{cell}})
\le u\alpha\sum_{k_{\mathrm{tail}}>L}\frac{(2MB)^{k_{\mathrm{tail}}}}{k_{\mathrm{tail}}!}.
\]

The tuning gives \(2MB=L/500\). For an integer \(k_{\mathrm{tail}}>L\),
\[
\frac{(L/500)^{k_{\mathrm{tail}}}}{k_{\mathrm{tail}}!}
\le\frac{(k_{\mathrm{tail}}/500)^{k_{\mathrm{tail}}}}{k_{\mathrm{tail}}!}
\le\left(\frac{e}{500}\right)^{k_{\mathrm{tail}}}
\le e^{-5k_{\mathrm{tail}}}.
\]
The middle inequality follows by retaining the \(k_{\mathrm{tail}}\)-th term of the positive series for \(e^{k_{\mathrm{tail}}}\). For the last inequality, \(e\le11/4\) gives
\(e^6\le(11/4)^6<500\). Since \(e^{-5}\le1/2\),
\[
\sum_{k_{\mathrm{tail}}>L}\frac{(L/500)^{k_{\mathrm{tail}}}}{k_{\mathrm{tail}}!}
\le\frac{e^{-5(L+1)}}{1-e^{-5}}
\le2e^{-5L}.
\]
Moreover,
\[
K_*u\alpha=128S(K_*b_0)\le\frac{32}{25}S,
\]
and
\[
Se^{-5L}
\le e^{\ell-5L}
\le e^{-39\ell}
\le e^{-8}\le\frac1{256}.
\]
Here \(S\le e^\ell\), \(L\ge8\ell\), \(\ell\ge1\), and
\(e^8\ge2^8=256\). Hence
\[
K_*\operatorname{TV}(\nu_1^{\mathrm{cell}},\nu_0^{\mathrm{cell}})
\le K_*u\alpha\sum_{k_{\mathrm{tail}}>L}\frac{(2MB)^{k_{\mathrm{tail}}}}{k_{\mathrm{tail}}!}
\le\frac1{100}<\frac1{64}.
\]

We now restore the full observation channels. Given a combined control count \(\kappa\), its complete-channel count has conditional distribution
\[
\Pr(\kappa_{\mathrm L}=k\mid\kappa)
=\binom{\kappa}{k}\left(\frac uM\right)^k
                    \left(\frac wM\right)^{\kappa-k},
\qquad 0\le k\le\kappa.
\]
This follows directly from
\[
\begin{aligned}
&\operatorname{Poi}(uS_0(v))\{k\}
 \operatorname{Poi}(wS_0(v))\{\kappa-k\}\\
&\hspace{1cm}=
\operatorname{Poi}(MS_0(v))\{\kappa\}
\binom{\kappa}{k}\left(\frac uM\right)^k
                    \left(\frac wM\right)^{\kappa-k}.
\end{aligned}
\]
The splitting law is independent of the latent coordinate and hypothesis. Applying it cannot increase total variation. Independence of the latent coordinates makes the rare-cell mixtures a product, whose distance is bounded by the sum of the cell distances; this follows by telescoping the difference of the product mass functions. The reservoir counts form a common independent factor, and the remaining cells have zero counts. Consequently the full two-channel histogram mixtures have distance below \(1/64\).

Uniformly ordering each channel's histogram is another common probability kernel. Its conditional law is the ordered Poisson experiment from the normalized table. To see this, for a probability mass function \(p_{\mathrm{mark}}\) on a finite alphabet \(\mathscr Z\), put
\[
\kappa_{\mathrm{tot}}=\sum_{z\in\mathscr Z}\kappa_z.
\]
The independent-count probabilities factor as
\[
\prod_{z\in\mathscr Z}
e^{-\lambda_{\mathrm P}p_{\mathrm{mark}}(z)}
\frac{(\lambda_{\mathrm P}p_{\mathrm{mark}}(z))^{\kappa_z}}{\kappa_z!}
=
e^{-\lambda_{\mathrm P}}
\frac{\lambda_{\mathrm P}^{\kappa_{\mathrm{tot}}}}{\kappa_{\mathrm{tot}}!}
\frac{\kappa_{\mathrm{tot}}!}{\prod_z\kappa_z!}
\prod_zp_{\mathrm{mark}}(z)^{\kappa_z}.
\]
Thus, conditional on the latent table, the ordered channels have the following law:
\[
\begin{aligned}
J_{\mathrm L}&\sim\operatorname{Poi}(uQ),&
\mathcal L_{\mathrm{raw}}\mid J_{\mathrm L}=k
&\sim P_\iota(\boldsymbol\lambda)^{\otimes k},\\
J_{\mathrm A}&\sim\operatorname{Poi}(wQ),&
\mathcal X_{\mathrm{raw}}\mid J_{\mathrm A}=k
&\sim(P_\iota(\boldsymbol\lambda))_{XA}^{\otimes k}.
\end{aligned}
\]
Both channels are independent conditional on the latent table. Their domains are
\[
\mathcal L_{\mathrm{raw}}\in
\bigsqcup_{k\ge0}
\bigl(\{1,\ldots,d\}\times\{0,1\}^2\bigr)^k,
\qquad
\mathcal X_{\mathrm{raw}}\in
\bigsqcup_{k\ge0}
\bigl(\{1,\ldots,d\}\times\{0,1\}\bigr)^k,
\]
with the empty tuple at length zero.

Append the independent seed \(U\sim\mathrm{Uniform}[0,1]\), and define
\[
\mathcal Q_\iota^{\mathrm{raw}}
=\operatorname{Law}(\mathcal L_{\mathrm{raw}},
                    \mathcal X_{\mathrm{raw}},U\mid\iota),
\]
where the latent table is mixed according to \(\nu_{\mathrm{lat}}\). Common kernels and a common independent seed preserve the distance bound, so
\[
\operatorname{TV}(\mathcal Q_1^{\mathrm{raw}},
                  \mathcal Q_0^{\mathrm{raw}})<\frac1{64}.
\]
% lean: CausalSmith.Stat.AnnotationRarearmFrontier.affineOrderedRawKernel_randomized_tv_small

\item
We next bound the normalized target's deviation from its raw mean. Since \(0\le V\le B\),
\[
\mathbb EV^2\le B\mathbb EV
\le\frac{B\alpha}{1-\epsilon}
\le\frac43B\alpha.
\]
For either hypothesis,
\[
0\le w_*(V)\mu_1^{(\iota)}(V,c_i)\le w_*(V),
\qquad
\mathbb Ew_*(V)^2
\le2b_0^2+2\mathbb EV^2
\le2b_0^2+\frac83B\alpha.
\]
Independence and the inequality
\(\operatorname{Var}(Z)\le\mathbb EZ^2\) consequently give
\[
\operatorname{Var}(Z_\iota^{\mathrm{tar}})
\le K_*\left(2b_0^2+\frac83B\alpha\right),
\qquad
\mathbb EQ=1,
\qquad
\operatorname{Var}(Q)\le\frac43K_*B\alpha.
\]
For the last two identities, use \(p_*=1-K_*\bar w\) and
\(\operatorname{Var}(b_0+V)=\operatorname{Var}(V)\).

The normalization identity is
\[
\tau_\iota-t_\iota^{\mathrm{tar}}
=(Z_\iota^{\mathrm{tar}}-t_\iota^{\mathrm{tar}})
+\tau_\iota(1-Q).
\]
Since \(0\le\tau_\iota\le1\), squaring and integrating yields
\[
\begin{aligned}
\mathbb E_{\nu_{\mathrm{lat}}}
(\tau_\iota-t_\iota^{\mathrm{tar}})^2
&\le2\operatorname{Var}(Z_\iota^{\mathrm{tar}})
   +2\operatorname{Var}(Q)\\
&\le4K_*b_0^2+8K_*B\alpha.
\end{aligned}
\]
Markov's inequality at the positive radius \(\Delta_{\mathrm{tar}}/8\), together with
\(\Delta_{\mathrm{tar}}\ge K_*b_0/12\), gives
\[
\begin{aligned}
\nu_{\mathrm{lat}}\left\{
|\tau_\iota-t_\iota^{\mathrm{tar}}|
\ge\frac{\Delta_{\mathrm{tar}}}{8}\right\}
&\le
\frac{64(4K_*b_0^2+8K_*B\alpha)}
     {\Delta_{\mathrm{tar}}^2}\\
&\le\frac{36864+7372800\epsilon^2L^2}{K_*}\\
&\le\frac{2^{23}(1+\epsilon^2L^2)}{K_*}.
\end{aligned}
\]
The middle estimate uses the tuning identity
\[
B\alpha=100\epsilon^2L^2b_0^2.
\]
% lean: CausalSmith.Stat.AnnotationRarearmFrontier.affine_normalized_target_gap_tail

\item
To calibrate this tail bound, define the ordinary logarithmic scale
\[
\ell_{\mathrm{plain}}=\log S.
\]
The information-scale assumption implies
\[
\ell_{\mathrm{plain}}\ge4096,\qquad
\ell\le\ell_{\mathrm{plain}}+1,\qquad
L\le10\ell_{\mathrm{plain}}.
\]
For the second inequality, \(S\ge2\) gives \(e+S\le eS\); the third follows from \(L\le9\ell\).

The floor-mass bound already established implies
\[
\left\lfloor\frac1{100b_0}\right\rfloor
\ge\frac1{200b_0}=500M\epsilon L
\ge64000SL.
\]
Also \(x^2>S^{-1}\) and \(x>0\) give
\[
d>\frac{N\epsilon\ell}{\sqrt S}\ge\sqrt S\,\ell.
\]
Thus the two caps defining \(K_*\) satisfy
\[
d-1\ge\frac d2>\frac{\sqrt S\,\ell}{2}
\ge\frac{\sqrt S\,L}{18},
\qquad
\left\lfloor\frac1{100b_0}\right\rfloor
\ge64000SL\ge\frac{\sqrt S\,L}{18},
\]
and hence
\[
K_*\ge\frac{\sqrt S\,L}{18}.
\]

The required exponential comparison follows from
\[
2^{50}\le e^{\ell_{\mathrm{plain}}/4},
\qquad
\frac{\ell_{\mathrm{plain}}}{4}
\le e^{\ell_{\mathrm{plain}}/4}.
\]
The first uses \(e>2\) and
\(50\le\ell_{\mathrm{plain}}/4\); the second follows from
\(1+t\le e^t\). Multiplying gives
\[
e^{\ell_{\mathrm{plain}}/2}
\ge2^{48}\ell_{\mathrm{plain}}
\ge180\cdot2^{30}\ell_{\mathrm{plain}}.
\]
Since \(\sqrt S=e^{\ell_{\mathrm{plain}}/2}\),
\[
K_*\ge10\cdot2^{30}\ell_{\mathrm{plain}}L.
\]
Finally,
\[
1\le\ell_{\mathrm{plain}}L,\qquad
\epsilon^2L^2\le\frac{L^2}{16}
\le\frac{10}{16}\ell_{\mathrm{plain}}L,
\]
so
\[
1+\epsilon^2L^2\le10\ell_{\mathrm{plain}}L,
\qquad
K_*\ge2^{30}(1+\epsilon^2L^2).
\]
Substituting into the preceding tail bound proves, for both hypotheses,
\[
\nu_{\mathrm{lat}}\left\{
|\tau_\iota-t_\iota^{\mathrm{tar}}|
\ge\frac{\Delta_{\mathrm{tar}}}{8}\right\}
\le\frac1{128}.
\]
% lean: CausalSmith.Stat.AnnotationRarearmFrontier.affine_normalized_target_tail_small

\item
Let \(\widehat\tau_{\mathrm{raw}}\) be any Borel rule on the ordered raw channels and seed. Define its two Bayes squared risks by
\[
\mathcal A_\iota(\widehat\tau_{\mathrm{raw}})
=
\int
\mathbb E\left[
\bigl(\widehat\tau_{\mathrm{raw}}
(\mathcal L_{\mathrm{raw}},\mathcal X_{\mathrm{raw}},U)
-\tau_\iota(\boldsymbol\lambda)\bigr)^2
\mid\iota,\boldsymbol\lambda
\right]d\nu_{\mathrm{lat}}(\boldsymbol\lambda).
\]
Threshold the rule at
\[
t_{\mathrm{mid}}^{\mathrm{tar}}
=\frac{t_0^{\mathrm{tar}}+t_1^{\mathrm{tar}}}{2},
\]
and define the test errors
\[
\begin{aligned}
e_0^{\mathrm{test}}
&=\mathcal Q_0^{\mathrm{raw}}
  \{\widehat\tau_{\mathrm{raw}}\ge t_{\mathrm{mid}}^{\mathrm{tar}}\},\\
e_1^{\mathrm{test}}
&=\mathcal Q_1^{\mathrm{raw}}
  \{\widehat\tau_{\mathrm{raw}}<t_{\mathrm{mid}}^{\mathrm{tar}}\}.
\end{aligned}
\]
The definition of total variation gives
\[
e_0^{\mathrm{test}}+e_1^{\mathrm{test}}
\ge1-\operatorname{TV}(\mathcal Q_0^{\mathrm{raw}},
                       \mathcal Q_1^{\mathrm{raw}})
\ge1-\frac1{64}.
\]
On a test error with
\(|\tau_\iota-t_\iota^{\mathrm{tar}}|\le\Delta_{\mathrm{tar}}/8\),
the estimation error is at least \(3\Delta_{\mathrm{tar}}/8\). The two target-tail bounds therefore imply
\[
\mathcal A_0(\widehat\tau_{\mathrm{raw}})
+\mathcal A_1(\widehat\tau_{\mathrm{raw}})
\ge
\left(\frac{3\Delta_{\mathrm{tar}}}{8}\right)^2
\left(e_0^{\mathrm{test}}+e_1^{\mathrm{test}}-\frac2{128}\right).
\]
Consequently,
\[
\max_{\iota\in\{0,1\}}\mathcal A_\iota(\widehat\tau_{\mathrm{raw}})
\ge
\frac12\left(\frac{3\Delta_{\mathrm{tar}}}{8}\right)^2
\left(1-\frac1{64}-\frac2{128}\right)
=\frac{279}{4096}\Delta_{\mathrm{tar}}^2
\ge\frac{\Delta_{\mathrm{tar}}^2}{64}.
\]
If either Bayes risk is infinite, this conclusion holds directly. The common seed is already part of the experiment, so this argument includes randomized rules.
% lean: CausalSmith.Stat.AnnotationRarearmFrontier.affineRawTaggedKernel_bayes_lower

\item
We transfer this bound to an arbitrary clipped Borel rule \(T\) in the original experiment. First we control prefix failures. If
\(J\sim\operatorname{Poi}(\lambda_{\mathrm P})\) and
\(\lambda_{\mathrm P}\ge64k\), exponential Markov gives
\[
\Pr(J<k)
\le\Pr(J\le k)
\le4^k\mathbb E4^{-J}
=\exp\left(k\log4-\frac34\lambda_{\mathrm P}\right)
\le e^{-31k},
\]
using \(\log4\le3\). The event \(J<0\) is empty.

Since \(Q>1/2\), the conditional means satisfy
\[
uQ\ge64n,\qquad wQ\ge64N.
\]
For the retention event
\[
\mathcal E_{\mathrm{ret}}=\{J_{\mathrm L}\ge n,\ J_{\mathrm A}\ge m\},
\]
the complete-channel tail is bounded at \(n\), while
\(\{J_{\mathrm A}<m\}\subseteq\{J_{\mathrm A}<N\}\). Hence, uniformly in the latent configuration and hypothesis,
\[
\Pr(\mathcal E_{\mathrm{ret}}^c\mid\iota,\boldsymbol\lambda)
\le e^{-31n}+e^{-31N}\le2e^{-31n}.
\]
When \(m=0\), the auxiliary failure event is empty.

We calibrate this last quantity against the target gap. Since
\(\epsilon\le1/4\), \(n\ge4S\). For \(S\ge1\),
\[
S\le e^{124(S-1)},
\qquad
Se^{-124S}\le e^{-124}.
\]
Also \(e^{124}\ge2^{124}\), so
\[
2e^{-31n}\le2e^{-124S}\le\frac{10^{-24}}{S}.
\]
If \(x\le1\), the affine-regime inequality gives
\((\min\{x,1\})^2=x^2>S^{-1}\). If \(x>1\), then
\((\min\{x,1\})^2=1\ge S^{-1}\). The calibrated mean gap consequently satisfies
\[
\Delta_{\mathrm{tar}}^2
\ge10^{-20}(\min\{x,1\})^2
\ge\frac{10^{-20}}S.
\]
It follows that
\[
\Pr(\mathcal E_{\mathrm{ret}}^c\mid\iota,\boldsymbol\lambda)
\le2e^{-31n}
\le\frac{\Delta_{\mathrm{tar}}^2}{128}.
\]

Define the original rule's risk and its worst-case risk by
\[
\begin{aligned}
\mathcal R_P(T)
&=\mathbb E_{P^{\otimes n}\otimes P_{XA}^{\otimes m}
                  \otimes\mathrm{Uniform}[0,1]}
  [(T(\mathcal D,U)-\tau(P))^2],\\
\mathcal R_{\max}(T)
&=\sup_{P\in\mathcal M_{d,\epsilon}}\mathcal R_P(T).
\end{aligned}
\]
Clipping and \(\tau(P)\in[-1,1]\) give
\[
0\le\mathcal R_P(T)\le4.
\]
Thus this supremum is finite, and every constructed table's risk is at most \(\mathcal R_{\max}(T)\).

The transferred rule is
\[
\widetilde T(\mathcal L_{\mathrm{raw}},\mathcal X_{\mathrm{raw}},U)
=
\begin{cases}
T(\mathcal L_{\mathrm{raw}}[1:n],
  \mathcal X_{\mathrm{raw}}[1:m],U),
&\mathcal E_{\mathrm{ret}},\\
0,&\mathcal E_{\mathrm{ret}}^c.
\end{cases}
\]
The brackets denote the first indicated number of ordered records, with the empty prefix when that number is zero. This is a Borel rule.

Conditional on the latent table and the channel totals, the records are independent draws from the normalized laws. Thus, on retention, their prefixes have exactly the original two-channel product law, independently of the seed. In particular,
\[
\begin{aligned}
&\mathbb E\left[
(\widetilde T-\tau_\iota)^2\mathbf1_{\mathcal E_{\mathrm{ret}}}
\mid\iota,\boldsymbol\lambda\right]\\
&\qquad=
\Pr(\mathcal E_{\mathrm{ret}}\mid\iota,\boldsymbol\lambda)
\,\mathcal R_{P_\iota(\boldsymbol\lambda)}(T)
\le\mathcal R_{\max}(T).
\end{aligned}
\]
On failure, the output is zero and \(0\le\tau_\iota\le1\), so
\[
\mathbb E\left[
(\widetilde T-\tau_\iota)^2\mathbf1_{\mathcal E_{\mathrm{ret}}^c}
\mid\iota,\boldsymbol\lambda\right]
\le\Pr(\mathcal E_{\mathrm{ret}}^c\mid\iota,\boldsymbol\lambda)
\le\frac{\Delta_{\mathrm{tar}}^2}{128}.
\]
Averaging under either latent prior gives
\[
\mathcal A_\iota(\widetilde T)
\le\mathcal R_{\max}(T)+\frac{\Delta_{\mathrm{tar}}^2}{128}.
\]
Combining this with the raw testing lower bound yields
\[
\frac{\Delta_{\mathrm{tar}}^2}{64}
\le\max_\iota\mathcal A_\iota(\widetilde T)
\le\mathcal R_{\max}(T)+\frac{\Delta_{\mathrm{tar}}^2}{128},
\qquad
\mathcal R_{\max}(T)\ge\frac{\Delta_{\mathrm{tar}}^2}{128}.
\]
% lean: CausalSmith.Stat.AnnotationRarearmFrontier.affine_fixed_rule_gap_lower

\item
For \(x\ge0\), the two cases \(x\le1\) and \(x>1\) give
\[
(\min\{x,1\})^2
=
\begin{cases}
x^2,&x\le1,\\
1,&x>1
\end{cases}
=\min\{1,x^2\}.
\]
Therefore every original clipped Borel randomized rule satisfies
\[
\mathcal R_{\max}(T)
\ge\frac{\Delta_{\mathrm{tar}}^2}{128}
\ge\frac{10^{-20}}{128}\min\{1,x^2\}.
\]
The constant zero rule belongs to the rule class. Taking the infimum over that class, as specified in \cref{obj:def:risk}, proves
\[
R(n,m,d,\epsilon)
\ge\frac{10^{-20}}{128}\min\{1,x^2\}.
\]
Together with the support and marginal equalities already established, this proves the assertion with the stated numerical constant \(c\).
% lean: CausalSmith.Stat.AnnotationRarearmFrontier.affine_fixed_minimax_rate_lower
\end{enumerate}
\end{proof}

For a common latent configuration, the two hypotheses have identical normalized marginal tables because their arm masses and normalization agree. Pushing the common latent distribution through the two table maps consequently gives the same distribution of \(P_{XA}\) under both priors. Across latent configurations, the marginal table itself can vary. Thus the prior property concerns equality of distributions over marginal tables, together with equality under the specified common-latent coupling.

The risk conclusion in \cref{obj:lem:normalized-affine-testing} concerns the joint observation law of complete and auxiliary records in \cref{obj:def:risk}. Its observation comparison incorporates the complete-record outcomes and their association with the latent marginal table. The common marginal distribution describes the auxiliary channel's unconditional law under the priors; the stated lower bound accounts for both channels together. The resulting many-cell scale depends on total marginal information through \(N\), while the logarithmic degree calibration depends on rare-arm information through \(S\).

Finally, \cref{obj:thm:rare-label-floor} supplies the capped inverse-\(S\) lower bound throughout the public parameter range. It controls the branch \(x^2\le S^{-1}\) and, with the fixed numerical threshold absorbed into the comparison constant, the bounded-information branch \(S<\exp(4096)\). In the remaining branch, \cref{obj:lem:normalized-affine-testing} supplies the capped \(x^2\) lower bound. Together these comparisons give the lower order
\[
\min\{1,S^{-1}+x^2\}
\]
specified by \cref{obj:def:rate-handle}. Combined with the attaining-estimator upper comparison in \cref{obj:thm:uniform-annotation-frontier}, they establish the uniform minimax comparison for the original two-channel experiment.


\section{Resource characterizations and verification note}

The uniform comparison in \cref{obj:thm:uniform-annotation-frontier} supplies the common basis for the resource conclusions in \cref{obj:thm:annotation-resource-phases}. For every allowed experiment sequence in \cref{obj:synth_8}, the rate definition in \cref{obj:def:rate-handle} gives
\[
R(n_k,m_k,d_k,\epsilon_k)
\asymp
\min\left\{
1,\frac{1}{S_k}
+\left(\frac{d_k}{N_k\epsilon_k\ell_k}\right)^2
\right\}.
\]
Here the sequence-specific scales are those introduced in \cref{obj:def:phases}. The comparison constants are independent of the public indices. Consequently, the resource characterizations apply to arbitrary allowed sequences, including sequences with oscillating budgets, expanding alphabets, and diminishing overlap.

\subsection{Consistency and oracle order}

The consistency characterization in \cref{obj:thm:annotation-resource-phases} assigns a separate requirement to each component of the rate:
\[
\mathbf v\in\mathcal C
\quad\Longleftrightarrow\quad
S_k\longrightarrow\infty,
\qquad
\frac{d_k}{N_k\epsilon_k\ell_k}\longrightarrow0.
\]
The first condition concerns the information supplied by rare-arm outcomes in the complete records. The second concerns marginal information relative to the alphabet size. Auxiliary records contribute through the second denominator, while the complete-record budget and overlap floor determine the first scale. These two requirements characterize vanishing minimax risk throughout the sequence domain.

Oracle order compares precision with the rare-label benchmark in \cref{obj:prop:benchmark-recovery}. Under \(S_k\to\infty\), that benchmark eventually equals \(S_k^{-1}\). The oracle characterization therefore places the many-cell component at or below this inverse-information order:
\[
\mathbf v\in\mathcal O
\quad\Longleftrightarrow\quad
S_k\longrightarrow\infty,
\qquad
d_k=O\left(
\frac{N_k\epsilon_k\ell_k}{\sqrt{S_k}}
\right).
\]
Equivalently, the total marginal-information budget satisfies
\[
\frac{d_k\sqrt{S_k}}{\epsilon_k\ell_k}=O(N_k).
\]
As stated in \cref{obj:thm:annotation-resource-phases}, this budget characterization is conditional on divergence of the rare-arm information scale. Its order constant may depend on the sequence, whereas the underlying risk comparison is uniform over public indices.

For the supervised specialization \(m_k=0\), the identity \(N_k\epsilon_k=S_k\) gives the consistency threshold \(d_k=o(S_k\ell_k)\) and the oracle threshold \(d_k=O(\sqrt{S_k}\ell_k)\), together with \(S_k\to\infty\). These specializations identify the marginal-information requirements supplied entirely by complete records.

\subsection{Strict relative improvement}

The improvement phase in \cref{obj:def:phases} compares the auxiliary-record experiment with the supervised experiment at the same complete-record budget, alphabet size, and overlap floor. Applying the uniform risk comparison to both experiments gives
\[
\frac{R(n_k,m_k,d_k,\epsilon_k)}
     {R(n_k,0,d_k,\epsilon_k)}
\asymp
\frac{
\min\left\{1,S_k^{-1}
+\left(d_k/(N_k\epsilon_k\ell_k)\right)^2\right\}
}{
\min\left\{1,S_k^{-1}
+\left(d_k/(S_k\ell_k)\right)^2\right\}
}.
\]
Both caps enter this comparison, accommodating supervised risks ranging from vanishing precision to saturation.

The exact characterization in \cref{obj:thm:annotation-resource-phases} combines three requirements:
\[
S_k\longrightarrow\infty,
\qquad
\frac{d_k}{\sqrt{S_k}\ell_k}\longrightarrow\infty,
\qquad
\frac{N_k/n_k}
{\max\{1,d_k/(S_k\ell_k)\}}
\longrightarrow\infty.
\]
The first supports increasing outcome precision. The second places the supervised experiment above rare-label oracle order by a diverging factor. The third measures the marginal expansion relative to the larger of a constant and the supervised many-cell scale. Together they characterize a vanishing relative risk, including for sequences that move between different supervised resource regions. Within the supervised specialization, the relative risk ratio equals one.

\subsection{Saturation and the precision plateau}

The finite-experiment conclusions in \cref{obj:thm:annotation-resource-phases} retain the cap in the rate. When
\[
N\epsilon\ell\le d,
\]
the many-cell component reaches the cap, and the minimax risk is bounded above and below by positive universal constants. This saturation condition identifies a region of constant-order risk.

At the larger marginal budget
\[
N\ge \frac{d\sqrt{S}}{\epsilon\ell},
\]
the same result establishes
\[
R(n,m,d,\epsilon)\asymp b(n,\epsilon)
=\min\{1,S^{-1}\}.
\]
The cap preserves the comparison for small rare-arm information scales. Along every oracle sequence, \(S_k\to\infty\), so the risk eventually has order \(S_k^{-1}\).

The plateau conclusion also preserves the quantification over auxiliary budgets. By \cref{obj:thm:annotation-resource-phases}, on every oracle sequence there is a positive lower comparison constant such that, for all sufficiently large indices and every replacement budget \(m'\in\mathbb N\),
\[
R(n_k,m',d_k,\epsilon_k)\ge \frac{c}{S_k}.
\]
The rare-label floor in \cref{obj:thm:rare-label-floor} supplies this bound uniformly over the auxiliary channel. Thus oracle attainment identifies a precision order governed by complete-record outcome information, with further marginal information operating within that order.

\subsection{Verification note}

The mathematical scope includes the canonical randomized minimax risks, the explicit estimator and its risk guarantee, the testing and approximation bounds, the benchmark comparisons, the resource characterizations, and causal realization and identification. The statistical results use the finite observable model in \cref{obj:def:model}, the overlap and product-sampling conditions in \cref{obj:ass:overlap,obj:ass:data-product}, and the rule domains in \cref{obj:def:risk,obj:def:known-risk}. Causal identification concerns compatible potential-outcome extensions satisfying \cref{obj:ass:consistency,obj:ass:exchangeability}, as specified in \cref{obj:lem:causal-realization}. Auxiliary results retain their displayed independence, calibration, and approximation hypotheses. These conditions are assumed mathematical inputs to the stated implications.


\section{Proofs of the main results}\label{sec:deferred-proofs}

\begin{proof}[Proof of \cref{obj:thm:uniform-annotation-frontier}]
\begin{enumerate}
\item Choose numerical constants \(c_{\mathrm{lab}}>0\) and \(c_{\mathrm{aff}}>0\) from \cref{obj:thm:rare-label-floor,obj:lem:normalized-affine-testing}, respectively. For any public indices in the stated range, write
\[
N=n+m,\qquad
S=n\epsilon,\qquad
\ell=\log(\exp(1)+S),\qquad
x=\frac{d}{N\epsilon\ell}.
\]
Then \(N>0\), \(S>0\), and \(\ell>0\). By \cref{obj:def:rate-handle},
\[
r(n,m,d,\epsilon)=\min\{1,S^{-1}+x^2\},
\qquad
0<r(n,m,d,\epsilon)\le1.
\]
The two cited results supply
\[
c_{\mathrm{lab}}\,b(n,\epsilon)
\le R(n,m,d,\epsilon),
\qquad
b(n,\epsilon)=\min\{1,S^{-1}\},
\]
throughout the public index range, and
\[
c_{\mathrm{aff}}\min\{1,x^2\}
\le R(n,m,d,\epsilon)
\quad\text{when}\quad
S\ge\exp(4096),\qquad x^2>S^{-1}.
\]
Both constants are independent of the public indices. We next establish a uniform upper constant for the specified estimator.
% lean: CausalSmith.Stat.AnnotationRarearmFrontier.rare_label_floor
% lean: CausalSmith.Stat.AnnotationRarearmFrontier.normalized_affine_testing

\item First suppose \(S\ge\exp(4096)\), and use the pool lengths and tuning parameters from \cref{obj:def:estimator-handle}. Since \(\epsilon\le1/4\),
\[
12\le4\exp(4096)\le n.
\]
The integer splits satisfy
\[
h_0\ge\frac n3-1\ge\frac n4,\qquad
h_p\ge h_0,\qquad
h_p\le h_f\le h_p+1.
\]
Indeed, \(n-3h_0\in\{0,1,2\}\), and splitting the remaining records into two parts gives the last two inequalities. Consequently,
\[
M_p
=h_p+\lfloor m/2\rfloor
\ge\frac n4+\frac{m-1}{2}
\ge\frac N8,
\]
where the last inequality uses \(n\ge12\) and \(m\ge0\). Moreover,
\[
0\le M_f-M_p
=(h_f-h_p)+(m-2\lfloor m/2\rfloor)\le2.
\]
Thus
\[
h_0,M_p,M_f\ge1,\qquad
\frac n4\le h_0,\qquad
\frac N8\le M_p,M_f,\qquad
M_p\le M_f\le3M_p.
\]
For the prescribed intensities this yields
\[
u=\frac{h_0}{8}\ge\frac n{32},\qquad
t_p=\frac{M_p}{8}\ge\frac N{64},\qquad
t=\frac{M_f}{8}\ge\frac N{64},\qquad
\frac13\le\frac t{t_p}\le3.
\]
These bounds include \(m=0\).
% lean: CausalSmith.Stat.AnnotationRarearmFrontier.hybrid_finite_pool_sizes
% lean: CausalSmith.Stat.AnnotationRarearmFrontier.hybrid_finite_pool_intensities

\item Under \cref{obj:ass:data-product}, the three disjoint pools are independent. Their observation laws are \(P\), \(P_{XA}\), and \(P_{XA}\), respectively: projecting a complete record gives the same treatment--covariate law as an auxiliary record. Extend each pool by an independent iid continuation and introduce independent prefix lengths
\[
J_{\mathrm o}\sim\operatorname{Poi}(u),\qquad
J_{\mathrm p}\sim\operatorname{Poi}(t_p),\qquad
J_{\mathrm f}\sim\operatorname{Poi}(t),
\qquad
J_{\mathrm o},J_{\mathrm p},J_{\mathrm f}\in\mathbb N_0,
\]
independent of the observations.

Here is the count calculation for each pool. Let its finite alphabet be \(\mathsf Z_{\mathrm{pool}}\), its atom probabilities be \(\nu_{\mathrm{pool}}(z)\), and its prefix intensity be \(\lambda_{\mathrm{pool}}>0\), so that
\[
\nu_{\mathrm{pool}}(z)\ge0,\qquad
\sum_{z\in\mathsf Z_{\mathrm{pool}}}\nu_{\mathrm{pool}}(z)=1.
\]
For counts \(\kappa_z\in\mathbb N_0\), define
\[
\kappa_{\mathrm{tot}}
=\sum_{z\in\mathsf Z_{\mathrm{pool}}}\kappa_z.
\]
Conditioning on the prefix length and using the multinomial count formula gives
\[
\begin{aligned}
\Pr\{\text{every atom }z\text{ has count }\kappa_z\}
&=
e^{-\lambda_{\mathrm{pool}}}
\frac{\lambda_{\mathrm{pool}}^{\kappa_{\mathrm{tot}}}}
     {\kappa_{\mathrm{tot}}!}
\frac{\kappa_{\mathrm{tot}}!}{\prod_z\kappa_z!}
\prod_z\nu_{\mathrm{pool}}(z)^{\kappa_z}\\
&=
\prod_z
e^{-\lambda_{\mathrm{pool}}\nu_{\mathrm{pool}}(z)}
\frac{(\lambda_{\mathrm{pool}}\nu_{\mathrm{pool}}(z))^{\kappa_z}}
     {\kappa_z!}.
\end{aligned}
\]
We use \(0^0=1\); an atom of zero probability has count zero almost surely. This factorization, together with independence of the pools, shows that the prefix counts from \cref{obj:def:estimator-handle} are jointly independent and satisfy
\[
Z_{aj}\sim\operatorname{Poi}(uq_{aj}(P)),\qquad
J_{aj}\sim\operatorname{Poi}(t_ps_{aj}(P)),\qquad
K_{aj}\sim\operatorname{Poi}(ts_{aj}(P)).
\]
Thus \cref{obj:lem:uniform-hybrid-arm-risk} applies with precisely the prescribed \(L,B,k_0\).

Define, for \(a\in\{0,1\}\),
\[
\widehat\psi_a=\sum_{j=1}^dV_{aj},
\qquad
\psi_a=\sum_{j=1}^dp_j(P)\mu_{aj}(P),
\]
on these ideal prefixes. The cited result supplies a numerical constant \(C_{\mathrm{hyb}}>0\) such that
\[
\mathbb E[(\widehat\psi_a-\psi_a)^2]
\le
C_{\mathrm{hyb}}
\left\{S^{-1}+\left(\frac d{N\epsilon L}\right)^2\right\}.
\]
The binary conditional means and the covariate probabilities satisfy
\[
0\le\mu_{aj}(P)\le1,\qquad
p_j(P)\ge0,\qquad
\sum_{j=1}^dp_j(P)=1,
\]
including the zero convention for null cells. Hence
\[
-1
=\sum_jp_j(P)(-1)
\le\tau(P)
=\psi_1-\psi_0
\le\sum_jp_j(P)=1.
\]
Define the clipped ideal-prefix contrast by
\[
F_{\mathrm{Poi}}
=\operatorname{clip}_{[-1,1]}(\widehat\psi_1-\widehat\psi_0).
\]
Clipping toward an interval containing the target decreases squared loss, and
\[
\begin{aligned}
\mathbb E[(F_{\mathrm{Poi}}-\tau(P))^2]
&\le
\mathbb E[(\widehat\psi_1-\widehat\psi_0-\tau(P))^2]\\
&\le
2\mathbb E[(\widehat\psi_1-\psi_1)^2]
+2\mathbb E[(\widehat\psi_0-\psi_0)^2]\\
&\le
4C_{\mathrm{hyb}}
\left\{S^{-1}+\left(\frac d{N\epsilon L}\right)^2\right\}.
\end{aligned}
\]
The second inequality is pointwise, so it applies with the shared counts used by the two arm statistics.
% lean: CausalSmith.Stat.AnnotationRarearmFrontier.hybrid_ordered_prefix_arm_risk
% lean: CausalSmith.Stat.AnnotationRarearmFrontier.hybrid_ordered_prefix_contrast_risk

\item We calibrate the last bound to the numerical rate and bound the prefix overflow cost. From \(S\ge\exp(4096)\) and the definition of \(L\),
\[
L=\left\lfloor\frac{\log S}{1024}\right\rfloor\ge4,
\qquad
\log S<1024(L+1)\le1280L.
\]
Also \(S\ge2\) and \(\exp(1)\ge2\), which imply
\[
\exp(1)+S\le S\exp(1),
\qquad
\ell\le\log S+1\le1281L.
\]
All denominators below are positive, so
\[
0\le\frac d{N\epsilon L}\le1281x,
\qquad
S^{-1}+\left(\frac d{N\epsilon L}\right)^2
\le1281^2(S^{-1}+x^2).
\]

For every positive pool length \(h\), the inequality \(e^h\ge1+h\ge h\) gives
\[
e^{-h}\le h^{-1}.
\]
Applying this to the three pool lengths and using the bounds already established,
\[
\begin{aligned}
e^{-h_0}+e^{-M_p}+e^{-M_f}
&\le\frac4n+\frac8n+\frac8n\\
&=\frac{20}{n}
\le\frac5S,
\end{aligned}
\]
where the last inequality uses \(\epsilon\le1/4\). Define the numerical constant
\[
C_{\mathrm{large}}=4C_{\mathrm{hyb}}1281^2+5>0.
\]
Combining the ideal-prefix risk, degree comparison, and overflow bound yields
\[
\begin{aligned}
\mathbb E[(F_{\mathrm{Poi}}-\tau(P))^2]
+e^{-h_0}+e^{-M_p}+e^{-M_f}
&\le
4C_{\mathrm{hyb}}1281^2(S^{-1}+x^2)+5S^{-1}\\
&=
C_{\mathrm{large}}(S^{-1}+x^2)-5x^2\\
&\le C_{\mathrm{large}}(S^{-1}+x^2).
\end{aligned}
\]
% lean: CausalSmith.Stat.AnnotationRarearmFrontier.hybrid_finite_degree_rate
% lean: CausalSmith.Stat.AnnotationRarearmFrontier.hybrid_finite_overflow_rate
% lean: CausalSmith.Stat.AnnotationRarearmFrontier.hybrid_tuned_prefix_and_overflow_rate

\item Apply \cref{obj:lem:finite-prefix-transfer} to the three positive pool lengths \(h_0,M_p,M_f\), the clipped prefix statistic, and the target \(\tau(P)\in[-1,1]\). Its finite average is exactly the estimator in \cref{obj:def:estimator-handle}, since
\[
\Pr(J_{\mathrm o}=h',J_{\mathrm p}=h'',J_{\mathrm f}=h''')
=
e^{-u-t_p-t}
\frac{u^{h'}}{h'!}
\frac{t_p^{h''}}{h''!}
\frac{t^{h'''}}{h'''!}.
\]
The zero output on overflow agrees with the construction. Independence and the stated pool laws therefore give
\[
\begin{aligned}
\mathbb E_P[
(\widehat\tau(\mathcal D,U)-\tau(P))^2]
&\le
\mathbb E[(F_{\mathrm{Poi}}-\tau(P))^2]
+e^{-h_0}+e^{-M_p}+e^{-M_f}\\
&\le C_{\mathrm{large}}(S^{-1}+x^2).
\end{aligned}
\]
This proves the uncapped upper bound for every model law when \(S\ge\exp(4096)\).
% lean: CausalSmith.Stat.AnnotationRarearmFrontier.hybrid_large_information_risk

\item Define
\[
C_{\mathrm{up}}
=\max\{\max\{C_{\mathrm{large}},4\},\exp(4096)\}>0.
\]
For fixed public indices, the observed-data space is finite, so the deterministic estimator is Borel measurable. Its definition is independent of \(U\). In the averaging branch, every prefix statistic lies in \([-1,1]\), the displayed prefix probabilities are nonnegative, and
\[
\sum_{h'=0}^{h_0}\sum_{h''=0}^{M_p}\sum_{h'''=0}^{M_f}
\Pr(J_{\mathrm o}=h',J_{\mathrm p}=h'',J_{\mathrm f}=h''')
\le1.
\]
Thus, with the remaining probability assigned output zero,
\[
-1\le\widehat\tau(\mathcal D,U)\le1.
\]
The zero branch has the same range. Since \(\tau(P)\in[-1,1]\), this also gives the universal bound
\[
\mathbb E_P[
(\widehat\tau(\mathcal D,U)-\tau(P))^2]\le4.
\]

We now follow the three upper-bound regimes. If \(S<\exp(4096)\), the estimator is zero, and
\[
\mathbb E_P[
(\widehat\tau(\mathcal D,U)-\tau(P))^2]
=\tau(P)^2\le1.
\]
Moreover,
\[
e^{-4096}\le1,\qquad e^{-4096}\le S^{-1},
\qquad
e^{-4096}\le r(n,m,d,\epsilon),
\]
so
\[
\mathbb E_P[
(\widehat\tau(\mathcal D,U)-\tau(P))^2]
\le1
\le\exp(4096)r(n,m,d,\epsilon)
\le C_{\mathrm{up}}r(n,m,d,\epsilon).
\]
If \(S\ge\exp(4096)\) and \(1\le S^{-1}+x^2\), then
\[
r(n,m,d,\epsilon)=1,
\qquad
\mathbb E_P[
(\widehat\tau(\mathcal D,U)-\tau(P))^2]
\le4\le C_{\mathrm{up}}r(n,m,d,\epsilon).
\]
Finally, if \(S\ge\exp(4096)\) and \(S^{-1}+x^2<1\), the preceding uncapped bound gives
\[
\mathbb E_P[
(\widehat\tau(\mathcal D,U)-\tau(P))^2]
\le C_{\mathrm{large}}(S^{-1}+x^2)
\le C_{\mathrm{up}}r(n,m,d,\epsilon).
\]
Taking the supremum over the model class proves
\[
\sup_{P\in\mathcal M_{d,\epsilon}}
\mathbb E_P[
(\widehat\tau(\mathcal D,U)-\tau(P))^2]
\le C_{\mathrm{up}}r(n,m,d,\epsilon).
\]
% lean: CausalSmith.Stat.AnnotationRarearmFrontier.hybrid_upper

\item Choose the final numerical constants by
\[
E_{\mathrm{cut}}=\exp(4096),\qquad
c=\min\left\{
\frac{c_{\mathrm{lab}}}{2E_{\mathrm{cut}}},
\frac{c_{\mathrm{aff}}}{2}
\right\},
\qquad
C=\max\{C_{\mathrm{up}},c\}.
\]
They satisfy
\[
1\le E_{\mathrm{cut}},\qquad
0<c\le C<\infty,
\]
and the calibrations needed below are
\[
c\le\frac{c_{\mathrm{lab}}}{2E_{\mathrm{cut}}}
\le\frac{c_{\mathrm{lab}}}{E_{\mathrm{cut}}},
\qquad
2c\le\frac{c_{\mathrm{lab}}}{E_{\mathrm{cut}}}
\le c_{\mathrm{lab}},
\qquad
2c\le c_{\mathrm{aff}}.
\]
All these constants were chosen independently of the public indices.

The lifted rule \((\mathcal D,U)\mapsto\widehat\tau(\mathcal D,U)\) is Borel measurable by composition with the data projection and takes values in \([-1,1]\). It is therefore admissible in \cref{obj:def:risk}. Squared-error risks are nonnegative, and the defining minimax infimum yields
\[
R(n,m,d,\epsilon)
\le
\sup_{P\in\mathcal M_{d,\epsilon}}
\mathbb E_P[
(\widehat\tau(\mathcal D,U)-\tau(P))^2]
\le C_{\mathrm{up}}r(n,m,d,\epsilon)
\le Cr(n,m,d,\epsilon),
\]
where the last inequality uses \(r(n,m,d,\epsilon)\ge0\).
% lean: CausalSmith.Stat.AnnotationRarearmFrontier.uniform_annotation_frontier

\item For the lower bound, first suppose \(S<E_{\mathrm{cut}}\). Positivity of \(S\) and \(E_{\mathrm{cut}}\ge1\) gives
\[
E_{\mathrm{cut}}^{-1}\le S^{-1},
\qquad
E_{\mathrm{cut}}^{-1}\le1,
\qquad
E_{\mathrm{cut}}^{-1}
\le b(n,\epsilon).
\]
Using \(r(n,m,d,\epsilon)\le1\), the constant calibration, and the rare-label bound,
\[
\begin{aligned}
c\,r(n,m,d,\epsilon)
&\le c\\
&\le c_{\mathrm{lab}}E_{\mathrm{cut}}^{-1}\\
&\le c_{\mathrm{lab}}b(n,\epsilon)\\
&\le R(n,m,d,\epsilon).
\end{aligned}
\]
% lean: CausalSmith.Stat.AnnotationRarearmFrontier.uniform_annotation_frontier

\item Now suppose \(S\ge E_{\mathrm{cut}}\). Then \(S\ge1\), and hence
\[
b(n,\epsilon)=S^{-1}.
\]
If \(x^2\le S^{-1}\), the rate satisfies
\[
r(n,m,d,\epsilon)
\le S^{-1}+x^2\le2S^{-1}.
\]
It follows that
\[
c\,r(n,m,d,\epsilon)
\le2cS^{-1}
\le c_{\mathrm{lab}}S^{-1}
\le R(n,m,d,\epsilon).
\]

If \(x^2>S^{-1}\), all hypotheses of \cref{obj:lem:normalized-affine-testing} hold, and it supplies
\[
c_{\mathrm{aff}}\min\{1,x^2\}\le R(n,m,d,\epsilon).
\]
To compare this quantity with the rate, split according to whether \(x^2\le1\). In that case,
\[
r(n,m,d,\epsilon)
\le S^{-1}+x^2
<2x^2
=2\min\{1,x^2\}.
\]
When \(x^2>1\),
\[
r(n,m,d,\epsilon)\le1\le2=2\min\{1,x^2\}.
\]
Thus in both cases,
\[
\begin{aligned}
c\,r(n,m,d,\epsilon)
&\le2c\min\{1,x^2\}\\
&\le c_{\mathrm{aff}}\min\{1,x^2\}\\
&\le R(n,m,d,\epsilon).
\end{aligned}
\]
Together with the preceding branches, this proves the lower bound throughout the public index range.
% lean: CausalSmith.Stat.AnnotationRarearmFrontier.uniform_annotation_frontier

\item Finally, fix \(d\ge2\), \(0<\epsilon\le1/4\), and \(P\in\mathcal M_{d,\epsilon}\). Choose the independent conditional extension specified in \cref{obj:lem:causal-realization}: \(H\) is a joint law of \((X,A,Y,Y_0,Y_1)\), with binary potential outcomes, satisfying
\[
\mathcal L_H(X,A,Y)=P,\qquad
Y=AY_1+(1-A)Y_0\quad H\text{-almost surely},
\qquad
(Y_0,Y_1)\perp\!\!\!\perp A\mid X.
\]
The cited result supplies both this extension and, for every extension with the same observed marginal satisfying \cref{obj:ass:consistency,obj:ass:exchangeability}, the identity
\[
\mathbb E_H[Y_1-Y_0]=\tau(P).
\]
This establishes the causal realization and identification assertions.
% lean: CausalSmith.Stat.AnnotationRarearmFrontier.causal_realization
\end{enumerate}
\end{proof}

\begin{proof}[Proof of \cref{obj:thm:rare-label-floor}]
\begin{enumerate}
\item \textit{Choice of the two laws.}
Fix \(n,m,d,\epsilon\) satisfying the hypotheses of the risk bounds, and set
\[
c:=\frac1{1024},
\qquad
S:=n\epsilon>0,
\qquad
\delta_n:=\frac1{16\sqrt{\max\{1,S\}}}.
\]
Since \(\max\{1,S\}\ge1\), we have
\[
0<\delta_n\le\frac1{16}\le\frac18.
\]
Write
\[
\mathsf Z_d:=\{1,\ldots,d\}\times\{0,1\}^2.
\]
For \(\xi\in\{-1,1\}\), define the law \(P_{\xi,n}\) on \(\mathsf Z_d\) by
\[
\begin{aligned}
P_{\xi,n}(j,1,1)&=\frac{\epsilon}{2}\left(\frac12+\xi\delta_n\right),&
P_{\xi,n}(j,1,0)&=\frac{\epsilon}{2}\left(\frac12-\xi\delta_n\right),\\
P_{\xi,n}(j,0,1)&=\frac{1-\epsilon}{4},&
P_{\xi,n}(j,0,0)&=\frac{1-\epsilon}{4},
&&j=1,2,
\end{aligned}
\]
and
\[
P_{\xi,n}(j,a,y)=0,
\qquad j>2,\quad a,y\in\{0,1\}.
\]
Also set
\[
P_{+,n}:=P_{1,n},
\qquad
P_{-,n}:=P_{-1,n}.
\]
Every atom in the first two covariate cells is positive. Summing within each such cell gives
\[
\sum_{a,y\in\{0,1\}}P_{\xi,n}(j,a,y)
=
\frac{\epsilon}{2}+\frac{1-\epsilon}{2}
=\frac12,
\qquad j=1,2.
\]
Thus the total mass is one. Summing over outcomes and taking the conditional ratios gives
\[
\begin{aligned}
p_j(P_{\xi,n})&=\frac12,
&
e_j(P_{\xi,n})&=\frac{\epsilon/2}{1/2}=\epsilon,\\
\mu_{0j}(P_{\xi,n})&=
\frac{(1-\epsilon)/4}{(1-\epsilon)/2}=\frac12,
&
\mu_{1j}(P_{\xi,n})&=
\frac{(\epsilon/2)(1/2+\xi\delta_n)}{\epsilon/2}
=\frac12+\xi\delta_n,
&&j=1,2.
\end{aligned}
\]
The remaining covariate masses are zero. Because
\[
\epsilon\le1-\epsilon,
\]
the occupied cells satisfy \cref{obj:ass:overlap}, so \cref{obj:def:model} gives
\[
P_{+,n},P_{-,n}\in\mathcal M_{d,\epsilon}.
\]
By \cref{obj:def:target},
\[
\tau(P_{\xi,n})
=
\sum_{j=1}^{2}\frac12\,\xi\delta_n
=\xi\delta_n.
\]
Define their common treatment--covariate marginal by
\[
\pi_{XA}(j,a):=
\begin{cases}
\epsilon/2,&j\in\{1,2\},\ a=1,\\
(1-\epsilon)/2,&j\in\{1,2\},\ a=0,\\
0,&j>2.
\end{cases}
\]
The outcome sums above show that
\[
(P_{+,n})_{XA}=(P_{-,n})_{XA}=\pi_{XA}.
\]
When \(d=2\), all statements concerning \(j>2\) have an empty index set.
% lean: CausalSmith.Stat.AnnotationRarearmFrontier.labelFloorAmplitude_scaled
% lean: CausalSmith.Stat.AnnotationRarearmFrontier.labelFloor_jointMass
% lean: CausalSmith.Stat.AnnotationRarearmFrontier.labelFloor_model
% lean: CausalSmith.Stat.AnnotationRarearmFrontier.labelFloor_target
% lean: CausalSmith.Stat.AnnotationRarearmFrontier.labelFloor_auxMarginal

\item \textit{Closeness of the complete-record experiments.}
Use the natural-logarithm divergence and finite total variation distance of
\cref{obj:lem:pinsker-finite}. The two laws have common positive support,
namely \(\{1,2\}\times\{0,1\}^2\), so their divergence is finite.
For either occupied cell, cancellation of the common control terms gives
\[
\begin{aligned}
&\sum_{a,y\in\{0,1\}}
P_{+,n}(j,a,y)
\log\frac{P_{+,n}(j,a,y)}{P_{-,n}(j,a,y)}\\
&\quad=
\frac{\epsilon}{2}\left[
\left(\frac12+\delta_n\right)
\log\frac{1/2+\delta_n}{1/2-\delta_n}
+
\left(\frac12-\delta_n\right)
\log\frac{1/2-\delta_n}{1/2+\delta_n}
\right]\\
&\quad=
\epsilon\delta_n
\log\frac{1/2+\delta_n}{1/2-\delta_n},
\qquad j=1,2,
\end{aligned}
\]
where the last equality uses \(\log(x^{-1})=-\log x\).
The null cells contribute zero. Hence
\[
D(P_{+,n}\Vert P_{-,n})
=
2\epsilon\delta_n
\log\frac{1/2+\delta_n}{1/2-\delta_n}.
\]
The ratio is positive, and \(\log x\le x-1\) yields
\[
\log\frac{1/2+\delta_n}{1/2-\delta_n}
\le
\frac{1/2+\delta_n}{1/2-\delta_n}-1
=
\frac{2\delta_n}{1/2-\delta_n}
\le8\delta_n.
\]
Indeed, \(1/2-\delta_n\ge3/8\ge1/4\). Multiplication by
\(2\epsilon\delta_n>0\) therefore gives
\[
D(P_{+,n}\Vert P_{-,n})
\le16\epsilon\delta_n^2.
\]

Define the complete-record array laws by
\[
\mathsf P_{\mathrm{lab},\pm}:=P_{\pm,n}^{\otimes n}.
\]
They too have common support. For an array
\(\boldsymbol z=(z_1,\ldots,z_n)\) in that support,
\[
\log
\frac{\mathsf P_{\mathrm{lab},+}(\boldsymbol z)}
     {\mathsf P_{\mathrm{lab},-}(\boldsymbol z)}
=
\sum_{i=1}^{n}
\log\frac{P_{+,n}(z_i)}{P_{-,n}(z_i)}.
\]
Taking expectations, each coordinate has law \(P_{+,n}\), and thus
\[
\mathbb E_{\mathsf P_{\mathrm{lab},+}}
\left[
\sum_{i=1}^{n}
\log\frac{P_{+,n}(z_i)}{P_{-,n}(z_i)}
\right]
=
nD(P_{+,n}\Vert P_{-,n}).
\]
Consequently the product divergence is finite and satisfies
\[
D(\mathsf P_{\mathrm{lab},+}\Vert\mathsf P_{\mathrm{lab},-})
\le
nD(P_{+,n}\Vert P_{-,n})
\le16S\delta_n^2.
\]
Squaring the chosen amplitude gives
\[
\delta_n^2
=
\frac1{256\max\{1,S\}},
\qquad
16S\delta_n^2
=
\frac{S}{16\max\{1,S\}}
\le\frac1{16}.
\]
Applying \cref{obj:lem:pinsker-finite} to these finite array laws now gives
\[
\operatorname{TV}
(\mathsf P_{\mathrm{lab},+},\mathsf P_{\mathrm{lab},-})
\le
\sqrt{
\frac{D(\mathsf P_{\mathrm{lab},+}\Vert
        \mathsf P_{\mathrm{lab},-})}{2}}
\le\frac14,
\]
since the squared expression under the final comparison is at most
\(1/32\le1/16\).
% lean: CausalSmith.Stat.AnnotationRarearmFrontier.labelFloor_KL_formula
% lean: CausalSmith.Stat.AnnotationRarearmFrontier.labelFloor_KL_bound
% lean: CausalSmith.Stat.AnnotationRarearmFrontier.labelFloor_absolutelyContinuous
% lean: CausalSmith.Stat.AnnotationRarearmFrontier.labelFloor_labeled_KL_bound
% lean: CausalSmith.Stat.AnnotationRarearmFrontier.labelFloor_scaled_tv

\item \textit{Calibration to the benchmark.}
The square identity gives
\[
\frac{\delta_n^2}{4}
=
\frac1{1024\max\{1,S\}}.
\]
If \(S\le1\), positivity of \(S\) gives \(S^{-1}\ge1\), and therefore
\[
\frac1{\max\{1,S\}}=1=\min\{1,S^{-1}\}.
\]
If \(S>1\), then \(S^{-1}<1\), and
\[
\frac1{\max\{1,S\}}=S^{-1}=\min\{1,S^{-1}\}.
\]
In both cases,
\[
\frac{\delta_n^2}{4}
=
\frac1{1024}\min\{1,(n\epsilon)^{-1}\}
=
c\,b(n,\epsilon).
\]
% lean: CausalSmith.Stat.AnnotationRarearmFrontier.labelFloor_scaled_rate

\item \textit{Adjoining the auxiliary observations and randomization.}
For probability measures on a general measurable space, use the extension
\[
\operatorname{TV}(\nu,\nu')
:=
\sup_{\mathcal E\ \mathrm{measurable}}
|\nu(\mathcal E)-\nu'(\mathcal E)|,
\]
which agrees with the finite formula in \cref{obj:lem:pinsker-finite}.

We first establish the common-factor identity needed here. Let
\(\nu,\nu'\) be probability laws on a finite set
\(\Omega_{\mathrm{fin}}\), and let \(\rho\) be a probability law on a
measurable space \(\Omega_{\mathrm{common}}\). For any measurable
\(\mathcal E\subseteq
\Omega_{\mathrm{fin}}\times\Omega_{\mathrm{common}}\), define its sections by
\[
\mathcal E_z:=\{w\in\Omega_{\mathrm{common}}:(z,w)\in\mathcal E\},
\qquad z\in\Omega_{\mathrm{fin}}.
\]
The signed masses sum to zero, so
\[
\sum_z\max\{\nu(z)-\nu'(z),0\}
=
\sum_z\max\{\nu'(z)-\nu(z),0\}
=
\frac12\sum_z|\nu(z)-\nu'(z)|.
\]
Since \(0\le\rho(\mathcal E_z)\le1\), it follows that
\[
\begin{aligned}
|(\nu\otimes\rho)(\mathcal E)-(\nu'\otimes\rho)(\mathcal E)|
&=
\left|
\sum_z(\nu(z)-\nu'(z))\rho(\mathcal E_z)
\right|\\
&\le\frac12\sum_z|\nu(z)-\nu'(z)|
=\operatorname{TV}(\nu,\nu').
\end{aligned}
\]
Conversely, for every subset
\(\mathcal E_{\mathrm{fin}}\subseteq\Omega_{\mathrm{fin}}\),
\[
(\nu\otimes\rho)
(\mathcal E_{\mathrm{fin}}\times\Omega_{\mathrm{common}})
=
\nu(\mathcal E_{\mathrm{fin}}),
\]
and the same identity holds for \(\nu'\). Taking suprema in the two
directions proves
\[
\operatorname{TV}(\nu\otimes\rho,\nu'\otimes\rho)
=
\operatorname{TV}(\nu,\nu').
\]

By \cref{obj:ass:data-product}, the augmented observation laws for the
original-record experiment are
\[
\mathsf P_{\mathrm{ann},\pm}
:=
\bigl(\mathsf P_{\mathrm{lab},\pm}
      \otimes\pi_{XA}^{\otimes m}\bigr)
\otimes\mathrm{Uniform}[0,1].
\]
First removing the common seed factor and then the common auxiliary
factor gives
\[
\begin{aligned}
\operatorname{TV}
(\mathsf P_{\mathrm{ann},+},\mathsf P_{\mathrm{ann},-})
&=
\operatorname{TV}
\bigl(
\mathsf P_{\mathrm{lab},+}\otimes\pi_{XA}^{\otimes m},
\mathsf P_{\mathrm{lab},-}\otimes\pi_{XA}^{\otimes m}
\bigr)\\
&=
\operatorname{TV}
(\mathsf P_{\mathrm{lab},+},\mathsf P_{\mathrm{lab},-})
\le\frac14.
\end{aligned}
\]
For \(m=0\), the auxiliary array space consists of the empty array and its
law is a point mass, so the same argument applies.
% lean: CausalSmith.Stat.AnnotationRarearmFrontier.labelFloor_tv_common_product

\item \textit{Bounded risks and the testing reduction.}
For any observable law \(P\), the conditional outcome means satisfy
\begingroup\small
\[
\mu_{aj}(P)=
\begin{cases}
\displaystyle
\frac{P(j,a,1)}{P(j,a,0)+P(j,a,1)},
& P(j,a,0)+P(j,a,1)>0,\\[6pt]
0,&P(j,a,0)+P(j,a,1)=0,
\end{cases}
\qquad a\in\{0,1\},\quad j\in\{1,\ldots,d\}.
\]
\endgroup
The numerator lies between zero and the denominator, so
\[
0\le\mu_{aj}(P)\le1,
\qquad
-1\le\mu_{1j}(P)-\mu_{0j}(P)\le1.
\]
Since \(p_j(P)\ge0\) and \(\sum_jp_j(P)=1\),
\cref{obj:def:target} yields
\[
-1
=
\sum_jp_j(P)(-1)
\le\tau(P)
\le
\sum_jp_j(P)
=1.
\]
Consequently, every clipped rule obeys
\[
|T(\mathcal D,U)-\tau(P)|\le2,
\qquad
0\le(T(\mathcal D,U)-\tau(P))^2\le4,
\]
and hence its risk is at most four under every model law. The same
pointwise bound applies to clipped exact-marginal rules. The constant
zero rule belongs to both rule classes.

To establish the testing inequality, let
\[
\nu_\pm\in\mathsf{Prob}(\Omega_{\mathrm{test}}),
\qquad
T_{\mathrm{test}}:\Omega_{\mathrm{test}}\longrightarrow[-1,1]
\quad\text{be measurable},
\qquad
\operatorname{TV}(\nu_+,\nu_-)\le\frac14,
\]
where \(\Omega_{\mathrm{test}}\) is any measurable space. Define
\[
\mathcal E_+
:=
\{\omega:|T_{\mathrm{test}}(\omega)-\delta_n|\ge\delta_n\},
\qquad
\mathcal E_-
:=
\{\omega:|T_{\mathrm{test}}(\omega)+\delta_n|\ge\delta_n\}.
\]
The targets are separated by
\[
|\delta_n-(-\delta_n)|=2\delta_n.
\]
If \(\omega\in\mathcal E_+^c\), the triangle inequality gives
\[
2\delta_n
\le
|T_{\mathrm{test}}(\omega)-\delta_n|
+
|T_{\mathrm{test}}(\omega)+\delta_n|
<
\delta_n+|T_{\mathrm{test}}(\omega)+\delta_n|.
\]
Thus \(\mathcal E_+^c\subseteq\mathcal E_-\), and the defining event bound
for total variation implies
\[
\begin{aligned}
\nu_+(\mathcal E_+)+\nu_-(\mathcal E_-)
&\ge
\nu_+(\mathcal E_+)+\nu_-(\mathcal E_+^c)\\
&=
1+\nu_+(\mathcal E_+)-\nu_-(\mathcal E_+)\\
&\ge1-\operatorname{TV}(\nu_+,\nu_-).
\end{aligned}
\]
Therefore
\[
\frac14
\le
\frac{1-\operatorname{TV}(\nu_+,\nu_-)}2
\le
\max\{\nu_+(\mathcal E_+),\nu_-(\mathcal E_-)\}.
\]
For each sign, the pointwise loss bound on the corresponding error
region is
\[
\delta_n^2\mathbf1_{\mathcal E_\pm}(\omega)
\le
\bigl(T_{\mathrm{test}}(\omega)\mp\delta_n\bigr)^2.
\]
The losses are integrable because the rule is bounded. Integrating and
taking the maximum yields
\[
\begin{aligned}
\frac{\delta_n^2}{4}
&\le
\delta_n^2
\max\{\nu_+(\mathcal E_+),\nu_-(\mathcal E_-)\}\\
&\le
\max\left\{
\int (T_{\mathrm{test}}-\delta_n)^2\,d\nu_+,
\int (T_{\mathrm{test}}+\delta_n)^2\,d\nu_-
\right\}.
\end{aligned}
\]
% lean: CausalSmith.Stat.AnnotationRarearmFrontier.ateFunctional_mem_Icc
% lean: CausalSmith.Stat.AnnotationRarearmFrontier.labelFloor_bounded_mse
% lean: Causalean.Stat.half_one_sub_tvDist_le_max_error
% lean: CausalSmith.Stat.AnnotationRarearmFrontier.labelFloor_testing_mse

\item \textit{The original-record minimax bound.}
Fix any clipped Borel randomized rule \(T\) admitted by
\cref{obj:def:risk}. Apply the testing inequality with
\[
\nu_\pm:=\mathsf P_{\mathrm{ann},\pm},
\qquad
T_{\mathrm{test}}(\mathcal D,U):=T(\mathcal D,U).
\]
Using the two target identities, we obtain
\[
\max\left\{
\mathbb E_{\mathsf P_{\mathrm{ann},+}}
[(T(\mathcal D,U)-\tau(P_{+,n}))^2],
\mathbb E_{\mathsf P_{\mathrm{ann},-}}
[(T(\mathcal D,U)-\tau(P_{-,n}))^2]
\right\}
\ge\frac{\delta_n^2}{4}
=c\,b(n,\epsilon).
\]
Both laws belong to \(\mathcal M_{d,\epsilon}\), so
\[
\sup_{P\in\mathcal M_{d,\epsilon}}
\mathbb E_{P^{\otimes n}\otimes P_{XA}^{\otimes m}
                   \otimes\mathrm{Uniform}[0,1]}
[(T(\mathcal D,U)-\tau(P))^2]
\ge c\,b(n,\epsilon).
\]
The supremum is finite by the bound of four, and the rule class is
nonempty. Taking the infimum over all such \(T\), as specified in
\cref{obj:def:risk}, proves
\[
R(n,m,d,\epsilon)\ge c\,b(n,\epsilon).
\]
% lean: Causalean.Stat.le_minimaxValue_of_two_point
% lean: CausalSmith.Stat.AnnotationRarearmFrontier.rare_label_floor

\item \textit{The exact-marginal minimax bound.}
Define the complete-record laws augmented by the seed by
\[
\mathsf P_{\mathrm K,\pm}
:=
\mathsf P_{\mathrm{lab},\pm}
\otimes\mathrm{Uniform}[0,1].
\]
The common-factor identity gives
\[
\operatorname{TV}
(\mathsf P_{\mathrm K,+},\mathsf P_{\mathrm K,-})
=
\operatorname{TV}
(\mathsf P_{\mathrm{lab},+},\mathsf P_{\mathrm{lab},-})
\le\frac14.
\]
Fix a clipped Borel rule \(T^{\mathrm K}\) admitted by
\cref{obj:def:known-risk}, and freeze its marginal input at the common
table:
\[
\widetilde T^{\mathrm K}(\mathcal L,U)
:=
T^{\mathrm K}(\mathcal L,\pi_{XA},U).
\]
This is a measurable clipped rule on the complete-record arrays and the
seed. Applying the testing inequality to this rule gives
\[
\max\left\{
\mathbb E_{\mathsf P_{\mathrm K,+}}
[(\widetilde T^{\mathrm K}(\mathcal L,U)-\delta_n)^2],
\mathbb E_{\mathsf P_{\mathrm K,-}}
[(\widetilde T^{\mathrm K}(\mathcal L,U)+\delta_n)^2]
\right\}
\ge\frac{\delta_n^2}{4}
=c\,b(n,\epsilon).
\]
Because \((P_{\pm,n})_{XA}=\pi_{XA}\), these are precisely the two risks
of \(T^{\mathrm K}\) at \(P_{+,n}\) and \(P_{-,n}\). Consequently,
\[
\sup_{P\in\mathcal M_{d,\epsilon}}
\mathbb E_{P^{\otimes n}\otimes\mathrm{Uniform}[0,1]}
\left[
\bigl(T^{\mathrm K}(\mathcal L,P_{XA},U)-\tau(P)\bigr)^2
\right]
\ge c\,b(n,\epsilon).
\]
Again the supremum is bounded by four and the rule class contains the
constant zero rule. The squared-error clipping reduction in
\cref{obj:def:risk} also applies when the exact marginal is supplied as
an input: projecting the rule's output onto \([-1,1]\) preserves Borel
measurability and weakly decreases squared loss, since
\(\tau(P)\in[-1,1]\). Thus the infimum over all Borel exact-marginal
rules equals the infimum over clipped Borel exact-marginal rules.
Taking the infimum in \cref{obj:def:known-risk}
proves
\[
R^{\mathrm K}(n,d,\epsilon)\ge c\,b(n,\epsilon).
\]
The chosen \(c=1/1024>0\) is independent of all public indices.
% lean: CausalSmith.Stat.AnnotationRarearmFrontier.labelFloor_tv_common_product
% lean: Causalean.Stat.le_minimaxValue_of_two_point
% lean: CausalSmith.Stat.AnnotationRarearmFrontier.rare_label_floor

\item \textit{The pair at an arbitrary admissible separation.}
Now fix \(d\ge2\), \(0<\epsilon\le1/4\), and \(0<\delta\le1/8\).
Define the asserted pair by
\[
\begin{aligned}
P_\pm(j,1,1)&=\frac{\epsilon}{2}\left(\frac12\pm\delta\right),&
P_\pm(j,1,0)&=\frac{\epsilon}{2}\left(\frac12\mp\delta\right),\\
P_\pm(j,0,1)&=\frac{1-\epsilon}{4},&
P_\pm(j,0,0)&=\frac{1-\epsilon}{4},
&&j=1,2,
\end{aligned}
\]
and
\[
P_\pm(j,a,y)=0,
\qquad j>2,\quad a,y\in\{0,1\}.
\]
All occupied-cell entries are positive, and their sums are
\[
\sum_{a,y}P_\pm(j,a,y)
=
\frac{\epsilon}{2}+\frac{1-\epsilon}{2}
=\frac12,
\qquad j=1,2.
\]
Thus each table is a probability law. Its occupied-cell quantities are
\[
p_j(P_\pm)=\frac12,
\qquad
e_j(P_\pm)=\frac{\epsilon/2}{1/2}=\epsilon,
\qquad
\mu_{0j}(P_\pm)=\frac12,
\qquad
\mu_{1j}(P_\pm)=\frac12\pm\delta.
\]
Every other covariate mass is zero. Since
\(\epsilon\in[\epsilon,1-\epsilon]\), \cref{obj:def:model} gives
\[
P_+,P_-\in\mathcal M_{d,\epsilon}.
\]
The two occupied contributions in \cref{obj:def:target} give
\[
\tau(P_\pm)
=
\frac12(\pm\delta)+\frac12(\pm\delta)
=\pm\delta.
\]
Finally, summing over outcomes gives
\[
(P_\pm)_{XA}(j,a)=
\begin{cases}
\epsilon/2,&j\in\{1,2\},\ a=1,\\
(1-\epsilon)/2,&j\in\{1,2\},\ a=0,\\
0,&j>2,
\end{cases}
\]
which proves \((P_+)_{XA}=(P_-)_{XA}\). The null-cell assertions are
vacuous when \(d=2\).
% lean: CausalSmith.Stat.AnnotationRarearmFrontier.labelFloor_model
% lean: CausalSmith.Stat.AnnotationRarearmFrontier.labelFloor_target
% lean: CausalSmith.Stat.AnnotationRarearmFrontier.labelFloor_auxMarginal
% lean: CausalSmith.Stat.AnnotationRarearmFrontier.labelFloor_cellMass
% lean: CausalSmith.Stat.AnnotationRarearmFrontier.labelFloor_propensity
% lean: CausalSmith.Stat.AnnotationRarearmFrontier.labelFloor_outcomeMean

\item \textit{Exact divergence at separation \(1/8\).}
At \(\delta=1/8\), the treated success and failure masses in each
occupied cell under \(P_+\) are \(5\epsilon/16\) and
\(3\epsilon/16\), respectively; under \(P_-\) they are exchanged.
The control terms have likelihood ratio one, and the null cells have
zero contribution. Therefore, for each \(j=1,2\),
\[
\begin{aligned}
\sum_{a,y}
P_+(j,a,y)\log\frac{P_+(j,a,y)}{P_-(j,a,y)}
&=
\frac{5\epsilon}{16}\log\frac53
+
\frac{3\epsilon}{16}\log\frac35\\
&=
\frac{\epsilon}{8}\log\frac53,
\end{aligned}
\]
using \(\log(3/5)=-\log(5/3)\). Summing the two occupied-cell
contributions proves
\[
D(P_+\Vert P_-)
=
\sum_{z:\,P_+(z)>0}
P_+(z)\log\frac{P_+(z)}{P_-(z)}
=
\frac{\epsilon}{4}\log\frac53.
\]
% lean: CausalSmith.Stat.AnnotationRarearmFrontier.labelFloor_KL_eighth
\end{enumerate}
\end{proof}

\begin{proof}[Proof of \cref{obj:prop:benchmark-recovery}]
We first establish the quantitative experiment comparison, then prove the remaining assertions in their stated order. Throughout, the public indices satisfy the ranges in the statement. We use the cell quantities
\begingroup\scriptsize
\[
\begin{aligned}
P(j,a,y)&=P(X=j,A=a,Y=y),\\
s_{aj}(P)&=\sum_{y=0}^1P(j,a,y),
&
q_{aj}(P)&=P(j,a,1),\\
p_j(P)&=s_{0j}(P)+s_{1j}(P),
&
\mu_{aj}(P)&=
\begin{cases}
q_{aj}(P)/s_{aj}(P),&s_{aj}(P)>0,\\
0,&s_{aj}(P)=0,
\end{cases}
\end{aligned}
\qquad
j\in\{1,\ldots,d\},\quad a,y\in\{0,1\}.
\]
\endgroup
Since the conditional means belong to \([0,1]\), \cref{obj:def:target} gives
\[
-1=\sum_jp_j(P)(-1)
\le \tau(P)
\le \sum_jp_j(P)=1.
\]
The laws supplied by \cref{obj:thm:rare-label-floor} make the model class nonempty. Both rule classes contain the zero rule, and every clipped rule has squared loss between zero and four. Thus the suprema and infima used below are finite.
% lean: CausalSmith.Stat.AnnotationRarearmFrontier.ateFunctional_mem_Icc

\begin{enumerate}
\item \emph{Comparison of experiments.}
Fix an original-record rule \(T\). For every real table
\(s'=(s'_{aj})_{j,a}\in\mathbb R^{2d}\) and auxiliary array
\(\mathcal V=((X_i^{\mathrm{aux}},A_i^{\mathrm{aux}}))_{i=1}^m
\in(\{1,\ldots,d\}\times\{0,1\})^m\), define
\[
\omega_{s'}(\mathcal V)
=\prod_{i=1}^m s'_{A_i^{\mathrm{aux}},X_i^{\mathrm{aux}}},
\]
and, for every complete-record array \(\mathcal L\) and every
\(u_{\mathrm{seed}}\in\mathbb R\), define
\[
T^{\mathrm{av}}(\mathcal L,s',u_{\mathrm{seed}})
=
\operatorname{clip}_{[-1,1]}
\left(
\sum_{\mathcal V}
\omega_{s'}(\mathcal V)
T((\mathcal L,\mathcal V),u_{\mathrm{seed}})
\right).
\]
Here clipping is as in \cref{obj:synth_4}. The sum is finite, its weights are polynomials in the supplied entries, and clipping makes this a Borel rule with values in \([-1,1]\) on the entire real table space.

At the true table \(s(P)=(s_{aj}(P))_{j,a}\), product sampling gives
\[
\omega_{s(P)}(\mathcal V)
=P_{XA}^{\otimes m}(\{\mathcal V\}),
\qquad
\omega_{s(P)}(\mathcal V)\ge0,
\qquad
\sum_{\mathcal V}\omega_{s(P)}(\mathcal V)=1.
\]
Consequently the average already lies in \([-1,1]\), and convexity of squared loss yields
\[
\bigl(T^{\mathrm{av}}(\mathcal L,s(P),u_{\mathrm{seed}})
-\tau(P)\bigr)^2
\le
\sum_{\mathcal V}\omega_{s(P)}(\mathcal V)
\bigl(T((\mathcal L,\mathcal V),u_{\mathrm{seed}})
-\tau(P)\bigr)^2.
\]
Integrating over \(\mathcal L\) and the independent seed gives
\[
\begin{aligned}
&\mathbb E_{P^{\otimes n}\otimes\mathrm{Uniform}[0,1]}
\bigl[(T^{\mathrm{av}}(\mathcal L,s(P),U)-\tau(P))^2\bigr]\\
&\qquad\le
\mathbb E_{P^{\otimes n}\otimes P_{XA}^{\otimes m}
\otimes\mathrm{Uniform}[0,1]}
\bigl[(T((\mathcal L,\mathcal V),U)-\tau(P))^2\bigr].
\end{aligned}
\]
For \(m=0\), the auxiliary-array space contains the single empty array, whose weight is the empty product \(1\); the same argument applies. Taking suprema over \(P\) and infima over \(T\) proves
\[
R^{\mathrm K}(n,d,\epsilon)\le R(n,m,d,\epsilon).
\]
The zero rule and nonnegativity of squared loss also give
\[
0\le R^{\mathrm K}(n,d,\epsilon)
\le R(n,m,d,\epsilon)\le1,
\]
and hence
\[
0\le R(n,m,d,\epsilon)-R^{\mathrm K}(n,d,\epsilon)\le1.
\]
% lean: CausalSmith.Stat.AnnotationRarearmFrontier.comparison_difference_bounds

For the reverse quantitative bound, assume \(m\ge1\). Define the empirical and adjusted auxiliary tables by
\[
\begin{aligned}
\widetilde s_{aj}
&=\frac1m\sum_{i=1}^m
\mathbf1\{X_i^{\mathrm{aux}}=j,\ A_i^{\mathrm{aux}}=a\},\\
\widetilde p_j&=\widetilde s_{0j}+\widetilde s_{1j},\\
\widehat s_{1j}
&=\min\{(1-\epsilon)\widetilde p_j,
          \max\{\epsilon\widetilde p_j,\widetilde s_{1j}\}\},\\
\widehat s_{0j}&=\widetilde p_j-\widehat s_{1j}.
\end{aligned}
\]
Because \(0<\epsilon\le1/4\), the two clipping endpoints are ordered. Thus
\[
\widehat s_{aj}\ge0,\qquad
\sum_{j,a}\widehat s_{aj}=1,\qquad
\epsilon(\widehat s_{0j}+\widehat s_{1j})
\le\widehat s_{1j}
\le(1-\epsilon)(\widehat s_{0j}+\widehat s_{1j}).
\]
In an empty empirical cell, all four empirical and adjusted arm entries are zero. The adjustment is a Borel function of the auxiliary array.
% lean: CausalSmith.Stat.AnnotationRarearmFrontier.comparisonEmpiricalTable_probability
% lean: CausalSmith.Stat.AnnotationRarearmFrontier.comparisonEmpiricalTable_overlap

Fix any clipped Borel supplied-table rule \(T^{\mathrm K}\), and define
\[
T^{\mathrm{plug}}((\mathcal L,\mathcal V),u_{\mathrm{seed}})
=
T^{\mathrm K}(\mathcal L,\widehat s(\mathcal V),u_{\mathrm{seed}}).
\]
This is an admissible original-record rule, so
\[
R(n,m,d,\epsilon)
\le
\sup_{P\in\mathcal M_{d,\epsilon}}
\mathbb E_{P^{\otimes n}\otimes P_{XA}^{\otimes m}
\otimes\mathrm{Uniform}[0,1]}
\bigl[(T^{\mathrm{plug}}-\tau(P))^2\bigr].
\]
For later use, define its supplied-table counterpart's maximal risk by
\[
\mathscr R^{\mathrm K}(T^{\mathrm K})
=
\sup_{P'\in\mathcal M_{d,\epsilon}}
\mathbb E_{(P')^{\otimes n}\otimes\mathrm{Uniform}[0,1]}
\bigl[
(T^{\mathrm K}(\mathcal L,s(P'),U)-\tau(P'))^2
\bigr],
\qquad
0\le\mathscr R^{\mathrm K}(T^{\mathrm K})\le4.
\]
% lean: CausalSmith.Stat.AnnotationRarearmFrontier.comparison_minimax_le_known_add

We next bound the risk at a fixed supplied table. For real tables define
\[
\|s'-s''\|_1=\sum_{j=1}^d\sum_{a=0}^1|s'_{aj}-s''_{aj}|,
\qquad s',s''\in\mathbb R^{2d}.
\]
Fix \(P\in\mathcal M_{d,\epsilon}\) and any table \(s'\) satisfying
\[
s'_{aj}\ge0,\qquad \sum_{j,a}s'_{aj}=1,\qquad
\epsilon(s'_{0j}+s'_{1j})\le s'_{1j}
\le(1-\epsilon)(s'_{0j}+s'_{1j}).
\]
Define a replacement law by the atoms
\[
P^{s'}(j,a,1)=s'_{aj}\mu_{aj}(P),
\qquad
P^{s'}(j,a,0)=s'_{aj}(1-\mu_{aj}(P)).
\]
These atoms are nonnegative and sum to one. Its marginal and cell masses are
\[
s_{aj}(P^{s'})=s'_{aj},
\qquad
p_j(P^{s'})=s'_{0j}+s'_{1j}.
\]
On every occupied replacement cell, division of the displayed legality inequalities by its cell mass proves the overlap restriction. Thus
\[
P^{s'}\in\mathcal M_{d,\epsilon}.
\]
Moreover, both replacement arm masses are positive in an occupied cell, so
\[
\mu_{aj}(P^{s'})=\mu_{aj}(P)
\qquad\text{whenever }p_j(P^{s'})>0.
\]
In a null replacement cell the cell-weighted means vanish. Originally null arms have inherited mean zero by the stated convention.
% lean: CausalSmith.Stat.AnnotationRarearmFrontier.comparisonReplacementLaw_model
% lean: CausalSmith.Stat.AnnotationRarearmFrontier.comparisonReplacementLaw_weighted_mean

The original atoms factor in the same way:
\[
P(j,a,1)=s_{aj}(P)\mu_{aj}(P),
\qquad
P(j,a,0)=s_{aj}(P)(1-\mu_{aj}(P)).
\]
For a null original arm, both atoms are zero, which verifies these identities there as well. It follows that
\[
\begin{aligned}
\sum_{j,a,y}|P(j,a,y)-P^{s'}(j,a,y)|
&=\sum_{j,a}|s_{aj}(P)-s'_{aj}|
   \{\mu_{aj}(P)+1-\mu_{aj}(P)\}\\
&=\|s(P)-s'\|_1.
\end{aligned}
\]
For finite probability laws, use
\[
\operatorname{TV}(\nu_1,\nu_2)
=\frac12\sum_z|\nu_1(z)-\nu_2(z)|,
\]
where the sum ranges over their common finite atom space. Therefore
\[
\operatorname{TV}(P,P^{s'})
\le\frac12\|s(P)-s'\|_1.
\]
The replacement target is
\[
\tau(P^{s'})
=\sum_j(s'_{0j}+s'_{1j})
       \{\mu_{1j}(P)-\mu_{0j}(P)\},
\]
and the contrast bound \(|\mu_{1j}(P)-\mu_{0j}(P)|\le1\) gives
\[
\begin{aligned}
|\tau(P)-\tau(P^{s'})|
&\le
\sum_j|p_j(P)-(s'_{0j}+s'_{1j})|\\
&\le\sum_{j,a}|s_{aj}(P)-s'_{aj}|
=\|s(P)-s'\|_1.
\end{aligned}
\]
% lean: CausalSmith.Stat.AnnotationRarearmFrontier.comparisonReplacementLaw_tv
% lean: CausalSmith.Stat.AnnotationRarearmFrontier.comparisonReplacementLaw_target_error

The product-distance bound is
\[
\operatorname{TV}(\nu_1^{\otimes n},\nu_2^{\otimes n})
\le n\operatorname{TV}(\nu_1,\nu_2).
\]
Indeed, the empty products have distance zero, and the triangle inequality, inserting
\(\nu_2\otimes\nu_1^{\otimes k}\), gives
\[
\operatorname{TV}(\nu_1^{\otimes(k+1)},\nu_2^{\otimes(k+1)})
\le
\operatorname{TV}(\nu_1,\nu_2)
+\operatorname{TV}(\nu_1^{\otimes k},\nu_2^{\otimes k}),
\qquad k\ge0.
\]
Here a common independent probability factor preserves total variation: integration over its sections proves the upper inequality, and cylinder events prove the reverse inequality. In particular,
\[
\begin{aligned}
&\operatorname{TV}\!\left(
P^{\otimes n}\otimes\mathrm{Uniform}[0,1],
(P^{s'})^{\otimes n}\otimes\mathrm{Uniform}[0,1]\right)\\
&\qquad=
\operatorname{TV}(P^{\otimes n},(P^{s'})^{\otimes n})
\le \frac n2\|s(P)-s'\|_1.
\end{aligned}
\]
% lean: Causalean.Stat.Minimax.MomentMatchedMixture.tvDist_pi_iid_le
% lean: CausalSmith.Stat.AnnotationRarearmFrontier.labelFloor_tv_common_product

Squared loss is in \([0,4]\). After integrating over the common seed, summing the loss against the positive and negative atom-mass differences bounds its change in expectation by four times total variation. Hence
\[
\begin{aligned}
&\mathbb E_{P^{\otimes n}\otimes\mathrm{Uniform}[0,1]}
\bigl[(T^{\mathrm K}(\mathcal L,s',U)-\tau(P))^2\bigr]\\
&\quad\le
\mathbb E_{(P^{s'})^{\otimes n}\otimes\mathrm{Uniform}[0,1]}
\bigl[(T^{\mathrm K}(\mathcal L,s',U)-\tau(P))^2\bigr]
+2n\|s(P)-s'\|_1.
\end{aligned}
\]
For every realization, the rule and both targets lie in \([-1,1]\), so
\[
\begin{aligned}
&\left|
(T^{\mathrm K}(\mathcal L,s',U)-\tau(P))^2
-(T^{\mathrm K}(\mathcal L,s',U)-\tau(P^{s'}))^2
\right|\\
&\quad=
|\tau(P)-\tau(P^{s'})|
\left|2T^{\mathrm K}(\mathcal L,s',U)-\tau(P)-\tau(P^{s'})\right|\\
&\quad\le4|\tau(P)-\tau(P^{s'})|
\le4\|s(P)-s'\|_1.
\end{aligned}
\]
Integrating this inequality under the replacement law, whose true marginal is \(s'\), and combining the two bounds proves
\[
\begin{aligned}
&\mathbb E_{P^{\otimes n}\otimes\mathrm{Uniform}[0,1]}
\bigl[(T^{\mathrm K}(\mathcal L,s',U)-\tau(P))^2\bigr]\\
&\qquad\le
\mathscr R^{\mathrm K}(T^{\mathrm K})
+(2n+4)\|s(P)-s'\|_1.
\end{aligned}
\]
% lean: CausalSmith.Stat.AnnotationRarearmFrontier.comparisonFixedTableRisk_transport

It remains to control the expected adjusted-table error. Define
\[
\Delta_j^{\mathrm{aux}}
=
|\widetilde s_{0j}-s_{0j}(P)|
+|\widetilde s_{1j}-s_{1j}(P)|.
\]
The population overlap restriction gives
\[
\epsilon s_{0j}(P)-(1-\epsilon)s_{1j}(P)\le0,
\qquad
\epsilon s_{1j}(P)-(1-\epsilon)s_{0j}(P)\le0.
\]
If \(\widetilde s_{1j}\le\epsilon\widetilde p_j\), the treated entry moves to the lower endpoint and
\[
\begin{aligned}
|\widehat s_{1j}-\widetilde s_{1j}|
&=\epsilon\widetilde s_{0j}-(1-\epsilon)\widetilde s_{1j}\\
&\le
\epsilon|\widetilde s_{0j}-s_{0j}(P)|
+(1-\epsilon)|\widetilde s_{1j}-s_{1j}(P)|\\
&\le\Delta_j^{\mathrm{aux}}.
\end{aligned}
\]
Otherwise, if \((1-\epsilon)\widetilde p_j\le\widetilde s_{1j}\), it moves to the upper endpoint and
\[
\begin{aligned}
|\widehat s_{1j}-\widetilde s_{1j}|
&=\epsilon\widetilde s_{1j}-(1-\epsilon)\widetilde s_{0j}\\
&\le
\epsilon|\widetilde s_{1j}-s_{1j}(P)|
+(1-\epsilon)|\widetilde s_{0j}-s_{0j}(P)|\\
&\le\Delta_j^{\mathrm{aux}}.
\end{aligned}
\]
In the remaining case the entry is unchanged. The control entry moves by the opposite amount, and the triangle inequality consequently gives
\[
\begin{aligned}
&|\widehat s_{0j}-s_{0j}(P)|
+|\widehat s_{1j}-s_{1j}(P)|\\
&\qquad\le
\Delta_j^{\mathrm{aux}}
+2|\widehat s_{1j}-\widetilde s_{1j}|
\le3\Delta_j^{\mathrm{aux}}.
\end{aligned}
\]
Summing over cells proves
\[
\|\widehat s-s(P)\|_1
\le3\|\widetilde s-s(P)\|_1.
\]
% lean: CausalSmith.Stat.AnnotationRarearmFrontier.comparisonEmpiricalTable_l1_error

Under \(P_{XA}^{\otimes m}\), expanding the empirical-coordinate square into its diagonal and off-diagonal terms gives
\[
\begin{aligned}
\mathbb E\widetilde s_{aj}&=s_{aj}(P),\\
\mathbb E\widetilde s_{aj}^{\,2}
&=\frac{s_{aj}(P)}m+\frac{m-1}{m}s_{aj}(P)^2,\\
\mathbb E(\widetilde s_{aj}-s_{aj}(P))^2
&=\frac{s_{aj}(P)-s_{aj}(P)^2}{m}.
\end{aligned}
\]
Cauchy--Schwarz for the \(2d\) coordinates yields
\[
\|\widetilde s-s(P)\|_1^2
\le2d\sum_{j,a}(\widetilde s_{aj}-s_{aj}(P))^2.
\]
Using the coordinate identity and \(\sum_{j,a}s_{aj}(P)=1\), we obtain
\[
\begin{aligned}
\mathbb E\|\widetilde s-s(P)\|_1^2
&\le\frac{2d}{m}
\sum_{j,a}\{s_{aj}(P)-s_{aj}(P)^2\}\\
&\le\frac{2d}{m}.
\end{aligned}
\]
Nonnegative variance then gives
\[
\bigl(\mathbb E\|\widetilde s-s(P)\|_1\bigr)^2
\le\mathbb E\|\widetilde s-s(P)\|_1^2
\le\frac{2d}{m},
\]
and therefore
\[
\mathbb E\|\widehat s-s(P)\|_1
\le3\mathbb E\|\widetilde s-s(P)\|_1
\le3\sqrt{\frac{2d}{m}}.
\]
% lean: CausalSmith.Stat.AnnotationRarearmFrontier.comparisonEmpiricalTable_expected_l1

Product sampling expresses the plug-in risk as the auxiliary average of its fixed-table risk. Applying the fixed-table bound at \(s'=\widehat s(\mathcal V)\) gives
\[
\begin{aligned}
&\mathbb E_{P^{\otimes n}\otimes P_{XA}^{\otimes m}
\otimes\mathrm{Uniform}[0,1]}
\bigl[(T^{\mathrm{plug}}-\tau(P))^2\bigr]\\
&\qquad\le
\mathscr R^{\mathrm K}(T^{\mathrm K})
+(2n+4)\mathbb E_{P_{XA}^{\otimes m}}\|\widehat s-s(P)\|_1\\
&\qquad\le
\mathscr R^{\mathrm K}(T^{\mathrm K})
+6(n+2)\sqrt{\frac{2d}{m}}.
\end{aligned}
\]
This holds uniformly in \(P\). Consequently, for every admissible \(T^{\mathrm K}\),
\[
R(n,m,d,\epsilon)-6(n+2)\sqrt{\frac{2d}{m}}
\le\mathscr R^{\mathrm K}(T^{\mathrm K}).
\]
Taking the infimum over all supplied-table rules, as in \cref{obj:def:known-risk}, proves
\[
R(n,m,d,\epsilon)
\le R^{\mathrm K}(n,d,\epsilon)
+6(n+2)\sqrt{\frac{2d}{m}}.
\]
Combining this with the earlier difference bounds establishes
\[
0\le R(n,m,d,\epsilon)-R^{\mathrm K}(n,d,\epsilon)
\le\min\left\{1,6(n+2)\sqrt{\frac{2d}{m}}\right\}.
\]
% lean: CausalSmith.Stat.AnnotationRarearmFrontier.comparisonPluginRule_risk_le
% lean: CausalSmith.Stat.AnnotationRarearmFrontier.comparison_minimax_le_known_add

\item \emph{Fixed-overlap benchmark.}
Fix \(0<\epsilon\le1/4\), and define
\[
N=n+m,\qquad
\ell=\log(\exp(1)+n\epsilon),\qquad
\Lambda_n=\log(\exp(1)n)=1+\log n,
\]
\[
K_{\log}=2+\log(1/\epsilon),
\qquad
B_c=\max\left\{1,
\left(\max\{\epsilon^{-1},K_{\log}/\epsilon\}\right)^2\right\}.
\]
These are positive, with \(B_c\ge1\). We first compare the logarithms. Since \(n\ge1\) and \(n\epsilon>0\),
\[
1\le\ell.
\]
The elementary inequality \(\exp(1)>2\) implies
\(\exp(1)+1\le\exp(2)\). Hence
\[
\exp(1)+n\epsilon
\le(\exp(1)+1)n
\le\exp(2)n,
\]
so
\[
\ell\le2+\log n\le2\Lambda_n.
\]
Also,
\[
n\le\frac{\exp(1)+n\epsilon}{\epsilon}
\quad\Longrightarrow\quad
\log n\le\ell+\log(1/\epsilon).
\]
Since \(\ell\ge1\) and \(\log(1/\epsilon)\ge0\), this gives
\[
\Lambda_n
\le1+\ell+\log(1/\epsilon)
\le K_{\log}\ell.
\]
% lean: CausalSmith.Stat.AnnotationRarearmFrontier.benchmark_log_comparison

The upper logarithmic bound and \(\epsilon\le1/4\) imply
\[
\epsilon\ell\le2\epsilon\Lambda_n\le\Lambda_n.
\]
Together with \(\Lambda_n\le K_{\log}\ell\), positivity of the denominators yields
\[
\frac d{N\Lambda_n}
\le\frac d{N\epsilon\ell}
\le\frac{K_{\log}}{\epsilon}\frac d{N\Lambda_n}.
\]
The chosen constant satisfies
\[
\epsilon^{-1}
\le\max\{\epsilon^{-1},K_{\log}/\epsilon\}
\le\left(\max\{\epsilon^{-1},K_{\log}/\epsilon\}\right)^2
\le B_c,
\qquad
(K_{\log}/\epsilon)^2\le B_c,
\]
where the middle inequality uses \(\epsilon^{-1}\ge1\). Define the uncapped quantities
\[
z_{\mathrm{fix}}
=\frac1n+\left(\frac d{N\Lambda_n}\right)^2,
\qquad
z_{\mathrm{rate}}
=\frac1{n\epsilon}+\left(\frac d{N\epsilon\ell}\right)^2.
\]
The preceding bounds give
\[
z_{\mathrm{fix}}\le z_{\mathrm{rate}}
\le
B_c\left\{\frac1n+\left(\frac d{N\Lambda_n}\right)^2\right\}
=B_cz_{\mathrm{fix}}.
\]
Finally, because \(B_c\ge1\),
\[
\begin{aligned}
f(n,m,d)
&=\min\{1,z_{\mathrm{fix}}\}
\le\min\{1,z_{\mathrm{rate}}\}
=r(n,m,d,\epsilon),\\
r(n,m,d,\epsilon)
&\le\min\{1,B_cz_{\mathrm{fix}}\}
\le B_c\min\{1,z_{\mathrm{fix}}\}
=B_cf(n,m,d).
\end{aligned}
\]
Thus the first assertion holds with \(a=1\) and the displayed \(B_c\), which depends only on \(\epsilon\).
% lean: CausalSmith.Stat.AnnotationRarearmFrontier.benchmark_fixed_overlap_comparison

\item \emph{Rare-label benchmarks.}
Let \(c_{\mathrm{lab}}>0\) be the numerical constant supplied by
\cref{obj:thm:rare-label-floor}, and let \(C_{\mathrm B}>0\) be the numerical constant supplied by \cref{obj:thm:uniform-baseline}. Choose
\[
c=c_{\mathrm{lab}},
\qquad
C=\max\{5C_{\mathrm B},2\}.
\]
The rare-label bounds give
\[
c\,b(n,\epsilon)\le R(n,m,2,\epsilon),
\qquad
c\,b(n,\epsilon)\le R^{\mathrm K}(n,d,\epsilon).
\]
The baseline is an admissible clipped Borel rule by
\cref{obj:thm:uniform-baseline}. Taking its worst-case risk and then the minimax infimum gives
\[
R(n,m,2,\epsilon)
\le
C_{\mathrm B}\min\left\{
1,\frac1{n\epsilon}
+\left(\frac2{(n+m)\epsilon}\right)^2
\right\}.
\]
% lean: CausalSmith.Stat.AnnotationRarearmFrontier.benchmark_minimax_le_baseline

Set
\[
S=n\epsilon.
\]
Since \((n+m)\epsilon\ge S>0\), if \(S\ge1\), then
\[
\frac1S+\left(\frac2{(n+m)\epsilon}\right)^2
\le\frac1S+\frac4{S^2}
\le\frac5S,
\qquad
b(n,\epsilon)=S^{-1}.
\]
If \(S<1\), then \(b(n,\epsilon)=1\), and the capped baseline quantity is at most \(1\le5b(n,\epsilon)\). Thus in both cases
\[
R(n,m,2,\epsilon)
\le5C_{\mathrm B}b(n,\epsilon)
\le Cb(n,\epsilon).
\]
% lean: CausalSmith.Stat.AnnotationRarearmFrontier.benchmark_binary_baseline_rate

For the supplied-table upper bound, define, for every real table \(s'\) and every observation \((j,a,y)\),
\[
G^{\mathrm K}_{s'}(j,a,y)
=(2a-1)y\,\frac{s'_{0j}+s'_{1j}}{s'_{aj}},
\qquad
s'\in\mathbb R^{2d},\quad
j\in\{1,\ldots,d\},\quad a,y\in\{0,1\}.
\]
Every quotient with zero denominator in this construction is assigned zero. Define the total rule
\[
T^{\mathrm{IP}}(\mathcal L,s',u_{\mathrm{seed}})
=
\begin{cases}
0,&S<1,\\[2pt]
\displaystyle
\operatorname{clip}_{[-1,1]}
\left(\frac1n\sum_{i=1}^n
G^{\mathrm K}_{s'}(X_i,A_i,Y_i)\right),&S\ge1,
\end{cases}
\qquad u_{\mathrm{seed}}\in\mathbb R.
\]
The totalized quotients are Borel functions, and clipping makes the rule admissible on every ambient table.

At the true table, finite summation gives
\[
\mathbb E_PG^{\mathrm K}_{s(P)}(X,A,Y)
=\sum_jp_j(P)\{\mu_{1j}(P)-\mu_{0j}(P)\}
=\tau(P).
\]
For each arm, the second-moment contribution obeys
\[
q_{aj}(P)\left(\frac{p_j(P)}{s_{aj}(P)}\right)^2
\le\frac{p_j(P)}{\epsilon}.
\]
Indeed, when \(s_{aj}(P)=0\), the left side is zero. Otherwise,
\(q_{aj}(P)\le s_{aj}(P)\) and overlap gives
\(s_{aj}(P)\ge\epsilon p_j(P)\), so
\[
\begin{aligned}
q_{aj}(P)\left(\frac{p_j(P)}{s_{aj}(P)}\right)^2
&\le s_{aj}(P)\left(\frac{p_j(P)}{s_{aj}(P)}\right)^2\\
&=p_j(P)\frac{p_j(P)}{s_{aj}(P)}
\le\frac{p_j(P)}{\epsilon}.
\end{aligned}
\]
Summing both arms and all cells therefore yields
\[
\mathbb E_P\bigl[G^{\mathrm K}_{s(P)}(X,A,Y)^2\bigr]
=
\sum_{j,a}q_{aj}(P)
\left(\frac{p_j(P)}{s_{aj}(P)}\right)^2
\le\frac2\epsilon.
\]
% lean: CausalSmith.Stat.AnnotationRarearmFrontier.knownMarginalScore_mean
% lean: CausalSmith.Stat.AnnotationRarearmFrontier.knownMarginalScore_second_moment

Define the un-clipped sample mean by
\[
\overline G^{\mathrm K}_P(\mathcal L)
=\frac1n\sum_{i=1}^nG^{\mathrm K}_{s(P)}(X_i,A_i,Y_i).
\]
Independence of the complete records gives
\[
\begin{aligned}
\mathbb E_{P^{\otimes n}}
\bigl[(\overline G^{\mathrm K}_P-\tau(P))^2\bigr]
&=\operatorname{Var}_{P^{\otimes n}}(\overline G^{\mathrm K}_P)\\
&=\frac1n\operatorname{Var}_P(G^{\mathrm K}_{s(P)})\\
&\le\frac1n\mathbb E_P[(G^{\mathrm K}_{s(P)})^2]
\le\frac2{n\epsilon}.
\end{aligned}
\]
Clipping satisfies
\[
\bigl(\operatorname{clip}_{[-1,1]}(x_{\mathrm{mean}})
-x_{\mathrm{target}}\bigr)^2
\le(x_{\mathrm{mean}}-x_{\mathrm{target}})^2,
\qquad
x_{\mathrm{mean}}\in\mathbb R,\quad x_{\mathrm{target}}\in[-1,1].
\]
For \(x_{\mathrm{mean}}<-1\), moving it to \(-1\) decreases its distance to the target; for \(x_{\mathrm{mean}}>1\), moving it to \(1\) does the same; and on \([-1,1]\) clipping is the identity. Consequently, when \(S\ge1\), the supplied-table rule has risk at most \(2/S=2b(n,\epsilon)\). When \(S<1\), its zero branch has risk
\[
\tau(P)^2\le1\le2b(n,\epsilon).
\]
Taking suprema and the supplied-table rule infimum proves
\[
R^{\mathrm K}(n,d,\epsilon)
\le2b(n,\epsilon)\le Cb(n,\epsilon).
\]
The chosen \(c,C\) are positive numerical constants common to both benchmark comparisons.
% lean: CausalSmith.Stat.AnnotationRarearmFrontier.knownMarginalMean_risk
% lean: CausalSmith.Stat.AnnotationRarearmFrontier.knownMarginal_clip_loss
% lean: CausalSmith.Stat.AnnotationRarearmFrontier.knownMarginalRisk_le_labelBenchmark

\item \emph{Limit along auxiliary sample sizes.}
Fix \(n,d,\epsilon\) in the stated range. Along integer \(m\to\infty\),
\[
m^{-1}\longrightarrow0,
\qquad
\frac{2d}{m}\longrightarrow0,
\qquad
6(n+2)\sqrt{\frac{2d}{m}}\longrightarrow0,
\]
where the last implication uses continuity of the square root at zero. For every \(m\ge1\), the comparison already established gives
\[
0\le R(n,m,d,\epsilon)-R^{\mathrm K}(n,d,\epsilon)
\le6(n+2)\sqrt{\frac{2d}{m}}.
\]
The squeeze theorem therefore yields
\[
R(n,m,d,\epsilon)-R^{\mathrm K}(n,d,\epsilon)\longrightarrow0.
\]
Adding the fixed value \(R^{\mathrm K}(n,d,\epsilon)\) proves the asserted limit.
% lean: CausalSmith.Stat.AnnotationRarearmFrontier.benchmark_comparison_error_tendsto
% lean: CausalSmith.Stat.AnnotationRarearmFrontier.benchmark_limit_of_comparison

\item \emph{Bounded rare-label scale.}
Fix \(M_0>0\), and use the numerical constant \(c_{\mathrm{lab}}>0\) from
\cref{obj:thm:rare-label-floor}. Define
\[
\kappa=\frac{c_{\mathrm{lab}}}{\max\{1,M_0\}}>0.
\]
If \(S=n\epsilon\le M_0\), then \(S>0\) and
\[
\frac1{\max\{1,M_0\}}\le1,
\qquad
\frac1{\max\{1,M_0\}}\le\frac1S.
\]
Taking the minimum of the two right sides gives
\[
\frac1{\max\{1,M_0\}}
\le\min\{1,S^{-1}\}=b(n,\epsilon).
\]
The rare-label risk floor now implies
\[
\kappa
\le c_{\mathrm{lab}}b(n,\epsilon)
\le R(n,m,d,\epsilon).
\]
This constant depends only on \(M_0\), as required.
% lean: CausalSmith.Stat.AnnotationRarearmFrontier.benchmark_label_floor_on_bounded_scale
\end{enumerate}
\end{proof}

\begin{proof}[Proof of \cref{obj:thm:uniform-baseline}]
For admissible \(n,m,d,\epsilon\) and \(P\in\mathcal M_{d,\epsilon}\), define
\[
\rho_{\mathrm B}
=
\frac{1}{n\epsilon}
+
\left(\frac{d}{(n+m)\epsilon}\right)^2,
\qquad
\mathcal R_{\mathrm B}(P)
=
\mathbb E_P\!\left[
\bigl(T^{\mathrm B}(\mathcal D)-\tau(P)\bigr)^2
\right].
\]
The conditional outcome means belong to \([0,1]\), while the covariate masses are nonnegative and sum to one. Thus \cref{obj:def:target} gives
\[
-1
=
\sum_{j=1}^d p_j(P)(-1)
\le \tau(P)
\le
\sum_{j=1}^d p_j(P)
=
1.
\]
We first establish an uncapped risk bound for \(n\ge6\).

\begin{enumerate}
\item
Set, as in \cref{obj:def:baseline},
\[
h_{\mathrm o}=\lfloor n/2\rfloor,
\qquad
h_{\mathrm m}=n-h_{\mathrm o}+m,
\qquad
u_{\mathrm B}=h_{\mathrm o}/8,
\qquad
t_{\mathrm B}=h_{\mathrm m}/8.
\]
Both pool lengths and both prefix means are positive. By
\cref{obj:ass:data-product}, the outcome pool and the marginal pool are independent, with respective laws
\[
P^{\otimes h_{\mathrm o}}
\quad\text{and}\quad
P_{XA}^{\otimes h_{\mathrm m}}.
\]
Indeed, the marginal pool uses projections of complete records disjoint from the outcome pool, followed by the independent auxiliary records.

Extend these pools by independent iid continuations, and write
\[
\mathcal S_{\mathrm o}
=
\bigl((X_i^{\mathrm o},A_i^{\mathrm o},Y_i^{\mathrm o})\bigr)_{i\ge1}
\sim P^{\otimes\mathbb N},
\qquad
\mathcal S_{\mathrm m}
=
\bigl((X_i^{\mathrm m},A_i^{\mathrm m})\bigr)_{i\ge1}
\sim P_{XA}^{\otimes\mathbb N}.
\]
Independently of these sequences, take independent prefix lengths
\[
J_{\mathrm o}\sim\operatorname{Poi}(u_{\mathrm B}),
\qquad
J_{\mathrm m}\sim\operatorname{Poi}(t_{\mathrm B}),
\]
where \(\operatorname{Poi}(\lambda)\) is the Poisson probability law with mean \(\lambda\). Define the complete-atom and marginal-atom counts by
\[
N_{jay}^{\mathrm o}
=
\sum_{i=1}^{J_{\mathrm o}}
\mathbf1\{X_i^{\mathrm o}=j,A_i^{\mathrm o}=a,Y_i^{\mathrm o}=y\},
\qquad
N_{ja}^{\mathrm m}
=
\sum_{i=1}^{J_{\mathrm m}}
\mathbf1\{X_i^{\mathrm m}=j,A_i^{\mathrm m}=a\},
\]
for \(j\in\{1,\ldots,d\}\) and \(a,y\in\{0,1\}\); empty sums equal zero.

For a nonnegative integer array
\(\boldsymbol c=(c_{jay})\), put
\[
|\boldsymbol c|=\sum_{j,a,y}c_{jay}.
\]
Conditioning on the outcome-prefix length gives its Poisson probability multiplied by the multinomial probability:
\[
\begin{aligned}
\Pr\!\left((N_{jay}^{\mathrm o})=\boldsymbol c\right)
&=
e^{-u_{\mathrm B}}
\frac{u_{\mathrm B}^{|\boldsymbol c|}}{|\boldsymbol c|!}
\frac{|\boldsymbol c|!}{\prod_{j,a,y}c_{jay}!}
\prod_{j,a,y}
P(X=j,A=a,Y=y)^{c_{jay}}\\
&=
\prod_{j,a,y}
e^{-u_{\mathrm B}P(X=j,A=a,Y=y)}
\frac{
\bigl(u_{\mathrm B}P(X=j,A=a,Y=y)\bigr)^{c_{jay}}
}{c_{jay}!}.
\end{aligned}
\]
Here \(0^0=1\). The same calculation for a nonnegative integer array
\(\boldsymbol k=(k_{ja})\) gives
\[
\Pr\!\left((N_{ja}^{\mathrm m})=\boldsymbol k\right)
=
\prod_{j,a}
e^{-t_{\mathrm B}s_{aj}(P)}
\frac{\bigl(t_{\mathrm B}s_{aj}(P)\bigr)^{k_{ja}}}{k_{ja}!}.
\]
Consequently, the atom counts within each pool are independent Poisson variables, and independence of the pools gives independence across the two collections.

In particular, define
\[
Z_{aj}^{\mathrm B}=N_{ja1}^{\mathrm o},
\qquad
K_{aj}^{\mathrm B}=N_{ja}^{\mathrm m}.
\]
Then
\[
Z_{aj}^{\mathrm B}\sim
\operatorname{Poi}\!\bigl(u_{\mathrm B}q_{aj}(P)\bigr),
\qquad
K_{aj}^{\mathrm B}\sim
\operatorname{Poi}\!\bigl(t_{\mathrm B}s_{aj}(P)\bigr).
\]
For each fixed arm \(a\), the collection
\[
\bigl(Z_{aj}^{\mathrm B},K_{aj}^{\mathrm B},
K_{1-a,j}^{\mathrm B}\bigr)_{j=1}^d
\]
is mutually independent. These conclusions include zero-mass cells, whose counts equal zero almost surely.
% lean: CausalSmith.Stat.AnnotationRarearmFrontier.baseline_prefix_counts_law
% lean: CausalSmith.Stat.AnnotationRarearmFrontier.baseline_prefix_arm_counts_independent

\item
Define the arm totals, their population targets, and the clipped Poisson-prefix contrast by
\[
H_{\mathrm{sum},a}^{\mathrm B}
=
\sum_{j=1}^d
\frac{Z_{aj}^{\mathrm B}}{u_{\mathrm B}}
\left(1+\frac{K_{1-a,j}^{\mathrm B}}{K_{aj}^{\mathrm B}+1}\right),
\qquad
\psi_a
=
\sum_{j=1}^d p_j(P)\mu_{aj}(P),
\]
\[
T_{\mathrm{Pois}}^{\mathrm B}
=
\max\!\left\{-1,\,
\min\!\left\{1,\,
H_{\mathrm{sum},1}^{\mathrm B}
-
H_{\mathrm{sum},0}^{\mathrm B}
\right\}\right\}.
\]
The count laws and independence just established discharge the hypotheses of
\cref{obj:lem:inverse-count-arm-risk}. Let \(C_{\mathrm{arm}}>0\) be its universal variance constant. For each \(a\in\{0,1\}\), it supplies
\[
\left|
\mathbb E H_{\mathrm{sum},a}^{\mathrm B}-\psi_a
\right|
\le
\min\!\left\{
1,\frac{d}{\exp(1)t_{\mathrm B}\epsilon}
\right\},
\]
\[
\operatorname{Var}\!\left(H_{\mathrm{sum},a}^{\mathrm B}\right)
\le
C_{\mathrm{arm}}
\left(
\frac{1}{u_{\mathrm B}\epsilon}
+
\frac{1}{t_{\mathrm B}\epsilon}
\right).
\]
The finite arm sums are square-integrable, using Poisson moments and the positive denominators. Expanding squared loss around the mean therefore yields
\[
\begin{aligned}
\mathbb E\!\left[
\bigl(H_{\mathrm{sum},a}^{\mathrm B}-\psi_a\bigr)^2
\right]
&=
\operatorname{Var}\!\left(H_{\mathrm{sum},a}^{\mathrm B}\right)
+
\left(\mathbb E H_{\mathrm{sum},a}^{\mathrm B}-\psi_a\right)^2\\
&\le
C_{\mathrm{arm}}
\left(
\frac{1}{u_{\mathrm B}\epsilon}
+
\frac{1}{t_{\mathrm B}\epsilon}
\right)
+
\left(\frac{d}{\exp(1)t_{\mathrm B}\epsilon}\right)^2.
\end{aligned}
\]
By \cref{obj:def:target},
\[
\tau(P)=\psi_1-\psi_0.
\]
Projection onto \([-1,1]\) decreases squared distance to \(\tau(P)\). Moreover,
\[
\begin{aligned}
\bigl(T_{\mathrm{Pois}}^{\mathrm B}-\tau(P)\bigr)^2
&\le
\bigl(H_{\mathrm{sum},1}^{\mathrm B}
-H_{\mathrm{sum},0}^{\mathrm B}-\tau(P)\bigr)^2\\
&\le
2\bigl(H_{\mathrm{sum},1}^{\mathrm B}-\psi_1\bigr)^2
+
2\bigl(H_{\mathrm{sum},0}^{\mathrm B}-\psi_0\bigr)^2,
\end{aligned}
\]
where the second inequality follows by expanding the square and using the nonnegativity of
\[
\left[
\bigl(H_{\mathrm{sum},1}^{\mathrm B}-\psi_1\bigr)
+
\bigl(H_{\mathrm{sum},0}^{\mathrm B}-\psi_0\bigr)
\right]^2.
\]
Set
\[
C_{\mathrm{prefix}}=4C_{\mathrm{arm}}>0.
\]
Taking expectations gives
\[
\mathbb E\!\left[
\bigl(T_{\mathrm{Pois}}^{\mathrm B}-\tau(P)\bigr)^2
\right]
\le
C_{\mathrm{prefix}}
\left(
\frac{1}{u_{\mathrm B}\epsilon}
+
\frac{1}{t_{\mathrm B}\epsilon}
\right)
+
4\left(\frac{d}{\exp(1)t_{\mathrm B}\epsilon}\right)^2.
\]
% lean: CausalSmith.Stat.AnnotationRarearmFrontier.baseline_poisson_arm_mse
% lean: CausalSmith.Stat.AnnotationRarearmFrontier.baseline_poisson_contrast_risk
% lean: CausalSmith.Stat.AnnotationRarearmFrontier.baseline_ordered_prefix_risk

\item
Define the Poisson risk allowance and the exponential overflow allowance by
\[
\mathcal A_{\mathrm B}
=
C_{\mathrm{prefix}}
\left(
\frac{1}{u_{\mathrm B}\epsilon}
+
\frac{1}{t_{\mathrm B}\epsilon}
\right)
+
4\left(\frac{d}{\exp(1)t_{\mathrm B}\epsilon}\right)^2,
\qquad
\mathcal E_{\mathrm B}
=
e^{-h_{\mathrm o}}+e^{-h_{\mathrm m}}.
\]
Both are nonnegative. The clipped prefix statistic is measurable and takes values in \([-1,1]\), including on empty prefixes. The target also belongs to this interval. Thus the two positive, independent iid pools satisfy the hypotheses of
\cref{obj:lem:finite-prefix-transfer}.

More explicitly, writing \(\mathcal S^{1:h}\) for the first \(h\) records of a sequence, define
\[
T_{\mathrm{fin}}^{\mathrm B}
=
\mathbb E\!\left[
T_{\mathrm{Pois}}^{\mathrm B}
\mathbf1\{J_{\mathrm o}\le h_{\mathrm o},\,
J_{\mathrm m}\le h_{\mathrm m}\}
\,\middle|\,
\mathcal S_{\mathrm o}^{1:h_{\mathrm o}},
\mathcal S_{\mathrm m}^{1:h_{\mathrm m}}
\right].
\]
Expanding this conditional expectation over the independent Poisson lengths gives exactly the finite average in \cref{obj:def:baseline}, with zero output on overflow. Hence, after forming the pools from the original records,
\[
T^{\mathrm B}(\mathcal D)=T_{\mathrm{fin}}^{\mathrm B},
\qquad
\mathcal R_{\mathrm B}(P)
=
\mathbb E\!\left[
\bigl(T_{\mathrm{fin}}^{\mathrm B}-\tau(P)\bigr)^2
\right].
\]
Integration over the independent seed leaves this risk unchanged.

The transfer bound in \cref{obj:lem:finite-prefix-transfer}, with its nonnegative overflow allowance enlarged by a factor of four, now gives
\[
\begin{aligned}
\mathcal R_{\mathrm B}(P)
&\le
\mathbb E\!\left[
\bigl(T_{\mathrm{Pois}}^{\mathrm B}-\tau(P)\bigr)^2
\right]
+
4\mathcal E_{\mathrm B}\\
&\le
C_{\mathrm{prefix}}
\left(
\frac{1}{u_{\mathrm B}\epsilon}
+
\frac{1}{t_{\mathrm B}\epsilon}
\right)
+
4\left(\frac{d}{\exp(1)t_{\mathrm B}\epsilon}\right)^2
+
4\left(e^{-h_{\mathrm o}}+e^{-h_{\mathrm m}}\right)\\
&=
\mathcal A_{\mathrm B}+4\mathcal E_{\mathrm B}.
\end{aligned}
\]
% lean: CausalSmith.Stat.AnnotationRarearmFrontier.baseline_ruleRisk_eq_fixed_average
% lean: Causalean.Stat.FiniteRaoBlackwell.IndependentPoissonPrefix.sqRisk_independentPrefixRaoBlackwellStatistic_le

\item
The integer pool sizes satisfy
\[
n\le3\lfloor n/2\rfloor,
\qquad
n+m\le2\bigl(n-\lfloor n/2\rfloor+m\bigr),
\]
because \(n\ge6\) and \(m\ge0\). Consequently,
\[
h_{\mathrm o}\ge\frac n3,
\qquad
h_{\mathrm m}\ge\frac{n+m}{2},
\]
and therefore
\[
\frac{1}{u_{\mathrm B}\epsilon}
\le\frac{24}{n\epsilon},
\qquad
\frac{1}{t_{\mathrm B}\epsilon}
\le\frac{16}{n\epsilon}.
\]
Also, \(\exp(1)\ge1\) and \(t_{\mathrm B}\ge(n+m)/16\), so
\[
\frac{d}{\exp(1)t_{\mathrm B}\epsilon}
\le
16\frac{d}{(n+m)\epsilon}.
\]

For the overflow allowance, both pool sizes are at least \(n/3\), giving
\[
\mathcal E_{\mathrm B}\le2e^{-n/3}.
\]
The elementary bound \(x e^{-x}\le e^{-1}\), obtained by maximizing \(x e^{-x}\), together with \(\epsilon\le1/4\), yields
\[
(n\epsilon)\,2e^{-n/3}
\le
\frac n2e^{-n/3}
=
\frac32\left(\frac n3e^{-n/3}\right)
\le
\frac32e^{-1}
\le2.
\]
Since \(n\epsilon>0\), this proves
\[
\mathcal E_{\mathrm B}\le\frac{2}{n\epsilon}.
\]
Combining these estimates gives
\[
\begin{aligned}
\mathcal A_{\mathrm B}+\mathcal E_{\mathrm B}
&\le
\frac{40C_{\mathrm{prefix}}+2}{n\epsilon}
+
1024\left(\frac{d}{(n+m)\epsilon}\right)^2\\
&\le
(40C_{\mathrm{prefix}}+1026)\rho_{\mathrm B}.
\end{aligned}
\]
% lean: CausalSmith.Stat.AnnotationRarearmFrontier.baseline_pool_sizes
% lean: CausalSmith.Stat.AnnotationRarearmFrontier.baseline_pool_overflow_bound
% lean: CausalSmith.Stat.AnnotationRarearmFrontier.baseline_pool_rate_bound

\item
Define
\[
C_{\mathrm{large}}
=
4(40C_{\mathrm{prefix}}+1026)>0.
\]
Using \(\mathcal A_{\mathrm B}\ge0\), the preceding bounds imply, for every \(n\ge6\),
\[
\begin{aligned}
\mathcal R_{\mathrm B}(P)
&\le
\mathcal A_{\mathrm B}+4\mathcal E_{\mathrm B}\\
&\le
4(\mathcal A_{\mathrm B}+\mathcal E_{\mathrm B})\\
&\le
C_{\mathrm{large}}\rho_{\mathrm B}.
\end{aligned}
\]
The constants used here are universal.
% lean: CausalSmith.Stat.AnnotationRarearmFrontier.uniform_baseline

\item
Choose the final constant
\[
C=\max\{C_{\mathrm{large}},4\}>0.
\]
For all admissible sample sizes, the rule in \cref{obj:def:baseline} is measurable: the small-sample branch is constant, and the remaining branch is a finite sum of measurable count functions with positive denominators.

To check its range directly, for \(n\ge6\) define
\[
w_{\mathrm B}(h',h'')
=
e^{-u_{\mathrm B}-t_{\mathrm B}}
\frac{u_{\mathrm B}^{h'}}{h'!}
\frac{t_{\mathrm B}^{h''}}{h''!},
\qquad
0\le h'\le h_{\mathrm o},\quad
0\le h''\le h_{\mathrm m}.
\]
These weights are nonnegative and satisfy
\[
\sum_{h'=0}^{h_{\mathrm o}}
\sum_{h''=0}^{h_{\mathrm m}}
w_{\mathrm B}(h',h'')
=
\Pr(J_{\mathrm o}\le h_{\mathrm o})\,
\Pr(J_{\mathrm m}\le h_{\mathrm m})
\le1.
\]
Every clipped prefix value belongs to \([-1,1]\). Thus
\[
-1
\le
-\sum_{h',h''}w_{\mathrm B}(h',h'')
\le
T^{\mathrm B}(\mathcal D)
\le
\sum_{h',h''}w_{\mathrm B}(h',h'')
\le1.
\]
For \(n<6\), the rule equals zero, so the same range conclusion holds.
% lean: CausalSmith.Stat.AnnotationRarearmFrontier.baselineEstimator_measurable
% lean: CausalSmith.Stat.AnnotationRarearmFrontier.baselineEstimator_mem_Icc

\item
Suppose \(1\le n<6\). The zero rule has risk
\[
\mathcal R_{\mathrm B}(P)=\tau(P)^2\le1.
\]
Moreover,
\[
0<n\epsilon\le\frac54,
\qquad
\frac{1}{n\epsilon}\ge\frac45,
\qquad
\min\{1,\rho_{\mathrm B}\}\ge\frac45.
\]
It follows that
\[
\mathcal R_{\mathrm B}(P)
\le1
\le2\min\{1,\rho_{\mathrm B}\}
\le C\min\{1,\rho_{\mathrm B}\},
\]
using \(C\ge4\) and the nonnegativity of the capped rate.
% lean: CausalSmith.Stat.AnnotationRarearmFrontier.baseline_small_sample_risk
% lean: CausalSmith.Stat.AnnotationRarearmFrontier.baseline_small_sample_rate

\item
Finally, suppose \(n\ge6\). If \(\rho_{\mathrm B}\ge1\), the estimator and target both belong to \([-1,1]\), so
\[
\bigl(T^{\mathrm B}(\mathcal D)-\tau(P)\bigr)^2\le4
\]
for every realization. Hence
\[
\mathcal R_{\mathrm B}(P)
\le4
\le C
=
C\min\{1,\rho_{\mathrm B}\}.
\]
If \(\rho_{\mathrm B}<1\), the uncapped bound established above gives
\[
\mathcal R_{\mathrm B}(P)
\le C_{\mathrm{large}}\rho_{\mathrm B}
\le C\rho_{\mathrm B}
=
C\min\{1,\rho_{\mathrm B}\}.
\]
These cases cover every admissible experiment and every
\(P\in\mathcal M_{d,\epsilon}\), with the same positive universal constant \(C\).
% lean: CausalSmith.Stat.AnnotationRarearmFrontier.baseline_risk_le_four
% lean: CausalSmith.Stat.AnnotationRarearmFrontier.uniform_baseline
\end{enumerate}
\end{proof}

\begin{proof}[Proof of \cref{obj:thm:annotation-resource-phases}]
\begin{enumerate}
\item \emph{Uniform risk comparison.}
By \cref{obj:thm:uniform-annotation-frontier}, there are numerical constants \(0<c_{\mathrm F}\le C_{\mathrm F}<\infty\) such that, throughout the allowed public index range,
\[
c_{\mathrm F}r(n,m,d,\epsilon)
\le R(n,m,d,\epsilon)
\le C_{\mathrm F}r(n,m,d,\epsilon),
\]
where \(r\) is defined in \cref{obj:def:rate-handle}.

Fix an allowed sequence \(\mathbf v\), and write
\[
\begin{gathered}
S_k=n_k\epsilon_k,\qquad
N_k=n_k+m_k,\qquad
\ell_k=\log(\exp(1)+S_k),\\
R_k=R(n_k,m_k,d_k,\epsilon_k),\qquad
r_k=r(n_k,m_k,d_k,\epsilon_k),\\
b_k=b(n_k,\epsilon_k)=\min\{1,S_k^{-1}\},\qquad
z_k=\frac{d_k}{N_k\epsilon_k\ell_k}.
\end{gathered}
\]
The public restrictions give \(S_k>0\), \(N_k>0\), and \(\ell_k>0\). Consequently,
\[
z_k>0,\qquad
r_k=\min\{1,S_k^{-1}+z_k^2\}>0,\qquad
0<c_{\mathrm F}r_k\le R_k\le C_{\mathrm F}r_k.
\]
The same comparison and positivity apply at every auxiliary budget, including zero.
% lean: CausalSmith.Stat.AnnotationRarearmFrontier.annotation_resource_phases

\item \emph{Consistency.}
The comparison gives
\[
0\le r_k\le \frac{R_k}{c_{\mathrm F}},
\qquad
0\le R_k\le C_{\mathrm F}r_k,
\]
so \(R_k\to0\) if and only if \(r_k\to0\). If \(r_k\to0\), then eventually \(r_k<1\), and hence
\[
r_k=S_k^{-1}+z_k^2.
\]
Both summands are nonnegative, so
\[
0<S_k^{-1}\le r_k\longrightarrow0,
\qquad
0\le z_k^2\le r_k\longrightarrow0.
\]
Taking reciprocals in the first limit and nonnegative square roots in the second yields \(S_k\to\infty\) and \(z_k\to0\). Conversely, these two limits imply
\[
0\le r_k\le S_k^{-1}+z_k^2\longrightarrow0.
\]
Since \(N_k\epsilon_k\ell_k>0\),
\[
z_k\longrightarrow0
\quad\Longleftrightarrow\quad
d_k=o(N_k\epsilon_k\ell_k).
\]
Together with \cref{obj:def:phases}, this proves the consistency characterization.
% lean: CausalSmith.Stat.AnnotationRarearmFrontier.consistency_of_frontier_comparison

\item \emph{Oracle membership.}
Suppose \(S_k\to\infty\). Then eventually \(S_k\ge1\), and
\[
b_k=S_k^{-1}.
\]
By \cref{obj:def:phases}, the remaining oracle condition is eventual boundedness above of \(R_k/b_k\).

First suppose that, for some real constant \(M_{\mathrm{or}}\),
\[
\frac{R_k}{b_k}\le M_{\mathrm{or}}
\]
eventually. Define
\[
K_{\mathrm{or}}=\frac{\max\{M_{\mathrm{or}},1\}}{c_{\mathrm F}}>0.
\]
The lower comparison then gives, eventually,
\[
c_{\mathrm F}r_k\le R_k
\le \frac{\max\{M_{\mathrm{or}},1\}}{S_k},
\qquad
r_k\le\frac{K_{\mathrm{or}}}{S_k}.
\]
Eventually \(S_k>K_{\mathrm{or}}\), so \(r_k<1\) and its cap is inactive. Thus
\[
1+S_kz_k^2=S_kr_k\le K_{\mathrm{or}},
\qquad
(z_k\sqrt{S_k})^2\le K_{\mathrm{or}}.
\]
All factors are nonnegative, and therefore
\[
z_k\sqrt{S_k}\le\sqrt{K_{\mathrm{or}}},
\qquad
d_k\le
\sqrt{K_{\mathrm{or}}}\,
\frac{N_k\epsilon_k\ell_k}{\sqrt{S_k}}
\]
eventually. This is the required big-\(O\) condition.

Conversely, that big-\(O\) condition supplies a constant \(K_{\mathrm{dim}}>0\) such that eventually
\[
d_k\le K_{\mathrm{dim}}
\frac{N_k\epsilon_k\ell_k}{\sqrt{S_k}}.
\]
Multiplication by the positive factor
\(\sqrt{S_k}/(N_k\epsilon_k\ell_k)\) gives
\[
z_k\sqrt{S_k}\le K_{\mathrm{dim}},
\qquad
S_kz_k^2\le K_{\mathrm{dim}}^2.
\]
Consequently,
\[
r_k\le S_k^{-1}+z_k^2
\le\frac{1+K_{\mathrm{dim}}^2}{S_k},
\qquad
\frac{R_k}{b_k}
\le C_{\mathrm F}(1+K_{\mathrm{dim}}^2)
\]
eventually. This proves the oracle equivalence.
% lean: CausalSmith.Stat.AnnotationRarearmFrontier.oracle_of_frontier_comparison

\item \emph{Necessary conditions for improvement.}
Define the supervised risk, supervised rate, and rate ratio by
\[
R_{\mathrm{sup},k}=R(n_k,0,d_k,\epsilon_k),\qquad
r_{\mathrm{sup},k}=r(n_k,0,d_k,\epsilon_k),\qquad
\rho_k=\frac{r_k}{r_{\mathrm{sup},k}}.
\]
Their denominators are positive. Applying the uniform comparison at both auxiliary budgets gives
\[
\frac{c_{\mathrm F}}{C_{\mathrm F}}\rho_k
\le \frac{R_k}{R_{\mathrm{sup},k}}
\le \frac{C_{\mathrm F}}{c_{\mathrm F}}\rho_k.
\]
Indeed, the left inequality follows by dividing
\(R_k\ge c_{\mathrm F}r_k\) by
\(R_{\mathrm{sup},k}\le C_{\mathrm F}r_{\mathrm{sup},k}\);
the right inequality follows from the opposite two bounds. Squeezing therefore yields
\[
\mathbf v\in\mathcal I
\quad\Longleftrightarrow\quad
\rho_k\longrightarrow0.
\]
% lean: CausalSmith.Stat.AnnotationRarearmFrontier.improvement_ratio_of_frontier_comparison

Assume \(\rho_k\to0\). Since \(0<r_{\mathrm{sup},k}\le1\),
\[
0\le r_k=\rho_kr_{\mathrm{sup},k}\le\rho_k\longrightarrow0.
\]
The consistency calculation gives \(S_k\to\infty\).
% lean: CausalSmith.Stat.AnnotationRarearmFrontier.phase_improvement_label_divergence

Introduce, for every \(k\),
\[
\xi_k=\frac{d_k}{S_k\ell_k}>0,\qquad
\eta_k=\frac{N_k}{n_k}>0,\qquad
a_{\mathrm{lab},k}=S_k^{-1}>0,\qquad
t_{\mathrm{cap},k}=\min\{1,\xi_k^2\}>0.
\]
The identity \(N_k\epsilon_k=\eta_kS_k\) gives
\[
r_k=\min\{1,a_{\mathrm{lab},k}+(\xi_k/\eta_k)^2\},
\qquad
r_{\mathrm{sup},k}=\min\{1,a_{\mathrm{lab},k}+\xi_k^2\}.
\]

Eventually \(S_k\ge1\) and \(0\le\rho_k\le1/2\). For such \(k\),
\[
a_{\mathrm{lab},k}\le1,\qquad
r_k=\rho_kr_{\mathrm{sup},k}\le\rho_k\le\frac12.
\]
Hence the cap in \(r_k\) is inactive:
\[
r_k=a_{\mathrm{lab},k}+(\xi_k/\eta_k)^2.
\]
In particular,
\[
a_{\mathrm{lab},k}
\le r_k
=\rho_kr_{\mathrm{sup},k}
\le\rho_k(a_{\mathrm{lab},k}+\xi_k^2)
\le\frac12(a_{\mathrm{lab},k}+\xi_k^2).
\]
It follows that
\[
a_{\mathrm{lab},k}\le\xi_k^2,
\qquad
\frac12a_{\mathrm{lab},k}
\le(1-\rho_k)a_{\mathrm{lab},k}
\le\rho_k\xi_k^2.
\]
Using \((\sqrt{S_k})^2=S_k\), we obtain the first inverse-square bound:
\[
(\sqrt{S_k}\xi_k)^{-2}
=\frac{a_{\mathrm{lab},k}}{\xi_k^2}
\le2\rho_k.
\]

For the second bound, the cap satisfies, for every \(k\),
\[
t_{\mathrm{cap},k}\le r_{\mathrm{sup},k}
\le t_{\mathrm{cap},k}+a_{\mathrm{lab},k}.
\]
The lower bound follows by increasing the second argument of the minimum. For the upper bound, if \(\xi_k^2\le1\), then
\[
r_{\mathrm{sup},k}\le a_{\mathrm{lab},k}+\xi_k^2
=a_{\mathrm{lab},k}+t_{\mathrm{cap},k};
\]
if \(\xi_k^2>1\), then
\[
r_{\mathrm{sup},k}\le1=t_{\mathrm{cap},k}
\le t_{\mathrm{cap},k}+a_{\mathrm{lab},k}.
\]
% lean: CausalSmith.Stat.AnnotationRarearmFrontier.phase_supervised_cap_bounds

For the eventual indices under consideration,
\(a_{\mathrm{lab},k}\le1\) and
\(a_{\mathrm{lab},k}\le\xi_k^2\), so
\[
a_{\mathrm{lab},k}\le t_{\mathrm{cap},k},
\qquad
r_{\mathrm{sup},k}\le2t_{\mathrm{cap},k}.
\]
Thus
\[
(\xi_k/\eta_k)^2
\le r_k
=\rho_kr_{\mathrm{sup},k}
\le2\rho_kt_{\mathrm{cap},k}.
\]
The exact budget identity is
\[
\frac{(\xi_k/\eta_k)^2}{t_{\mathrm{cap},k}}
=
\left(\frac{\max\{1,\xi_k\}}{\eta_k}\right)^2
=
\left(\frac{\eta_k}{\max\{1,\xi_k\}}\right)^{-2}.
\]
If \(\xi_k\le1\), the first two expressions both equal
\(\eta_k^{-2}\); if \(\xi_k>1\), they both equal
\((\xi_k/\eta_k)^2\). Therefore,
\[
\left(\frac{\eta_k}{\max\{1,\xi_k\}}\right)^{-2}
\le2\rho_k.
\]
% lean: CausalSmith.Stat.AnnotationRarearmFrontier.phase_improvement_necessary_bounds

Both inverse squares tend to zero. Taking positive square roots and then reciprocals gives
\[
\sqrt{S_k}\xi_k\longrightarrow\infty,
\qquad
\frac{\eta_k}{\max\{1,\xi_k\}}\longrightarrow\infty.
\]
Finally,
\[
\sqrt{S_k}\xi_k
=\frac{d_k}{\sqrt{S_k}\ell_k},
\qquad
\frac{\eta_k}{\max\{1,\xi_k\}}
=
\frac{N_k/n_k}{\max\{1,d_k/(S_k\ell_k)\}},
\]
which proves all three necessary limits.
% lean: CausalSmith.Stat.AnnotationRarearmFrontier.phase_improvement_ratio_necessary

\item \emph{Sufficient conditions for improvement.}
Suppose the three limits in the claimed improvement characterization hold. In the notation just defined, they are
\[
S_k\longrightarrow\infty,\qquad
\sqrt{S_k}\xi_k\longrightarrow\infty,\qquad
\frac{\eta_k}{\max\{1,\xi_k\}}\longrightarrow\infty.
\]
For every \(k\), positivity and
\(t_{\mathrm{cap},k}\le r_{\mathrm{sup},k}\) give
\[
0\le\rho_k
\le
\frac{a_{\mathrm{lab},k}+(\xi_k/\eta_k)^2}
     {t_{\mathrm{cap},k}}
=
\frac{a_{\mathrm{lab},k}}{t_{\mathrm{cap},k}}
+
\frac{(\xi_k/\eta_k)^2}{t_{\mathrm{cap},k}}.
\]
The first quotient satisfies
\[
\frac{a_{\mathrm{lab},k}}{t_{\mathrm{cap},k}}
=
\begin{cases}
(\sqrt{S_k}\xi_k)^{-2},&\xi_k^2\le1,\\
S_k^{-1},&\xi_k^2>1,
\end{cases}
\le S_k^{-1}+(\sqrt{S_k}\xi_k)^{-2}.
\]
Combining this with the budget identity yields
\[
0\le\rho_k
\le
S_k^{-1}
+(\sqrt{S_k}\xi_k)^{-2}
+\left(\frac{\eta_k}{\max\{1,\xi_k\}}\right)^{-2}
\longrightarrow0.
\]
The risk-ratio comparison then proves
\(\mathbf v\in\mathcal I\).
% lean: CausalSmith.Stat.AnnotationRarearmFrontier.phase_improvement_ratio_sufficient

\item \emph{The supervised slice.}
If \(m_k=0\) for every \(k\), then
\[
N_k\epsilon_k=S_k,\qquad
\frac{N_k\epsilon_k\ell_k}{\sqrt{S_k}}
=\frac{S_k\ell_k}{\sqrt{S_k}}
=\sqrt{S_k}\ell_k,
\]
where the final equality uses \(S_k=(\sqrt{S_k})^2\) and
\(\sqrt{S_k}>0\). Substitution into the consistency and oracle characterizations gives the two stated supervised conditions.

Moreover, \(R_{\mathrm{sup},k}=R_k>0\), so
\[
\frac{R_k}{R_{\mathrm{sup},k}}=1
\qquad\text{for every }k.
\]
Its limit is \(1\), and hence \(\mathbf v\notin\mathcal I\).
% lean: CausalSmith.Stat.AnnotationRarearmFrontier.annotation_resource_phases

\item \emph{The equivalent oracle budget.}
For every \(k\) and every constant \(K_{\mathrm{bud}}\ge0\), multiplication or division by the positive quantities
\(\sqrt{S_k}\) and \(\epsilon_k\ell_k\) gives
\[
d_k\le K_{\mathrm{bud}}
\frac{N_k\epsilon_k\ell_k}{\sqrt{S_k}}
\quad\Longleftrightarrow\quad
\frac{d_k\sqrt{S_k}}{\epsilon_k\ell_k}
\le K_{\mathrm{bud}}N_k.
\]
Thus the two eventual inequalities use the same comparison constant, and
\[
d_k=O\!\left(\frac{N_k\epsilon_k\ell_k}{\sqrt{S_k}}\right)
\quad\Longleftrightarrow\quad
\frac{d_k\sqrt{S_k}}{\epsilon_k\ell_k}=O(N_k).
\]
Under \(S_k\to\infty\), the oracle characterization proves the asserted budget equivalence.
% lean: CausalSmith.Stat.AnnotationRarearmFrontier.phase_dimension_budget_iff

\item \emph{Risk order and the uniform auxiliary-budget floor.}
For any allowed experiment, increasing the second argument of the minimum gives
\[
b(n,\epsilon)
=\min\{1,(n\epsilon)^{-1}\}
\le
\min\left\{1,(n\epsilon)^{-1}
+\left(\frac{d}{(n+m)\epsilon\ell}\right)^2\right\}
=r(n,m,d,\epsilon).
\]
Hence, for every \(k\) and every \(m'\in\mathbb N\),
\[
c_{\mathrm F}b_k
\le c_{\mathrm F}r(n_k,m',d_k,\epsilon_k)
\le R(n_k,m',d_k,\epsilon_k).
\]
% lean: CausalSmith.Stat.AnnotationRarearmFrontier.phase_labelBenchmark_le_frontier

Now let \(\mathbf v\in\mathcal O\). Choose a real constant
\(M_{\mathrm{risk}}\) giving an eventual upper bound for \(R_k/b_k\), and define
\[
C_{\mathrm{or}}=\max\{M_{\mathrm{risk}},1\}>0.
\]
Since \(S_k\to\infty\), eventually \(b_k=S_k^{-1}\). On that same tail,
\[
\frac{c_{\mathrm F}}{S_k}\le R_k
\le\frac{M_{\mathrm{risk}}}{S_k}
\le\frac{C_{\mathrm{or}}}{S_k}.
\]
The lower comparison also gives, on a tail independent of \(m'\),
\[
\frac{c_{\mathrm F}}{S_k}
\le R(n_k,m',d_k,\epsilon_k)
\qquad\text{for every }m'\in\mathbb N.
\]
This proves both oracle risk assertions with a positive lower constant independent of the auxiliary budget.
% lean: CausalSmith.Stat.AnnotationRarearmFrontier.annotation_resource_phases

\item \emph{Single-experiment saturation.}
Fix an experiment satisfying the stated public restrictions, and use
\[
N=n+m,\qquad
S=n\epsilon,\qquad
\ell=\log(\exp(1)+S).
\]
All three denominator factors are positive. If \(N\epsilon\ell\le d\), then
\[
\frac{d}{N\epsilon\ell}\ge1,
\qquad
S^{-1}+\left(\frac{d}{N\epsilon\ell}\right)^2\ge1,
\qquad
r(n,m,d,\epsilon)=1.
\]
The uniform comparison therefore yields
\[
c_{\mathrm F}\le R(n,m,d,\epsilon).
\]
% lean: CausalSmith.Stat.AnnotationRarearmFrontier.phase_frontier_saturation

For the upper bound, the laws in
\cref{obj:thm:rare-label-floor} ensure that the model class is nonempty. For each law \(P\) in that class, the binary outcome conditional means lie in \([0,1]\). The target formula in \cref{obj:def:target} thus gives
\[
|\tau(P)|
\le\sum_{j=1}^{d}p_j(P)
       |\mu_{1j}(P)-\mu_{0j}(P)|
\le\sum_{j=1}^{d}p_j(P)=1.
\]
Define the constant original-record rule by
\[
T_{\mathrm{zero}}(\mathcal D,U)=0
\]
for every data realization and seed. This rule is Borel measurable and takes values in \([-1,1]\). Its squared error is the constant \(\tau(P)^2\), so \cref{obj:def:risk} gives
\[
R(n,m,d,\epsilon)
\le
\sup_{P\in\mathcal M_{d,\epsilon}}
\mathbb E_P\!\left[
(T_{\mathrm{zero}}(\mathcal D,U)-\tau(P))^2
\right]
=
\sup_{P\in\mathcal M_{d,\epsilon}}\tau(P)^2
\le1.
\]
This proves the saturation bounds with the universal constant \(c_{\mathrm F}\).
% lean: CausalSmith.Stat.AnnotationRarearmFrontier.phase_minimaxRisk_le_one

\item \emph{Single-experiment oracle plateau.}
Suppose
\[
\frac{d\sqrt S}{\epsilon\ell}\le N.
\]
Multiplication by \(\epsilon\ell>0\) and division by
\(N\epsilon\ell>0\) give
\[
d\sqrt S\le N\epsilon\ell,
\qquad
\frac{d}{N\epsilon\ell}\sqrt S\le1.
\]
Squaring, using nonnegativity and \((\sqrt S)^2=S\), and dividing by \(S>0\), we obtain
\[
\left(\frac{d}{N\epsilon\ell}\right)^2\le S^{-1}.
\]
If \(S^{-1}\le1\), then
\[
r(n,m,d,\epsilon)
\le S^{-1}+\left(\frac{d}{N\epsilon\ell}\right)^2
\le2S^{-1}
=2b(n,\epsilon).
\]
If \(S^{-1}>1\), then
\[
b(n,\epsilon)=1,\qquad
r(n,m,d,\epsilon)\le1\le2b(n,\epsilon).
\]
Together with the universal lower comparison
\(b(n,\epsilon)\le r(n,m,d,\epsilon)\), these cases give
\[
c_{\mathrm F}b(n,\epsilon)
\le R(n,m,d,\epsilon)
\le C_{\mathrm F}r(n,m,d,\epsilon)
\le2C_{\mathrm F}b(n,\epsilon).
\]
The constants \(c_{\mathrm F}\) and \(2C_{\mathrm F}\) are positive and independent of all public indices.
% lean: CausalSmith.Stat.AnnotationRarearmFrontier.phase_frontier_plateau
\end{enumerate}
\end{proof}

\bibliographystyle{plainnat}

\bibliography{references}

\end{document}
